% Generated mechanically from frozen input inputs/p03/ja/tex/sites-modules.tex.
% Do not edit this generated file; edit the bound locale source instead.

% OK, start here.
%

\setcounter{chapter}{17}
\chapter{サイト上の加群}



\phantomsection
\label{sites-modules-section-phantom}


\section{序論}
\label{sites-modules-section-introduction}

\noindent
本章では，環付きトポスまたは環付きサイト上の加群の層に関する
基本概念を展開する。まずアーベル群の層に関するいくつかの基本事実を扱い，
その後，環付きサイトと環付きトポスを導入する。さらに，空間上の層の章に
おける $\mathcal{O}$-加群の（前）層の議論と同様に，
$\mathcal{O}$-加群の（前）層に関するごく基本的な概念を調べる。
これらを準備した後，加群の層の章における議論の多くをここで再現する
（《加群の層》第 \ref{modules-section-introduction} 節参照）。
基本的な参考文献は \cite{FAC}, \cite{EGA}, \cite{SGA4} である。






\section{アーベル群の前層}
\label{sites-modules-section-abelian-pre-sheaves}

\noindent
$\mathcal{C}$ を圏とする。
アーベル群の前層は《サイト》第 \ref{sites-section-presheaves} 節および
第 \ref{sites-section-sheaves} 節で導入され，《サイト》第
\ref{sites-section-sheaves-algebraic-structures} 節でさらに論じられた。
最後の参照先の規約に従い，アーベル群の前層を，各切断集合に制限写像と
両立する加法則を備えた集合の前層とみなす。
$\mathcal{C}$ 上のアーベル群の前層の圏を
$\textit{PAb}(\mathcal{C})$ と書くことを思い出そう。

\medskip\noindent
圏 $\textit{PAb}(\mathcal{C})$ は，《ホモロジー》定義
\ref{homology-definition-abelian-category} の意味でアーベル圏である。
前層の射 $\varphi : \mathcal{G}_1 \to \mathcal{G}_2$ が与えられたとき，
$\varphi$ の核はアーベル群の前層
$U \mapsto \Ker(\mathcal{G}_1(U) \to \mathcal{G}_2(U))$ であり，
$\varphi$ の余核は前層
$U \mapsto \Coker(\mathcal{G}_1(U) \to \mathcal{G}_2(U))$ である。
アーベル群の圏はアーベル圏であり，各 $U$ 上で $\Coim = \Im$ が
成り立つから，ここでも $\Coim = \Im$ である。
アーベル群の前層の列
$$
\mathcal{G}_1 \longrightarrow
\mathcal{G}_2 \longrightarrow
\mathcal{G}_3
$$
が完全であることと，
$\mathcal{G}_1(U) \to \mathcal{G}_2(U) \to \mathcal{G}_3(U)$
がすべての $U \in \Ob(\mathcal{C})$ に対してアーベル群の完全列で
あることとは同値である。検証は読者に委ねる。

\begin{lemma}
\label{sites-modules-lemma-limits-colimits-abelian-presheaves}
$\mathcal{C}$ を圏とする。
\begin{enumerate}
\item $\textit{PAb}(\mathcal{C})$ にはすべての極限と余極限が存在する。
\item すべての極限と余極限は，$\mathcal{C}$ の対象上で切断を取る操作と
可換である。
\end{enumerate}
\end{lemma}

\begin{proof}
$\mathcal{I} \to \textit{PAb}(\mathcal{C})$，
$i \mapsto \mathcal{F}_i$ を図式とする。アーベル群の前層 $L$ と $C$ を
単に次の規則で定義すればよい：
$$
L : U \longmapsto \lim_i \mathcal{F}_i(U)
$$
および
$$
C : U \longmapsto \colim_i \mathcal{F}_i(U).
$$
各 $U$ 上の切断の群における対応する射を用いれば，アーベル群の前層の射
$L \to \mathcal{F}_i$ および $\mathcal{F}_i \to C$ があることは明らかである。
これらの射を備えた $L$ と $C$ が，$\textit{PAb}(\mathcal{C})$ における
この図式の極限と余極限であることは直ちに確認できる。
これで (1) と (2) が従う。詳細は省略する。
\end{proof}


\section{アーベル群の層}
\label{sites-modules-section-abelian-sheaves}

\noindent
$\mathcal{C}$ をサイトとする。
$\mathcal{C}$ 上のアーベル群の層の圏を
$\textit{Ab}(\mathcal{C})$ と書く。これは，台となる集合の前層が層であるような
アーベル群の前層からなる $\textit{PAb}(\mathcal{C})$ の充満部分圏である。
《サイト》第 \ref{sites-section-sheaves-algebraic-structures} 節の性質
($\alpha$) -- ($\zeta$) が成り立つ。《サイト》命題
\ref{sites-proposition-functoriality-algebraic-structures-topoi} を参照せよ。
特に，包含関手
$\textit{Ab}(\mathcal{C}) \to \textit{PAb}(\mathcal{C})$
は左随伴，すなわち層化関手
$\mathcal{G} \mapsto \mathcal{G}^\#$
をもつ。

\medskip\noindent
以下の証明を読むよりも，まず補題を手元で証明してみることを読者に勧める。

\begin{lemma}
\label{sites-modules-lemma-abelian-abelian}
$\mathcal{C}$ をサイトとし，
$\varphi : \mathcal{F} \to \mathcal{G}$ を $\mathcal{C}$ 上の
アーベル群の層の射とする。
\begin{enumerate}
\item 圏 $\textit{Ab}(\mathcal{C})$ はアーベル圏である。
\item $\varphi$ の核 $\Ker(\varphi)$ は，前層の射としての
$\varphi$ の核と同じである。
\item 射 $\varphi$ が単射であること
（《ホモロジー》定義 \ref{homology-definition-injective-surjective}）と，
$\varphi$ が前層の写像として単射であること
（《サイト》定義 \ref{sites-definition-presheaves-injective-surjective}）と，
$\varphi$ が層の写像として単射であること
（《サイト》定義 \ref{sites-definition-sheaves-injective-surjective}）とは
同値である。
\item $\varphi$ の余核 $\Coker(\varphi)$ は，前層の射としての
$\varphi$ の余核を層化したものである。
\item 射 $\varphi$ が全射であること
（《ホモロジー》定義 \ref{homology-definition-injective-surjective}）と，
$\varphi$ が層の写像として全射であること
（《サイト》定義 \ref{sites-definition-sheaves-injective-surjective}）とは
同値である。
\item アーベル群の層の複体
$$
\mathcal{F} \to \mathcal{G} \to \mathcal{H}
$$
が $\mathcal{G}$ において完全であることと，すべての
$U \in \Ob(\mathcal{C})$ および $\mathcal{H}(U)$ で零に写るすべての
$s \in \mathcal{G}(U)$ に対し，$\mathcal{C}$ における被覆
$\{U_i \to U\}_{i \in I}$ が存在して，各 $s|_{U_i}$ が
$\mathcal{F}(U_i) \to \mathcal{G}(U_i)$ の像に属することとは同値である。
\end{enumerate}
\end{lemma}

\begin{proof}
《ホモロジー》補題 \ref{homology-lemma-adjoint-get-abelian} を，
圏 $\mathcal{A} = \textit{Ab}(\mathcal{C})$，
$\mathcal{B} = \textit{PAb}(\mathcal{C})$，および関手
$a : \mathcal{A} \to \mathcal{B}$（包含）と
$b : \mathcal{B} \to \mathcal{A}$（層化）に適用できると主張する。
《ホモロジー》補題 \ref{homology-lemma-adjoint-get-abelian} の仮定を確認しよう。
仮定 (1) は，$\mathcal{A}$ と $\mathcal{B}$ が
加法圏であり，$a$ と $b$ が加法関手で，$a$ が $b$ の右随伴であるという
ものである。初めの二つは明らかであり，随伴性は《サイト》第
\ref{sites-section-sheaves-algebraic-structures} 節 ($\epsilon$) である。
仮定 (2) は，$\textit{PAb}(\mathcal{C})$ がアーベル圏であることと，
層化が左完全であることをいう。前者は第
\ref{sites-modules-section-abelian-pre-sheaves} 節で見たとおりであり，後者は
《サイト》第 \ref{sites-section-sheaves-algebraic-structures} 節
($\zeta$) である。最後の仮定 $ba \cong \text{id}_\mathcal{A}$ は
《サイト》第 \ref{sites-section-sheaves-algebraic-structures} 節
($\delta$) である。したがって《ホモロジー》補題
\ref{homology-lemma-adjoint-get-abelian} が適用でき，
$\textit{Ab}(\mathcal{C})$ はアーベル圏である。

\medskip\noindent
《ホモロジー》補題 \ref{homology-lemma-adjoint-get-abelian} の証明では，
$\Ker(\varphi)$ と $\Coker(\varphi)$ が，アーベル群の前層の射としての
$\varphi$ の核と余核を層化したものに等しいことが示されている。
これで (4) が従う。核は等化子，すなわち極限であり，層化は有限極限と
可換するから，(2) も従う。

\medskip\noindent
(2) から (3) が従う。層化の記述により，(4) から (5) が従う。
(6) の完全性の特徴づけは，(2)，(5)，および列が完全であることと
$\Im(\mathcal{F} \to \mathcal{G}) =
\Ker(\mathcal{G} \to \mathcal{H})$
が成り立つこととの同値性から従う。
\end{proof}

\noindent
補題の (6) は，アーベル群の層の列
$$
\mathcal{F}_1 \longrightarrow
\mathcal{F}_2 \longrightarrow
\mathcal{F}_3
$$
が完全であることと，
$U \mapsto \Im(\mathcal{F}_1(U) \to \mathcal{F}_2(U))$
を層化したものが $\mathcal{F}_2 \to \mathcal{F}_3$ の核に等しいこととの
同値性として言い換えることもできる。

\begin{lemma}
\label{sites-modules-lemma-limits-colimits-abelian-sheaves}
$\mathcal{C}$ をサイトとする。
\begin{enumerate}
\item $\textit{Ab}(\mathcal{C})$ にはすべての極限と余極限が存在する。
\item 極限は，$\mathcal{C}$ 上のアーベル群の前層として取った対応する
極限と同じである（すなわち，$\mathcal{C}$ の対象上で切断を取る操作と
可換する）。
\item 有限直和は，$\mathcal{C}$ 上のアーベル群の前層の圏で取った
対応する有限直和と同じである。
\item 余極限は，アーベル群の前層の圏で取った対応する余極限を層化した
ものである。
\item フィルター付き余極限は完全である。
\end{enumerate}
\end{lemma}

\begin{proof}
補題 \ref{sites-modules-lemma-limits-colimits-abelian-presheaves} により，
アーベル群の前層の極限と余極限は存在し，各対象上の切断のレベルで
極限と余極限を取ることによって記述される。

\medskip\noindent
$\mathcal{I} \to \textit{Ab}(\mathcal{C})$，
$i \mapsto \mathcal{F}_i$ を図式とする。
$\lim_i \mathcal{F}_i$ を，アーベル群の前層として取ったこの図式の
極限とする。《サイト》補題 \ref{sites-lemma-limit-sheaf} により，
これはアーベル群の層である。すると $\lim_i \mathcal{F}_i$ が
$\textit{Ab}(\mathcal{C})$ における図式の極限であることは容易に分かる。
これで極限の存在と (2) が従う。

\medskip\noindent
《圏論》補題 \ref{categories-lemma-adjoint-exact} と，層化が包含関手の
左随伴であることにより，$\colim_i \mathcal{F}$ は存在し，
$\textit{PAb}(\mathcal{C})$ における余極限を層化したものになる。
これで余極限の存在と (4) が従う。

\medskip\noindent
任意のアーベル圏において有限直和は有限直積と同じである。
したがって (2) から (3) が従う。

\medskip\noindent
最後に (5) を証明する。この主張は，有向集合 $I$ 上のアーベル群の層の
完全列の系
$0 \to \mathcal{F}_i \to \mathcal{G}_i \to \mathcal{H}_i \to 0$
が与えられたとき，列
$0 \to \colim \mathcal{F}_i \to \colim \mathcal{G}_i \to
\colim \mathcal{H}_i \to 0$
も完全であるという意味である。《ホモロジー》補題
\ref{homology-lemma-check-exactness} と余極限の定義を用いる形式的な議論により，
列
$\colim \mathcal{F}_i \to \colim \mathcal{G}_i \to \colim \mathcal{H}_i \to 0$
は完全である。各写像 $\mathcal{F}_i \to \mathcal{G}_i$ が単射なので，
前層の余極限の間の写像は単射であり，
$\colim \mathcal{F}_i \to \colim \mathcal{G}_i$ はその写像を層化したものに
なることに注意せよ。層化は完全だから，結論が従う。
\end{proof}







\section{自由アーベル群の前層}
\label{sites-modules-section-free-abelian-presheaf}

\noindent
以下の定義で用いる記法を準備するため，集合 $S$ 上の自由アーベル群を
\footnote{他の章では，記法 $\mathbf{Z}[S]$ が $S$ を変数とする
$\mathbf{Z}$ 上の多項式環を表すこともある。}
$\mathbf{Z}[S] = \bigoplus_{s \in S} \mathbf{Z}$ と書くことにする。
これは性質
$$
\Mor_{\textit{Ab}}(\mathbf{Z}[S], A)
=
\Mor_{\textit{Sets}}(S, A)
$$
によって特徴づけられる。言い換えれば，構成
$S \mapsto \mathbf{Z}[S]$ は忘却関手
$\textit{Ab} \to \textit{Sets}$ の左随伴である。

\begin{definition}
\label{sites-modules-definition-free-abelian-presheaf-on}
$\mathcal{C}$ を圏とし，$\mathcal{G}$ を集合の前層とする。
$\mathcal{G}$ 上の {\it 自由アーベル群の前層}
$\mathbf{Z}_\mathcal{G}$ とは，規則
$$
U \longmapsto \mathbf{Z}[\mathcal{G}(U)].
$$
で定義されるアーベル群の前層である。$\mathcal{C}$ の対象 $X$ に付随する
表現可能前層 $\mathcal{G}=h_X$ の場合には，
$\mathbf{Z}_X = \mathbf{Z}_{h_X}$ と書く。すなわち，
$$
\mathbf{Z}_X(U) = \mathbf{Z}[\Mor_\mathcal{C}(U, X)].
$$
\end{definition}

\noindent
この構成が前層 $\mathcal{G}$ に関して関手的であることは明らかである。
実際，これは忘却関手
$\textit{PAb}(\mathcal{C}) \to \textit{PSh}(\mathcal{C})$
の左随伴である。正確には次が成り立つ。

\begin{lemma}
\label{sites-modules-lemma-obvious-adjointness}
$\mathcal{C}$ を圏とする。$\mathcal{G}$ と $\mathcal{F}$ を集合の前層，
$\mathcal{A}$ をアーベル群の前層，$U$ を $\mathcal{C}$ の対象とする。
このとき，
\begin{align*}
\Mor_{\textit{PSh}(\mathcal{C})}(h_U, \mathcal{F})
& =
\mathcal{F}(U), \\
\Mor_{\textit{PAb}(\mathcal{C})}(\mathbf{Z}_\mathcal{G}, \mathcal{A})
& =
\Mor_{\textit{PSh}(\mathcal{C})}(\mathcal{G}, \mathcal{A}), \\
\Mor_{\textit{PAb}(\mathcal{C})}(\mathbf{Z}_U, \mathcal{A})
& =
\mathcal{A}(U).
\end{align*}
これらの等式はすべて関手的である。
\end{lemma}

\begin{proof}
省略する。
\end{proof}

\begin{lemma}
\label{sites-modules-lemma-coproduct-sum-free-abelian-presheaf}
$\mathcal{C}$ を圏，$I$ を集合とする。各 $i \in I$ に対し，
$\mathcal{G}_i$ を集合の前層とする。このとき
$$
\mathbf{Z}_{\coprod_i \mathcal{G}_i}
=
\bigoplus\nolimits_{i \in I} \mathbf{Z}_{\mathcal{G}_i}
$$
が $\textit{PAb}(\mathcal{C})$ において成り立つ。
\end{lemma}

\begin{proof}
省略する。
\end{proof}



\section{自由アーベル群の層}
\label{sites-modules-section-free-abelian-sheaf}

\noindent
集合の層上の自由アーベル群の層を次のように定義する。

\begin{definition}
\label{sites-modules-definition-free-abelian-sheaf-on}
$\mathcal{C}$ をサイトとし，$\mathcal{G}$ を集合の前層とする。
$\mathcal{G}$ 上の {\it 自由アーベル群の層}
$\mathbf{Z}_\mathcal{G}^\#$ とは，$\mathcal{G}$ 上の自由アーベル群の
前層を層化して得られるアーベル群の層
$\mathbf{Z}_\mathcal{G}^\#$ のことである。$\mathcal{C}$ の対象 $X$ に
付随する表現可能前層 $\mathcal{G}=h_X$ の場合には
$\mathbf{Z}_X^\#$ と書く。
\end{definition}

\noindent
この構成が前層 $\mathcal{G}$ に関して関手的であることは明らかである。
実際，これは忘却関手
$\textit{Ab}(\mathcal{C}) \to \Sh(\mathcal{C})$
の左随伴を与える。正確には次が成り立つ。

\begin{lemma}
\label{sites-modules-lemma-obvious-adjointness-sheaves}
$\mathcal{C}$ をサイトとする。$\mathcal{G}$ と $\mathcal{F}$ を集合の層，
$\mathcal{A}$ をアーベル群の層，$U$ を $\mathcal{C}$ の対象とする。
このとき，
\begin{align*}
\Mor_{\Sh(\mathcal{C})}(h_U^\#, \mathcal{F})
& =
\mathcal{F}(U), \\
\Mor_{\textit{Ab}(\mathcal{C})}(\mathbf{Z}_\mathcal{G}^\#,
\mathcal{A})
& =
\Mor_{\Sh(\mathcal{C})}(\mathcal{G}, \mathcal{A}), \\
\Mor_{\textit{Ab}(\mathcal{C})}(\mathbf{Z}_U^\#, \mathcal{A})
& =
\mathcal{A}(U).
\end{align*}
これらの等式はすべて関手的である。
\end{lemma}

\begin{proof}
省略する。
\end{proof}

\begin{lemma}
\label{sites-modules-lemma-may-sheafify-before-abelianize}
$\mathcal{C}$ をサイトとし，$\mathcal{G}$ を集合の前層とする。このとき
$\mathbf{Z}_\mathcal{G}^\# = (\mathbf{Z}_{\mathcal{G}^\#})^\#$ である。
\end{lemma}

\begin{proof}
省略する。
\end{proof}








\section{環付きサイト}
\label{sites-modules-section-ringed-sites}

\noindent
本章では主として，環付きサイト上の加群の層を扱う。
そこで，まずこの概念を定義する。

\begin{definition}
\label{sites-modules-definition-ringed-site}
環付きサイトについて次のように定義する。
\begin{enumerate}
\item {\it 環付きサイト}とは，$\mathcal{C}$ がサイトで，
$\mathcal{O}$ が $\mathcal{C}$ 上の環の層であるような組
$(\mathcal{C}, \mathcal{O})$ のことである。層 $\mathcal{O}$ を
環付きサイトの {\it 構造層}という。
\item $(\mathcal{C}, \mathcal{O})$ と
$(\mathcal{C}', \mathcal{O}')$ を環付きサイトとする。
{\it 環付きサイトの射}
$$
(f, f^\sharp) :
(\mathcal{C}, \mathcal{O})
\longrightarrow
(\mathcal{C}', \mathcal{O}')
$$
は，サイトの射 $f : \mathcal{C} \to \mathcal{C}'$
（《サイト》定義 \ref{sites-definition-morphism-sites} 参照）と，
環の層の射
$f^\sharp : f^{-1}\mathcal{O}' \to \mathcal{O}$ の組として与えられる。
随伴により，後者は環の層の射
$f^\sharp : \mathcal{O}' \to f_*\mathcal{O}$
と同じことである。
\item
$(f, f^\sharp) :
(\mathcal{C}_1, \mathcal{O}_1) \to (\mathcal{C}_2, \mathcal{O}_2)$ と
$(g, g^\sharp) :
(\mathcal{C}_2, \mathcal{O}_2) \to (\mathcal{C}_3, \mathcal{O}_3)$
を環付きサイトの射とする。このとき {\it 環付きサイトの射の合成}を
規則
$$
(g, g^\sharp) \circ (f, f^\sharp) = (g \circ f, f^\sharp \circ g^\sharp).
$$
で定義する。ここで《サイト》定義
\ref{sites-definition-composition-morphisms-sites} で定義された
サイトの射の合成を用い，$f^\sharp \circ g^\sharp$ は環の層の射
$$
\mathcal{O}_3 \xrightarrow{g^\sharp} g_*\mathcal{O}_2
\xrightarrow{g_*f^\sharp} g_*f_*\mathcal{O}_1 = (g \circ f)_*\mathcal{O}_1
$$
を表す。
\end{enumerate}
\end{definition}






\section{環付きトポス}
\label{sites-modules-section-ringed-topoi}

\noindent
環付きトポスは環付きサイトと同じ種類の対象であるが，
環付きトポスの射の概念は環付きサイトの射の概念とは異なる。

\begin{definition}
\label{sites-modules-definition-ringed-topos}
環付きトポスについて次のように定義する。
\begin{enumerate}
\item {\it 環付きトポス}とは，$\mathcal{C}$ がサイトで，
$\mathcal{O}$ が $\mathcal{C}$ 上の環の層であるような組
$(\Sh(\mathcal{C}), \mathcal{O})$
のことである。層 $\mathcal{O}$ を環付きトポスの {\it 構造層}という。
\item $(\Sh(\mathcal{C}), \mathcal{O})$ と
$(\Sh(\mathcal{C}'), \mathcal{O}')$ を環付きトポスとする。
{\it 環付きトポスの射}
$$
(f, f^\sharp) :
(\Sh(\mathcal{C}), \mathcal{O})
\longrightarrow
(\Sh(\mathcal{C}'), \mathcal{O}')
$$
は，トポスの射 $f : \Sh(\mathcal{C}) \to \Sh(\mathcal{C}')$
（《サイト》定義 \ref{sites-definition-topos} 参照）と，環の層の射
$f^\sharp : f^{-1}\mathcal{O}' \to \mathcal{O}$ の組として与えられる。
随伴により，後者は環の層の射
$f^\sharp : \mathcal{O}' \to f_*\mathcal{O}$
と同じことである。
\item
$(f, f^\sharp) :
(\Sh(\mathcal{C}_1), \mathcal{O}_1)
\to (\Sh(\mathcal{C}_2), \mathcal{O}_2)$ と
$(g, g^\sharp) :
(\Sh(\mathcal{C}_2), \mathcal{O}_2) \to
(\Sh(\mathcal{C}_3), \mathcal{O}_3)$
を環付きトポスの射とする。このとき {\it 環付きトポスの射の合成}を
規則
$$
(g, g^\sharp) \circ (f, f^\sharp) = (g \circ f, f^\sharp \circ g^\sharp).
$$
で定義する。ここで《サイト》定義 \ref{sites-definition-topos} で
定義されたトポスの射の合成を用い，$f^\sharp \circ g^\sharp$ は環の層の射
$$
\mathcal{O}_3 \xrightarrow{g^\sharp} g_*\mathcal{O}_2
\xrightarrow{g_*f^\sharp} g_*f_*\mathcal{O}_1 = (g \circ f)_*\mathcal{O}_1
$$
を表す。
\end{enumerate}
\end{definition}

\noindent
任意の環付きトポスの射は，環付きトポスの同値と，環付きサイトの射に
付随する環付きトポスの射との合成である。正確には次が成り立つ。

\begin{lemma}
\label{sites-modules-lemma-morphism-ringed-topoi-comes-from-morphism-ringed-sites}
$(f, f^\sharp) :
(\Sh(\mathcal{C}), \mathcal{O}_\mathcal{C})
\to (\Sh(\mathcal{D}), \mathcal{O}_\mathcal{D})$
を環付きトポスの射とする。このとき，次の分解が存在する：
$$
\xymatrix{
(\Sh(\mathcal{C}), \mathcal{O}_\mathcal{C})
\ar[rr]_{(f, f^\sharp)}
\ar[d]_{(g, g^\sharp)}
& &
(\Sh(\mathcal{D}), \mathcal{O}_\mathcal{D}) \ar[d]^{(e, e^\sharp)}
\\
(\Sh(\mathcal{C}'), \mathcal{O}_{\mathcal{C}'})
\ar[rr]^{(h, h^\sharp)} & &
(\Sh(\mathcal{D}'), \mathcal{O}_{\mathcal{D}'})
}
$$
ここで
\begin{enumerate}
\item $g : \Sh(\mathcal{C}) \to \Sh(\mathcal{C}')$
は特殊余連続関手 $\mathcal{C} \to \mathcal{C}'$ によって誘導される
トポスの同値である（《サイト》定義
\ref{sites-definition-special-cocontinuous-functor} 参照）。
\item $e : \Sh(\mathcal{D}) \to \Sh(\mathcal{D}')$
は特殊余連続関手 $\mathcal{D} \to \mathcal{D}'$ によって誘導される
トポスの同値である（《サイト》定義
\ref{sites-definition-special-cocontinuous-functor} 参照）。
\item $\mathcal{O}_{\mathcal{C}'} = g_*\mathcal{O}_\mathcal{C}$
であり，$g^\sharp$ は明らかな写像である。
\item $\mathcal{O}_{\mathcal{D}'} = e_*\mathcal{O}_\mathcal{D}$
であり，$e^\sharp$ は明らかな写像である。
\item サイト $\mathcal{C}'$ と $\mathcal{D}'$ は終対象とファイバー積
（すなわち，すべての有限極限）をもつ。
\item $h$ は，すべての有限極限と可換する連続関手
$u : \mathcal{D}' \to \mathcal{C}'$ によって誘導されるサイトの射である
（すなわち，《サイト》命題 \ref{sites-proposition-get-morphism} の仮定を
満たす）。
\item $\mathcal{C}$ 上の任意の層の族 $\mathcal{F}_i$
（それぞれ $\mathcal{D}$ 上の任意の層の族 $\mathcal{G}_j$）が
与えられたとき，それぞれを $\mathcal{C}'$
（それぞれ $\mathcal{D}'$）上の表現可能層と仮定できる。
\end{enumerate}
さらに，$(f, f^\sharp)$ が環付きトポスの同値ならば，
$\mathcal{C}' = \mathcal{D}'$,
$\mathcal{O}_{\mathcal{C}'} = \mathcal{O}_{\mathcal{D}'}$
かつ $(h, h^\sharp)$ が恒等射となるように図式を選べる。
\end{lemma}

\begin{proof}
これは《サイト》補題
\ref{sites-lemma-morphism-topoi-comes-from-morphism-sites}，および注
\ref{sites-remark-morphism-topoi-comes-from-morphism-sites} と
\ref{sites-remark-equivalence-topoi-comes-from-morphism-sites}
から従う。環の層も併せて追跡すればよい。詳細は省略する。
\end{proof}







\section{環付きトポスの 2-射}
\label{sites-modules-section-2-category}

\noindent
本節では，環付きトポスの $2$-射の概念を簡潔に扱う。

\begin{definition}
\label{sites-modules-definition-2-morphism-ringed-topoi}
$f, g :
(\Sh(\mathcal{C}), \mathcal{O}_\mathcal{C})
\to
(\Sh(\mathcal{D}), \mathcal{O}_\mathcal{D})$
を環付きトポスの二つの射とする。{\it $f$ から $g$ への 2-射}とは，
図式
$$
\xymatrix{
& \mathcal{O}_\mathcal{D}
\ar[ld]_{f^\sharp}
\ar[rd]^{g^\sharp} \\
f_*\mathcal{O}_\mathcal{C} \ar[rr]^t & &
g_*\mathcal{O}_\mathcal{C}
}
$$
が可換となるような関手の変換 $t : f_* \to g_*$ のことである。
\end{definition}

\noindent
$t$ を図式的に次のように表すことがある：
$$
\xymatrix{
(\Sh(\mathcal{C}), \mathcal{O}_\mathcal{C})
\rrtwocell^f_g{t}
&
&
(\Sh(\mathcal{D}), \mathcal{O}_\mathcal{D})
}
$$
《サイト》第 \ref{sites-section-2-category} 節と同様に，
2-射 $t : f_* \to g_*$ を与えることは，図式
$$
\xymatrix{
f^{-1}\mathcal{O}_\mathcal{D}
\ar[rd]_{f^\sharp}  & &
g^{-1}\mathcal{O}_\mathcal{D} \ar[ll]^t \ar[ld]^{g^\sharp} \\
& \mathcal{O}_\mathcal{C}
}
$$
が可換となるような $t : g^{-1} \to f^{-1}$（通常，同じ記号で表す）を
与えることと同値である。また《サイト》第
\ref{sites-section-2-category} 節と同様に，前掲箇所で説明した水平合成と
垂直合成により，厳密 $2$-圏の公理が成り立つ。










\section{加群の前層}
\label{sites-modules-section-presheaves-modules}

\noindent
$\mathcal{C}$ を圏とし，$\mathcal{O}$ を $\mathcal{C}$ 上の環の前層とする。
ここまで $\mathcal{O}$-加群の前層をまだ定義していなかったので，
ここで定義する。

\begin{definition}
\label{sites-modules-definition-presheaf-modules}
$\mathcal{C}$ を圏とし，$\mathcal{O}$ を $\mathcal{C}$ 上の
環の前層とする。
\begin{enumerate}
\item {\it $\mathcal{O}$-加群の前層}とは，アーベル群の前層
$\mathcal{F}$ と，集合の前層の写像
$$
\mathcal{O} \times \mathcal{F} \longrightarrow \mathcal{F}
$$
の組であって，$\mathcal{C}$ の各対象 $U$ に対して写像
$\mathcal{O}(U) \times \mathcal{F}(U) \to \mathcal{F}(U)$
がアーベル群 $\mathcal{F}(U)$ 上に $\mathcal{O}(U)$-加群構造を
定めるものをいう。
\item {\it $\mathcal{O}$-加群の前層の射
$\varphi : \mathcal{F} \to \mathcal{G}$}とは，図式
$$
\xymatrix{
\mathcal{O} \times \mathcal{F} \ar[r] \ar[d]_{\text{id} \times \varphi} &
\mathcal{F} \ar[d]^{\varphi} \\
\mathcal{O} \times \mathcal{G} \ar[r] &
\mathcal{G}
}
$$
が可換となるようなアーベル群の前層の射
$\varphi : \mathcal{F} \to \mathcal{G}$ のことである。
\item 上のような $\mathcal{O}$-加群の射全体の集合を
$\Hom_\mathcal{O}(\mathcal{F}, \mathcal{G})$ と書く。
\item $\mathcal{O}$-加群の前層の圏を
$\textit{PMod}(\mathcal{O})$ と書く。
\end{enumerate}
\end{definition}

\noindent
$\mathcal{O}_1 \to \mathcal{O}_2$ を圏 $\mathcal{C}$ 上の環の前層の射と
する。このとき，$\mathcal{F}$ が $\mathcal{O}_2$-加群の前層ならば，
合成
$$
\mathcal{O}_1 \times \mathcal{F}
\to
\mathcal{O}_2 \times \mathcal{F}
\to
\mathcal{F}.
$$
を用いることにより，$\mathcal{F}$ を $\mathcal{O}_1$-加群の前層と
みなせる。このスカラー制限を明示するため，
$\mathcal{F}_{\mathcal{O}_1}$ と書くことがある。これを
{\it $\mathcal{F}$ の制限}と呼ぶ。制限関手
$$
\textit{PMod}(\mathcal{O}_2)
\longrightarrow
\textit{PMod}(\mathcal{O}_1)
$$

\medskip\noindent
を得る。一方，$\mathcal{O}_1$-加群の前層 $\mathcal{G}$ が
与えられたとき，$\mathcal{O}_2$-加群の前層
$\mathcal{O}_2 \otimes_{p, \mathcal{O}_1} \mathcal{G}$
を規則
$$
U \longmapsto
\left(\mathcal{O}_2 \otimes_{p, \mathcal{O}_1} \mathcal{G}\right)(U)
=
\mathcal{O}_2(U) \otimes_{\mathcal{O}_1(U)} \mathcal{G}(U)
$$
で構成できる。ここで $U \in \Ob(\mathcal{C})$ であり，制限写像は
明らかなものを用いる。添字 $p$ は「点」ではなく「前層」を表す。
この前層をテンソル積前層と呼ぶ。{\it 環の変更}関手
$$
\textit{PMod}(\mathcal{O}_1)
\longrightarrow
\textit{PMod}(\mathcal{O}_2)
$$
を得る。

\begin{lemma}
\label{sites-modules-lemma-adjointness-tensor-restrict-presheaves}
上の $\mathcal{C}$，$\mathcal{O}_1 \to \mathcal{O}_2$，
$\mathcal{F}$，$\mathcal{G}$ に対し，標準的な全単射
$$
\Hom_{\mathcal{O}_1}(\mathcal{G}, \mathcal{F}_{\mathcal{O}_1})
=
\Hom_{\mathcal{O}_2}(
\mathcal{O}_2 \otimes_{p, \mathcal{O}_1} \mathcal{G},
\mathcal{F}
)
$$
が存在する。言い換えれば，上で定義した制限関手と環の変更関手は
互いに随伴である。
\end{lemma}

\begin{proof}
環の射 $A \to B$ に対する制限関手と環の変更関手が互いに随伴であることから
従う。
\end{proof}


\section{加群の層}
\label{sites-modules-section-sheaves-modules}

\begin{definition}
\label{sites-modules-definition-sheaf-modules}
$\mathcal{C}$ をサイトとし，$\mathcal{O}$ を $\mathcal{C}$ 上の
環の層とする。
\begin{enumerate}
\item {\it $\mathcal{O}$-加群の層}とは，アーベル群の台となる前層
$\mathcal{F}$ が層であるような $\mathcal{O}$-加群の前層
$\mathcal{F}$ のことである。定義 \ref{sites-modules-definition-presheaf-modules} を
参照せよ。
\item {\it $\mathcal{O}$-加群の層の射}とは，
$\mathcal{O}$-加群の前層の射のことである。
\item $\mathcal{F}$ と $\mathcal{G}$ を $\mathcal{O}$-加群の層とするとき，
$\mathcal{O}$-加群の層の射全体の集合を
$\Hom_\mathcal{O}(\mathcal{F}, \mathcal{G})$ と書く。
\item $\mathcal{O}$-加群の層の圏を
$\textit{Mod}(\mathcal{O})$ と書く。
\end{enumerate}
\end{definition}

\noindent
$\mathcal{O}$ が単なる環の前層である場合にも，この定義はある程度
意味をもつ。しかし，それが有用となる例を知らないので，
$\mathcal{O}$ が環の層でない場合には「$\mathcal{O}$-加群の層」という
用語を用いない。



\section{加群の前層の層化}
\label{sites-modules-section-sheafification-presheaves-modules}

\begin{lemma}
\label{sites-modules-lemma-sheafification-presheaf-modules}
$\mathcal{C}$ をサイト，$\mathcal{O}$ を $\mathcal{C}$ 上の環の前層，
$\mathcal{F}$ を $\mathcal{O}$-加群の前層とする。
$\mathcal{O}^\#$ を環の前層としての $\mathcal{O}$ の層化とする
（《サイト》第 \ref{sites-section-sheaves-algebraic-structures} 節参照）。
$\mathcal{F}^\#$ をアーベル群の前層としての $\mathcal{F}$ の層化とする。
集合の層の写像
$$
\mathcal{O}^\# \times \mathcal{F}^\#
\longrightarrow
\mathcal{F}^\#
$$
であって，図式
$$
\xymatrix{
\mathcal{O} \times \mathcal{F} \ar[r] \ar[d] &
\mathcal{F} \ar[d] \\
\mathcal{O}^\# \times \mathcal{F}^\# \ar[r] &
\mathcal{F}^\#
}
$$
を可換にし，かつ $\mathcal{F}^\#$ を $\mathcal{O}^\#$-加群の層に
するものが一意に存在する。さらに，$\mathcal{G}$ が
$\mathcal{O}^\#$-加群の層ならば，$\mathcal{O}$-加群の前層の任意の射
$\mathcal{F} \to \mathcal{G}$（$\mathcal{G}$ を $\mathcal{O}$-加群へ
制限したものへの射）は，一意に
$\mathcal{F} \to \mathcal{F}^\# \to \mathcal{G}$ と分解する。
ここで $\mathcal{F}^\# \to \mathcal{G}$ は
$\mathcal{O}^\#$-加群の射である。
\end{lemma}

\begin{proof}
省略する。
\end{proof}

\noindent
これは，関手
$i : \textit{Mod}(\mathcal{O}^\#) \to \textit{PMod}(\mathcal{O})$
（制限と，層を前層へ包含する操作との合成）と，上の補題の層化関手
${}^\# : \textit{PMod}(\mathcal{O}) \to \textit{Mod}(\mathcal{O}^\#)$
が随伴であるという意味である。式で書けば
$$
\Mor_{\textit{PMod}(\mathcal{O})}(\mathcal{F}, i\mathcal{G})
=
\Mor_{\textit{Mod}(\mathcal{O}^\#)}(\mathcal{F}^\#, \mathcal{G})
$$
である。特に重要なのは $\mathcal{O}$ がすでに環の層である場合である。
このとき式は
$$
\Mor_{\textit{PMod}(\mathcal{O})}(\mathcal{F}, i\mathcal{G})
=
\Mor_{\textit{Mod}(\mathcal{O})}(\mathcal{F}^\#, \mathcal{G})
$$
となる。実際，この場合 $\mathcal{O} = \mathcal{O}^\#$ である。

\begin{lemma}
\label{sites-modules-lemma-sheafification-exact}
$\mathcal{C}$ をサイトとし，$\mathcal{O}$ を $\mathcal{C}$ 上の
環の前層とする。層化関手
$$
\textit{PMod}(\mathcal{O}) \longrightarrow \textit{Mod}(\mathcal{O}^\#), \quad
\mathcal{F} \longmapsto \mathcal{F}^\#
$$
は完全である。
\end{lemma}

\begin{proof}
層化 $\textit{PAb}(\mathcal{C}) \to \textit{Ab}(\mathcal{C})$ に対して
同じ主張が成り立つからである。第 \ref{sites-modules-section-abelian-sheaves} 節の
議論を参照せよ。
\end{proof}

\noindent
$\mathcal{C}$ をサイトとし，
$\mathcal{O}_1 \to \mathcal{O}_2$ を $\mathcal{C}$ 上の環の層の射とする。
第 \ref{sites-modules-section-presheaves-modules} 節では，この状況に付随する
加群の前層上の制限関手と環の変更関手を定義した。

\medskip\noindent
$\mathcal{F}$ が $\mathcal{O}_2$-加群の層ならば，
$\mathcal{F}$ の制限 $\mathcal{F}_{\mathcal{O}_1}$ は明らかに
$\mathcal{O}_1$-加群の層である。したがって制限関手
$$
\textit{Mod}(\mathcal{O}_2)
\longrightarrow
\textit{Mod}(\mathcal{O}_1)
$$
を得る。

\medskip\noindent
一方，$\mathcal{O}_1$-加群の層 $\mathcal{G}$ が与えられたとき，
$\mathcal{O}_2$-加群の前層
$\mathcal{O}_2 \otimes_{p, \mathcal{O}_1} \mathcal{G}$
は一般には層ではない。そこで，{\it テンソル積層}
$\mathcal{O}_2 \otimes_{\mathcal{O}_1} \mathcal{G}$
を，前層に対する構成を層化したものとして式
$$
\mathcal{O}_2 \otimes_{\mathcal{O}_1} \mathcal{G}
=
(\mathcal{O}_2 \otimes_{p, \mathcal{O}_1} \mathcal{G})^\#
$$
で定義する。{\it 環の変更}関手
$$
\textit{Mod}(\mathcal{O}_1)
\longrightarrow
\textit{Mod}(\mathcal{O}_2)
$$
を得る。

\begin{lemma}
\label{sites-modules-lemma-adjointness-tensor-restrict}
上の $X$，$\mathcal{O}_1$，$\mathcal{O}_2$，$\mathcal{F}$，
$\mathcal{G}$ に対し，標準的な全単射
$$
\Hom_{\mathcal{O}_1}(\mathcal{G}, \mathcal{F}_{\mathcal{O}_1})
=
\Hom_{\mathcal{O}_2}(
\mathcal{O}_2 \otimes_{\mathcal{O}_1} \mathcal{G},
\mathcal{F}
)
$$
が存在する。言い換えれば，制限関手と環の変更関手は互いに随伴である。
\end{lemma}

\begin{proof}
補題 \ref{sites-modules-lemma-adjointness-tensor-restrict-presheaves} と，等式
$\Hom_{\mathcal{O}_2}(
\mathcal{O}_2 \otimes_{\mathcal{O}_1} \mathcal{G},
\mathcal{F}
)
=
\Hom_{\mathcal{O}_2}(
\mathcal{O}_2 \otimes_{p, \mathcal{O}_1} \mathcal{G},
\mathcal{F}
)$
から従う。この等式は $\mathcal{F}$ が層であるために成り立つ。
\end{proof}

\begin{lemma}
\label{sites-modules-lemma-epimorphism-modules}
$\mathcal{C}$ をサイトとし，$\mathcal{O} \to \mathcal{O}'$ を
環の層のエピ射とする。$\mathcal{G}_1,\mathcal{G}_2$ を
$\mathcal{O}'$-加群とする。このとき
$$
\Hom_{\mathcal{O}'}(\mathcal{G}_1, \mathcal{G}_2) =
\Hom_\mathcal{O}(\mathcal{G}_1, \mathcal{G}_2).
$$
言い換えれば，制限関手
$\textit{Mod}(\mathcal{O}') \to \textit{Mod}(\mathcal{O})$ は充満忠実である。
\end{lemma}

\begin{proof}
これは《代数》補題 \ref{algebra-lemma-epimorphism-modules} の
層に対する版であり，まったく同じ方法で証明される。
\end{proof}




\section{トポスの射と加群の層}
\label{sites-modules-section-sheaves-modules-functorial}

\noindent
以下の内容はすべて完全に形式的である。
ここではトポスの射の場合にすべてを定式化するが，
もちろん結果はサイトの射の場合にも成り立つ。

\begin{lemma}
\label{sites-modules-lemma-pushforward-module}
$\mathcal{C}$ と $\mathcal{D}$ をサイトとし，
$f : \Sh(\mathcal{C}) \to \Sh(\mathcal{D})$
をトポスの射とする。
$\mathcal{O}$ を $\mathcal{C}$ 上の環の層，
$\mathcal{F}$ を $\mathcal{O}$-加群の層とする。
集合の層の間の自然な写像
$$
f_*\mathcal{O} \times f_*\mathcal{F}
\longrightarrow
f_*\mathcal{F}
$$
が存在し，これによって $f_*\mathcal{F}$ は
$f_*\mathcal{O}$-加群の層となる。
この構成は $\mathcal{F}$ に関して関手的である。
\end{lemma}

\begin{proof}
乗法写像を
$\mu : \mathcal{O} \times \mathcal{F} \to \mathcal{F}$ と書く。
集合の層に対する $f_*$ は左完全であり，したがって積と可換することを
想起せよ。それゆえ $f_*\mu$ が上に表示した写像である。
これで補題が証明された。
\end{proof}

\begin{lemma}
\label{sites-modules-lemma-pullback-module}
$\mathcal{C}$ と $\mathcal{D}$ をサイトとし，
$f : \Sh(\mathcal{C}) \to \Sh(\mathcal{D})$
をトポスの射とする。
$\mathcal{O}$ を $\mathcal{D}$ 上の環の層，
$\mathcal{G}$ を $\mathcal{O}$-加群の層とする。
集合の層の間の自然な写像
$$
f^{-1}\mathcal{O} \times f^{-1}\mathcal{G}
\longrightarrow
f^{-1}\mathcal{G}
$$
が存在し，これによって $f^{-1}\mathcal{G}$ は
$f^{-1}\mathcal{O}$-加群の層となる。
この構成は $\mathcal{G}$ に関して関手的である。
\end{lemma}

\begin{proof}
乗法写像を
$\mu : \mathcal{O} \times \mathcal{G} \to \mathcal{G}$ と書く。
集合の層に対する $f^{-1}$ は完全であり，したがって積と可換することを
想起せよ。それゆえ $f^{-1}\mu$ が上に表示した写像である。
これで補題が証明された。
\end{proof}

\begin{lemma}
\label{sites-modules-lemma-adjoint-push-pull-modules}
$\mathcal{C}$ と $\mathcal{D}$ をサイトとし，
$f : \Sh(\mathcal{C}) \to \Sh(\mathcal{D})$
をトポスの射とする。
$\mathcal{O}$ を $\mathcal{D}$ 上の環の層，
$\mathcal{G}$ を $\mathcal{O}$-加群の層，
$\mathcal{F}$ を $f^{-1}\mathcal{O}$-加群の層とする。
このとき
$$
\Mor_{\textit{Mod}(f^{-1}\mathcal{O})}(f^{-1}\mathcal{G}, \mathcal{F})
=
\Mor_{\textit{Mod}(\mathcal{O})}(\mathcal{G}, f_*\mathcal{F}).
$$
ここでは補題 \ref{sites-modules-lemma-pullback-module} と
\ref{sites-modules-lemma-pushforward-module} を用い，
$\mathcal{O} \to f_*f^{-1}\mathcal{O}$ に沿って制限することにより
$f_*\mathcal{F}$ を $\mathcal{O}$-加群とみなしている。
\end{lemma}

\begin{proof}
まず，
$$
\Mor_{\textit{Ab}(\mathcal{C})}(f^{-1}\mathcal{G}, \mathcal{F})
=
\Mor_{\textit{Ab}(\mathcal{D})}(\mathcal{G}, f_*\mathcal{F}).
$$
が成り立つことに注意する。これは《サイト》命題
\ref{sites-proposition-functoriality-algebraic-structures-topoi}
による。アーベル群の層の射
$\alpha : f^{-1}\mathcal{G} \to \mathcal{F}$ と
$\beta : \mathcal{G} \to f_*\mathcal{F}$ が，上の式によって
対応していると仮定する。$\alpha$ が
$f^{-1}\mathcal{O}$-線形であることと，$\beta$ が
$\mathcal{O}$-線形であることが同値であることを示さなければならない。
例えば $\alpha$ が $f^{-1}\mathcal{O}$-線形であると仮定すると，
$f_*\alpha$ は明らかに $f_*f^{-1}\mathcal{O}$-線形であり，
したがって（制限は関手なので）$\mathcal{O}$-線形である。
ゆえに，随伴写像
$\mathcal{G} \to f_*f^{-1}\mathcal{G}$ が
$\mathcal{O}$-線形であることを示せば十分である。
集合の層に対する $f_*$ と $f^{-1}$ はともに積と可換するので，
これは図式
$$
\xymatrix{
\mathcal{O} \times \mathcal{G} \ar[r] \ar[d] &
f_*f^{-1}(\mathcal{O} \times \mathcal{G}) \ar[d] \\
\mathcal{G} \ar[r] & f_*f^{-1}\mathcal{G}
}
$$
が可換であることに帰着する。これは随伴写像
$\text{id}_{\Sh(\mathcal{D})} \to f_*f^{-1}$ が
関手の自然変換であることから成り立つ。
「$\beta$ が線形ならば $\alpha$ が線形である」という向きの証明は
省略する。
\end{proof}

\begin{lemma}
\label{sites-modules-lemma-adjoint-pull-push-modules}
$\mathcal{C}$ と $\mathcal{D}$ をサイトとし，
$f : \Sh(\mathcal{C}) \to \Sh(\mathcal{D})$
をトポスの射とする。
$\mathcal{O}$ を $\mathcal{C}$ 上の環の層，
$\mathcal{F}$ を $\mathcal{O}$-加群の層，
$\mathcal{G}$ を $f_*\mathcal{O}$-加群の層とする。
このとき
$$
\Mor_{\textit{Mod}(\mathcal{O})}(
\mathcal{O} \otimes_{f^{-1}f_*\mathcal{O}} f^{-1}\mathcal{G}, \mathcal{F})
=
\Mor_{\textit{Mod}(f_*\mathcal{O})}(\mathcal{G}, f_*\mathcal{F}).
$$
ここでは補題 \ref{sites-modules-lemma-pullback-module} と
\ref{sites-modules-lemma-pushforward-module} を用い，テンソル積の定義には標準写像
$f^{-1}f_*\mathcal{O} \to \mathcal{O}$ を用いている。
\end{lemma}

\begin{proof}
補題 \ref{sites-modules-lemma-adjointness-tensor-restrict} により，
$$
\Mor_{\textit{Mod}(\mathcal{O})}(
\mathcal{O} \otimes_{f^{-1}f_*\mathcal{O}} f^{-1}\mathcal{G}, \mathcal{F})
=
\Mor_{\textit{Mod}(f^{-1}f_*\mathcal{O})}(
f^{-1}\mathcal{G}, \mathcal{F}_{f^{-1}f_*\mathcal{O}})
$$
である。したがって結論は補題
\ref{sites-modules-lemma-adjoint-push-pull-modules} から従う。
\end{proof}






\section{環付きトポスの射と加群}
\label{sites-modules-section-functoriality-modules}

\noindent
これで，環付きトポスの射に沿う加群の引き戻しと順像を定義するために
十分な記法を導入した。

\begin{definition}
\label{sites-modules-definition-pushforward}
写像
$(f, f^\sharp) :
(\Sh(\mathcal{C}), \mathcal{O}_\mathcal{C})
\to
(\Sh(\mathcal{D}), \mathcal{O}_\mathcal{D})$
を環付きトポスまたは環付きサイトの射とする。
\begin{enumerate}
\item $\mathcal{F}$ を $\mathcal{O}_\mathcal{C}$-加群の層とする。
$\mathcal{F}$ の {\it 順像}を，アーベル群の層としては
$f_*\mathcal{F}$ に等しく，加群構造が
$f^\sharp : \mathcal{O}_\mathcal{D} \to f_*\mathcal{O}_\mathcal{C}$
に沿って加群構造
$$
f_*\mathcal{O}_\mathcal{C} \times f_*\mathcal{F}
\longrightarrow
f_*\mathcal{F}
$$
を制限して得られる $\mathcal{O}_\mathcal{D}$-加群の層として定義する。
ここでこの加群構造は補題 \ref{sites-modules-lemma-pushforward-module} による。
\item $\mathcal{G}$ を $\mathcal{O}_\mathcal{D}$-加群の層とする。
$\mathcal{G}$ の {\it 引き戻し} $f^*\mathcal{G}$ を，式
$$
f^*\mathcal{G}
=
\mathcal{O}_\mathcal{C} \otimes_{f^{-1}\mathcal{O}_\mathcal{D}}
f^{-1}\mathcal{G}
$$
で定義される $\mathcal{O}_\mathcal{C}$-加群の層とする。
ここで環の写像
$f^{-1}\mathcal{O}_\mathcal{D} \to \mathcal{O}_\mathcal{C}$
は $f^\sharp$ であり，加群構造は補題
\ref{sites-modules-lemma-pullback-module} によって与えられる。
\end{enumerate}
\end{definition}

\noindent
こうして関手
\begin{eqnarray*}
f_* : \textit{Mod}(\mathcal{O}_\mathcal{C})
& \longrightarrow &
\textit{Mod}(\mathcal{O}_\mathcal{D}) \\
f^* : \textit{Mod}(\mathcal{O}_\mathcal{D})
& \longrightarrow &
\textit{Mod}(\mathcal{O}_\mathcal{C})
\end{eqnarray*}
を定義した。これらの関手に関する最後の結果は，期待どおり実際に
互いに随伴であるというものである。

\begin{lemma}
\label{sites-modules-lemma-adjoint-pullback-pushforward-modules}
写像
$(f, f^\sharp) :
(\Sh(\mathcal{C}), \mathcal{O}_\mathcal{C})
\to
(\Sh(\mathcal{D}), \mathcal{O}_\mathcal{D})$
を環付きトポスまたは環付きサイトの射とする。
$\mathcal{F}$ を $\mathcal{O}_\mathcal{C}$-加群の層，
$\mathcal{G}$ を $\mathcal{O}_\mathcal{D}$-加群の層とする。
標準的な全単射
$$
\Hom_{\mathcal{O}_\mathcal{C}}(f^*\mathcal{G}, \mathcal{F})
=
\Hom_{\mathcal{O}_\mathcal{D}}(\mathcal{G}, f_*\mathcal{F}).
$$
が存在する。言い換えれば，関手 $f^*$ は $f_*$ の左随伴である。
\end{lemma}

\begin{proof}
これはこれまでの結果から従う：
\begin{eqnarray*}
\Hom_{\mathcal{O}_\mathcal{C}}(f^*\mathcal{G}, \mathcal{F})
& = &
\Mor_{\textit{Mod}(\mathcal{O}_\mathcal{C})}(
\mathcal{O}_\mathcal{C}
\otimes_{f^{-1}\mathcal{O}_\mathcal{D}} f^{-1}\mathcal{G},
\mathcal{F}) \\
& = &
\Mor_{\textit{Mod}(f^{-1}\mathcal{O}_\mathcal{D})}(
f^{-1}\mathcal{G}, \mathcal{F}_{f^{-1}\mathcal{O}_\mathcal{D}}) \\
& = &
\Hom_{\mathcal{O}_\mathcal{D}}(\mathcal{G}, f_*\mathcal{F}).
\end{eqnarray*}
ここでは補題 \ref{sites-modules-lemma-adjointness-tensor-restrict} と
\ref{sites-modules-lemma-adjoint-push-pull-modules} を用いた。
\end{proof}

\begin{lemma}
\label{sites-modules-lemma-push-pull-composition-modules}
$(f, f^\sharp) :
(\Sh(\mathcal{C}_1), \mathcal{O}_1)
\to (\Sh(\mathcal{C}_2), \mathcal{O}_2)$ と
$(g, g^\sharp) :
(\Sh(\mathcal{C}_2), \mathcal{O}_2) \to
(\Sh(\mathcal{C}_3), \mathcal{O}_3)$
を環付きトポスの射とする。関手の標準同型
$(g \circ f)_* \cong g_* \circ f_*$ と
$(g \circ f)^* \cong f^* \circ g^*$ が存在する。
\end{lemma}

\begin{proof}
定義から明らかである。
\end{proof}





\section{加群の層のアーベル圏}
\label{sites-modules-section-kernels}

\noindent
$(\Sh(\mathcal{C}), \mathcal{O})$ を環付きトポスとする。
$\mathcal{F}$ と $\mathcal{G}$ を $\mathcal{O}$-加群の層とする
（《層》定義 \ref{sheaves-definition-sheaf-modules} 参照）。
$\varphi, \psi : \mathcal{F} \to \mathcal{G}$ を
$\mathcal{O}$-加群の層の射とする。
$\varphi + \psi : \mathcal{F} \to \mathcal{G}$ を，アーベル群の層の
射としての $\varphi$ と $\psi$ の和として定義する。
これは明らかに再び $\mathcal{O}$-加群の写像である。
また，$\mathcal{O}$-加群の写像の合成がこの加法に関して双線形であることも
明らかである。したがって $\textit{Mod}(\mathcal{O})$ は前加法圏である
（《ホモロジー》定義 \ref{homology-definition-preadditive} 参照）。

\medskip\noindent
$\mathcal{C}$ のすべての対象 $U$ において一定値 $\{0\}$ をとる
$\mathcal{O}$-加群の層を $0$ と書く。
これは明らかに $\textit{Mod}(\mathcal{O})$ の終対象であると同時に
始対象でもある。$\mathcal{O}$-加群の射
$\varphi : \mathcal{F} \to \mathcal{G}$ に対し，次は同値である：
(a) $\varphi$ は零である，(b) $\varphi$ は $0$ を経由して分解する，
(c) 各対象 $U$ 上の切断において $\varphi$ は零である。

\medskip\noindent
さらに，二つの $\mathcal{O}$-加群の層
$\mathcal{F}$ と $\mathcal{G}$ に対し，直和を
$$
\mathcal{F} \oplus \mathcal{G} = \mathcal{F} \times \mathcal{G}
$$
として定義できる。ここで写像 $(i, j, p, q)$ は《ホモロジー》定義
\ref{homology-definition-direct-sum} における明らかな写像である。
したがって $\textit{Mod}(\mathcal{O})$ は加法圏である
（《ホモロジー》定義
\ref{homology-definition-additive-category} 参照）。

\medskip\noindent
$\varphi : \mathcal{F} \to \mathcal{G}$ を
$\mathcal{O}$-加群の射とする。$\Ker(\varphi)$ を，アーベル群の層の
写像としての $\varphi$ の核として定義できる。
第 \ref{sites-modules-section-abelian-sheaves} 節により，これは各
$\mathcal{C}$ の対象 $U$ における切断が
$$
\Ker(\varphi)(U) =
\{ s \in \mathcal{F}(U) \mid \varphi(s) = 0 \text{ in } \mathcal{G}(U)\}
$$
である $\mathcal{F}$ の部分層である。これが実際に
$\mathcal{O}$-加群の圏における核であることは容易に分かる。
言い換えれば，射 $\alpha : \mathcal{H} \to \mathcal{F}$ が
$\Ker(\varphi)$ を経由して分解することと
$\varphi \circ \alpha = 0$ であることとは同値である。

\medskip\noindent
同様に，$\Coker(\varphi)$ を，アーベル群の層の写像としての
$\varphi$ の余核として定義する。写像
$$
\mathcal{O} \times \Coker(\varphi) \longrightarrow \Coker(\varphi)
$$
であって，$\mathcal{G} \to \Coker(\varphi)$ を
$\mathcal{O}$-加群の射にするものが一意に存在する
（確認は省略する）。写像 $\mathcal{G} \to \Coker(\varphi)$ は全射である
（集合の層の写像として。第 \ref{sites-modules-section-abelian-sheaves} 節参照）。
$\Coker(\varphi)$ が $\textit{Mod}(\mathcal{O})$ における余核であることを
示すため，$\beta : \mathcal{G} \to \mathcal{H}$ を
$\beta \circ \varphi$ が零となる $\mathcal{O}$-加群の射とする。
すると $\mathcal{C}$ の各対象 $U$ に対し，$\beta$ が誘導する写像
$\mathcal{G}(U)/\varphi(\mathcal{F}(U)) \to \mathcal{H}(U)$ を得る。
層化の普遍性（《層》補題
\ref{sheaves-lemma-sheafification-presheaf-modules} 参照）により，
元の $\beta$ が合成
$\mathcal{G} \to \Coker(\varphi) \to \mathcal{H}$ に等しくなるような
標準写像 $\Coker(\varphi) \to \mathcal{H}$ を得る。
この射 $\Coker(\varphi) \to \mathcal{H}$ は，上で述べた全射性により
一意である。

\begin{lemma}
\label{sites-modules-lemma-abelian}
$(\Sh(\mathcal{C}), \mathcal{O})$ を環付きトポスとする。
圏 $\textit{Mod}(\mathcal{O})$ はアーベル圏である。忘却関手
$\textit{Mod}(\mathcal{O}) \to \textit{Ab}(\mathcal{C})$
は完全である。したがって $\mathcal{O}$-加群の核，余核，および完全性は，
アーベル群の層に対する対応する概念と一致する。
\end{lemma}

\begin{proof}
上で，$\textit{Mod}(\mathcal{O})$ は核と余核をもつ加法圏であり，
$\textit{Mod}(\mathcal{O}) \to \textit{Ab}(\mathcal{C})$ が核と余核を
保つことを見た。《ホモロジー》定義
\ref{homology-definition-abelian-category} により，像と余像が一致することを
示せばよい。これは $\textit{Ab}(\mathcal{C})$ で成り立つので明らかで
ある。補題が従う。
\end{proof}

\begin{lemma}
\label{sites-modules-lemma-limits-colimits}
$(\Sh(\mathcal{C}), \mathcal{O})$ を環付きトポスとする。
$\textit{Mod}(\mathcal{O})$ にはすべての極限と余極限が存在し，忘却関手
$\textit{Mod}(\mathcal{O}) \to \textit{Ab}(\mathcal{C})$
はそれらと可換する。さらに，フィルター付き余極限は完全である。
\end{lemma}

\begin{proof}
最後の主張は最初の主張から従う。実際，補題
\ref{sites-modules-lemma-limits-colimits-abelian-sheaves} により，
$\textit{Ab}(\mathcal{C})$ においてフィルター付き余極限は完全である。
$\mathcal{I} \to \textit{Mod}(\mathcal{C})$，
$i \mapsto \mathcal{F}_i$ を図式とする。
$\lim_i \mathcal{F}_i$ を $\textit{Ab}(\mathcal{C})$ におけるこの図式の
極限とする。補題 \ref{sites-modules-lemma-limits-colimits-abelian-sheaves} における
この極限の記述から，乗法
$$
\mathcal{O} \times \lim_i \mathcal{F}_i
\longrightarrow
\lim_i \mathcal{F}_i
$$
であって，$\lim_i \mathcal{F}_i$ を $\mathcal{O}$-加群の層にするものが
存在することが直ちに分かる。これが
$\textit{Mod}(\mathcal{C})$ における図式の極限であることは容易に分かる。
$\colim_i \mathcal{F}_i$ を $\textit{PAb}(\mathcal{C})$ における図式の
余極限とする。補題
\ref{sites-modules-lemma-limits-colimits-abelian-presheaves} の証明における
この余極限の記述から，乗法
$$
\mathcal{O} \times \colim_i \mathcal{F}_i
\longrightarrow
\colim_i \mathcal{F}_i
$$
であって，$\colim_i \mathcal{F}_i$ を $\mathcal{O}$-加群の前層に
するものが存在することが直ちに分かる。層化を施すことにより，
$\mathcal{O}$-加群の層 $(\colim_i \mathcal{F}_i)^\#$ を得る
（補題 \ref{sites-modules-lemma-sheafification-presheaf-modules} 参照）。
$(\colim_i \mathcal{F}_i)^\#$ が
$\textit{Mod}(\mathcal{O})$ における図式の余極限であることは容易に分かる。
また補題 \ref{sites-modules-lemma-limits-colimits-abelian-sheaves} により，
$\mathcal{O}$-加群構造を忘れたものは
$\textit{Ab}(\mathcal{C})$ における余極限である。
\end{proof}

\noindent
極限と余極限が存在することにより，$\mathcal{O}$-加群の圏上で定義された
関手の完全性を，《圏論》第
\ref{categories-section-exact-functor} 節と同様に，
極限と余極限を用いて考察できる。
短完全列による完全性の記述については《ホモロジー》補題
\ref{homology-lemma-exact-functor} を参照せよ。

\begin{lemma}
\label{sites-modules-lemma-exactness-pushforward-pullback}
$f : (\Sh(\mathcal{C}), \mathcal{O}_\mathcal{C})
\to (\Sh(\mathcal{D}), \mathcal{O}_\mathcal{D})$
を環付きトポスの射とする。
\begin{enumerate}
\item 関手 $f_*$ は左完全である。実際，すべての極限と可換する。
\item 関手 $f^*$ は右完全である。実際，すべての余極限と可換する。
\end{enumerate}
\end{lemma}

\begin{proof}
これは $(f^*, f_*)$ が随伴関手対であることから成り立つ
（補題 \ref{sites-modules-lemma-adjoint-pullback-pushforward-modules} 参照）。
《圏論》第 \ref{categories-section-adjoint} 節を参照せよ。
\end{proof}

\begin{lemma}
\label{sites-modules-lemma-check-exactness-stalks}
$\mathcal{C}$ をサイトとする。$\{p_i\}_{i \in I}$ が保存的な点の族ならば，
アーベル群の層の列の完全性は，各 $i \in I$ に対する点 $p_i$ での
茎上で確認できる。$\mathcal{C}$ が十分多くの点をもつならば，
アーベル群の層の列の完全性は茎上で確認できる。
\end{lemma}

\begin{proof}
《サイト》補題 \ref{sites-lemma-exactness-stalks} から直ちに従う。
\end{proof}





\section{順像の完全性}
\label{sites-modules-section-pushforward}

\noindent
順像の完全性に関するいくつかの技術的補題を述べる。

\begin{lemma}
\label{sites-modules-lemma-reflect-surjections}
$f : \Sh(\mathcal{C}) \to \Sh(\mathcal{D})$ をトポスの射とする。
次は同値である：
\begin{enumerate}
\item $\textit{Ab}(\mathcal{C})$ のすべての $\mathcal{F}$ に対し，
$f^{-1}f_*\mathcal{F} \to \mathcal{F}$ は全射である。
\item $f_* : \textit{Ab}(\mathcal{C}) \to \textit{Ab}(\mathcal{D})$
は全射を反映する。
\end{enumerate}
この場合，関手
$f_* : \textit{Ab}(\mathcal{C}) \to \textit{Ab}(\mathcal{D})$
は忠実である。
\end{lemma}

\begin{proof}
(1) を仮定する。$a : \mathcal{F} \to \mathcal{F}'$ を
$\mathcal{C}$ 上のアーベル群の層の写像で，$f_*a$ が全射であるものとする。
$f^{-1}$ は完全なので，
$f^{-1}f_*a : f^{-1}f_*\mathcal{F} \to f^{-1}f_*\mathcal{F}'$
も全射である。(1) と合わせると，$a$ が全射であることが従う。
したがって (2) が成り立つ。

\medskip\noindent
(2) を仮定する。$\mathcal{F}$ を $\mathcal{C}$ 上のアーベル群の層とする。
写像 $f^{-1}f_*\mathcal{F} \to \mathcal{F}$ が全射であることを
示さなければならない。(2) により，
$f_*f^{-1}f_*\mathcal{F} \to f_*\mathcal{F}$ が全射であることを
示せば十分である。これは，その片側逆である標準写像
$f_*\mathcal{F} \to f_*f^{-1}f_*\mathcal{F}$ が存在することから成り立つ。

\medskip\noindent
最後の主張の証明は省略する。
\end{proof}

\begin{lemma}
\label{sites-modules-lemma-exactness}
$f : \Sh(\mathcal{C}) \to \Sh(\mathcal{D})$ をトポスの射とする。
次の性質のうち少なくとも一つが成り立つと仮定する：
\begin{enumerate}
\item $f_*$ は集合の層の全射を全射に写す。
\item $f_*$ はアーベル群の層の全射を全射に写す。
\item $f_*$ は集合の層における余等化子と可換する。
\item $f_*$ は集合の層における押し出しと可換する。
\end{enumerate}
このとき
$f_* : \textit{Ab}(\mathcal{C}) \to \textit{Ab}(\mathcal{D})$ は完全である。
\end{lemma}

\begin{proof}
$f_* : \textit{Ab}(\mathcal{C}) \to \textit{Ab}(\mathcal{D})$ は
右随伴なので，$\mathcal{C}$ 上のアーベル群の層の短完全列
$0 \to \mathcal{F}_1 \to \mathcal{F}_2 \to \mathcal{F}_3 \to 0$
を完全列
$$
0 \to f_*\mathcal{F}_1 \to f_*\mathcal{F}_2 \to f_*\mathcal{F}_3
$$
に写すことはすでに分かっている。《圏論》第
\ref{categories-section-exact-functor} 節および
\ref{categories-section-adjoint} 節，ならびに《ホモロジー》第
\ref{homology-section-functors} 節を参照せよ。
したがって写像 $f_*\mathcal{F}_2 \to f_*\mathcal{F}_3$ が
全射であることを示せば十分である。(1) または (2) が成り立つ場合，
これは定義から明らかである。《サイト》補題
\ref{sites-lemma-exactness-properties} により，(3) または (4) の
いずれからも形式的に (1) が従うので，これらの場合も証明が完了する。
\end{proof}

\begin{lemma}
\label{sites-modules-lemma-morphism-ringed-sites-almost-cocontinuous}
$f : \mathcal{D} \to \mathcal{C}$ を，連続関手
$u : \mathcal{C} \to \mathcal{D}$ に付随するサイトの射とする。
$u$ がほとんど余連続であると仮定する。このとき
\begin{enumerate}
\item $f_* : \textit{Ab}(\mathcal{D}) \to \textit{Ab}(\mathcal{C})$ は
完全である。
\item $f$ が環付きサイトの射となるように
$f^\sharp : f^{-1}\mathcal{O}_\mathcal{C} \to \mathcal{O}_\mathcal{D}$
が与えられているならば，
$f_* : \textit{Mod}(\mathcal{O}_\mathcal{D}) \to
\textit{Mod}(\mathcal{O}_\mathcal{C})$ は完全である。
\end{enumerate}
\end{lemma}

\begin{proof}
(2) は補題 \ref{sites-modules-lemma-limits-colimits} により (1) から従う。
(1) は《サイト》補題
\ref{sites-lemma-morphism-of-sites-almost-cocontinuous} および
\ref{sites-lemma-exactness-properties} から従う。
\end{proof}





\section{関手 $g_!$ の完全性}
\label{sites-modules-section-exactness-lower-shriek}

\noindent
$u : \mathcal{C} \to \mathcal{D}$ をサイト間の関手とする。
次を仮定する：
\begin{enumerate}
\item[(a)] $u$ は余連続である。
\item[(b)] $u$ は連続である。
\end{enumerate}
$g : \Sh(\mathcal{C}) \to \Sh(\mathcal{D})$ を $u$ に付随する
トポスの射とする（《サイト》補題
\ref{sites-lemma-cocontinuous-morphism-topoi} 参照）。
$g^{-1} = u^p$，すなわち $g^{-1}$ は簡単な式
$(g^{-1}\mathcal{G})(U) = \mathcal{G}(u(U))$ で与えられることを
想起せよ（《サイト》補題 \ref{sites-lemma-when-shriek} 参照）。
関手
$g^{-1} : \textit{Ab}(\mathcal{D}) \to \textit{Ab}(\mathcal{C})$
が左随伴 $g_!$ をもつことを示したい。《サイト》補題
\ref{sites-lemma-when-shriek} により，
関手 $g^{Sh}_! = (u_p\ )^\#$ は集合の層における左随伴である。
さらに，$g^{Sh}_!\mathcal{F}$ は前層
$$
V \longmapsto \colim_{V \to u(U)} \mathcal{F}(U)
$$
に付随する層であることが分かっている。ここで余極限は
$(\mathcal{I}_V^u)^{opp}$ 上で，集合の圏において取る。
したがって次の定義は自然である。

\begin{definition}
\label{sites-modules-definition-g-shriek}
$u : \mathcal{C} \to \mathcal{D}$ は上の (a)，(b) を満たすとする。
$\mathcal{F} \in \textit{PAb}(\mathcal{C})$ に対し，
前層 {\it $g_{p!}\mathcal{F}$} を
$$
V \longmapsto \colim_{V \to u(U)} \mathcal{F}(U)
$$
によって定義する。ここで $(\mathcal{I}_V^u)^{opp}$ 上の余極限は
$\textit{Ab}$ において取る。
$\mathcal{F} \in \textit{PAb}(\mathcal{C})$ に対し，
{\it $g_!\mathcal{F} = (g_{p!}\mathcal{F})^\#$} と置く。
\end{definition}

\noindent
ここでこれほど明示的に述べたのは，関手 $g^{Sh}_!$ と $g_!$ が
異なるからである。両方を用いるときには，区別が明確になるよう
注意しなければならない。

\begin{lemma}
\label{sites-modules-lemma-g-shriek-adjoint}
関手 $g_{p!}$ は関手 $u^p$ の左随伴である。
関手 $g_!$ は関手 $g^{-1}$ の左随伴である。
言い換えれば，式
\begin{align*}
\Mor_{\textit{PAb}(\mathcal{C})}(\mathcal{F}, u^p\mathcal{G})
& =
\Mor_{\textit{PAb}(\mathcal{D})}(g_{p!}\mathcal{F}, \mathcal{G}), \\
\Mor_{\textit{Ab}(\mathcal{C})}(\mathcal{F}, g^{-1}\mathcal{G})
& =
\Mor_{\textit{Ab}(\mathcal{D})}(g_!\mathcal{F}, \mathcal{G})
\end{align*}
は $\mathcal{F}$ と $\mathcal{G}$ に関して双関手的に成り立つ。
\end{lemma}

\begin{proof}
第二の式は第一の式から形式的に従う。実際，
$\mathcal{F}$ と $\mathcal{G}$ がアーベル群の層ならば，層化の普遍性により
\begin{align*}
\Mor_{\textit{Ab}(\mathcal{C})}(\mathcal{F}, g^{-1}\mathcal{G})
& =
\Mor_{\textit{PAb}(\mathcal{D})}(g_{p!}\mathcal{F}, \mathcal{G}) \\
& =
\Mor_{\textit{Ab}(\mathcal{D})}(g_!\mathcal{F}, \mathcal{G})
\end{align*}
である。

\medskip\noindent
第一の式を証明するため，$\mathcal{F}$ と $\mathcal{G}$ を
アーベル群の前層とする。左辺の群から右辺の群への写像と，
右辺の群から左辺の群への写像を構成する。
これらが互いに逆であることの確認は省略する。

\medskip\noindent
アーベル群の前層の標準写像
$\mathcal{F} \to u^pg_{p!}\mathcal{F}$ が存在し，$U$ 上の切断では
自然な写像
$\mathcal{F}(U) \to \colim_{u(U) \to u(U')} \mathcal{F}(U')$ であることに
注意せよ。《サイト》補題 \ref{sites-lemma-recover} を参照せよ。
写像 $\alpha : g_{p!}\mathcal{F} \to \mathcal{G}$ が与えられると，
$u^p\alpha : u^pg_{p!}\mathcal{F} \to u^p\mathcal{G}$ を得る。
これに写像 $\mathcal{F} \to u^pg_{p!}\mathcal{F}$ を前合成できる。

\medskip\noindent
アーベル群の前層の標準写像
$g_{p!}u^p\mathcal{G} \to \mathcal{G}$ が存在し，$V$ 上の切断では
自然な写像
$\colim_{V \to u(U)} \mathcal{G}(u(U)) \to \mathcal{G}(V)$ であることに
注意せよ。これは $t : V \to u(U)$ に対応する成分の切断
$s \in u(U)$ を $t^*s \in \mathcal{G}(V)$ に写す。
したがって，写像 $\beta : \mathcal{F} \to u^p\mathcal{G}$ が
与えられると，写像
$g_{p!}\beta : g_{p!}\mathcal{F} \to g_{p!}u^p\mathcal{G}$ を得る。
これに上の写像 $g_{p!}u^p\mathcal{G} \to \mathcal{G}$ を後合成できる。
\end{proof}

\begin{lemma}
\label{sites-modules-lemma-exactness-lower-shriek}
$\mathcal{C}$ と $\mathcal{D}$ をサイトとし，
$u : \mathcal{C} \to \mathcal{D}$ を関手とする。
次を仮定する：
\begin{enumerate}
\item[(a)] $u$ は余連続である。
\item[(b)] $u$ は連続である。
\item[(c)] $\mathcal{C}$ にはファイバー積と等化子が存在し，
$u$ はそれらと可換する。
\end{enumerate}
この場合，関手
$g_! : \textit{Ab}(\mathcal{C}) \to \textit{Ab}(\mathcal{D})$
は完全である。
\end{lemma}

\begin{proof}
《サイト》補題 \ref{sites-lemma-preserve-equalizers} と比較せよ。
(a)，(b)，(c) を仮定する。
$g_!$ は左随伴なので右完全であることはすでに分かっている。
《圏論》補題 \ref{categories-lemma-exact-adjoint} および
《ホモロジー》第 \ref{homology-section-functors} 節を参照せよ。
$g_! = (g_{p!}\ )^\#$ である。$g_!$ がアーベル群の層の単射を
アーベル群の前層の単射に写すことを示さなければならない。
アーベル群の前層の層化は完全であることを想起せよ
（補題 \ref{sites-modules-lemma-limits-colimits-abelian-sheaves} 参照）。
したがって，$g_{p!}$ がアーベル群の前層の単射をアーベル群の前層の
単射に写すことを示せば十分である。そのためには，《サイト》第
\ref{sites-section-functoriality-PSh} 節の圏
$(\mathcal{I}_V^u)^{opp}$ 上の
上の余極限が，図式間の単射を単射に写せば十分である。
これは《サイト》補題 \ref{sites-lemma-almost-directed} および
《代数》補題 \ref{algebra-lemma-almost-directed-colimit-exact} から従う。
\end{proof}

\begin{lemma}
\label{sites-modules-lemma-back-and-forth}
$\mathcal{C}$ と $\mathcal{D}$ をサイトとし，
$u : \mathcal{C} \to \mathcal{D}$ を関手とする。
次を仮定する：
\begin{enumerate}
\item[(a)] $u$ は余連続である。
\item[(b)] $u$ は連続である。
\item[(c)] $u$ は充満忠実である。
\end{enumerate}
上の $g_!, g^{-1}, g_*$ に対し，標準写像
$\mathcal{F} \to g^{-1}g_!\mathcal{F}$ と
$g^{-1}g_*\mathcal{F} \to \mathcal{F}$
は，$\mathcal{C}$ 上のすべてのアーベル群の層 $\mathcal{F}$ に対して
同型である。
\end{lemma}

\begin{proof}
写像 $g^{-1}g_*\mathcal{F} \to \mathcal{F}$ は，《サイト》補題
\ref{sites-lemma-back-and-forth}，およびアーベル群の層の引き戻しと順像が
台となる集合の層の引き戻しと順像に一致することにより，同型である。

\medskip\noindent
$U \in \Ob(\mathcal{C})$ を取る。
$g^{-1}g_!\mathcal{F}(U) = \mathcal{F}(U)$ を示そう。まず，
$g^{-1}g_!\mathcal{F}(U) = g_!\mathcal{F}(u(U))$ であることに注意する。
したがって $g_!\mathcal{F}(u(U)) = \mathcal{F}(U)$ を示せば十分である。
$g_!\mathcal{F}$ は，規則
$$
V \longmapsto \colim_{V \to u(U')} \mathcal{F}(U')
$$
によって定義される前層 $g_{p!}\mathcal{F}$ に付随する
アーベル群の層であり，余極限は $\textit{Ab}$ において取る。
$V = u(U)$ ならば，$u$ は充満忠実なので，この余極限は
$U \to U'$ 上の余極限である。したがって
$g_{p!}\mathcal{F}(u(U) = \mathcal{F}(U)$ を得る。
$u$ は余連続かつ連続なので，$\mathcal{D}$ における $u(U)$ の任意の被覆は，
$\{U_i \to U\}$ が $\mathcal{C}$ における被覆となるような
$\mathcal{D}$ の被覆 (!) $\{u(U_i) \to u(U)\}$ によって細分できる。
したがって $(g_{p!}\mathcal{F})^+(u(U)) = \mathcal{F}(U)$ も成り立つ。
実際，$u(U)$ における $(g_{p!}\mathcal{F})^+$ の値を定義する余極限では，
上の形の被覆 $\{u(U_i) \to u(U)\}$ からなる共終系に制限できる。
ゆえに，$\mathcal{C}$ のすべての対象 $U$ に対しても
$(g_{p!}\mathcal{F})^+(u(U)) = \mathcal{F}(U)$ である。
この議論をもう一度繰り返すと，$\mathcal{C}$ のすべての対象 $U$ に対して
$(g_{p!}\mathcal{F})^\#(u(U)) = \mathcal{F}(U)$ を得る。
これから所望の等式 $g^{-1}g_!\mathcal{F} = \mathcal{F}$ を得る。
\end{proof}

\begin{remark}
\label{sites-modules-remark-no-extension}
一般に，$g$ が環付きトポスの射（の一部）である場合，関手 $g_!$ を
加群の圏へ拡張することはできない。実際，任意の環の写像 $A \to B$ に
対して，関手 $M \mapsto B \otimes_A M$ は右随伴（制限）をもつが，
一般には左随伴をもたない（その存在は $A \to B$ が平坦であることを
含意するからである）。後の第 \ref{sites-modules-section-localize} 節で，
サイトの局所化の場合には加群の層上で $j_!$ を定義できることを見る。
より一般的な形は後の第 \ref{sites-modules-section-lower-shriek-modules} 節で論じる。
\end{remark}

\begin{lemma}
\label{sites-modules-lemma-have-left-adjoint}
$\mathcal{C}$ と $\mathcal{D}$ をサイトとする。
$g : \Sh(\mathcal{C}) \to \Sh(\mathcal{D})$ を，連続かつ余連続な関手
$u : \mathcal{C} \to \mathcal{D}$ に付随するトポスの射とする。
\begin{enumerate}
\item $u$ が左随伴 $w$ をもつならば，台となる集合の層上で
$g_!$ は $g_!^{\Sh}$ と一致し，$g_!$ は完全である。
\item さらに $w$ が余連続ならば，$g_! = h^{-1}$ かつ
$g^{-1} = h_*$ である。ここで
$h : \Sh(\mathcal{D}) \to \Sh(\mathcal{C})$ は $w$ に付随する
トポスの射である。
\end{enumerate}
\end{lemma}

\begin{proof}
この補題は《サイト》補題 \ref{sites-lemma-have-left-adjoint} の類似である。
《サイト》補題 \ref{sites-lemma-adjoint-functors} により，
圏 $\mathcal{I}_V^u$ は始対象をもつ。したがって
$(\mathcal{I}_V^u)^{opp}$ 上のアーベル群の系の余極限の台集合は，
台集合の余極限である。ゆえに，定義 \ref{sites-modules-definition-g-shriek} における
$g_!$ の構成から $g_!^{\Sh}$ と $g_!$ の一致が従う。
完全性と (2) は，《サイト》補題
\ref{sites-lemma-have-left-adjoint} の対応する主張から直ちに従う。
\end{proof}





\section{加群の大域的な型}
\label{sites-modules-section-global}

\begin{definition}
\label{sites-modules-definition-global}
$(\Sh(\mathcal{C}), \mathcal{O})$ を環付きトポスとし，
$\mathcal{F}$ を $\mathcal{O}$-加群の層とする。
\begin{enumerate}
\item $\mathcal{F}$ が $\mathcal{O}$-加群として
$\bigoplus_{i \in I} \mathcal{O}$ の形の層と同型であるとき，
$\mathcal{F}$ は {\it 自由 $\mathcal{O}$-加群}であるという。
\item $\mathcal{F}$ が $\mathcal{O}$-加群として，
有限添字集合 $I$ に対する
$\bigoplus_{i \in I} \mathcal{O}$ の形の層と同型であるとき，
$\mathcal{F}$ は {\it 有限自由}であるという。
\item 自由 $\mathcal{O}$-加群から $\mathcal{F}$ への全射
$$
\bigoplus\nolimits_{i \in I} \mathcal{O} \longrightarrow \mathcal{F}
$$
が存在するとき，$\mathcal{F}$ は {\it 大域切断で生成される}という。
\item $r \geq 0$ とする。全射
$\mathcal{O}^{\oplus r} \to \mathcal{F}$ が存在するとき，
$\mathcal{F}$ は {\it $r$ 個の大域切断で生成される}という。
\item ある $r \geq 0$ に対して $\mathcal{F}$ が $r$ 個の大域切断で
生成されるとき，$\mathcal{F}$ は
{\it 有限個の大域切断で生成される}という。
\item $\mathcal{O}$-加群の完全列
$$
\bigoplus\nolimits_{j \in J} \mathcal{O} \longrightarrow
\bigoplus\nolimits_{i \in I} \mathcal{O} \longrightarrow
\mathcal{F} \longrightarrow 0
$$
が存在するとき，$\mathcal{F}$ は {\it 大域表示}をもつという。
\item $I$ と $J$ が有限集合であるような $\mathcal{O}$-加群の完全列
$$
\bigoplus\nolimits_{j \in J} \mathcal{O} \longrightarrow
\bigoplus\nolimits_{i \in I} \mathcal{O} \longrightarrow
\mathcal{F} \longrightarrow 0
$$
が存在するとき，$\mathcal{F}$ は {\it 大域有限表示}をもつという。
\end{enumerate}
\end{definition}

\noindent
任意の集合 $I$ に対して直和
$\bigoplus_{i \in I} \mathcal{O}$ が存在し
（補題 \ref{sites-modules-lemma-limits-colimits}），前層
$U \mapsto \bigoplus_{i \in I} \mathcal{O}(U)$ の層化であることに
注意せよ。この加群を {\it 集合 $I$ 上の自由 $\mathcal{O}$-加群}という。

\begin{lemma}
\label{sites-modules-lemma-global-pullback}
写像
$(f, f^\sharp) :
(\Sh(\mathcal{C}), \mathcal{O}_\mathcal{C})
\to
(\Sh(\mathcal{D}), \mathcal{O}_\mathcal{D})$
を環付きトポスの射とし，
$\mathcal{F}$ を $\mathcal{O}_\mathcal{D}$-加群とする。
\begin{enumerate}
\item $\mathcal{F}$ が自由ならば，$f^*\mathcal{F}$ も自由である。
\item $\mathcal{F}$ が有限自由ならば，$f^*\mathcal{F}$ も有限自由である。
\item $\mathcal{F}$ が大域切断で生成されるならば，
$f^*\mathcal{F}$ も大域切断で生成される。
\item $r \geq 0$ とする。$\mathcal{F}$ が $r$ 個の大域切断で
生成されるならば，$f^*\mathcal{F}$ も $r$ 個の大域切断で生成される。
\item $\mathcal{F}$ が有限個の大域切断で生成されるならば，
$f^*\mathcal{F}$ も有限個の大域切断で生成される。
\item $\mathcal{F}$ が大域表示をもつならば，
$f^*\mathcal{F}$ も大域表示をもつ。
\item $\mathcal{F}$ が大域有限表示をもつならば，
$f^*\mathcal{F}$ も大域有限表示をもつ。
\end{enumerate}
\end{lemma}

\begin{proof}
$f^*$ は任意の余極限と可換し
（補題 \ref{sites-modules-lemma-exactness-pushforward-pullback}），かつ
$f^*\mathcal{O}_\mathcal{D} = \mathcal{O}_\mathcal{C}$ だからである。
\end{proof}






\section{加群の内在的性質}
\label{sites-modules-section-intrinsic}

\noindent
$\mathcal{P}$ を環付きトポス上の加群の層の性質とする。
環付きトポスの同値
$(f, f^\sharp) :
(\Sh(\mathcal{C}'), \mathcal{O}')
\to
(\Sh(\mathcal{C}), \mathcal{O})$
に対して常に
$\mathcal{P}(\mathcal{F}) \Leftrightarrow \mathcal{P}(f^*\mathcal{F})$
が成り立つとき，$\mathcal{P}$ を {\it 内在的性質}という。
例えば，自由であるという性質は内在的である。実際，集合 $I$ 上の
自由 $\mathcal{O}$-加群は，可変な
$\mathcal{F} \in \textit{Mod}(\mathcal{O})$ に対する性質
$$
\Mor_{\textit{Mod}(\mathcal{O})}(
\bigoplus\nolimits_{i \in I} \mathcal{O},
\mathcal{F})
=
\prod\nolimits_{i \in I} \Mor_{\Sh(\mathcal{C})}(\{*\},
\mathcal{F})
$$
によって特徴づけられる。別の方法として，補題
\ref{sites-modules-lemma-global-pullback} を用いても，自由であることが内在的であると
分かる。実際，同じ理由により，定義 \ref{sites-modules-definition-global} で
定義した各性質は内在的である。
では，$\mathcal{O}$-加群の他の内在的性質をどのように定義すればよいだろうか。

\medskip\noindent
補題 \ref{sites-modules-lemma-morphism-ringed-topoi-comes-from-morphism-ringed-sites}
の要点は次のとおりである。トポス上の $\mathcal{O}$-加群の
内在的性質 $\mathcal{P}$ を定義したいとする。このとき，次のように
進めることができる：
\begin{enumerate}
\item 任意のサイト $\mathcal{C}$，$\mathcal{C}$ 上の任意の環の層
$\mathcal{O}$，および任意の $\mathcal{O}$-加群 $\mathcal{F}$ に対し，
対応する性質
$\mathcal{P}(\mathcal{C}, \mathcal{O}, \mathcal{F})$ を定義する。
\item 任意の二つのサイト $\mathcal{C}$，$\mathcal{C}'$，任意の
特殊余連続関手 $u : \mathcal{C} \to \mathcal{C}'$，
$\mathcal{C}$ 上の任意の環の層 $\mathcal{O}$，および任意の
$\mathcal{O}$-加群 $\mathcal{F}$ に対して
$$
\mathcal{P}(\mathcal{C}, \mathcal{O}, \mathcal{F})
\Leftrightarrow
\mathcal{P}(\mathcal{C}', g_*\mathcal{O}, g_*\mathcal{F})
$$
を示す。ここで $g : \Sh(\mathcal{C}) \to \Sh(\mathcal{C}')$ は
$u$ に付随するトポスの同値である。
\end{enumerate}
この場合，任意の環付きトポス $(\Sh(\mathcal{C}), \mathcal{O})$ と
任意の $\mathcal{O}$-加群の層 $\mathcal{F}$ に対し，
$\mathcal{P}(\mathcal{C}, \mathcal{O}, \mathcal{F})$ が真であるとき，
単に $\mathcal{F}$ は性質 $\mathcal{P}$ をもつという。
補題 \ref{sites-modules-lemma-morphism-ringed-topoi-comes-from-morphism-ringed-sites} と
上の (2) を合わせると，この定義が適切であることが保証される。

\medskip\noindent
さらに，同じ補題
\ref{sites-modules-lemma-morphism-ringed-topoi-comes-from-morphism-ringed-sites} により，
加えて次が成り立つならば，
\begin{enumerate}
\item[(3)] 任意の環付きサイトの射
$(f, f^\sharp) :
(\mathcal{C}, \mathcal{O}_\mathcal{C})
\to
(\mathcal{D}, \mathcal{O}_\mathcal{D})$
であって，$f$ が《サイト》命題
\ref{sites-proposition-get-morphism} の仮定を満たす関手
$u : \mathcal{D} \to \mathcal{C}$ によって与えられるもの，および任意の
$\mathcal{O}_\mathcal{D}$-加群 $\mathcal{G}$ に対し，
$$
\mathcal{P}(\mathcal{D}, \mathcal{O}_\mathcal{D}, \mathcal{F})
\Rightarrow
\mathcal{P}(\mathcal{C}, \mathcal{O}_\mathcal{C}, f^*\mathcal{F})
$$
が成り立つ。
\end{enumerate}
性質 $\mathcal{P}$ は任意の環付きトポスの射に関する加群の引き戻しで
保たれることも保証される。

\medskip\noindent
以下の節ではこの方法を用いて，
局所自由，
局所的に切断で生成されること，
局所的に $r$ 個の切断で生成されること，
有限型，
有限表示，
準連接，および
連接
が加群の内在的性質であることを見る。

\medskip\noindent
ここで概説した手間のかかる手順よりも，これらの概念の内在的な定義を
見つける方が満足のいく方法かもしれない。
他方，これらの定義を適用したい多くの幾何学的状況では，具体的な
環付きサイトと具体的な加群の層が与えられているので，この言語に
あらかじめ適合した定義があることには利点がある。




\section{環付きサイトの局所化}
\label{sites-modules-section-localize}

\noindent
$(\mathcal{C}, \mathcal{O})$ を環付きサイトとし，
$U \in \Ob(\mathcal{C})$ とする。
この設定における《サイト》第 \ref{sites-section-localize} 節の
結果に対応するものを説明する。

\medskip\noindent
$\mathcal{O}$ のサイト $\mathcal{C}/U$ への制限
$\mathcal{O}_U = j_U^{-1}\mathcal{O}$ を $\mathcal{O}_U$ と書く。
これは簡単な規則 $\mathcal{O}_U(V/U) = \mathcal{O}(V)$ で記述される。
この記法の下で，局所化射 $j_U$ は環付きトポスの射
$$
(j_U, j_U^\sharp) :
(\Sh(\mathcal{C}/U), \mathcal{O}_U)
\longrightarrow
(\Sh(\mathcal{C}), \mathcal{O})
$$
となる。すなわち，
$j_U^\sharp : j_U^{-1}\mathcal{O} \to \mathcal{O}_U$ を恒等写像とする。
さらに，加群の順像と引き戻しについて次の記述を得る。

\begin{definition}
\label{sites-modules-definition-localize-ringed-site}
$(\mathcal{C}, \mathcal{O})$ を環付きサイトとし，
$U \in \Ob(\mathcal{C})$ とする。
\begin{enumerate}
\item 環付きサイト $(\mathcal{C}/U, \mathcal{O}_U)$ を
{\it 環付きサイト $(\mathcal{C}, \mathcal{O})$ の対象 $U$ における
局所化}という。
\item 環付きトポスの射
$(j_U, j_U^\sharp) :
(\Sh(\mathcal{C}/U), \mathcal{O}_U)
\to
(\Sh(\mathcal{C}), \mathcal{O})$
を {\it 局所化射}という。
\item 関手
$j_{U*} : \textit{Mod}(\mathcal{O}_U) \to \textit{Mod}(\mathcal{O})$
を {\it 順像関手}という。
\item $\mathcal{C}$ 上の $\mathcal{O}$-加群の層 $\mathcal{F}$ に対し，
層 $j_U^*\mathcal{F}$ を {\it $\mathcal{F}$ の
$\mathcal{C}/U$ への制限}という。これを
$\mathcal{F}|_{\mathcal{C}/U}$，あるいは単に $\mathcal{F}|_U$ と
書くこともある。これは簡単な規則
$j_U^*(\mathcal{F})(X/U) = \mathcal{F}(X)$ で記述される。
\item 制限の左随伴
$j_{U!} : \textit{Mod}(\mathcal{O}_U) \to \textit{Mod}(\mathcal{O})$
を {\it 零による延長}という。補題 \ref{sites-modules-lemma-extension-by-zero} と
\ref{sites-modules-lemma-extension-by-zero-exact} により，これは存在して完全である。
\end{enumerate}
\end{definition}

\noindent
位相的な場合と同様に（《層》第
\ref{sheaves-section-open-immersions} 節参照），零による延長関手
$j_{U!}$ は，集合の層上で定義される空集合による延長 $j_{U!}$ とは
異なる。零による延長を定義する補題は次のとおりである。

\begin{lemma}
\label{sites-modules-lemma-extension-by-zero}
$(\mathcal{C}, \mathcal{O})$ を環付きサイトとし，
$U \in \Ob(\mathcal{C})$ とする。制限関手
$j_U^* : \textit{Mod}(\mathcal{O}) \to \textit{Mod}(\mathcal{O}_U)$
は左随伴
$j_{U!} : \textit{Mod}(\mathcal{O}_U) \to \textit{Mod}(\mathcal{O})$.
をもつ。したがって
$$
\Mor_{\textit{Mod}(\mathcal{O}_U)}(\mathcal{G}, j_U^*\mathcal{F})
=
\Mor_{\textit{Mod}(\mathcal{O})}(j_{U!}\mathcal{G}, \mathcal{F})
$$
が，$\mathcal{F} \in \Ob(\textit{Mod}(\mathcal{O}))$ および
$\mathcal{G} \in \Ob(\textit{Mod}(\mathcal{O}_U))$ に対して成り立つ。
さらに，$\mathcal{G}$ の零による延長 $j_{U!}\mathcal{G}$ は，
前層
$$
V
\longmapsto
\bigoplus\nolimits_{\varphi \in \Mor_\mathcal{C}(V, U)}
\mathcal{G}(V \xrightarrow{\varphi} U)
$$
に付随する層である。ここで制限写像と $\mathcal{O}$-加群構造は
明らかなものを用いる。
\end{lemma}

\begin{proof}
前層上の $\mathcal{O}$-加群構造を次のように定義する。
$f \in \mathcal{O}(V)$ および
$s \in \mathcal{G}(V \xrightarrow{\varphi} U)$ に対し，
$f \cdot s = fs$ と定義する。ここで
$f \in \mathcal{O}_U(\varphi : V \to U) = \mathcal{O}(V)$
である（$\mathcal{O}_U$ は $\mathcal{O}$ の $\mathcal{C}/U$ への
制限だからである）。

\medskip\noindent
同様に，$\alpha : \mathcal{G} \to \mathcal{F}|_U$ を
$\mathcal{O}_U$-加群の射とする。この場合，補題の前層から
$\mathcal{F}$ への写像
$$
\bigoplus\nolimits_{\varphi \in \Mor_\mathcal{C}(V, U)}
\mathcal{G}(V \xrightarrow{\varphi} U)
\longrightarrow
\mathcal{F}(V)
$$
を，$s \in \mathcal{G}(V \xrightarrow{\varphi} U)$ を
$\alpha(s) \in \mathcal{F}(V)$ に写す規則によって定義できる。
これは明らかに $\mathcal{O}$-線形である。したがって加群の層化の性質
（補題 \ref{sites-modules-lemma-sheafification-presheaf-modules}）により，
$\mathcal{O}$-加群の射
$\alpha' : j_{U!}\mathcal{G} \to \mathcal{F}$ を誘導する。

\medskip\noindent
逆に，$\beta : j_{U!}\mathcal{G} \to \mathcal{F}$ を
$\mathcal{O}$-加群の写像とする。《サイト》第
\ref{sites-section-localize} 節により，集合の層上には
$j^{Sh}_{U!} : \Sh(\mathcal{C}/U) \to \Sh(\mathcal{C})$
という空集合による延長が存在し，$j_U^{-1}$ の左随伴であることを
想起せよ。さらに，$j^{Sh}_{U!}\mathcal{G}$ は前層
$$
V
\longmapsto
\coprod\nolimits_{\varphi \in \Mor_\mathcal{C}(V, U)}
\mathcal{G}(V \xrightarrow{\varphi} U)
$$
に付随する層である。したがって集合の層の自然な写像
$j^{Sh}_{U!}\mathcal{G} \to j_{U!}\mathcal{G}$ が存在する。
この写像を $\beta$ に前合成すると，集合の層の写像
$j^{Sh}_{U!}\mathcal{G} \to \mathcal{F}$ を得る。これは随伴により，
集合の層の写像 $\beta' : \mathcal{G} \to \mathcal{F}|_U$ に対応する。
$\beta'$ が $\mathcal{O}_U$-線形であることを主張する。
実際，$\varphi : V \to U$ を $\mathcal{C}/U$ の対象とし，
$s, s' \in \mathcal{G}(\varphi : V \to U)$，および
$f \in \mathcal{O}(V) = \mathcal{O}_U(\varphi : V \to U)$ とする。
上の議論により，$\mathcal{F}|_U(\varphi : V \to U)$ における
$\beta'(s + s')$ と $\beta'(fs)$ は，$\mathcal{F}(V)$ における
$\beta(s + s')$ と $\beta(fs)$ にそれぞれ対応する。
$\beta$ は準同型なので，主張が従う。

\medskip\noindent
補題の証明を完了するには，構成
$\alpha \mapsto \alpha'$ と $\beta \mapsto \beta'$ が互いに逆であることを
示さなければならない。確認は省略する。
\end{proof}

\noindent
定義 \ref{sites-modules-definition-localize-ringed-site} の状況では，
\begin{equation}
\label{sites-modules-equation-map-lower-shriek-OU-into-module}
\Hom_\mathcal{O}(j_{U!}\mathcal{O}_U, \mathcal{F}) =
\Hom_{\mathcal{O}_U}(\mathcal{O}_U, j_U^*\mathcal{F}) =
\mathcal{F}(U)
\end{equation}
$\mathcal{O}$-加群 $\mathcal{F}$ のすべてに対して成り立つ。
実際，第一の等式は $j_{U!}$ と $j_U^*$ の随伴性により成り立ち，
第二の等式は
$\Hom_{\mathcal{O}_U}(\mathcal{O}_U, j_U^*\mathcal{F}) =
j_U^*\mathcal{F}(U/U) = \mathcal{F}|_U(U/U) = \mathcal{F}(U)$
だからである。

\begin{lemma}
\label{sites-modules-lemma-extension-by-zero-exact}
$(\mathcal{C}, \mathcal{O})$ を環付きサイトとし，
$U \in \Ob(\mathcal{C})$ とする。関手
$j_{U!} : \textit{Mod}(\mathcal{O}_U) \to \textit{Mod}(\mathcal{O})$
は完全である。
\end{lemma}

\begin{proof}
$j_{U!}$ は $j_U^*$ の左随伴なので右完全である
（《圏論》補題 \ref{categories-lemma-exact-adjoint} および
《ホモロジー》第 \ref{homology-section-functors} 節参照）。
したがって，$\mathcal{G}_1 \to \mathcal{G}_2$ が
$\mathcal{O}_U$-加群の単射ならば，
$j_{U!}\mathcal{G}_1 \to j_{U!}\mathcal{G}_2$ も単射であることを
示せば十分である。対象 $V$ 上の前層の切断における写像
（補題 \ref{sites-modules-lemma-extension-by-zero} と同様）は
$$
\bigoplus\nolimits_{\varphi \in \Mor_\mathcal{C}(V, U)}
\mathcal{G}_1(V \xrightarrow{\varphi} U)
\longrightarrow
\bigoplus\nolimits_{\varphi \in \Mor_\mathcal{C}(V, U)}
\mathcal{G}_2(V \xrightarrow{\varphi} U)
$$
であり，仮定により単射である。補題
\ref{sites-modules-lemma-sheafification-exact} により層化は完全なので，
$j_{U!}\mathcal{G}_1 \to j_{U!}\mathcal{G}_2$ は単射であり，結論を得る。
\end{proof}

\begin{lemma}
\label{sites-modules-lemma-j-shriek-reflects-exactness}
$(\mathcal{C}, \mathcal{O})$ を環付きサイトとし，
$U \in \Ob(\mathcal{C})$ とする。$\mathcal{O}_U$-加群の複体
$\mathcal{G}_1 \to \mathcal{G}_2 \to \mathcal{G}_3$ が完全であることと，
$j_{U!}\mathcal{G}_1 \to j_{U!}\mathcal{G}_2 \to j_{U!}\mathcal{G}_3$
が $\mathcal{O}$-加群の列として完全であることとは同値である。
\end{lemma}

\begin{proof}
補題 \ref{sites-modules-lemma-extension-by-zero-exact} により，
$j_{U!}$ が完全であることはすでに分かっている。
したがって
$j_{U!} :  \textit{Mod}(\mathcal{O}_U) \to \textit{Mod}(\mathcal{O})$
が単射と全射を反映することを示せば十分である。

\medskip\noindent
$\textit{Mod}(\mathcal{O}_U)$ の各 $\mathcal{G}$ に対し，随伴の単位
$\mathcal{G} \to j_U^*j_{U!}\mathcal{G}$ がある。
この写像は層の単射であると主張する。実際，補題
\ref{sites-modules-lemma-extension-by-zero} の構成を見ると，この写像は
$\mathcal{C}/U$ の対象 $V/U$ を単射
$$
\mathcal{G}(V/U) \longrightarrow
\bigoplus\nolimits_{\varphi \in \Mor_\mathcal{C}(V, U)}
\mathcal{G}(V \xrightarrow{\varphi} U)
$$
に写す規則の層化である。ここでこの単射は，構造射 $V \to U$ に
対応する直和成分の包含によって与えられる。
層化は完全なので主張が従う。詳細の一部は省略する。

\medskip\noindent
$\mathcal{G} \to \mathcal{G}'$ を
$j_{U!}\mathcal{G} \to j_{U!}\mathcal{G}'$ が単射となる
$\mathcal{O}_U$-加群の写像とする。このとき
$j_U^*j_{U!}\mathcal{G} \to j_U^*j_{U!}\mathcal{G}'$ は単射である
（制限は完全である）。したがって
$\mathcal{G} \to j_U^*j_{U!}\mathcal{G}'$ は単射であり，ゆえに
$\mathcal{G} \to \mathcal{G}'$ は単射である。
$j_{U!}$ は単射を反映すると結論する。

\medskip\noindent
$a : \mathcal{G} \to \mathcal{G}'$ を，
$j_{U!}\mathcal{G} \to j_{U!}\mathcal{G}'$ が全射となる
$\mathcal{O}_U$-加群の写像とする。
$\mathcal{H}$ を $a$ の余核とする。$j_{U!}$ は完全なので
$j_{U!}\mathcal{H} = 0$ である。上により，写像
$\mathcal{H} \to j^*_U j_{U!}\mathcal{H}$ は単射である。
したがって所望のとおり $\mathcal{H} = 0$ である。
\end{proof}

\begin{lemma}
\label{sites-modules-lemma-relocalize}
$(\mathcal{C}, \mathcal{O})$ を環付きサイトとし，
$f : V \to U$ を $\mathcal{C}$ の射とする。このとき環付きトポスの
可換図式
$$
\xymatrix{
(\Sh(\mathcal{C}/V), \mathcal{O}_V)
\ar[rd]_{(j_V, j_V^\sharp)} \ar[rr]_{(j, j^\sharp)} & &
(\Sh(\mathcal{C}/U), \mathcal{O}_U)
\ar[ld]^{(j_U, j_U^\sharp)} \\
& (\Sh(\mathcal{C}), \mathcal{O}) &
}
$$
が存在する。ここで $(j, j^\sharp)$ は，環付きサイト
$(\mathcal{C}/V, \mathcal{O}_V)$ の対象 $V/U$ に付随する局所化射である。
\end{lemma}

\begin{proof}
確認すべきことは
$j_V^\sharp = j^\sharp \circ j^{-1}(j_U^\sharp)$,
だけである。他のすべては《サイト》補題
\ref{sites-lemma-relocalize} と《サイト》式
(\ref{sites-equation-relocalize}) から直ちに従う。
等式の確認は省略する。
\end{proof}

\begin{remark}
\label{sites-modules-remark-localize-shriek-equal}
補題 \ref{sites-modules-lemma-extension-by-zero} の状況では，図式
$$
\xymatrix{
\textit{Mod}(\mathcal{O}_U) \ar[r]_{j_{U!}} \ar[d]_{forget} &
\textit{Mod}(\mathcal{O}_\mathcal{C}) \ar[d]^{forget} \\
\textit{Ab}(\mathcal{C}/U) \ar[r]^{j^{Ab}_{U!}} &
\textit{Ab}(\mathcal{C})
}
$$
は可換である。これは補題における関手 $j_{U!}$ の明示的な記述から
明らかである。
\end{remark}

\begin{remark}
\label{sites-modules-remark-localize-presheaves}
局所化と加群の前層について述べる。《サイト》注
\ref{sites-remark-localize-presheaves} を参照せよ。
$\mathcal{C}$ を圏，$\mathcal{O}$ を環の前層，
$U$ を $\mathcal{C}$ の対象とする。
厳密にいえば，関手 $j_U^*$，$j_{U*}$，$j_{U!}$ は
$\mathcal{O}$-加群の前層に対して定義されていない。
しかしもちろん，前層を $\mathcal{C}$ 上のカオス位相に関する層と
みなせる（《サイト》例 \ref{sites-example-indiscrete} 参照）。
したがって関手
$$
j_U^* :
\textit{PMod}(\mathcal{O})
\longrightarrow
\textit{PMod}(\mathcal{O}_U)
$$
と関手
$$
j_{U*}, j_{U!} :
\textit{PMod}(\mathcal{O}_U)
\longrightarrow
\textit{PMod}(\mathcal{O})
$$
を得る。これらはそれぞれ $j_U^*$ の右随伴と左随伴である。
補題 \ref{sites-modules-lemma-extension-by-zero} の証明を見ると，
$j_{U!}\mathcal{G}$ は前層
$$
V \longmapsto
\bigoplus\nolimits_{\varphi \in \Mor_\mathcal{C}(V, U)}
\mathcal{G}(V \xrightarrow{\varphi} U)
$$
であることが分かる。さらに関手 $j_{U!}$ は完全である
（離散位相の場合の補題 \ref{sites-modules-lemma-extension-by-zero-exact} による）。
また，$\mathcal{C}$ が実際にサイトであり，$\mathcal{O}$ が実際に
環の層であるならば，図式
$$
\xymatrix{
\textit{Mod}(\mathcal{O}_U) \ar[r]_{j_{U!}} \ar[d]_{forget} &
\textit{Mod}(\mathcal{O}) \\
\textit{PMod}(\mathcal{O}_U) \ar[r]^{j_{U!}} &
\textit{PMod}(\mathcal{O}) \ar[u]_{(\ )^\#}
}
$$
は可換である。
\end{remark}

\begin{lemma}
\label{sites-modules-lemma-restrict-back}
$\mathcal{C}$ をサイトとし，$U \in \Ob(\mathcal{C})$ とする。
$\mathcal{C}$ の各 $X$ から $U$ への射が高々一つであると仮定する。
$\mathcal{F}$ を $\mathcal{C}/U$ 上のアーベル群の層とする。
標準写像 $\mathcal{F} \to j_U^{-1}j_{U!}\mathcal{F}$ と
$j_U^{-1}j_{U*}\mathcal{F} \to \mathcal{F}$ は同型である。
\end{lemma}

\begin{proof}
これは補題 \ref{sites-modules-lemma-back-and-forth} の特別な場合である。
実際，$U$ に関する仮定は，局所化関手
$\mathcal{C}/U \to \mathcal{C}$ の充満忠実性と同値である。
\end{proof}









\section{環付きサイトの射の局所化}
\label{sites-modules-section-localize-morphisms}

\noindent
本節は《サイト》第 \ref{sites-section-localize-morphisms} 節の類似である。

\begin{lemma}
\label{sites-modules-lemma-localize-morphism-ringed-sites}
写像
$(f, f^\sharp) :
(\mathcal{C}, \mathcal{O})
\longrightarrow
(\mathcal{D}, \mathcal{O}')$
を環付きサイトの射とし，$f$ は連続関手
$u : \mathcal{D} \to \mathcal{C}$ によって与えられるとする。
$V$ を $\mathcal{D}$ の対象とし，$U = u(V)$ と置く。
このとき，環の層の標準写像 $(f')^\sharp$ であって，
《サイト》補題 \ref{sites-lemma-localize-morphism} の図式を
環付きトポスの可換図式
$$
\xymatrix{
(\Sh(\mathcal{C}/U), \mathcal{O}_U)
\ar[rr]_{(j_U, j_U^\sharp)} \ar[d]_{(f', (f')^\sharp)} & &
(\Sh(\mathcal{C}), \mathcal{O})
\ar[d]^{(f, f^\sharp)} \\
(\Sh(\mathcal{D}/V), \mathcal{O}'_V)
\ar[rr]^{(j_V, j_V^\sharp)} & &
(\Sh(\mathcal{D}), \mathcal{O}').
}
$$
にするものが存在する。さらに，この状況では
$f'_*j_U^{-1} = j_V^{-1}f_*$ かつ $f'_*j_U^* = j_V^*f_*$ である。
\end{lemma}

\begin{proof}
$(f')^\sharp$ を
$$
(f')^{-1}\mathcal{O}'_V =
(f')^{-1}j_V^{-1}\mathcal{O}' =
j_U^{-1}f^{-1}\mathcal{O}' \xrightarrow{j_U^{-1}f^\sharp}
j_U^{-1}\mathcal{O} = \mathcal{O}_U
$$
と取ればよい。他のすべては《サイト》補題
\ref{sites-lemma-localize-morphism} から従う。
（$j$ が局所化射ならば，加群の層上で $j^{-1} = j^*$ であることに
注意せよ。したがって関手の第一の等式から第二の等式が従う。）
\end{proof}

\begin{lemma}
\label{sites-modules-lemma-relocalize-morphism-ringed-sites}
写像
$(f, f^\sharp) :
(\mathcal{C}, \mathcal{O})
\longrightarrow
(\mathcal{D}, \mathcal{O}')$
を環付きサイトの射とし，$f$ は連続関手
$u : \mathcal{D} \to \mathcal{C}$ によって与えられるとする。
$V \in \Ob(\mathcal{D})$，$U \in \Ob(\mathcal{C})$ とし，
$c : U \to u(V)$ を $\mathcal{C}$ の射とする。
環付きトポスの可換図式
$$
\xymatrix{
(\Sh(\mathcal{C}/U), \mathcal{O}_U)
\ar[rr]_{(j_U, j_U^\sharp)} \ar[d]_{(f_c, f_c^\sharp)} & &
(\Sh(\mathcal{C}), \mathcal{O}) \ar[d]^{(f, f^\sharp)} \\
(\Sh(\mathcal{D}/V), \mathcal{O}'_V)
\ar[rr]^{(j_V, j_V^\sharp)} & &
(\Sh(\mathcal{D}), \mathcal{O}').
}
$$
が存在する。射 $(f_c, f_c^\sharp)$ は，補題
\ref{sites-modules-lemma-localize-morphism-ringed-sites} の射
$$
(f', (f')^\sharp) :
(\Sh(\mathcal{C}/u(V)), \mathcal{O}_{u(V)})
\longrightarrow
(\Sh(\mathcal{D}/V), \mathcal{O}'_V)
$$
と，補題 \ref{sites-modules-lemma-relocalize} の射
$$
(j, j^\sharp) :
(\Sh(\mathcal{C}/U), \mathcal{O}_U)
\to
(\Sh(\mathcal{C}/u(V)), \mathcal{O}_{u(V)})
$$
の合成に等しい。射 $b : V' \to V$，$a : U' \to U$，および
$c' : U' \to u(V')$ が与えられ，図式
$$
\xymatrix{
U' \ar[r]_-{c'} \ar[d]_a & u(V') \ar[d]^{u(b)} \\
U \ar[r]^-c & u(V)
}
$$
が可換であるとする。このとき，環付きトポスの図式
$$
\xymatrix{
(\Sh(\mathcal{C}/U'), \mathcal{O}_{U'})
\ar[rr]_{(j_{U'/U}, j_{U'/U}^\sharp)} \ar[d]_{(f_{c'}, f_{c'}^\sharp)} & &
(\Sh(\mathcal{C}/U), \mathcal{O}_U)
\ar[d]^{(f_c, f_c^\sharp)} \\
(\Sh(\mathcal{D}/V'), \mathcal{O}'_{V'})
\ar[rr]^{(j_{V'/V}, j_{V'/V}^\sharp)} & &
(\Sh(\mathcal{D}/V), \mathcal{O}'_{V'})
}
$$
は可換である。
\end{lemma}

\begin{proof}
トポスの射のレベルでは，これは《サイト》補題
\ref{sites-lemma-relocalize-morphism} である。
図式が環付きトポスの射として可換であることを確認するには，補題
\ref{sites-modules-lemma-relocalize} と
\ref{sites-modules-lemma-localize-morphism-ringed-sites} を，《サイト》補題
\ref{sites-lemma-relocalize-morphism} の証明とまったく同様に用いればよい。
\end{proof}
















\section{環付きトポスの局所化}
\label{sites-modules-section-localize-ringed-topoi}

\noindent
本節は，環付きトポスの設定における《サイト》第
\ref{sites-section-localize-topoi} 節の類似である。

\begin{lemma}
\label{sites-modules-lemma-localize-ringed-topos}
$(\Sh(\mathcal{C}), \mathcal{O})$ を環付きトポスとし，
$\mathcal{F} \in \Sh(\mathcal{C})$ を層とする。
$\mathcal{C}$ 上の層 $\mathcal{H}$ に対し，層
$\mathcal{H} \times \mathcal{F}$ を圏
$\Sh(\mathcal{C})/\mathcal{F}$ の対象とみなしたものを
$\mathcal{H}_\mathcal{F}$ と書く。組
$(\Sh(\mathcal{C})/\mathcal{F}, \mathcal{O}_\mathcal{F})$
は環付きトポスであり，環付きトポスの標準射
$$
(j_\mathcal{F}, j_\mathcal{F}^\sharp) :
(\Sh(\mathcal{C})/\mathcal{F}, \mathcal{O}_\mathcal{F})
\longrightarrow
(\Sh(\mathcal{C}), \mathcal{O})
$$
であって，第 \ref{sites-modules-section-localize} 節の意味で局所化となり，
次を満たすものが存在する：
\begin{enumerate}
\item 関手 $j_\mathcal{F}^{-1}$ は関手
$\mathcal{H} \mapsto \mathcal{H}_\mathcal{F}$ である。
\item 関手 $j_\mathcal{F}^*$ は関手
$\mathcal{H} \mapsto \mathcal{H}_\mathcal{F}$ である。
\item 集合の層上の関手 $j_{\mathcal{F}!}$ は忘却関手
$\mathcal{G}/\mathcal{F} \mapsto \mathcal{G}$ である。
\item 加群の層上の関手 $j_{\mathcal{F}!}$ は，
$\mathcal{O}_\mathcal{F}$-加群
$\varphi : \mathcal{G} \to \mathcal{F}$ に，前層
$$
V \longmapsto
\bigoplus\nolimits_{s \in \mathcal{F}(V)}
\{\sigma \in \mathcal{G}(V) \mid \varphi(\sigma) = s \}
$$
の層化である $\mathcal{O}$-加群を対応させる。
ここで $V \in \Ob(\mathcal{C})$ である。
\end{enumerate}
\end{lemma}

\begin{proof}
《サイト》補題 \ref{sites-lemma-localize-topos} により，
$\Sh(\mathcal{C})/\mathcal{F}$ はトポスであり，(1) と (3) が成り立つ。
特に，これから
$j_\mathcal{F}^{-1}\mathcal{O} = \mathcal{O}_\mathcal{F}$
であり，$\mathcal{O}_\mathcal{F}$ が環の層であることが分かる。
したがって写像 $j_\mathcal{F}^\sharp$ を恒等写像に取ることができ，
特に (2) が成り立つ。また《サイト》補題
\ref{sites-lemma-localize-topos} の証明により，$\mathcal{C}$ は
すべての有限極限をもつサイトで，位相は劣標準的であり，
$\mathcal{F} = h_U$ が $\mathcal{C}$ のある対象 $U$ に対して成り立つと
仮定できる。このとき (4) は補題 \ref{sites-modules-lemma-extension-by-zero} における
$j_{\mathcal{F}!}$ の記述から従う。
別の方法として，(4) で記述した関手が $j_\mathcal{F}^*$ の左随伴である
ことを直接示すこともできる。
\end{proof}

\begin{definition}
\label{sites-modules-definition-localize-ringed-topos}
$(\Sh(\mathcal{C}), \mathcal{O})$ を環付きトポスとし，
$\mathcal{F} \in \Sh(\mathcal{C})$ とする。
\begin{enumerate}
\item 環付きトポス
$(\Sh(\mathcal{C})/\mathcal{F}, \mathcal{O}_\mathcal{F})$
を {\it 環付きトポス $(\Sh(\mathcal{C}), \mathcal{O})$ の
$\mathcal{F}$ における局所化}という。
\item 補題 \ref{sites-modules-lemma-localize-ringed-topos} の環付きトポスの射
$(j_\mathcal{F}, j_\mathcal{F}^\sharp) :
(\Sh(\mathcal{C})/\mathcal{F}, \mathcal{O}_\mathcal{F})
\to
(\Sh(\mathcal{C}), \mathcal{O})$
を {\it 局所化射}という。
\end{enumerate}
\end{definition}

\noindent
サイトの章で確立した慣例に従い，意味をなす場合には常に，
トポス上の局所化の構成がサイト上の局所化の構成と両立することを確認する。

\begin{lemma}
\label{sites-modules-lemma-localize-compare}
補題 \ref{sites-modules-lemma-localize-ringed-topos} の
$(\Sh(\mathcal{C}), \mathcal{O})$ と
$\mathcal{F} \in \Sh(\mathcal{C})$ を考える。
$\mathcal{C}$ のある対象 $U$ に対して $\mathcal{F} = h_U^\#$ ならば，
《サイト》補題 \ref{sites-lemma-essential-image-j-shriek} の同一視
$\Sh(\mathcal{C}/U) = \Sh(\mathcal{C})/h_U^\#$
の下で，次が成り立つ：
\begin{enumerate}
\item 標準的に $\mathcal{O}_U = \mathcal{O}_\mathcal{F}$ である。
\item これらの同一視の下で
$(j_\mathcal{F}, j_\mathcal{F}^\sharp) = (j_U, j_U^\sharp)$ である。
\end{enumerate}
\end{lemma}

\begin{proof}
台となるトポスに関する主張は《サイト》補題
\ref{sites-lemma-localize-compare} である。
$\mathcal{O}_U$ は $\mathcal{O}$ の制限であり，《サイト》補題
\ref{sites-lemma-compute-j-shriek-restrict} により，《サイト》補題
\ref{sites-lemma-essential-image-j-shriek} の同値の下で
$\mathcal{O} \times h_U^\#$ に対応することに注意せよ。
$\mathcal{O}_\mathcal{F}$ の定義から (1) を得る。
残るのは，この同一視の下で
$j_\mathcal{F}^\sharp = j_U^\sharp$ を示すことである。
確認は省略する。
\end{proof}

\noindent
局所化は次の二つの意味で関手的である。
局所化を「再局所化」することができる
（補題 \ref{sites-modules-lemma-relocalize-ringed-topos} 参照）。
また，環付きトポスの射が与えられたとき，下側の層の逆像において
上側を局所化して，環付きトポスの可換図式を得ることができる
（補題 \ref{sites-modules-lemma-localize-morphism-ringed-topoi} 参照）。

\begin{lemma}
\label{sites-modules-lemma-relocalize-ringed-topos}
$(\Sh(\mathcal{C}), \mathcal{O})$ を環付きトポスとする。
$s : \mathcal{G} \to \mathcal{F}$ が $\mathcal{C}$ 上の層の射ならば，
環付きトポスの射からなる自然な可換図式
$$
\xymatrix{
(\Sh(\mathcal{C})/\mathcal{G}, \mathcal{O}_\mathcal{G})
\ar[rd]_{(j_\mathcal{G}, j_\mathcal{G}^\sharp)} \ar[rr]_{(j, j^\sharp)} & &
(\Sh(\mathcal{C})/\mathcal{F}, \mathcal{O}_\mathcal{F})
\ar[ld]^{(j_\mathcal{F}, j_\mathcal{F}^\sharp)} \\
& (\Sh(\mathcal{C}), \mathcal{O}) &
}
$$
が存在する。ここで $(j, j^\sharp)$ は，環付きトポス
$(\Sh(\mathcal{C})/\mathcal{F}, \mathcal{O}_\mathcal{F})$
の対象 $\mathcal{G}/\mathcal{F}$ における局所化射である。
\end{lemma}

\begin{proof}
等式
$j_\mathcal{G}^\sharp = j^\sharp \circ j^{-1}(j_\mathcal{F}^\sharp)$。
を除くすべての主張は《サイト》補題
\ref{sites-lemma-relocalize-topos} から従う。この等式の確認は省略する。
\end{proof}

\begin{lemma}
\label{sites-modules-lemma-relocalize-compare}
補題 \ref{sites-modules-lemma-relocalize-ringed-topos} の
$(\Sh(\mathcal{C}), \mathcal{O})$ と
$s : \mathcal{G} \to \mathcal{F}$ を考える。
$\mathcal{G} = h_V^\#$，$\mathcal{F} = h_U^\#$ であり，
$s$ が $f$ によって誘導されるような $\mathcal{C}$ の射
$f : V \to U$ が存在するならば，補題 \ref{sites-modules-lemma-relocalize} と
補題 \ref{sites-modules-lemma-relocalize-ringed-topos} の図式は，補題
\ref{sites-modules-lemma-localize-compare} の同一視
$(j_\mathcal{F}, j_\mathcal{F}^\sharp) = (j_U, j_U^\sharp)$
および
$(j_\mathcal{G}, j_\mathcal{G}^\sharp) = (j_V, j_V^\sharp)$
の下で一致する。
\end{lemma}

\begin{proof}
二つの写像 $j^\sharp$ が一致するという主張を除き，すべての主張は
《サイト》補題 \ref{sites-lemma-relocalize-compare} から従う。
この主張は，どちらの場合にも写像 $j^\sharp$ が単に恒等写像であることから
成り立つ。詳細の一部は省略する。
\end{proof}










\section{環付きトポスの射の局所化}
\label{sites-modules-section-localize-morphisms-ringed-topoi}

\noindent
本節は《サイト》第
\ref{sites-section-localize-morphisms-topoi} 節の類似である。

\begin{lemma}
\label{sites-modules-lemma-localize-morphism-ringed-topoi}
写像
$$
f :
(\Sh(\mathcal{C}), \mathcal{O})
\longrightarrow
(\Sh(\mathcal{D}), \mathcal{O}')
$$
を環付きトポスの射とする。$\mathcal{G}$ を $\mathcal{D}$ 上の層とし，
$\mathcal{F} = f^{-1}\mathcal{G}$ と置く。
このとき，環付きトポスの可換図式
$$
\xymatrix{
(\Sh(\mathcal{C})/\mathcal{F}, \mathcal{O}_\mathcal{F})
\ar[rr]_{(j_\mathcal{F}, j_\mathcal{F}^\sharp)}
\ar[d]_{(f', (f')^\sharp)} & &
(\Sh(\mathcal{C}), \mathcal{O}) \ar[d]^{(f, f^\sharp)} \\
(\Sh(\mathcal{D})/\mathcal{G}, \mathcal{O}'_\mathcal{G})
\ar[rr]^{(j_\mathcal{G}, j_\mathcal{G}^\sharp)} & &
(\Sh(\mathcal{D}), \mathcal{O}')
}
$$
が存在する。さらに
$f'_*j_\mathcal{F}^{-1} = j_\mathcal{G}^{-1}f_*$ および
$f'_*j_\mathcal{F}^* = j_\mathcal{G}^*f_*$ が成り立ち，
射 $f'$ は次の規則によって特徴づけられる：
$$
(f')^{-1}(\mathcal{H} \xrightarrow{\varphi} \mathcal{G})
=
(f^{-1}\mathcal{H} \xrightarrow{f^{-1}\varphi} \mathcal{F}).
$$
\end{lemma}

\begin{proof}
《サイト》補題 \ref{sites-lemma-localize-morphism-topoi} により，
台となるトポスの図式，等式
$f'_*j_\mathcal{F}^{-1} = j_\mathcal{G}^{-1}f_*$，および
$(f')^{-1}$ の記述が得られる。
$(f')^\sharp$ の定義には写像
$$
(f')^\sharp :
\mathcal{O}'_\mathcal{G} =
j_\mathcal{G}^{-1} \mathcal{O}'
\xrightarrow{j_\mathcal{G}^{-1}f^\sharp}
j_\mathcal{G}^{-1} f_*\mathcal{O} =
f'_* j_\mathcal{F}^{-1}\mathcal{O} =
f'_* \mathcal{O}_\mathcal{F}
$$
を用いる。これは，同値な形では写像
$$
(f')^\sharp :
(f')^{-1}\mathcal{O}'_\mathcal{G} =
(f')^{-1}j_\mathcal{G}^{-1} \mathcal{O}' =
j_\mathcal{F}^{-1}f^{-1}\mathcal{O}'
\xrightarrow{j_\mathcal{F}^{-1}f^\sharp}
j_\mathcal{F}^{-1} \mathcal{O} =
\mathcal{O}_\mathcal{F}.
$$
である。この二つの写像が実際に互いに随伴することの確認は省略する。
$(f')^\sharp$ の第二の構成から，この図式は環付きトポスの $2$-圏で
可換であることが分かる（写像 $j_\mathcal{F}^\sharp$ と
$j_\mathcal{G}^\sharp$ は恒等写像だからである）。
最後に，等式 $f'_*j_\mathcal{F}^* = j_\mathcal{G}^*f_*$ は，等式
$f'_*j_\mathcal{F}^{-1} = j_\mathcal{G}^{-1}f_*$ と，
加群の層に対する引き戻しが集合の層に対する引き戻しと一致することから従う。
補題 \ref{sites-modules-lemma-localize-ringed-topos} を参照せよ。
\end{proof}

\begin{lemma}
\label{sites-modules-lemma-localize-morphism-compare}
写像
$$
f :
(\Sh(\mathcal{C}), \mathcal{O})
\longrightarrow
(\Sh(\mathcal{D}), \mathcal{O}')
$$
を環付きトポスの射とする。
$\mathcal{G}$ を $\mathcal{D}$ 上の層とし，
$\mathcal{F} = f^{-1}\mathcal{G}$ と置く。
$f$ が連続関手 $u : \mathcal{D} \to \mathcal{C}$ によって与えられ，
$\mathcal{G} = h_V^\#$ ならば，補題
\ref{sites-modules-lemma-localize-morphism-ringed-sites} と補題
\ref{sites-modules-lemma-localize-morphism-ringed-topoi} の可換図式は，補題
\ref{sites-modules-lemma-localize-compare} の同一視の下で一致する。
\end{lemma}

\begin{proof}
トポスの射のレベルでは，これは《サイト》補題
\ref{sites-lemma-localize-morphism-compare} である。
補題 \ref{sites-modules-lemma-localize-morphism-ringed-sites} と補題
\ref{sites-modules-lemma-localize-morphism-ringed-topoi} の証明で
$(f')^\sharp$ を定める公式は一致するので，環付きトポスの射のレベルでも
同じことが成り立つ。
\end{proof}

\begin{lemma}
\label{sites-modules-lemma-relocalize-morphism-ringed-topoi}
写像
$(f, f^\sharp) :
(\Sh(\mathcal{C}), \mathcal{O})
\to
(\Sh(\mathcal{D}), \mathcal{O}')$
を環付きトポスの射とする。
$\mathcal{G}$ を $\mathcal{D}$ 上の層，$\mathcal{F}$ を $\mathcal{C}$ 上の層とし，
$s : \mathcal{F} \to f^{-1}\mathcal{G}$ を層の射とする。
環付きトポスの可換図式
$$
\xymatrix{
(\Sh(\mathcal{C})/\mathcal{F}, \mathcal{O}_\mathcal{F})
\ar[rr]_{(j_\mathcal{F}, j_\mathcal{F}^\sharp)}
\ar[d]_{(f_c, f_c^\sharp)} & &
(\Sh(\mathcal{C}), \mathcal{O})
\ar[d]^{(f, f^\sharp)} \\
(\Sh(\mathcal{D})/\mathcal{G}, \mathcal{O}'_\mathcal{G})
\ar[rr]^{(j_\mathcal{G}, j_\mathcal{G}^\sharp)} & &
(\Sh(\mathcal{D}), \mathcal{O}').
}
$$
が存在する。射 $(f_s, f_s^\sharp)$ は，補題
\ref{sites-modules-lemma-localize-morphism-ringed-topoi} の射
$$
(f', (f')^\sharp) :
(\Sh(\mathcal{C})/f^{-1}\mathcal{G}, \mathcal{O}_{f^{-1}\mathcal{G}})
\longrightarrow
(\Sh(\mathcal{D})/{\mathcal{G}}, \mathcal{O}'_\mathcal{G})
$$
と，補題 \ref{sites-modules-lemma-relocalize-ringed-topos} の射
$$
(j, j^\sharp) :
(\Sh(\mathcal{C})/\mathcal{F}, \mathcal{O}_\mathcal{F})
\to
(\Sh(\mathcal{C})/f^{-1}\mathcal{G}, \mathcal{O}_{f^{-1}\mathcal{G}})
$$
の合成に等しい。射 $b : \mathcal{G}' \to \mathcal{G}$，
$a : \mathcal{F}' \to \mathcal{F}$，および
$s' : \mathcal{F}' \to f^{-1}\mathcal{G}'$ が与えられ，図式
$$
\xymatrix{
\mathcal{F}' \ar[r]_-{s'} \ar[d]_a &
f^{-1}\mathcal{G}' \ar[d]^{f^{-1}b} \\
\mathcal{F} \ar[r]^-s &
f^{-1}\mathcal{G}
}
$$
が可換であるとする。このとき，環付きトポスの図式
$$
\xymatrix{
(\Sh(\mathcal{C})/\mathcal{F}', \mathcal{O}_{\mathcal{F}'})
\ar[rr]_{(j_{\mathcal{F}'/\mathcal{F}}, j_{\mathcal{F}'/\mathcal{F}}^\sharp)}
\ar[d]_{(f_{s'}, f_{s'}^\sharp)} & &
(\Sh(\mathcal{C})/\mathcal{F}, \mathcal{O}_\mathcal{F})
\ar[d]^{(f_s, f_s^\sharp)} \\
(\Sh(\mathcal{D})/\mathcal{G}', \mathcal{O}'_{\mathcal{G}'})
\ar[rr]^{(j_{\mathcal{G}'/\mathcal{G}}, j_{\mathcal{G}'/\mathcal{G}}^\sharp)}
& &
(\Sh(\mathcal{D})/\mathcal{G}, \mathcal{O}'_{\mathcal{G}'})
}
$$
は可換である。
\end{lemma}

\begin{proof}
トポスの射のレベルでは，これは《サイト》補題
\ref{sites-lemma-relocalize-morphism-topoi} である。
図式が環付きトポスの射として可換であることを確認するには，補題
\ref{sites-modules-lemma-relocalize-ringed-topos} と
\ref{sites-modules-lemma-localize-morphism-ringed-topoi} の可換図式を用いればよい。
\end{proof}

\begin{lemma}
\label{sites-modules-lemma-relocalize-morphism-compare}
補題 \ref{sites-modules-lemma-relocalize-morphism-ringed-topoi} と同じく，写像
$(f, f^\sharp) :
(\Sh(\mathcal{C}), \mathcal{O})
\to
(\Sh(\mathcal{D}), \mathcal{O}')$,
$s : \mathcal{F} \to f^{-1}\mathcal{G}$ を取る。
$f$ が連続関手 $u : \mathcal{D} \to \mathcal{C}$ によって与えられ，
$\mathcal{G} = h_V^\#$，$\mathcal{F} = h_U^\#$ であり，
$s$ が射 $c : U \to u(V)$ から来るならば，補題
\ref{sites-modules-lemma-relocalize-morphism-ringed-sites} と補題
\ref{sites-modules-lemma-relocalize-morphism-ringed-topoi} の可換図式は，補題
\ref{sites-modules-lemma-localize-compare} の同一視の下で一致する。
\end{lemma}

\begin{proof}
これは補題 \ref{sites-modules-lemma-relocalize-compare} と
\ref{sites-modules-lemma-localize-morphism-compare} から形式的に従う。
\end{proof}

















\section{加群の局所的な型}
\label{sites-modules-section-local}

\noindent
第 \ref{sites-modules-section-intrinsic} 節で説明した一般方針に従い，まず環付きサイト上の
加群の層について局所的な型を定義し，次いで直ちに，これらの型が内在的で
あること，したがって環付きトポス上の加群の層にも意味をもつことを示す。

\begin{definition}
\label{sites-modules-definition-site-local}
$(\mathcal{C}, \mathcal{O})$ を環付きサイトとし，
$\mathcal{F}$ を $\mathcal{O}$-加群の層とする。
定義 \ref{sites-modules-definition-global} で定めた概念を自由に用いる。
\begin{enumerate}
\item $\mathcal{C}$ の任意の対象 $U$ に対し，各制限
$\mathcal{F}|_{\mathcal{C}/U_i}$ が自由 $\mathcal{O}_{U_i}$-加群となるような
$\mathcal{C}$ の被覆 $\{U_i \to U\}_{i \in I}$ が存在するとき，
$\mathcal{F}$ は {\it 局所自由}であるという。
\item $\mathcal{C}$ の任意の対象 $U$ に対し，各制限
$\mathcal{F}|_{\mathcal{C}/U_i}$ が有限自由
$\mathcal{O}_{U_i}$-加群となるような被覆
$\{U_i \to U\}_{i \in I}$ が存在するとき，
$\mathcal{F}$ は {\it 有限局所自由}であるという。
\item $\mathcal{C}$ の任意の対象 $U$ に対し，各制限
$\mathcal{F}|_{\mathcal{C}/U_i}$ が大域切断で生成される
$\mathcal{O}_{U_i}$-加群となるような被覆
$\{U_i \to U\}_{i \in I}$ が存在するとき，
$\mathcal{F}$ は {\it 局所的に切断で生成される}という。
\item $r \geq 0$ とする。$\mathcal{C}$ の任意の対象 $U$ に対し，各制限
$\mathcal{F}|_{\mathcal{C}/U_i}$ が $r$ 個の大域切断で生成される
$\mathcal{O}_{U_i}$-加群となるような被覆
$\{U_i \to U\}_{i \in I}$ が存在するとき，
$\mathcal{F}$ は {\it 局所的に $r$ 個の切断で生成される}という。
\item $\mathcal{C}$ の任意の対象 $U$ に対し，各制限
$\mathcal{F}|_{\mathcal{C}/U_i}$ が有限個の大域切断で生成される
$\mathcal{O}_{U_i}$-加群となるような被覆
$\{U_i \to U\}_{i \in I}$ が存在するとき，
$\mathcal{F}$ は {\it 有限型}であるという。
\item $\mathcal{C}$ の任意の対象 $U$ に対し，各制限
$\mathcal{F}|_{\mathcal{C}/U_i}$ が大域表示をもつ
$\mathcal{O}_{U_i}$-加群となるような被覆
$\{U_i \to U\}_{i \in I}$ が存在するとき，
$\mathcal{F}$ は {\it 準連接}であるという。
\item $\mathcal{C}$ の任意の対象 $U$ に対し，各制限
$\mathcal{F}|_{\mathcal{C}/U_i}$ が大域有限表示をもつ
$\mathcal{O}_{U_i}$-加群となるような被覆
$\{U_i \to U\}_{i \in I}$ が存在するとき，
$\mathcal{F}$ は {\it 有限表示}であるという。
\item $\mathcal{F}$ が有限型であり，かつ $\mathcal{C}$ の任意の対象
$U$ と任意の $s_1, \ldots, s_n \in \mathcal{F}(U)$ に対して写像
$\bigoplus_{i = 1, \ldots, n} \mathcal{O}_U \to \mathcal{F}|_U$
の核が $(\mathcal{C}/U, \mathcal{O}_U)$ 上で有限型であるとき，かつそのときに限り，
$\mathcal{F}$ は {\it 連接}であるという。
\end{enumerate}
\end{definition}

\begin{lemma}
\label{sites-modules-lemma-special-locally-free}
定義 \ref{sites-modules-definition-site-local} の性質 (1) -- (8) はいずれも内在的である
（第 \ref{sites-modules-section-intrinsic} 節の議論を参照）。
\end{lemma}

\begin{proof}
サイト $\mathcal{C}$，$\mathcal{D}$ と特殊余連続関手
$u : \mathcal{C} \to \mathcal{D}$ を取る。
$\mathcal{O}$ を $\mathcal{C}$ 上の環の層，
$\mathcal{F}$ を $\mathcal{C}$ 上の $\mathcal{O}$-加群の層とする。
$g : \Sh(\mathcal{C}) \to \Sh(\mathcal{D})$ を $u$ に付随するトポスの同値とする。
$\mathcal{O}' = g_*\mathcal{O}$ と置き，
$g^\sharp : \mathcal{O}' \to g_*\mathcal{O}$ を恒等写像とする。
最後に $\mathcal{F}' = g_*\mathcal{F}$ と置く。
$\mathcal{P}_l$ を定義 \ref{sites-modules-definition-site-local} に挙げた性質 (1) -- (7) の
いずれかとする（連接の場合は証明の末尾で論じる）。
$\mathcal{P}_g$ を定義 \ref{sites-modules-definition-global} に挙げた対応する性質とする。
$\mathcal{P}_g$ が内在的であることは既に見た。示すべきことは
$\mathcal{P}_l(\mathcal{C}, \mathcal{O}, \mathcal{F})$
が成り立つことと
$\mathcal{P}_l(\mathcal{D}, \mathcal{O}', \mathcal{F}')$
が成り立つことが同値であることである。

\medskip\noindent
$\mathcal{F}$ が $\mathcal{P}_l$ をもつと仮定する。
$V$ を $\mathcal{D}$ の対象とする。特殊余連続関手の性質の一つにより，
サイト $\mathcal{D}$ に被覆 $\{u(U_i) \to V\}_{i \in I}$ が存在する。
仮定により，各 $i$ に対し，各制限 $\mathcal{F}|_{U_{ij}}$ が
$\mathcal{P}_g$ をもつような $\mathcal{C}$ の被覆
$\{U_{ij} \to U_i\}_{j \in J_i}$ が存在する。《サイト》補題
\ref{sites-lemma-localize-special-cocontinuous} により，環付きトポスの可換図式
$$
\xymatrix{
(\Sh(\mathcal{C}/U_{ij}), \mathcal{O}_{U_{ij}}) \ar[r] \ar[d] &
(\Sh(\mathcal{C}), \mathcal{O}) \ar[d] \\
(\Sh(\mathcal{D}/u(U_{ij})), \mathcal{O}'_{u(U_{ij})}) \ar[r] &
(\Sh(\mathcal{D}), \mathcal{O}')
}
$$
があり，縦の矢印は同値である。したがって
$\mathcal{F}'|_{u(U_{ij})}$ も性質 $\mathcal{P}_g$ をもつ。
さらに $\{u(U_{ij}) \to V\}_{i \in I, j \in J_i}$ はサイト
$\mathcal{D}$ の被覆である。ゆえに $\mathcal{F}'$ は性質
$\mathcal{P}_l$ をもつ。

\medskip\noindent
$\mathcal{F}'$ が $\mathcal{P}_l$ をもつと仮定する。
$U$ を $\mathcal{C}$ の対象とする。仮定により，
$\mathcal{F}'|_{V_i}$ が性質 $\mathcal{P}_g$ をもつような被覆
$\{V_i \to u(U)\}_{i \in I}$ が存在する。$u$ は余連続なので，この被覆を
族 $\{u(U_j) \to u(U)\}_{j \in J}$ で細分できる。ここで
$\{U_j \to U\}_{j \in J}$ は $\mathcal{C}$ の被覆である。
この細分が $\alpha : J \to I$ と射
$u(U_j) \to V_{\alpha(j)}$ で与えられるとする。
制限は推移的，すなわち
$(\mathcal{F}'|_{V_{\alpha(j)}})|_{u(U_j)} = \mathcal{F}'|_{u(U_j)}$
である。したがって補題 \ref{sites-modules-lemma-global-pullback} により
$\mathcal{F}'|_{u(U_j)}$ は性質 $\mathcal{P}_g$ をもつ。ゆえに，縦の矢印が
同値である図式
$$
\xymatrix{
(\Sh(\mathcal{C}/U_j), \mathcal{O}_{U_j}) \ar[r] \ar[d] &
(\Sh(\mathcal{C}), \mathcal{O}) \ar[d] \\
(\Sh(\mathcal{D}/u(U_j)), \mathcal{O}'_{u(U_j)})
\ar[r] &
(\Sh(\mathcal{D}), \mathcal{O}')
}
$$
から $\mathcal{F}|_{U_j}$ も性質 $\mathcal{P}_g$ をもつことが分かる。
したがって望みどおり $\mathcal{F}$ は性質 $\mathcal{P}_l$ をもつ。

\medskip\noindent
最後に $\mathcal{P}_l = coherent$ の場合に補題を証明する\footnote{この議論は
やや扱いにくいので，読者には証明のこの部分を読み飛ばすことを勧める。}。
$\mathcal{F}$ が連接であると仮定する。このとき $\mathcal{F}$ は有限型であり，
証明の前半により $\mathcal{F}'$ も有限型である。
$V$ を $\mathcal{D}$ の対象とし，
$s_1, \ldots, s_n \in \mathcal{F}'(V)$ とする。写像
$\bigoplus_{j = 1, \ldots, n} \mathcal{O}_V \to \mathcal{F}'|_V$
の核 $\mathcal{K}'$ が $\mathcal{D}/V$ 上で有限型であることを示さなければ
ならない。すなわち，任意の $V'/V$ に対し，
$\mathcal{F}'|_{V'_i}$ が有限個の切断で生成されるような被覆
$\{V'_i \to V'\}$ が存在することを示す必要がある。
$V$ を $V'$ で置き換え（切断 $s_j$ も $V'$ に制限し），$V' = V$ の場合に帰着する。
$u$ は特殊余連続関手なので，サイト $\mathcal{D}$ に被覆
$\{u(U_i) \to V\}_{i \in I}$ が存在する。トポスの同型
$\Sh(\mathcal{C}/U_i) = \Sh(\mathcal{D}/u(U_i))$
を用いると，$\mathcal{K}'|_{u(U_i)}$ は写像
$\bigoplus_{j = 1, \ldots, n} \mathcal{O}_{U_i} \to \mathcal{F}|_{U_i}$
の核 $\mathcal{K}_i$ に対応する。$\mathcal{F}$ は連接なので
$\mathcal{K}_i$ は有限型である。したがって（再び証明の前半により）
$\mathcal{K}|_{u(U_i)}$ は有限型である。ゆえに，
$\mathcal{K}|_{V_{il}}$ が有限個の大域切断で生成されるような被覆
$\{V_{il} \to u(U_i)\}$ が存在する。
$\{V_{il} \to V\}$ は $\mathcal{D}$ の被覆だから，望みどおり
$\mathcal{K}$ は有限型である。

\medskip\noindent
$\mathcal{F}'$ が連接であると仮定する。このとき $\mathcal{F}'$ は有限型であり，
証明の前半により $\mathcal{F}$ も有限型である。
$U$ を $\mathcal{C}$ の対象とし，
$s_1, \ldots, s_n \in \mathcal{F}(U)$ とする。写像
$\bigoplus_{j = 1, \ldots, n} \mathcal{O}_U \to \mathcal{F}|_U$
の核 $\mathcal{K}$ が $\mathcal{C}/U$ 上で有限型であることを示さなければならない。
トポスの同型
$\Sh(\mathcal{C}/U) = \Sh(\mathcal{D}/u(U))$
を用いると，$\mathcal{K}|_U$ は写像
$\bigoplus_{j = 1, \ldots, n} \mathcal{O}_{u(U)} \to \mathcal{F}'|_{u(U)}$
の核 $\mathcal{K}'$ に対応する。$\mathcal{F}'$ は連接だから，
$\mathcal{K}'$ は有限型である。したがって再び証明の前半により，
$\mathcal{K}$ は有限型である。
\end{proof}

\noindent
したがって今後は，サイトを指定せずに，定義 \ref{sites-modules-definition-site-local} で
定めた $\mathcal{O}$-加群の性質に言及してよい。

\begin{lemma}
\label{sites-modules-lemma-local-final-object}
$(\Sh(\mathcal{C}), \mathcal{O})$ を環付きトポスとし，
$\mathcal{F}$ を $\mathcal{O}$-加群とする。
サイト $\mathcal{C}$ が終対象 $X$ をもつと仮定する。このとき：
\begin{enumerate}
\item 次は同値である：
\begin{enumerate}
\item $\mathcal{F}$ は局所自由である；
\item 各制限 $\mathcal{F}|_{\mathcal{C}/X_i}$ が局所自由
$\mathcal{O}_{X_i}$-加群となるような $\mathcal{C}$ の被覆
$\{X_i \to X\}$ が存在する；
\item 各制限 $\mathcal{F}|_{\mathcal{C}/X_i}$ が自由
$\mathcal{O}_{X_i}$-加群となるような $\mathcal{C}$ の被覆
$\{X_i \to X\}$ が存在する。
\end{enumerate}
\item 次は同値である：
\begin{enumerate}
\item $\mathcal{F}$ は有限局所自由である；
\item 各制限 $\mathcal{F}|_{\mathcal{C}/X_i}$ が有限局所自由
$\mathcal{O}_{X_i}$-加群となるような $\mathcal{C}$ の被覆
$\{X_i \to X\}$ が存在する；
\item 各制限 $\mathcal{F}|_{\mathcal{C}/X_i}$ が有限自由
$\mathcal{O}_{X_i}$-加群となるような $\mathcal{C}$ の被覆
$\{X_i \to X\}$ が存在する。
\end{enumerate}
\item 次は同値である：
\begin{enumerate}
\item $\mathcal{F}$ は局所的に切断で生成される；
\item 各制限 $\mathcal{F}|_{\mathcal{C}/X_i}$ が局所的に切断で生成される
$\mathcal{O}_{X_i}$-加群となるような $\mathcal{C}$ の被覆
$\{X_i \to X\}$ が存在する；
\item 各制限 $\mathcal{F}|_{\mathcal{C}/X_i}$ が大域切断で生成される
$\mathcal{O}_{X_i}$-加群となるような $\mathcal{C}$ の被覆
$\{X_i \to X\}$ が存在する。
\end{enumerate}
\item $r \geq 0$ とする。次は同値である：
\begin{enumerate}
\item $\mathcal{F}$ は局所的に $r$ 個の切断で生成される；
\item 各制限 $\mathcal{F}|_{\mathcal{C}/X_i}$ が局所的に $r$ 個の切断で
生成される $\mathcal{O}_{X_i}$-加群となるような $\mathcal{C}$ の被覆
$\{X_i \to X\}$ が存在する；
\item 各制限 $\mathcal{F}|_{\mathcal{C}/X_i}$ が $r$ 個の大域切断で
生成される $\mathcal{O}_{X_i}$-加群となるような $\mathcal{C}$ の被覆
$\{X_i \to X\}$ が存在する。
\end{enumerate}
\item 次は同値である：
\begin{enumerate}
\item $\mathcal{F}$ は有限型である；
\item 各制限 $\mathcal{F}|_{\mathcal{C}/X_i}$ が有限型の
$\mathcal{O}_{X_i}$-加群となるような $\mathcal{C}$ の被覆
$\{X_i \to X\}$ が存在する；
\item 各制限 $\mathcal{F}|_{\mathcal{C}/X_i}$ が有限個の大域切断で
生成される $\mathcal{O}_{X_i}$-加群となるような $\mathcal{C}$ の被覆
$\{X_i \to X\}$ が存在する。
\end{enumerate}
\item 次は同値である：
\begin{enumerate}
\item $\mathcal{F}$ は準連接である；
\item 各制限 $\mathcal{F}|_{\mathcal{C}/X_i}$ が準連接
$\mathcal{O}_{X_i}$-加群となるような $\mathcal{C}$ の被覆
$\{X_i \to X\}$ が存在する；
\item 各制限 $\mathcal{F}|_{\mathcal{C}/X_i}$ が大域表示をもつ
$\mathcal{O}_{X_i}$-加群となるような $\mathcal{C}$ の被覆
$\{X_i \to X\}$ が存在する。
\end{enumerate}
\item 次は同値である：
\begin{enumerate}
\item $\mathcal{F}$ は有限表示である；
\item 各制限 $\mathcal{F}|_{\mathcal{C}/X_i}$ が有限表示の
$\mathcal{O}_{X_i}$-加群となるような $\mathcal{C}$ の被覆
$\{X_i \to X\}$ が存在する；
\item 各制限 $\mathcal{F}|_{\mathcal{C}/X_i}$ が大域有限表示をもつ
$\mathcal{O}_{X_i}$-加群となるような $\mathcal{C}$ の被覆
$\{X_i \to X\}$ が存在する。
\end{enumerate}
\item 次は同値である：
\begin{enumerate}
\item $\mathcal{F}$ は連接である；
\item 各制限 $\mathcal{F}|_{\mathcal{C}/X_i}$ が連接
$\mathcal{O}_{X_i}$-加群となるような $\mathcal{C}$ の被覆
$\{X_i \to X\}$ が存在する。
\end{enumerate}
\end{enumerate}
\end{lemma}

\begin{proof}
各場合に (a) $\Rightarrow (b)$ が成り立つ。(1) -- (6) の各場合には，
サイトの公理 (2)（《サイト》定義 \ref{sites-definition-site} 参照）と
加群の局所的な型の定義により，条件 (b) は条件 (c) を含意する。
$\{X_i \to X\}$ を被覆とする。このとき $\mathcal{C}$ の任意の対象 $U$ に対し，
誘導された被覆 $\{X_i \times_X U \to U\}$ を得る。さらに，(c) における
$\mathcal{F}|_{\mathcal{C}/X_i}$ の大域的性質は，補題
\ref{sites-modules-lemma-global-pullback} により
$\mathcal{F}|_{\mathcal{C}/X_i \times_X U}$ の対応する大域的性質を含意する。
したがって定義により層は性質 (a) をもつ。(7) の場合の
(b) $\Rightarrow$ (a) の証明は省略する。
\end{proof}

\begin{lemma}
\label{sites-modules-lemma-local-pullback}
写像
$(f, f^\sharp) :
(\Sh(\mathcal{C}), \mathcal{O}_\mathcal{C})
\to
(\Sh(\mathcal{D}), \mathcal{O}_\mathcal{D})$
を環付きトポスの射とし，
$\mathcal{F}$ を $\mathcal{O}_\mathcal{D}$-加群とする。
\begin{enumerate}
\item $\mathcal{F}$ が局所自由ならば，$f^*\mathcal{F}$ も局所自由である。
\item $\mathcal{F}$ が有限局所自由ならば，
$f^*\mathcal{F}$ も有限局所自由である。
\item $\mathcal{F}$ が局所的に切断で生成されるならば，
$f^*\mathcal{F}$ も局所的に切断で生成される。
\item $\mathcal{F}$ が局所的に $r$ 個の切断で生成されるならば，
$f^*\mathcal{F}$ も局所的に $r$ 個の切断で生成される。
\item $\mathcal{F}$ が有限型ならば，$f^*\mathcal{F}$ も有限型である。
\item $\mathcal{F}$ が準連接ならば，$f^*\mathcal{F}$ も準連接である。
\item $\mathcal{F}$ が有限表示ならば，$f^*\mathcal{F}$ も有限表示である。
\end{enumerate}
\end{lemma}

\begin{proof}
第 \ref{sites-modules-section-intrinsic} 節の議論により，環付きサイトの射
$(f, f^\sharp) :
(\mathcal{C}, \mathcal{O}_\mathcal{C})
\to
(\mathcal{D}, \mathcal{O}_\mathcal{D})$
に対する引き戻しで保存されることだけを確認すればよい。ここで $f$ は，
すべての有限極限をもつサイト間の左完全連続関手
$u : \mathcal{D} \to \mathcal{C}$ によって与えられるとする。
$\mathcal{G}$ を，定義 \ref{sites-modules-definition-site-local} の性質 (1) -- (6) の
いずれかをもつ $\mathcal{O}_\mathcal{D}$-加群の層とする。
$\mathcal{D}$ は終対象 $Y$ をもち，$X = u(Y)$ は $\mathcal{C}$ の終対象である。
仮定により，$\mathcal{G}|_{\mathcal{D}/Y_i}$ が対応する大域的性質をもつような
被覆 $\{Y_i \to Y\}$ がある。$X_i = u(Y_i)$ と置けば，
$\{X_i \to X\}$ は $\mathcal{C}$ の被覆である。
環付きサイトの射の可換図式
$$
\xymatrix{
(\mathcal{C}/X_i, \mathcal{O}_\mathcal{C}|_{X_i}) \ar[r] \ar[d] &
(\mathcal{C}, \mathcal{O}_\mathcal{C}) \ar[d] \\
(\mathcal{D}/Y_i, \mathcal{O}_\mathcal{D}|_{Y_i}) \ar[r] &
(\mathcal{D}, \mathcal{O}_\mathcal{D})
}
$$
を《サイト》補題 \ref{sites-lemma-localize-morphism-strong} により得る。
したがって補題 \ref{sites-modules-lemma-global-pullback} により，
$f^*\mathcal{G}|_{X_i}$ は対応する大域的性質をもつ。
ゆえに補題 \ref{sites-modules-lemma-local-final-object} により，$\mathcal{G}$ は最初に仮定した
局所的性質をもつと結論できる。
\end{proof}







\section{加群の局所的な型に関する基本結果}
\label{sites-modules-section-basics}

\noindent
上で定めた概念に関する基本的な補題を述べる。

\begin{lemma}
\label{sites-modules-lemma-cokernel-finite-finite-presentation}
$(\mathcal{C},\mathcal{O})$ を環付きサイトとし，
$\mathcal{F}$ を有限表示の $\mathcal{O}$-加群とする。
$\varphi : \mathcal{G} \to \mathcal{F}$ を $\mathcal{O}$-加群の射とする。
$\mathcal{G}$ が有限型ならば，$\Coker(\varphi)$ は有限表示である。
\end{lemma}

\begin{proof}
省略する。ヒント：《加群》補題
\ref{modules-lemma-cokernel-finite-finite-presentation} を参照せよ。
\end{proof}

\begin{lemma}
\label{sites-modules-lemma-kernel-surjection-finite-onto-finite-presentation}
$(\mathcal{C}, \mathcal{O})$ を環付きサイトとする。
$\theta : \mathcal{G} \to \mathcal{F}$ を全射な $\mathcal{O}$-加群準同型とし，
$\mathcal{F}$ は有限表示，$\mathcal{G}$ は有限型であるとする。
このとき $\Ker(\theta)$ は有限型である。
\end{lemma}

\begin{proof}
省略する。ヒント：《加群》補題
\ref{modules-lemma-kernel-surjection-finite-free-onto-finite-presentation}
を参照せよ。
\end{proof}

\begin{lemma}
\label{sites-modules-lemma-chaotic-quasi-coherent}
$\mathcal{C}$ をカオス位相を備えたサイトとみなした圏とする
（《サイト》例 \ref{sites-example-indiscrete} 参照）。
$\mathcal{O}$ を $\mathcal{C}$ 上の環の層，
$\mathcal{F}$ を $\mathcal{O}$-加群の層とする。このとき，
$\mathcal{F}$ が準連接であることと，$\mathcal{C}$ のすべての射
$U \to V$ に対して標準写像
$$
\mathcal{F}(V) \otimes_{\mathcal{O}(V)} \mathcal{O}(U)
\longrightarrow
\mathcal{F}(U)
$$
が同型であることは同値である。
\end{lemma}

\begin{proof}
$\mathcal{F}$ が準連接であると仮定し，$U \to V$ を $\mathcal{C}$ の射とする。
$V$ の各被覆は同型によって与えられるので，定義
\ref{sites-modules-definition-site-local} から表示
$$
\bigoplus\nolimits_{j \in J} \mathcal{O}_V \longrightarrow
\bigoplus\nolimits_{i \in I} \mathcal{O}_V \longrightarrow
\mathcal{F}|_{\mathcal{C}/V} \longrightarrow 0
$$
が存在する。$\mathcal{C}$ の位相はカオス位相なので，
$\mathcal{C}$ の任意の対象上で切断を取る関手は完全である。
したがって，$\mathcal{F}(V)$ の $\mathcal{O}(V)$-加群としての表示
$$
\bigoplus\nolimits_{j \in J} \mathcal{O}(V) \longrightarrow
\bigoplus\nolimits_{i \in I} \mathcal{O}(V) \longrightarrow
\mathcal{F}(V) \longrightarrow 0
$$
を得る。$\mathcal{F}(U)$ についても同様である。
これにより，補題の主張に表示した写像が同型であることが直ちに分かる。

\medskip\noindent
補題の主張に表示した写像が $\mathcal{C}$ のすべての射 $U \to V$ に対して
同型であると仮定する。$V$ を固定し，$\mathcal{F}(V)$ の
$\mathcal{O}(V)$-加群としての表示
$$
\bigoplus\nolimits_{j \in J} \mathcal{O}(V) \longrightarrow
\bigoplus\nolimits_{i \in I} \mathcal{O}(V) \longrightarrow
\mathcal{F}(V) \longrightarrow 0
$$
を取る。このとき $\mathcal{F}$ に関する仮定は，対応する列
$$
\bigoplus\nolimits_{j \in J} \mathcal{O}_V \longrightarrow
\bigoplus\nolimits_{i \in I} \mathcal{O}_V \longrightarrow
\mathcal{F}|_{\mathcal{C}/V} \longrightarrow 0
$$
が完全であることとちょうど同値である。したがって $\mathcal{F}$ は準連接である。
\end{proof}

\begin{lemma}
\label{sites-modules-lemma-chaotic-flat-quasi-coherent}
$\mathcal{C}$ をカオス位相を備えたサイトとみなした圏とする
（《サイト》例 \ref{sites-example-indiscrete} 参照）。
$\mathcal{O}$ を $\mathcal{C}$ 上の環の層とする。
$\mathcal{C}$ のすべての射 $U \to V$ に対し，制限写像
$\mathcal{O}(V) \to \mathcal{O}(U)$ が平坦な環準同型であると仮定する。
このとき準連接 $\mathcal{O}$-加群の圏は
$\textit{Mod}(\mathcal{O})$ の弱 Serre 部分圏である。
\end{lemma}

\begin{proof}
弱 Serre 部分圏の定義を確認する（《ホモロジー》定義
\ref{homology-definition-serre-subcategory} 参照）。そのために補題
\ref{sites-modules-lemma-chaotic-quasi-coherent} で与えた準連接加群の特徴づけを用いる。
$\textit{Mod}(\mathcal{O})$ の完全列
$$
\mathcal{F}_0 \to \mathcal{F}_1 \to
\mathcal{F}_2 \to \mathcal{F}_3 \to
\mathcal{F}_4
$$
を考え，$\mathcal{F}_0$，$\mathcal{F}_1$，$\mathcal{F}_3$，
$\mathcal{F}_4$ は準連接であるとする。
$U \to V$ を $\mathcal{C}$ の射とし，可換図式
$$
\resizebox{\textwidth}{!}{\ensuremath{\xymatrix{
\mathcal{F}_0(V) \otimes_{\mathcal{O}(V)} \mathcal{O}(U) \ar[r] \ar[d] &
\mathcal{F}_1(V) \otimes_{\mathcal{O}(V)} \mathcal{O}(U) \ar[r] \ar[d] &
\mathcal{F}_2(V) \otimes_{\mathcal{O}(V)} \mathcal{O}(U) \ar[r] \ar[d] &
\mathcal{F}_3(V) \otimes_{\mathcal{O}(V)} \mathcal{O}(U) \ar[r] \ar[d] &
\mathcal{F}_4(V) \otimes_{\mathcal{O}(V)} \mathcal{O}(U) \ar[d] \\
\mathcal{F}_0(U) \ar[r] &
\mathcal{F}_1(U) \ar[r] &
\mathcal{F}_2(U) \ar[r] &
\mathcal{F}_3(U) \ar[r] &
\mathcal{F}_4(U)
}}}
$$
を考える。仮定により，添字 $0$，$1$，$3$，$4$ の縦の矢印は同型である。
$\mathcal{C}$ の位相はカオス位相なので，$\mathcal{C}$ の対象上で切断を取る関手は
完全であり，したがって下の行は完全である。
$\mathcal{O}(V) \to \mathcal{O}(U)$ は平坦なので上の行も完全である。
ゆえに五項補題（《ホモロジー》補題
\ref{homology-lemma-five-lemma}）により中央の矢印は同型である。
\end{proof}






\section{環付きトポスの閉埋め込み}
\label{sites-modules-section-closed-immersion}

\noindent
環付きトポスの射
$i : (\Sh(\mathcal{C}), \mathcal{O}) \to (\Sh(\mathcal{D}), \mathcal{O}')$
をいつ閉埋め込みと呼ぶべきだろうか。《加群》第
\ref{modules-section-closed-immersion} 節の議論との類推から，少なくとも次を
仮定するのが自然に思われる：
\begin{enumerate}
\item 関手 $i$ はトポスの閉埋め込みである
（《サイト》定義 \ref{sites-definition-immersion-topoi}）。
\item 付随する写像 $\mathcal{O}' \to i_*\mathcal{O}$ は全射である。
\end{enumerate}
これらの条件だけでもいくつかの好ましい結果が従うので，本節でそれを論じる。
しかし環付きトポスの閉埋め込みについては，採用しうる定義がいくつもあるため，
実際の定義をここで設けない方が慎重であろう。

\begin{lemma}
\label{sites-modules-lemma-i-star-equivalence}
$i : (\Sh(\mathcal{C}), \mathcal{O}) \to (\Sh(\mathcal{D}), \mathcal{O}')$
を環付きトポスの射とする。$i$ はトポスの閉埋め込みであり，
$i^\sharp : \mathcal{O}' \to i_*\mathcal{O}$ は全射であると仮定する。
$\mathcal{I} \subset \mathcal{O}'$ を $i^\sharp$ の核とする。関手
$$
i_* :
\textit{Mod}(\mathcal{O})
\longrightarrow
\textit{Mod}(\mathcal{O}')
$$
は完全かつ充満忠実であり，その本質像は
$\mathcal{I}\mathcal{G} = 0$ を満たす $\mathcal{O}'$-加群
$\mathcal{G}$ 全体である。
\end{lemma}

\begin{proof}
補題 \ref{sites-modules-lemma-exactness} と《サイト》補題
\ref{sites-lemma-closed-immersion} により $i_*$ は完全である。
$i_*$ が集合の層上で充満忠実であることと $i^\sharp$ が全射であることから，
$i_*$ は関手 $\textit{Mod}(\mathcal{O}) \to \textit{Mod}(\mathcal{O}')$
として充満忠実であることが従う。実際，
$\alpha : i_*\mathcal{F}_1 \to i_*\mathcal{F}_2$ を
$\mathcal{O}'$-加群準同型とする。$i_*$ の充満忠実性により，集合の層の写像
$\beta : \mathcal{F}_1 \to \mathcal{F}_2$ を得る。
$\beta$ が加群準同型であることを示すには，図式
$$
\xymatrix{
\mathcal{O} \times \mathcal{F}_1 \ar[r] \ar[d] &
\mathcal{F}_1 \ar[d] \\
\mathcal{O} \times \mathcal{F}_2 \ar[r] &
\mathcal{F}_2
}
$$
が可換であることを示せばよい。$i_*$ を適用した後で可換性を示せば十分である。
図式
$$
\xymatrix{
\mathcal{O}' \times i_*\mathcal{F}_1 \ar[r] \ar[d] &
i_*\mathcal{O} \times i_*\mathcal{F}_1 \ar[r] \ar[d] &
i_*\mathcal{F}_1 \ar[d] \\
\mathcal{O}' \times i_*\mathcal{F}_2 \ar[r] &
i_*\mathcal{O} \times i_*\mathcal{F}_2 \ar[r] &
i_*\mathcal{F}_2
}
$$
を考える。外側の長方形は可換である。$i^\sharp$ は全射なので，結論が従う。

\medskip\noindent
証明を終えるには，$i_*$ の本質像に関する主張を示す必要がある。
任意の $\mathcal{O}$-加群 $\mathcal{F}$ に対し，
$i_*\mathcal{F}$ が $\mathcal{I}$ で零化されることは明らかである。
逆に，$\mathcal{I}\mathcal{G} = 0$ を満たす $\mathcal{O}'$-加群
$\mathcal{G}$ を取る。閉部分トポスの定義により，$\mathcal{D}$ の終対象の
部分層 $\mathcal{U}$ であって，集合の層上の $i_*$ の本質像が，写像
$\mathcal{H} \times \mathcal{U} \to \mathcal{U}$ が同型となる集合の層
$\mathcal{H}$ 全体であるようなものが存在する。特に
$i_*\mathcal{O} \times \mathcal{U} = \mathcal{U}$ である。したがって
$\mathcal{I} \times \mathcal{U} = \mathcal{O} \times \mathcal{U}$。
これより，加群 $\mathcal{G}$ は
$\mathcal{G} \times \mathcal{U} = \{0\} \times \mathcal{U} = \mathcal{U}$
を満たす（零加群は集合の層の終対象と同型だからである）。
ゆえに $i_*\mathcal{F} = \mathcal{G}$ を満たす $\mathcal{C}$ 上の集合の層
$\mathcal{F}$ が存在する。$i_*$ は集合の層上で充満忠実なので，加法
$\mathcal{F} \times \mathcal{F} \to \mathcal{F}$ と乗法
$\mathcal{O} \times \mathcal{F} \to \mathcal{F}$ を定めるには，加法
$$
\mathcal{G} \times \mathcal{G} \longrightarrow \mathcal{G}
$$
（$\mathcal{G}$ は $\mathcal{O}'$-加群なので与えられている）と乗法
$$
i_*\mathcal{O} \times \mathcal{G} \to \mathcal{G}
$$
を用いれば十分である。後者は，仮定により $\mathcal{I}$ を零化する
$\mathcal{O}'$ の乗法作用と $i_*\mathcal{O} = \mathcal{O}'/\mathcal{I}$
から与えられる。構成により，$\mathcal{G}$ はこうして得た
$\mathcal{O}$-加群 $\mathcal{F}$ の順像と同型である。
\end{proof}












\section{テンソル積}
\label{sites-modules-section-tensor-product}

\noindent
第 \ref{sites-modules-section-presheaves-modules} 節と第
\ref{sites-modules-section-sheafification-presheaves-modules} 節では，テンソル積の構成により
環の変更関手を定義した。この構成が意味をもつことを確かめるためにも，
$\mathcal{O}$-加群の前層のテンソル積を定義する。
正確には，$\mathcal{C}$ を圏，$\mathcal{O}$ を環の前層，
$\mathcal{F}$，$\mathcal{G}$ を $\mathcal{O}$-加群の前層とする。この場合，
$\mathcal{F} \otimes_{p, \mathcal{O}} \mathcal{G}$ を前層
$$
U
\longmapsto
(\mathcal{F} \otimes_{p, \mathcal{O}} \mathcal{G})(U)
=
\mathcal{F}(U) \otimes_{\mathcal{O}(U)} \mathcal{G}(U)
$$
として定義する。$\mathcal{C}$ がサイト，$\mathcal{O}$ が環の層，
$\mathcal{F}$，$\mathcal{G}$ が $\mathcal{O}$-加群の層ならば，
$$
\mathcal{F} \otimes_\mathcal{O} \mathcal{G}
=
(\mathcal{F} \otimes_{p, \mathcal{O}} \mathcal{G})^\#
$$
を前層 $\mathcal{F} \otimes_{p, \mathcal{O}} \mathcal{G}$ に付随する
$\mathcal{O}$-加群の層として定義する。

\medskip\noindent
以下，断りなく用いるいくつかの公式を挙げる：
$$
(\mathcal{F}
\otimes_{p, \mathcal{O}} \mathcal{G})
\otimes_{p, \mathcal{O}} \mathcal{H}
=
\mathcal{F}
\otimes_{p, \mathcal{O}} (\mathcal{G}
\otimes_{p, \mathcal{O}} \mathcal{H}),
$$
層についても同様である。$\mathcal{O}_1 \to \mathcal{O}_2$ が環の前層の写像ならば，
$$
(\mathcal{F} \otimes_{p, \mathcal{O}_1} \mathcal{G})
\otimes_{p, \mathcal{O}_1} \mathcal{O}_2 =
(\mathcal{F} \otimes_{p, \mathcal{O}_1} \mathcal{O}_2)
\otimes_{p, \mathcal{O}_2}
(\mathcal{G} \otimes_{p, \mathcal{O}_1} \mathcal{O}_2),
$$
であり，層についても同様である。これらは代数における対応する公式と層化から従う。

\begin{lemma}
\label{sites-modules-lemma-sheafification-tensor}
$\mathcal{C}$ をサイト，$\mathcal{O}$ を環の前層とし，
$\mathcal{F}$，$\mathcal{G}$ を $\mathcal{O}$-加群の前層とする。このとき
$\mathcal{F}^\# \otimes_{\mathcal{O}^\#} \mathcal{G}^\#$
は
$(\mathcal{F} \otimes_{p, \mathcal{O}} \mathcal{G})^\#$
に等しい。
\end{lemma}

\begin{proof}
省略する。ヒント：以下の双線形写像によるテンソル積の特徴づけと，
層化の普遍性を用いよ。
\end{proof}

\noindent
$\mathcal{C}$ をサイト，$\mathcal{O}$ を環の層とし，
$\mathcal{F}$，$\mathcal{G}$，$\mathcal{H}$ を $\mathcal{O}$-加群の層とする。
この場合，
$$
\text{Bilin}_\mathcal{O}(\mathcal{F} \times \mathcal{G}, \mathcal{H})
=
\{\varphi \in
\Mor_{\Sh(\mathcal{C})}(
\mathcal{F} \times \mathcal{G}, \mathcal{H}) \mid
\varphi \text{ は }\mathcal{O}\text{-双線形}\}.
$$
この定義により
$$
\Hom_\mathcal{O}
(\mathcal{F} \otimes_\mathcal{O} \mathcal{G}, \mathcal{H})
=
\text{Bilin}_\mathcal{O}(\mathcal{F} \times \mathcal{G}, \mathcal{H}).
$$
が成り立つ。言い換えると，$\mathcal{F} \otimes_\mathcal{O} \mathcal{G}$ は，
$\mathcal{H}$ に双線形写像 $\mathcal{F} \times \mathcal{G} \to \mathcal{H}$
の集合を対応させる関手を表現する。特に，環付きトポス上の二つの加群に対して
双線形写像の概念は意味をもつので，加群の層のテンソル積はトポスに内在的である
（第 \ref{sites-modules-section-intrinsic} 節の議論と比較せよ）。実際，次が成り立つ。

\begin{lemma}
\label{sites-modules-lemma-tensor-product-pullback}
$f : (\Sh(\mathcal{C}), \mathcal{O}_\mathcal{C})
\to (\Sh(\mathcal{D}), \mathcal{O}_\mathcal{D})$ を環付きトポスの射とし，
$\mathcal{F}$，$\mathcal{G}$ を $\mathcal{O}_\mathcal{D}$-加群とする。このとき
$f^*(\mathcal{F} \otimes_{\mathcal{O}_\mathcal{D}} \mathcal{G})
= f^*\mathcal{F} \otimes_{\mathcal{O}_\mathcal{C}} f^*\mathcal{G}$
が $\mathcal{F}$，$\mathcal{G}$ に関して関手的に成り立つ。
\end{lemma}

\begin{proof}
$\mathcal{O}_\mathcal{C}$-加群の層 $\mathcal{H}$ に対して
\begin{align*}
\Hom_{\mathcal{O}_\mathcal{C}}(
f^*(\mathcal{F} \otimes_\mathcal{O} \mathcal{G}), \mathcal{H})
& =
\Hom_{\mathcal{O}_\mathcal{D}}(
\mathcal{F} \otimes_\mathcal{O} \mathcal{G}, f_*\mathcal{H}) \\
& =
\text{Bilin}_{\mathcal{O}_\mathcal{D}}(
\mathcal{F} \times \mathcal{G}, f_*\mathcal{H}) \\
& =
\text{Bilin}_{f^{-1}\mathcal{O}_\mathcal{D}}(
f^{-1}\mathcal{F} \times f^{-1}\mathcal{G}, \mathcal{H}) \\
& =
\Hom_{f^{-1}\mathcal{O}_\mathcal{D}}(
f^{-1}\mathcal{F} \otimes_{f^{-1}\mathcal{O}_\mathcal{D}} f^{-1}\mathcal{G},
\mathcal{H}) \\
& =
\Hom_{\mathcal{O}_\mathcal{C}}(
f^*\mathcal{F} \otimes_{f^*\mathcal{O}_\mathcal{D}} f^*\mathcal{G},
\mathcal{H})
\end{align*}
が成り立つ。この等式列で興味深い「$=$」は第三の等式である。
これは定義と，前節までに論じた $f_*$ と $f^{-1}$ の随伴性から直ちに従う。
\end{proof}

\begin{lemma}
\label{sites-modules-lemma-tensor-product-permanence}
$(\mathcal{C}, \mathcal{O})$ を環付きサイトとし，
$\mathcal{F}$，$\mathcal{G}$ を $\mathcal{O}$-加群の層とする。
\begin{enumerate}
\item $\mathcal{F}$，$\mathcal{G}$ が局所自由ならば，
$\mathcal{F} \otimes_\mathcal{O} \mathcal{G}$ も局所自由である。
\item $\mathcal{F}$，$\mathcal{G}$ が有限局所自由ならば，
$\mathcal{F} \otimes_\mathcal{O} \mathcal{G}$ も有限局所自由である。
\item $\mathcal{F}$，$\mathcal{G}$ が局所的に切断で生成されるならば，
$\mathcal{F} \otimes_\mathcal{O} \mathcal{G}$ も局所的に切断で生成される。
\item $\mathcal{F}$，$\mathcal{G}$ が有限型ならば，
$\mathcal{F} \otimes_\mathcal{O} \mathcal{G}$ も有限型である。
\item $\mathcal{F}$，$\mathcal{G}$ が準連接ならば，
$\mathcal{F} \otimes_\mathcal{O} \mathcal{G}$ も準連接である。
\item $\mathcal{F}$，$\mathcal{G}$ が有限表示ならば，
$\mathcal{F} \otimes_\mathcal{O} \mathcal{G}$ も有限表示である。
\item $\mathcal{F}$ が有限表示で $\mathcal{G}$ が連接ならば，
$\mathcal{F} \otimes_\mathcal{O} \mathcal{G}$ は連接である。
\item $\mathcal{F}$，$\mathcal{G}$ が連接ならば，
$\mathcal{F} \otimes_\mathcal{O} \mathcal{G}$ も連接である。
\end{enumerate}
\end{lemma}

\begin{proof}
省略する。ヒント：《加群の層》補題
\ref{modules-lemma-tensor-product-permanence} と比較せよ。
\end{proof}



\section{内部 Hom}
\label{sites-modules-section-internal-hom}

\noindent
$\mathcal{C}$ を圏とし，$\mathcal{O}$ を環の前層とする。
$\mathcal{F}$，$\mathcal{G}$ を $\mathcal{O}$-加群の前層とする。対応
$$
U \longmapsto \Hom_{\mathcal{O}_U}(\mathcal{F}|_U, \mathcal{G}|_U).
$$
$\mathcal{C}$ の射 $\varphi : V \to U$ に対し，制限写像
$$
\Hom_{\mathcal{O}_U}(\mathcal{F}|_U, \mathcal{G}|_U)
\longrightarrow
\Hom_{\mathcal{O}_V}(\mathcal{F}|_V, \mathcal{G}|_V)
$$
を，再局所化射 $j : \mathcal{C}/V \to \mathcal{C}/U$ によって制限することで
定める（《サイト》の補題 \ref{sites-lemma-relocalize} を参照）。したがって，
これは前層 $\SheafHom_\mathcal{O}(\mathcal{F}, \mathcal{G})$ を定める。
さらに，元
$\varphi \in \Hom_{\mathcal{O}|_U}(\mathcal{F}|_U, \mathcal{G}|_U)$
と切断 $f \in \mathcal{O}(U)$ が与えられたとき，
$f\varphi \in \Hom_{\mathcal{O}|_U}(\mathcal{F}|_U, \mathcal{G}|_U)$
を，$\mathcal{F}|_U$ 上の $f$ 倍写像を前合成することによっても，
$\mathcal{G}|_U$ 上の $f$ 倍写像を後合成することによっても定められる
（両者は同じ結果を与える）。したがって，実際には
$\mathcal{O}$-加群の前層が得られる。標準的な「評価」射
$$
\mathcal{F}
\otimes_{p, \mathcal{O}}
\SheafHom_\mathcal{O}(\mathcal{F}, \mathcal{G})
\longrightarrow
\mathcal{G}.
$$
がある。

\begin{lemma}
\label{sites-modules-lemma-internal-hom}
$\mathcal{C}$ がサイト，$\mathcal{O}$ が環の層，
$\mathcal{F}$ が $\mathcal{O}$-加群の前層，
$\mathcal{G}$ が $\mathcal{O}$-加群の層であるならば，
$\SheafHom_\mathcal{O}(\mathcal{F}, \mathcal{G})$
は $\mathcal{O}$-加群の層である。
\end{lemma}

\begin{proof}
省略する。ヒント：まず
$\SheafHom_\mathcal{O}(\mathcal{F}, \mathcal{G})
= \SheafHom_\mathcal{O}(\mathcal{F}^\#, \mathcal{G})$ であることに注意すれば，
問題は $\mathcal{F}$ と $\mathcal{G}$ がともに層である場合に帰着される。
集合の層に対する結果は《サイト》の補題 \ref{sites-lemma-glue-maps} である。
\end{proof}

\begin{lemma}
\label{sites-modules-lemma-internal-hom-restriction}
$(\mathcal{C}, \mathcal{O})$ を環付きサイトとし，
$\mathcal{F}, \mathcal{G}$ を $\mathcal{O}$-加群の層とする。このとき，
$\SheafHom_\mathcal{O}(\mathcal{F}, \mathcal{G})$ の構成は，
$U \in \Ob(\mathcal{C})$ への制限と可換である。
\end{lemma}

\begin{proof}
定義から直ちに従う。
\end{proof}

\begin{remark}
\label{sites-modules-remark-pullback-internal-hom}
$f : (\mathcal{C}, \mathcal{O}_\mathcal{C}) \to
(\mathcal{D}, \mathcal{O}_\mathcal{D})$ を環付きサイトの射とする。
$\mathcal{F}, \mathcal{G}$ を $\mathcal{O}_\mathcal{D}$-加群の層とする。
標準写像
$$
f^*\SheafHom_{\mathcal{O}_\mathcal{D}}(\mathcal{F}, \mathcal{G})
\longrightarrow
\SheafHom_{\mathcal{O}_\mathcal{C}}(f^*\mathcal{F}, f^*\mathcal{G})
$$
がある。すなわち，この写像は次の写像の随伴である：
$$
\SheafHom_{\mathcal{O}_\mathcal{D}}(\mathcal{F}, \mathcal{G})
\longrightarrow
f_*\SheafHom_{\mathcal{O}_\mathcal{C}}(f^*\mathcal{F}, f^*\mathcal{G})
$$
この写像を次のように定める。$f$ が連続関手
$u : \mathcal{D} \to \mathcal{C}$ によって与えられるとする。
$V \in \Ob(\mathcal{D})$ 上の切断に対して，写像
\begin{align*}
\Gamma(V, \SheafHom_{\mathcal{O}_\mathcal{D}}(\mathcal{F}, \mathcal{G}))
& =
\Hom_{\mathcal{O}_V}(\mathcal{F}|_V, \mathcal{G}|_V) \\
& \longrightarrow
\Hom_{\mathcal{O}_{u(V)}}(f^*\mathcal{F}|_{u(V)}, f^*\mathcal{G}|_{u(V)}) \\
& =
\Gamma(u(V), \SheafHom_{\mathcal{O}_\mathcal{C}}(f^*\mathcal{F},
f^*\mathcal{G})) \\
& =
\Gamma(V, f_*\SheafHom_{\mathcal{O}_\mathcal{C}}(f^*\mathcal{F},
f^*\mathcal{G}))
\end{align*}
を用いる。ここで矢印には，$f$ が誘導する射
$(\mathcal{C}/u(V), \mathcal{O}_{u(V)}) \to
(\mathcal{D}/V, \mathcal{O}_V)$ による引き戻しを用いる。
\end{remark}

\noindent
補題 \ref{sites-modules-lemma-internal-hom} の状況では，「評価」射は
加群の層のテンソル積を経由して
$$
\mathcal{F}
\otimes_\mathcal{O}
\SheafHom_\mathcal{O}(\mathcal{F}, \mathcal{G})
\longrightarrow
\mathcal{G}.
$$
と分解する。

\begin{lemma}
\label{sites-modules-lemma-internal-hom-commute-limits}
内部 Hom と（余）極限。
$\mathcal{C}$ を圏とし，$\mathcal{O}$ を環の前層とする。
\begin{enumerate}
\item 任意の $\mathcal{O}$-加群の前層 $\mathcal{F}$ に対し，関手
$$
\textit{PMod}(\mathcal{O}) \longrightarrow \textit{PMod}(\mathcal{O})
, \quad
\mathcal{G} \longmapsto \SheafHom_\mathcal{O}(\mathcal{F}, \mathcal{G})
$$
は任意の極限と可換である。
\item 任意の $\mathcal{O}$-加群の前層 $\mathcal{G}$ に対し，関手
$$
\textit{PMod}(\mathcal{O}) \longrightarrow \textit{PMod}(\mathcal{O})^{opp}
, \quad
\mathcal{F} \longmapsto \SheafHom_\mathcal{O}(\mathcal{F}, \mathcal{G})
$$
は任意の余極限と可換である。式で書けば
$$
\SheafHom_\mathcal{O}(\colim_i \mathcal{F}_i, \mathcal{G})
=
\lim_i \SheafHom_\mathcal{O}(\mathcal{F}_i, \mathcal{G}).
$$
\end{enumerate}
$\mathcal{C}$ がサイトであり，$\mathcal{O}$ が環の層であるとする。
\begin{enumerate}
\item[(3)] 任意の $\mathcal{O}$-加群の層 $\mathcal{F}$ に対し，関手
$$
\textit{Mod}(\mathcal{O}) \longrightarrow \textit{Mod}(\mathcal{O})
, \quad
\mathcal{G} \longmapsto \SheafHom_\mathcal{O}(\mathcal{F}, \mathcal{G})
$$
は任意の極限と可換である。
\item[(4)] 任意の $\mathcal{O}$-加群の層 $\mathcal{G}$ に対し，関手
$$
\textit{Mod}(\mathcal{O}) \longrightarrow \textit{Mod}(\mathcal{O})^{opp}
, \quad
\mathcal{F} \longmapsto \SheafHom_\mathcal{O}(\mathcal{F}, \mathcal{G})
$$
は任意の余極限と可換である。式で書けば
$$
\SheafHom_\mathcal{O}(\colim_i \mathcal{F}_i, \mathcal{G})
=
\lim_i \SheafHom_\mathcal{O}(\mathcal{F}_i, \mathcal{G}).
$$
\end{enumerate}
\end{lemma}

\begin{proof}
$\mathcal{I} \to \textit{PMod}(\mathcal{O})$，$i \mapsto \mathcal{G}_i$
を図式とし，$U$ を圏 $\mathcal{C}$ の対象とする。
$j_U^*$ は左随伴でも右随伴でもあるから，
$\lim_i j_U^*\mathcal{G}_i = j_U^* \lim_i \mathcal{G}_i$ である。したがって
\begin{align*}
\SheafHom_\mathcal{O}(\mathcal{F}, \lim_i \mathcal{G}_i)(U)
& =
\Hom_{\mathcal{O}_U}(\mathcal{F}|_U, \lim_i \mathcal{G}_i|_U) \\
& =
\lim_i \Hom_{\mathcal{O}_U}(\mathcal{F}|_U, \mathcal{G}_i|_U) \\
& = \lim_i \SheafHom_\mathcal{O}(\mathcal{F}, \mathcal{G}_i)(U)
\end{align*}
を得る。二つ目の等号は極限の定義による。これで (1) が示された。
(2) もまったく同様に示される。層の図式の極限は前層の圏における極限と
同じであるから，(3) は (1) から従う。最後に，式
$$
\Mor_{\textit{Mod}(\mathcal{O})}(
\colim_i \mathcal{F}_i, \mathcal{G})
=
\Mor_{\textit{PMod}(\mathcal{O})}(
\colim^{PSh}_i \mathcal{F}_i, \mathcal{G})
$$
が成り立つ。実際，余極限 $\colim_i \mathcal{F}_i$ は，
$\textit{PMod}(\mathcal{O})$ における余極限
$\colim^{PSh}_i \mathcal{F}_i$ の層化である。したがって，(4) は (2) から
従う（ここでも上の極限に関する注意を用いる）。
\end{proof}

\begin{lemma}
\label{sites-modules-lemma-internal-hom-exact}
$(\mathcal{C}, \mathcal{O})$ を環付きサイトとし，
$\mathcal{F}$，$\mathcal{G}$ を $\mathcal{O}$-加群とする。
\begin{enumerate}
\item $\mathcal{F}_2 \to \mathcal{F}_1 \to \mathcal{F} \to 0$ が
$\mathcal{O}$-加群の完全列ならば，
$$
0 \to
\SheafHom_\mathcal{O}(\mathcal{F}, \mathcal{G}) \to
\SheafHom_\mathcal{O}(\mathcal{F}_1, \mathcal{G}) \to
\SheafHom_\mathcal{O}(\mathcal{F}_2, \mathcal{G})
$$
は完全である。
\item $0 \to \mathcal{G} \to \mathcal{G}_1 \to \mathcal{G}_2$ が
$\mathcal{O}$-加群の完全列ならば，
$$
0 \to
\SheafHom_\mathcal{O}(\mathcal{F}, \mathcal{G}) \to
\SheafHom_\mathcal{O}(\mathcal{F}, \mathcal{G}_1) \to
\SheafHom_\mathcal{O}(\mathcal{F}, \mathcal{G}_2)
$$
は完全である。
\end{enumerate}
\end{lemma}

\begin{proof}
補題 \ref{sites-modules-lemma-internal-hom-commute-limits} と
《ホモロジー》の補題 \ref{homology-lemma-exact-functor} から従う。
\end{proof}

\begin{lemma}
\label{sites-modules-lemma-internal-hom-adjoint-tensor}
$\mathcal{C}$ を圏とし，$\mathcal{O}$ を環の前層とする。
\begin{enumerate}
\item $\mathcal{F}$，$\mathcal{G}$，$\mathcal{H}$ を
$\mathcal{O}$-加群の前層とする。標準同型
$$
\SheafHom_\mathcal{O}
(\mathcal{F} \otimes_{p, \mathcal{O}} \mathcal{G}, \mathcal{H})
\longrightarrow
\SheafHom_\mathcal{O}
(\mathcal{F}, \SheafHom_\mathcal{O}(\mathcal{G}, \mathcal{H}))
$$
があり，三つの変数すべてについて関手的である（三箇所とも層 Hom である）。
特に，
$$
\Mor_{\textit{PMod}(\mathcal{O})}(
\mathcal{F} \otimes_{p, \mathcal{O}} \mathcal{G}, \mathcal{H})
=
\Mor_{\textit{PMod}(\mathcal{O})}(
\mathcal{F}, \SheafHom_\mathcal{O}(\mathcal{G}, \mathcal{H}))
$$
\item
$\mathcal{C}$ がサイト，$\mathcal{O}$ が環の層であり，
$\mathcal{F}$，$\mathcal{G}$，$\mathcal{H}$ が
$\mathcal{O}$-加群の層であるとする。標準同型
$$
\SheafHom_\mathcal{O}
(\mathcal{F} \otimes_\mathcal{O} \mathcal{G}, \mathcal{H})
\longrightarrow
\SheafHom_\mathcal{O}
(\mathcal{F}, \SheafHom_\mathcal{O}(\mathcal{G}, \mathcal{H}))
$$
があり，三つの変数すべてについて関手的である（三箇所とも層 Hom である）。
特に，
$$
\Mor_{\textit{Mod}(\mathcal{O})}(
\mathcal{F} \otimes_\mathcal{O} \mathcal{G}, \mathcal{H})
=
\Mor_{\textit{Mod}(\mathcal{O})}(
\mathcal{F}, \SheafHom_\mathcal{O}(\mathcal{G}, \mathcal{H}))
$$
\end{enumerate}
\end{lemma}

\begin{proof}
これは《代数》の補題 \ref{algebra-lemma-hom-from-tensor-product} の類似である。
証明も同じなので省略する。
\end{proof}

\begin{lemma}
\label{sites-modules-lemma-tensor-commute-colimits}
テンソル積と余極限。
$\mathcal{C}$ を圏とし，$\mathcal{O}$ を環の前層とする。
\begin{enumerate}
\item 任意の $\mathcal{O}$-加群の前層 $\mathcal{F}$ に対し，関手
$$
\textit{PMod}(\mathcal{O}) \longrightarrow \textit{PMod}(\mathcal{O})
, \quad
\mathcal{G} \longmapsto \mathcal{F} \otimes_{p, \mathcal{O}} \mathcal{G}
$$
は任意の余極限と可換である。
\item
$\mathcal{C}$ がサイト，$\mathcal{O}$ が環の層であるとする。
任意の $\mathcal{O}$-加群の層 $\mathcal{F}$ に対し，関手
$$
\textit{Mod}(\mathcal{O}) \longrightarrow \textit{Mod}(\mathcal{O})
, \quad
\mathcal{G} \longmapsto \mathcal{F} \otimes_\mathcal{O} \mathcal{G}
$$
は任意の余極限と可換である。
\end{enumerate}
\end{lemma}

\begin{proof}
補題 \ref{sites-modules-lemma-internal-hom-adjoint-tensor} により，テンソル積が内部 Hom に
随伴することから従う。《圏》の補題 \ref{categories-lemma-adjoint-exact} を参照。
\end{proof}

\begin{lemma}
\label{sites-modules-lemma-adjoint-hom-restrict}
$\mathcal{C}$ を圏，それぞれサイトとする。
$\mathcal{O} \to \mathcal{O}'$ を環の前層，それぞれ環の層の写像とする。このとき
$$
\Hom_\mathcal{O}(\mathcal{G}, \mathcal{F}) =
\Hom_{\mathcal{O}'}(\mathcal{G},
\SheafHom_\mathcal{O}(\mathcal{O}', \mathcal{F}))
$$
は，任意の $\mathcal{O}'$-加群 $\mathcal{G}$ と
$\mathcal{O}$-加群 $\mathcal{F}$ に対して成り立つ。
\end{lemma}

\begin{proof}
これは《代数》の補題 \ref{algebra-lemma-adjoint-hom-restrict} の類似である。
証明も同じなので省略する。
\end{proof}

\begin{lemma}
\label{sites-modules-lemma-j-shriek-and-tensor}
$(\mathcal{C}, \mathcal{O})$ を環付きサイトとし，
$U \in \Ob(\mathcal{C})$ とする。
$\textit{Mod}(\mathcal{O}_U)$ の $\mathcal{G}$ と
$\textit{Mod}(\mathcal{O})$ の $\mathcal{F}$ に対して
$j_{U!}\mathcal{G} \otimes_\mathcal{O} \mathcal{F} =
j_{U!}(\mathcal{G} \otimes_{\mathcal{O}_U} \mathcal{F}|_U)$ である。
\end{lemma}

\begin{proof}
$\mathcal{H}$ を $\textit{Mod}(\mathcal{O})$ の対象とする。このとき
\begin{align*}
\Hom_\mathcal{O}(j_{U!}(\mathcal{G} \otimes_{\mathcal{O}_U} \mathcal{F}|_U),
\mathcal{H})
& =
\Hom_{\mathcal{O}_U}(\mathcal{G} \otimes_{\mathcal{O}_U} \mathcal{F}|_U,
\mathcal{H}|_U) \\
& =
\Hom_{\mathcal{O}_U}(\mathcal{G},
\SheafHom_{\mathcal{O}_U}(\mathcal{F}|_U, \mathcal{H}|_U)) \\
& =
\Hom_{\mathcal{O}_U}(\mathcal{G},
\SheafHom_\mathcal{O}(\mathcal{F}, \mathcal{H})|_U) \\
& =
\Hom_\mathcal{O}(j_{U!}\mathcal{G},
\SheafHom_\mathcal{O}(\mathcal{F}, \mathcal{H})) \\
& =
\Hom_\mathcal{O}(j_{U!}\mathcal{G} \otimes_\mathcal{O} \mathcal{F},
\mathcal{H})
\end{align*}
最初の等号は，$j_{U!}$ が加群の制限の左随伴であることによる。
二つ目は補題 \ref{sites-modules-lemma-internal-hom-adjoint-tensor} による。
三つ目は補題 \ref{sites-modules-lemma-internal-hom-restriction} による。
四つ目は，$j_{U!}$ が加群の制限の左随伴であることによる。
五つ目は補題 \ref{sites-modules-lemma-internal-hom-adjoint-tensor} による。
以上と Yoneda の補題から本補題が従う。
\end{proof}

\begin{remark}
\label{sites-modules-remark-j-shriek-tensor}
$\mathcal{C}$ をサイトとする。$\mathcal{F}$ を $\mathcal{C}$ 上の集合の層とし，
局所化射
$j : \Sh(\mathcal{C})/\mathcal{F} \to \Sh(\mathcal{C})$.
を考える。《サイト》の定義 \ref{sites-definition-localize-topos} を参照。
次を主張する：(a) $j_!\mathbf{Z} = \mathbf{Z}_\mathcal{F}^\#$；
(b) $\mathcal{C}$ 上の任意のアーベル群の層 $\mathcal{H}$ に対し，
$j_!(j^{-1}\mathcal{H}) = j_!\mathbf{Z} \otimes_\mathbf{Z} \mathcal{H}$。
$\mathcal{G}$ を $\mathcal{C}$ 上のアーベル群の層とする。
(a) は Yoneda の補題から従う。実際，
$$
\Hom(j_!\mathbf{Z}, \mathcal{G}) =
\Hom(\mathbf{Z}, j^{-1}\mathcal{G}) =
\Hom(\mathbf{Z}_\mathcal{F}^\#, \mathcal{G})
$$
であり，二つ目の等号は，等号の両辺が写像 $\mathcal{F} \to \mathcal{G}$ の集合を
アーベル群とみなしたものに一致することによる。(b) には Yoneda の補題と
\begin{align*}
\Hom(j_!(j^{-1}\mathcal{H}), \mathcal{G})
& =
\Hom(j^{-1}\mathcal{H}, j^{-1}\mathcal{G}) \\
& =
\Hom(\mathbf{Z}, \SheafHom(j^{-1}\mathcal{H}, j^{-1}\mathcal{G})) \\
& =
\Hom(\mathbf{Z}, j^{-1}\SheafHom(\mathcal{H}, \mathcal{G})) \\
& =
\Hom(j_!\mathbf{Z}, \SheafHom(\mathcal{H}, \mathcal{G})) \\
& =
\Hom(j_!\mathbf{Z} \otimes_\mathbf{Z} \mathcal{H}, \mathcal{G})
\end{align*}
を用いる。ここでは随伴性，$\SheafHom$ をとる操作が局所化と可換であること，
および補題 \ref{sites-modules-lemma-internal-hom-adjoint-tensor} を用いた。
\end{remark}

\begin{lemma}
\label{sites-modules-lemma-hom-finite-presentation-colimit}
$(\mathcal{C}, \mathcal{O})$ を環付きサイトとする。
$\mathcal{F}$ を有限表示の $\mathcal{O}$-加群とする。
$\mathcal{G} = \colim_{\lambda \in \Lambda} \mathcal{G}_\lambda$
を $\mathcal{O}$-加群のフィルター付き余極限とする。このとき標準写像
$$
\colim_\lambda \SheafHom_\mathcal{O}(\mathcal{F}, \mathcal{G}_\lambda)
\longrightarrow
\SheafHom_\mathcal{O}(\mathcal{F}, \mathcal{G})
$$
は同型である。
\end{lemma}

\begin{proof}
$\mathcal{C}$ のすべての $U$ への制限後にこの矢印が同型であることを示せば十分である。
加群の層の余極限をとる操作と内部 Hom をとる操作は，いずれも $U$ への制限と
可換である。例えば補題 \ref{sites-modules-lemma-exactness-pushforward-pullback} と
\ref{sites-modules-lemma-internal-hom-restriction} を参照。$U$ を固定する。
被覆 $\{U_i \to U\}_{i \in I}$ が与えられたとき，各 $U_i$ への制限が
同型であることを示せば十分である。したがって，$\mathcal{F}$ が大域表示
$$
\bigoplus\nolimits_{j = 1, \ldots, m}
\mathcal{O}
\longrightarrow
\bigoplus\nolimits_{i = 1, \ldots, n}
\mathcal{O}
\to
\mathcal{F}
\to
0
$$
をもつと仮定してよい。関手 $\SheafHom_\mathcal{O}(-, -)$ は，どちらの変数に
ついても有限直和と可換し，$\SheafHom_\mathcal{O}(\mathcal{O}, -)$ は恒等関手である。
これと補題 \ref{sites-modules-lemma-internal-hom-exact} により，完全列
$$
0 \to
\SheafHom_\mathcal{O}(\mathcal{F}, \mathcal{G}) \to
\bigoplus\nolimits_{i = 1, \ldots, n} \mathcal{G} \to
\bigoplus\nolimits_{j = 1, \ldots, m} \mathcal{G}
$$
を得る。補題 \ref{sites-modules-lemma-limits-colimits} により，
$\textit{Mod}(\mathcal{O})$ ではフィルター付き余極限が完全であるから，
次の可換図式の上段も完全である：
$$
\xymatrix{
0 \ar[r] &
\colim_\lambda \SheafHom_\mathcal{O}(\mathcal{F}, \mathcal{G}_\lambda)
\ar[r] \ar[d] &
\colim_\lambda \bigoplus\nolimits_{i = 1, \ldots, n} \mathcal{G}_\lambda
\ar[r] \ar[d] &
\colim_\lambda \bigoplus\nolimits_{j = 1, \ldots, m} \mathcal{G}_\lambda
\ar[d] \\
0 \ar[r] &
\SheafHom_\mathcal{O}(\mathcal{F}, \mathcal{G}) \ar[r] &
\bigoplus\nolimits_{i = 1, \ldots, n} \mathcal{G} \ar[r] &
\bigoplus\nolimits_{j = 1, \ldots, m} \mathcal{G}
}
$$
右側の二つの垂直矢印は同型なので，結論が従う。
\end{proof}

\begin{lemma}
\label{sites-modules-lemma-finite-presentation-quasi-compact-colimit}
$(\mathcal{C}, \mathcal{O})$ を環付きサイトとする。
$\mathcal{G} = \colim_{\lambda \in \Lambda} \mathcal{G}_\lambda$
を $\mathcal{O}$-加群のフィルター付き余極限とし，
$\mathcal{F}$ を有限表示の $\mathcal{O}$-加群とする。このとき
$$
\colim_\lambda \Hom_\mathcal{O}(\mathcal{F}, \mathcal{G}_\lambda)
=
\Hom_\mathcal{O}(\mathcal{F}, \mathcal{G}).
$$
が成り立つ。ただし，サイト $\mathcal{C}$ に対して《サイト》の補題
\ref{sites-lemma-directed-colimits-global-sections} の (4) の仮定が満たされるものとする。
《サイト》の注意 \ref{sites-remark-stronger-conditions} も参照されたい。
\end{lemma}

\begin{proof}
$\mathcal{H} =
\SheafHom_\mathcal{O}(\mathcal{F}, \colim \mathcal{G}_\lambda)$
および $\mathcal{H}_\lambda =
\SheafHom_\mathcal{O}(\mathcal{F}, \mathcal{G}_\lambda)$ とおく。構成により
$$
\Hom_\mathcal{O}(\mathcal{F}, \mathcal{G}) = \Gamma(\mathcal{C}, \mathcal{H})
\quad\text{and}\quad
\Hom_\mathcal{O}(\mathcal{F}, \mathcal{G}_\lambda) =
\Gamma(\mathcal{C}, \mathcal{H}_\lambda)
$$
である。補題 \ref{sites-modules-lemma-hom-finite-presentation-colimit} により
$\mathcal{H} = \colim \mathcal{H}_\lambda$ である。したがって，本補題は
《サイト》の補題 \ref{sites-lemma-directed-colimits-global-sections} から従う。
\end{proof}









\section{平坦加群}
\label{sites-modules-section-flat}

\noindent
平坦加群は，環上の加群の場合とまったく同様に定義できる。

\begin{definition}
\label{sites-modules-definition-flat}
$\mathcal{C}$ を圏とし，$\mathcal{O}$ を環の前層とする。
\begin{enumerate}
\item $\mathcal{O}$-加群の前層 $\mathcal{F}$ が {\it 平坦}であるとは，関手
$$
\textit{PMod}(\mathcal{O})
\longrightarrow
\textit{PMod}(\mathcal{O}), \quad
\mathcal{G} \mapsto \mathcal{G} \otimes_{p, \mathcal{O}} \mathcal{F}
$$
が完全であることをいう。
\item 環の前層の写像 $\mathcal{O} \to \mathcal{O}'$ が {\it 平坦}であるとは，
$\mathcal{O}'$ が $\mathcal{O}$-加群の前層として平坦であることをいう。
\item $\mathcal{C}$ がサイト，$\mathcal{O}$ が環の層，
$\mathcal{F}$ が $\mathcal{O}$-加群の層であるとする。このとき
$\mathcal{F}$ が {\it 平坦}であるとは，関手
$$
\textit{Mod}(\mathcal{O})
\longrightarrow
\textit{Mod}(\mathcal{O}), \quad
\mathcal{G} \mapsto \mathcal{G} \otimes_\mathcal{O} \mathcal{F}
$$
が完全であることをいう。
\item サイト上の環の層の写像 $\mathcal{O} \to \mathcal{O}'$ が
{\it 平坦}であるとは，$\mathcal{O}'$ が $\mathcal{O}$-加群の層として
平坦であることをいう。
\end{enumerate}
\end{definition}

\noindent
平坦加群および平坦な環写像の概念は内在的である
（第 \ref{sites-modules-section-intrinsic} 節）。

\begin{lemma}
\label{sites-modules-lemma-flatness-presheaves}
$\mathcal{C}$ を圏，$\mathcal{O}$ を環の前層，
$\mathcal{F}$ を $\mathcal{O}$-加群の前層とする。
各 $\mathcal{F}(U)$ が平坦な $\mathcal{O}(U)$-加群ならば，
$\mathcal{F}$ は平坦である。
\end{lemma}

\begin{proof}
定義から直ちに従う。
\end{proof}

\begin{lemma}
\label{sites-modules-lemma-flatness-sheafification}
$\mathcal{C}$ をサイト，$\mathcal{O}$ を環の前層，
$\mathcal{F}$ を $\mathcal{O}$-加群の前層とする。
$\mathcal{F}$ が平坦な $\mathcal{O}$-加群ならば，
$\mathcal{F}^\#$ は平坦な $\mathcal{O}^\#$-加群である。
\end{lemma}

\begin{proof}
省略する。（ヒント：層化は完全である。）
\end{proof}

\begin{lemma}
\label{sites-modules-lemma-flatness-sheafification-refined}
$\mathcal{C}$ をサイト，$\mathcal{O}$ を環の前層，
$\mathcal{F}$ を $\mathcal{O}$-加群の前層とする。
$\mathcal{C}$ の各対象 $U$ が，$\mathcal{F}(U_i)$ が平坦な
$\mathcal{O}(U_i)$-加群となる被覆 $\{U_i \to U\}_{i \in I}$ をもつと仮定する。
このとき $\mathcal{F}^\#$ は平坦な $\mathcal{O}^\#$-加群である。
\end{lemma}

\begin{proof}
$\mathcal{G} \subset \mathcal{G}'$ を $\mathcal{O}^\#$-加群の包含とする。
示すべきことは，
$$
\mathcal{G} \otimes_{\mathcal{O}^\#} \mathcal{F}^\#
\longrightarrow
\mathcal{G}' \otimes_{\mathcal{O}^\#} \mathcal{F}^\#
$$
が単射であることである。補題 \ref{sites-modules-lemma-sheafification-tensor} により，
この矢印の始域は前層 $\mathcal{G} \otimes_{p, \mathcal{O}} \mathcal{F}$ の
層化であり，終域についても同様である。$U$ が $\mathcal{C}$ の対象で，
$\mathcal{F}(U)$ が平坦な $\mathcal{O}(U)$-加群ならば，
$$
(\mathcal{G} \otimes_{p, \mathcal{O}} \mathcal{F})(U) =
\mathcal{G}(U) \otimes_{\mathcal{O}(U)} \mathcal{F}(U)
\longrightarrow
\mathcal{G}'(U) \otimes_{\mathcal{O}(U)} \mathcal{F}(U) =
(\mathcal{G}' \otimes_{p, \mathcal{O}} \mathcal{F})(U)
$$
は単射である。したがって，次を示すことに帰着する：
$\mathcal{C}$ 上の前層の写像 $f : \mathcal{H} \to \mathcal{H}'$ が与えられ，
$\mathcal{C}$ の各 $U$ が，$\mathcal{H}(U_i) \to \mathcal{H}'(U_i)$ が単射となる
被覆 $\{U_i \to U\}_{i \in I}$ をもつならば，$f^\#$ は単射である。
これは読者への演習として残す。
\end{proof}

\begin{lemma}
\label{sites-modules-lemma-colimits-flat}
余極限とテンソル積。
\begin{enumerate}
\item 平坦な加群の前層のフィルター付き余極限は平坦である。
平坦な加群の前層の直和は平坦である。
\item 平坦な加群の層のフィルター付き余極限は平坦である。
平坦な加群の層の直和は平坦である。
\end{enumerate}
\end{lemma}

\begin{proof}
各対象上の切断を見れば，(1) は補題 \ref{sites-modules-lemma-tensor-commute-colimits} と
《代数》の補題 \ref{algebra-lemma-directed-colimit-exact} から従う。
(2) には，補題 \ref{sites-modules-lemma-tensor-commute-colimits} と，完全複体の
フィルター付き余極限が完全複体であるという事実を用いる
（ここでは層化が完全で余極限と可換であることを使う）。細部は省略する。
\end{proof}

\begin{lemma}
\label{sites-modules-lemma-restriction-flat}
$(\mathcal{C}, \mathcal{O})$ を環付きサイトとし，
$U$ を $\mathcal{C}$ の対象とする。
$\mathcal{F}$ が平坦な $\mathcal{O}$-加群ならば，
$\mathcal{F}|_U$ は平坦な $\mathcal{O}_U$-加群である。
\end{lemma}

\begin{proof}
$\mathcal{G}_1 \to \mathcal{G}_2 \to \mathcal{G}_3$ を
$\mathcal{O}_U$-加群の完全複体とする。
$j_{U!}$ は完全であり（補題 \ref{sites-modules-lemma-extension-by-zero-exact}），
$\mathcal{F}$ は $\mathcal{O}$-加群として平坦なので，加群
$$
j_{U!}(\mathcal{G}_i \otimes_{\mathcal{O}_U} \mathcal{F}|_U) =
j_{U!}\mathcal{G}_i \otimes_\mathcal{O} \mathcal{F}
$$
からなる複体は完全である（補題 \ref{sites-modules-lemma-j-shriek-and-tensor}）。したがって，
$\mathcal{G}_1 \otimes_{\mathcal{O}_U} \mathcal{F}|_U \to
\mathcal{G}_2 \otimes_{\mathcal{O}_U} \mathcal{F}|_U \to
\mathcal{G}_3 \otimes_{\mathcal{O}_U} \mathcal{F}|_U$
は補題 \ref{sites-modules-lemma-j-shriek-reflects-exactness} により完全である。
\end{proof}

\begin{lemma}
\label{sites-modules-lemma-j-shriek-flat}
$\mathcal{C}$ を圏，$\mathcal{O}$ を環の前層，
$U$ を $\mathcal{C}$ の対象とする。関手
$j_U : \mathcal{C}/U \to \mathcal{C}$ を考える。
\begin{enumerate}
\item $\mathcal{O}$-加群の前層 $j_{U!}\mathcal{O}_U$ は平坦である
（注意 \ref{sites-modules-remark-localize-presheaves} を参照）。
\item $\mathcal{C}$ がサイトで $\mathcal{O}$ が環の層ならば，
$j_{U!}\mathcal{O}_U$ は平坦な $\mathcal{O}$-加群の層である。
\end{enumerate}
\end{lemma}

\begin{proof}
(1) を証明する。注意 \ref{sites-modules-remark-localize-presheaves} の議論により，
$$
j_{U!}\mathcal{O}_U(V)
=
\bigoplus\nolimits_{\varphi \in \Mor_\mathcal{C}(V, U)}
\mathcal{O}(V)
$$
であり，これは平坦な $\mathcal{O}(V)$-加群である。したがって，(1) は
補題 \ref{sites-modules-lemma-flatness-presheaves} から従う。さらに
$j_{U!}\mathcal{O}_U = (j_{U!}\mathcal{O}_U)^\#$
（最初の $j_{U!}$ は層上，二つ目は前層上）と補題
\ref{sites-modules-lemma-flatness-sheafification} から (2) が従う。
\end{proof}

\begin{lemma}
\label{sites-modules-lemma-module-quotient-flat}
$\mathcal{C}$ を圏とし，$\mathcal{O}$ を環の前層とする。
\begin{enumerate}
\item 任意の $\mathcal{O}$-加群の前層は，直和
$\bigoplus j_{U_i!}\mathcal{O}_{U_i}$ の商である。
\item 任意の $\mathcal{O}$-加群の前層は，平坦な
$\mathcal{O}$-加群の前層の商である。
\item $\mathcal{C}$ がサイトで $\mathcal{O}$ が環の層ならば，任意の
$\mathcal{O}$-加群の層は，直和 $\bigoplus j_{U_i!}\mathcal{O}_{U_i}$ の商である。
\item $\mathcal{C}$ がサイトで $\mathcal{O}$ が環の層ならば，任意の
$\mathcal{O}$-加群の層は，平坦な $\mathcal{O}$-加群の層の商である。
\end{enumerate}
\end{lemma}

\begin{proof}
(1) を証明する。$\mathcal{C}$ の各対象 $U$ と各
$s \in \mathcal{F}(U)$ に対し，射 $\mathcal{O}_U \to \mathcal{F}|_U$，
$1 \mapsto s$ の随伴として射 $j_{U!}\mathcal{O}_U \to \mathcal{F}$ を得る。
明らかに写像
$$
\bigoplus\nolimits_{(U, s)} j_{U!}\mathcal{O}_U
\longrightarrow
\mathcal{F}
$$
は全射である。補題 \ref{sites-modules-lemma-colimits-flat} と
\ref{sites-modules-lemma-j-shriek-flat} を合わせれば始域は平坦であり，(2) が従う。
層の場合は，これを層化するか，同じ議論を繰り返せば従う。
\end{proof}

\begin{lemma}
\label{sites-modules-lemma-flat-tor-zero}
$\mathcal{C}$ を圏とし，$\mathcal{O}$ を環の前層とする。
$$
0 \to \mathcal{F}'' \to \mathcal{F}' \to \mathcal{F} \to 0
$$
を $\mathcal{O}$-加群の前層の短完全列とする。
$\mathcal{G}$ を $\mathcal{O}$-加群の前層とする。
\begin{enumerate}
\item $\mathcal{F}$ が平坦な加群の前層ならば，列
$$
0 \to
\mathcal{F}'' \otimes_{p, \mathcal{O}} \mathcal{G} \to
\mathcal{F}' \otimes_{p, \mathcal{O}} \mathcal{G} \to
\mathcal{F} \otimes_{p, \mathcal{O}} \mathcal{G} \to 0
$$
は完全である。
\item $\mathcal{C}$ がサイトで，$\mathcal{O}$，
$\mathcal{F}$，$\mathcal{F}'$，$\mathcal{F}''$，$\mathcal{G}$ が層であり，
$\mathcal{F}$ が加群の層として平坦ならば，列
$$
0 \to
\mathcal{F}'' \otimes_\mathcal{O} \mathcal{G} \to
\mathcal{F}' \otimes_\mathcal{O} \mathcal{G} \to
\mathcal{F} \otimes_\mathcal{O} \mathcal{G} \to 0
$$
は完全である。
\end{enumerate}
\end{lemma}

\begin{proof}
$\mathcal{G}$ へ全射する平坦な $\mathcal{O}$-加群の前層
$\mathcal{G}'$ を選ぶ。これは補題 \ref{sites-modules-lemma-module-quotient-flat} により可能である。
さらに
$\mathcal{G}'' = \Ker(\mathcal{G}' \to \mathcal{G})$.
とおく。次の図式に蛇の補題を適用すれば本補題が従う：
$$
\begin{matrix}
 & & 0 & & 0 & & 0 & & \\
 & & \uparrow & & \uparrow & & \uparrow & & \\
 & & \mathcal{F}'' \otimes_{p, \mathcal{O}} \mathcal{G} & \to &
     \mathcal{F}' \otimes_{p, \mathcal{O}} \mathcal{G} & \to &
     \mathcal{F} \otimes_{p, \mathcal{O}} \mathcal{G} & \to & 0 \\
 & & \uparrow & & \uparrow & & \uparrow & & \\
0 & \to & \mathcal{F}'' \otimes_{p, \mathcal{O}} \mathcal{G}' & \to &
          \mathcal{F}' \otimes_{p, \mathcal{O}} \mathcal{G}' & \to &
          \mathcal{F} \otimes_{p, \mathcal{O}} \mathcal{G}' & \to & 0 \\
 & & \uparrow & & \uparrow & & \uparrow & & \\
 & & \mathcal{F}'' \otimes_{p, \mathcal{O}} \mathcal{G}'' & \to &
     \mathcal{F}' \otimes_{p, \mathcal{O}} \mathcal{G}'' & \to &
     \mathcal{F} \otimes_{p, \mathcal{O}} \mathcal{G}'' & \to & 0 \\
 & & & & & & \uparrow & & \\
 & & & & & & 0 & &
\end{matrix}
$$
ここですべての行と列は完全である。中段が完全なのは，平坦加群
$\mathcal{G}'$ とのテンソル積をとる操作が完全だからである。
層の場合の証明もまったく同じである。
\end{proof}

\begin{lemma}
\label{sites-modules-lemma-flat-ses}
$\mathcal{C}$ を圏とし，$\mathcal{O}$ を環の前層とする。
$$
0 \to
\mathcal{F}_2 \to
\mathcal{F}_1 \to
\mathcal{F}_0 \to 0
$$
を $\mathcal{O}$-加群の前層の短完全列とする。
\begin{enumerate}
\item $\mathcal{F}_2$ と $\mathcal{F}_0$ が平坦ならば，
$\mathcal{F}_1$ も平坦である。
\item $\mathcal{F}_1$ と $\mathcal{F}_0$ が平坦ならば，
$\mathcal{F}_2$ も平坦である。
\end{enumerate}
$\mathcal{C}$ がサイトで $\mathcal{O}$ が環の層ならば，
$\textit{Mod}(\mathcal{O})$ においても同じ結果が成り立つ。
\end{lemma}

\begin{proof}
$\mathcal{G}^\bullet$ を $\mathcal{O}$-加群の前層の任意の完全複体とし，
$\mathcal{F}_0$ が平坦であると仮定する。補題 \ref{sites-modules-lemma-flat-tor-zero} により，
$$
0 \to
\mathcal{G}^\bullet \otimes_{p, \mathcal{O}} \mathcal{F}_2 \to
\mathcal{G}^\bullet \otimes_{p, \mathcal{O}} \mathcal{F}_1 \to
\mathcal{G}^\bullet \otimes_{p, \mathcal{O}} \mathcal{F}_0 \to 0
$$
は $\mathcal{O}$-加群の前層の複体の短完全列である。
したがって，(1) と (2) は蛇の補題から従う。加群の層の場合も同様に示される。
\end{proof}

\begin{lemma}
\label{sites-modules-lemma-flat-resolution-of-flat}
$\mathcal{C}$ を圏とし，$\mathcal{O}$ を環の前層とする。
$$
\ldots \to
\mathcal{F}_2 \to
\mathcal{F}_1 \to
\mathcal{F}_0 \to
\mathcal{Q} \to 0
$$
を $\mathcal{O}$-加群の前層の完全複体とする。
$\mathcal{Q}$ とすべての $\mathcal{F}_i$ が平坦な $\mathcal{O}$-加群ならば，
任意の $\mathcal{O}$-加群の前層 $\mathcal{G}$ に対して，複体
$$
\ldots \to
\mathcal{F}_2 \otimes_{p, \mathcal{O}} \mathcal{G} \to
\mathcal{F}_1 \otimes_{p, \mathcal{O}} \mathcal{G} \to
\mathcal{F}_0 \otimes_{p, \mathcal{O}} \mathcal{G} \to
\mathcal{Q} \otimes_{p, \mathcal{O}} \mathcal{G} \to 0
$$
も完全である。$\mathcal{C}$ がサイトで $\mathcal{O}$ が環の層ならば，
$\textit{Mod}(\mathcal{O})$ においても同じ結果が成り立つ。
\end{lemma}

\begin{proof}
補題 \ref{sites-modules-lemma-flat-tor-zero} から従う。実際，複体を短完全列に分割し，
補題 \ref{sites-modules-lemma-flat-ses} を用いて
$\Im(\mathcal{F}_{i + 1} \to \mathcal{F}_i)$ が平坦であることを帰納的に示す。
\end{proof}

\begin{lemma}
\label{sites-modules-lemma-tensor-flats}
$(\mathcal{C}, \mathcal{O})$ を環付きサイトとする。
$\mathcal{G}$ と $\mathcal{F}$ が平坦な $\mathcal{O}$-加群ならば，
$\mathcal{G} \otimes_\mathcal{O} \mathcal{F}$ は平坦な $\mathcal{O}$-加群である。
\end{lemma}

\begin{proof}
実際，
$$
(\mathcal{G} \otimes_\mathcal{O} \mathcal{F}) \otimes_\mathcal{O} \mathcal{H}
=
\mathcal{G} \otimes_\mathcal{O} (\mathcal{F} \otimes_\mathcal{O} \mathcal{H})
$$
であり，完全関手の合成は完全だからである。
\end{proof}

\begin{lemma}
\label{sites-modules-lemma-flat-change-of-rings}
$\mathcal{O}_1 \to \mathcal{O}_2$ をサイト $\mathcal{C}$ 上の環の層の写像とする。
$\mathcal{G}$ が平坦な $\mathcal{O}_1$-加群ならば，
$\mathcal{G} \otimes_{\mathcal{O}_1} \mathcal{O}_2$
は平坦な $\mathcal{O}_2$-加群である。
\end{lemma}

\begin{proof}
実際，
$$
(\mathcal{G} \otimes_{\mathcal{O}_1} \mathcal{O}_2)
\otimes_{\mathcal{O}_2} \mathcal{H}
=
\mathcal{G} \otimes_{\mathcal{O}_1} \mathcal{F}
$$
である（例えばアーベル群の層として）。
\end{proof}

\noindent
次の補題は平坦性の方程式的判定法
（《代数》の補題 \ref{algebra-lemma-flat-eq}）の類似である。

\begin{lemma}
\label{sites-modules-lemma-flat-eq}
$(\mathcal{C}, \mathcal{O})$ を環付きサイトとし，
$\mathcal{F}$ を $\mathcal{O}$-加群とする。次は同値である：
\begin{enumerate}
\item $\mathcal{F}$ は平坦な $\mathcal{O}$-加群である。
\item $U$ を $\mathcal{C}$ の対象とし，
$$
\mathcal{O}_U \xrightarrow{(f_1, \ldots, f_n)}
\mathcal{O}_U^{\oplus n} \xrightarrow{(s_1, \ldots, s_n)}
\mathcal{F}|_U
$$
を $\mathcal{O}_U$-加群の複体とする。このとき被覆 $\{U_i \to U\}$ が存在し，
各 $i$ に対して $(s_1, \ldots, s_n)|_{U_i}$ の分解
$$
\mathcal{O}_{U_i}^{\oplus n}
\xrightarrow{B_i}
\mathcal{O}_{U_i}^{\oplus l_i} \xrightarrow{(t_{i1}, \ldots, t_{il_i})}
\mathcal{F}|_{U_i}
$$
が存在して，$B_i \circ (f_1, \ldots, f_n)|_{U_i} = 0$ を満たす。
\item $U$ を $\mathcal{C}$ の対象とし，
$$
\mathcal{O}_U^{\oplus m} \xrightarrow{A}
\mathcal{O}_U^{\oplus n} \xrightarrow{(s_1, \ldots, s_n)}
\mathcal{F}|_U
$$
を $\mathcal{O}_U$-加群の複体とする。このとき被覆 $\{U_i \to U\}$ が存在し，
各 $i$ に対して $(s_1, \ldots, s_n)|_{U_i}$ の分解
$$
\mathcal{O}_{U_i}^{\oplus n}
\xrightarrow{B_i}
\mathcal{O}_{U_i}^{\oplus l_i} \xrightarrow{(t_{i1}, \ldots, t_{il_i})}
\mathcal{F}|_{U_i}
$$
が存在して，$B_i \circ A|_{U_i} = 0$ を満たす。
\end{enumerate}
\end{lemma}

\begin{proof}
(1) を仮定する。$\mathcal{I} \subset \mathcal{O}_U$ を
$f_1, \ldots, f_n$ で生成されるイデアルの層とする。このとき
$\sum f_j \otimes s_j$ は
$\mathcal{I} \otimes_{\mathcal{O}_U} \mathcal{F}|_U$ の切断であり，
$\mathcal{F}|_U$ において零に写る。$\mathcal{F}|_U$ は平坦なので
（補題 \ref{sites-modules-lemma-restriction-flat}），写像
$\mathcal{I} \otimes_{\mathcal{O}_U} \mathcal{F}|_U \to \mathcal{F}|_U$
は単射である。$\mathcal{I} \otimes_{\mathcal{O}_U} \mathcal{F}|_U$ は
前層のテンソル積に付随する層なので，被覆 $\{U_i \to U\}$ が存在して，
$\sum f_j|_{U_i} \otimes s_j|_{U_i}$ は
$\mathcal{I}(U_i) \otimes_{\mathcal{O}(U_i)} \mathcal{F}(U_i)$ において零となる。
《代数》の補題 \ref{algebra-lemma-relations} を用いて定義を書き下せば，
$t_{i1}, \ldots, t_{i l_i} \in \mathcal{F}(U_i)$ と
$a_{ijk} \in \mathcal{O}(U_i)$ が存在して，
$\sum_j a_{ijk}f_j|_{U_i} = 0$ および
$s_j|_{U_i} = \sum_k a_{ijk}t_{ik}$ を満たすことが分かる。
したがって (2) が成り立つ。

\medskip\noindent
(2) を仮定する。(3) における $U$，$n$，$m$，$A$，
$s_1, \ldots, s_n$ が与えられたとする。$A$ は $m$ 個の列をもつ。
$m$ に関する帰納法で (3) の主張を示す。基底の場合 $m = 0$ には，
被覆 $\{U \to U\}$ をとり，分解には
$\mathcal{O}_U^n \to \mathcal{O}_U^n \to \mathcal{F}|_U$ を用いる。
ここで最初の写像は恒等写像，二つ目は $(s_1, \ldots, s_n)$ である。
$m > 0$ と仮定する。$(f_1, \ldots, f_n)$ を $A$ の最後の列として (2) を
適用すると，新しい図式
$$
\mathcal{O}_{U_i}^{\oplus m} \xrightarrow{B_i \circ A|_{U_i}}
\mathcal{O}_{U_i}^{\oplus l_i} \xrightarrow{(t_{i1}, \ldots, t_{il_i})}
\mathcal{F}|_{U_i}
$$
を得るが，$A_i = B_i \circ A|_{U_i}$ の最後の列は零である。
そこで $U_i$，$l_i$，$m - 1$，$A_i$ の最初の $m - 1$ 列からなる行列，
および $t_{i1}, \ldots, t_{il_i}$ に帰納法の仮定を適用すると，
被覆 $\{U_{ij} \to U_j\}$ と分解
$$
\mathcal{O}_{U_{ij}}^{\oplus l_i}
\xrightarrow{C_{ij}}
\mathcal{O}_{U_{ij}}^{\oplus k_{ij}}
\xrightarrow{(v_{ij1}, \ldots, v_{ij k_{ij}})}
\mathcal{F}|_{U_{ij}}
$$
を得る。これは $(t_{i1}, \ldots, t_{il_i})|_{U_{ij}}$ の分解で，
$C_i \circ B_i|_{U_{ij}} \circ A|_{U_{ij}} = 0$ を満たす。
このとき $\{U_{ij} \to U\}$ は被覆であり，
$B_{ij} = C_i \circ B_i|_{U_{ij}}$ と $v_{ija}$ を用いれば所望の分解を得る。
以上により (2) は (3) を含意する。

\medskip\noindent
(3) を仮定する。$\mathcal{G} \to \mathcal{H}$ を
$\mathcal{O}$-加群の単射準同型とする。示すべきことは，
$\mathcal{G} \otimes_\mathcal{O} \mathcal{F} \to
\mathcal{H} \otimes_\mathcal{O} \mathcal{F}$
が単射であることである。$U$ を $\mathcal{C}$ の対象とし，
$s \in (\mathcal{G} \otimes_\mathcal{O} \mathcal{F})(U)$ を，
$\mathcal{H} \otimes_\mathcal{O} \mathcal{F}$ において零へ写る切断とする。
$s$ が零であることを示せばよい。
$\mathcal{G} \otimes_\mathcal{O} \mathcal{F}$
は層なので，すべての $i \in I$ に対して $s|_{U_i}$ が零となる
$\mathcal{C}$ の被覆 $\{U_i \to U\}_{i \in I}$ を見つければ十分である。
したがって，いつでも $U$ を被覆の各元で置き換えてよい。
特に，$\mathcal{G} \otimes_\mathcal{O} \mathcal{F}$ は
$\mathcal{G} \otimes_{p, \mathcal{O}} \mathcal{F}$ の層化なので，
$s' \in \mathcal{G}(U) \otimes_{\mathcal{O}(U)} \mathcal{F}(U)$ の像であると
仮定してよい。$\mathcal{H} \otimes_\mathcal{O} \mathcal{F}$ についても
同様に論じれば，$s'$ は
$\mathcal{H}(U) \otimes_{\mathcal{O}(U)} \mathcal{F}(U)$ において零へ写ると
仮定してよい。
$\mathcal{F}(U) = \colim M_\alpha$ を，有限表示
$\mathcal{O}(U)$-加群 $M_\alpha$ のフィルター付き余極限として書く
（《代数》の補題 \ref{algebra-lemma-module-colimit-fp}）。
テンソル積はフィルター付き余極限と可換するので
（《代数》の補題 \ref{algebra-lemma-tensor-products-commute-with-limits}），
$s'$ がある $s'' \in \mathcal{G}(U) \otimes_{\mathcal{O}(U)} M_\alpha$ から来て，
$s''$ が $\mathcal{H}(U) \otimes_{\mathcal{O}(U)} M_\alpha$ において
零へ写るような $\alpha$ を選べる。$\alpha$ と $s''$ を固定し，表示
$$
\mathcal{O}(U)^{\oplus m} \xrightarrow{A} \mathcal{O}(U)^{\oplus n}
\to M_\alpha \to 0
$$
を選ぶ。対応する $\mathcal{O}_U$-加群の複体
$$
\mathcal{O}_U^{\oplus m} \xrightarrow{A}
\mathcal{O}_U^{\oplus n} \xrightarrow{(s_1, \ldots, s_n)}
\mathcal{F}|_U
$$
に (3) を適用する。$U$ を被覆の各元 $U_i$ で置き換えると，写像
$$
M_\alpha \to \mathcal{F}(U)
$$
は，ある $l$ に対する自由加群 $\mathcal{O}(U)^{\oplus l}$ を経由することが分かる。
$\mathcal{G}(U) \to \mathcal{H}(U)$ は単射なので，
$$
\mathcal{G}(U) \otimes_{\mathcal{O}(U)} \mathcal{O}(U)^{\oplus l}
\to
\mathcal{H}(U) \otimes_{\mathcal{O}(U)} \mathcal{O}(U)^{\oplus l}
$$
も単射である。したがって，$s''$ は右辺の加群において零へ写るので，
左辺の加群においても零へ写る。すなわち，所望のとおり $s$ は零である。
\end{proof}

\begin{lemma}
\label{sites-modules-lemma-flat-over-thickening}
$\mathcal{C}$ をサイトとする。$\mathcal{O}' \to \mathcal{O}$ を環の層の全射とし，
その核 $\mathcal{I}$ は二乗零イデアルであるとする。
$\mathcal{F}'$ を $\mathcal{O}'$-加群とし，
$\mathcal{F} = \mathcal{F}'/\mathcal{I}\mathcal{F}'$ とおく。次は同値である：
\begin{enumerate}
\item $\mathcal{F}'$ は平坦な $\mathcal{O}'$-加群である；
\item $\mathcal{F}$ は平坦な $\mathcal{O}$-加群であり，
$\mathcal{I} \otimes_\mathcal{O} \mathcal{F} \to \mathcal{F}'$
は単射である。
\end{enumerate}
\end{lemma}

\begin{proof}
(1) が成り立つとする。このとき，補題 \ref{sites-modules-lemma-flat-change-of-rings} により
$\mathcal{F} = \mathcal{F}' \otimes_{\mathcal{O}'} \mathcal{O}$ は
$\mathcal{O}$ 上平坦である。また，完全列
$0 \to \mathcal{I} \to \mathcal{O}' \to \mathcal{O} \to 0$ に
$- \otimes_{\mathcal{O}'} \mathcal{F}'$ を適用すれば，写像
$\mathcal{I} \otimes_\mathcal{O} \mathcal{F} \to \mathcal{F}'$
が単射であることが分かる。補題 \ref{sites-modules-lemma-flat-tor-zero} を参照。
(2) を仮定する。以下の証明では，$\mathcal{I}$ で零化される任意の
$\mathcal{O}'$-加群 $\mathcal{K}$ に対し
$\mathcal{K} \otimes_{\mathcal{O}'} \mathcal{F}' =
\mathcal{K} \otimes_\mathcal{O} \mathcal{F}$
であることを断りなく用いる。
$\alpha : \mathcal{G}' \to \mathcal{H}'$ を $\mathcal{O}'$-加群の単射とする。
$\mathcal{G} \subset \mathcal{G}'$，それぞれ
$\mathcal{H} \subset \mathcal{H}'$ を，$\mathcal{I}$ で零化される切断の部分層とする。
図式
$$
\xymatrix{
\mathcal{G} \otimes_{\mathcal{O}'} \mathcal{F}' \ar[r] \ar[d] &
\mathcal{G}' \otimes_{\mathcal{O}'} \mathcal{F}' \ar[r] \ar[d] &
\mathcal{G}'/\mathcal{G} \otimes_{\mathcal{O}'} \mathcal{F}' \ar[r] \ar[d] &
0 \\
\mathcal{H} \otimes_{\mathcal{O}'} \mathcal{F}' \ar[r] &
\mathcal{H}' \otimes_{\mathcal{O}'} \mathcal{F}' \ar[r] &
\mathcal{H}'/\mathcal{H} \otimes_{\mathcal{O}'} \mathcal{F}' \ar[r] & 0
}
$$
を考える。$\mathcal{G}'/\mathcal{G}$ と $\mathcal{H}'/\mathcal{H}$ は
$\mathcal{I}$ で零化され，$\mathcal{G}'/\mathcal{G} \to
\mathcal{H}'/\mathcal{H}$ は単射であることに注意する。
$\mathcal{F}$ は $\mathcal{O}$ 上平坦なので，右側の垂直矢印は単射である。
左側の垂直矢印についても同様である。したがって中央の垂直矢印は単射であり，
$\mathcal{F}'$ は平坦である。
\end{proof}

\begin{lemma}
\label{sites-modules-lemma-neighbourhood-isomorphism}
$\mathcal{C}$ をサイトとする。$\mathcal{O} \to \mathcal{O}'$ を
環の層の平坦準同型とする。$\mathcal{I} \subset \mathcal{O}$ をイデアルの層とし，
誘導写像
$\mathcal{O}/\mathcal{I} \to \mathcal{O}'/\mathcal{I}\mathcal{O}'$
が同型であるとする。ある $n \geq 0$ に対して $\mathcal{I}^n$ で零化される
任意の $\mathcal{O}$-加群 $\mathcal{F}$ に対し，写像
$\text{id} \otimes 1 :
\mathcal{F} \to \mathcal{F} \otimes_\mathcal{O} \mathcal{O}'$
は同型である。
\end{lemma}

\begin{proof}
省略する。ヒント：《代数詳論》の補題
\ref{more-algebra-lemma-neighbourhood-isomorphism} を参照。
\end{proof}





\section{双対}
\label{sites-modules-section-duals}

\noindent
$(\mathcal{C}, \mathcal{O})$ を環付きサイトとする。
第 \ref{sites-modules-section-tensor-product} 節で構成したテンソル積を備えた
$\mathcal{O}$-加群の圏は対称モノイド圏である。
$\mathcal{O}$-加群 $\mathcal{F}$ に対し，次は同値である：
\begin{enumerate}
\item $\mathcal{O}$-加群のモノイド圏において $\mathcal{F}$ は左双対をもつ；
\item $\mathcal{C}$ の各対象 $U$ に対し，被覆 $\{U_i \to U\}$ が存在して，
$\mathcal{F}|_{U_i}$ は有限自由 $\mathcal{O}|_{U_i}$-加群の直和因子となる；
\item $\mathcal{F}$ は有限表示かつ平坦な $\mathcal{O}$-加群である。
\end{enumerate}
これは本節の例 \ref{sites-modules-example-dual} と補題 \ref{sites-modules-lemma-left-dual-module}，
\ref{sites-modules-lemma-flat-locally-finite-presentation} で証明される。

\begin{example}
\label{sites-modules-example-dual}
$(\mathcal{C}, \mathcal{O})$ を環付きサイトとする。
$\mathcal{F}$ を $\mathcal{O}$-加群とし，$\mathcal{C}$ の各対象 $U$ に対して，
$\mathcal{F}|_{U_i}$ が有限自由 $\mathcal{O}|_{U_i}$-加群の直和因子となるような
被覆 $\{U_i \to U\}$ が存在するとする。このとき写像
$$
\mathcal{F} \otimes_\mathcal{O}
\SheafHom_\mathcal{O}(\mathcal{F}, \mathcal{O})
\longrightarrow
\SheafHom_\mathcal{O}(\mathcal{F}, \mathcal{F})
$$
は同型である。実際，これは局所的な問題であり，$\mathcal{F}$ が有限自由ならば
成り立ち，さらに，これが成り立つ加群の任意の直和因子に対しても成り立つ
（細部は省略する）。写像
$$
\eta :
\mathcal{O}
\longrightarrow
\mathcal{F} \otimes_\mathcal{O}
\SheafHom_\mathcal{O}(\mathcal{F}, \mathcal{O})
$$
を $\eta$ と書く。これは $1$ を，上の同型のもとで
$\text{id}_\mathcal{F}$ に対応する切断へ送る写像である。また，評価写像
$$
\epsilon : 
\SheafHom_\mathcal{O}(\mathcal{F}, \mathcal{O})
\otimes_\mathcal{O} \mathcal{F}
\longrightarrow
\mathcal{O}
$$
を $\epsilon$ と書く。このとき，
$\SheafHom_\mathcal{O}(\mathcal{F}, \mathcal{O}), \eta, \epsilon$
は《圏》の定義 \ref{categories-definition-dual} の意味で
$\mathcal{F}$ の左双対である。次の等式の確認は省略する：
$(1 \otimes \epsilon) \circ (\eta \otimes 1) = \text{id}_\mathcal{F}$
および
$(\epsilon \otimes 1) \circ (1 \otimes \eta) =
\text{id}_{\SheafHom_\mathcal{O}(\mathcal{F}, \mathcal{O})}$.
\end{example}

\begin{lemma}
\label{sites-modules-lemma-left-dual-module}
$(\mathcal{C}, \mathcal{O})$ を環付きサイトとし，
$\mathcal{F}$ を $\mathcal{O}$-加群とする。
$\mathcal{G}, \eta, \epsilon$ を $\mathcal{O}$-加群のモノイド圏における
$\mathcal{F}$ の左双対とする。《圏》の定義
\ref{categories-definition-dual} を参照。このとき，
\begin{enumerate}
\item $\mathcal{C}$ の各対象 $U$ に対し，被覆 $\{U_i \to U\}$ が存在して，
$\mathcal{F}|_{U_i}$ は有限自由 $\mathcal{O}|_{U_i}$-加群の直和因子となる；
\item 局所切断 $\lambda$ を $(\lambda \otimes 1)(\eta)$ へ送る写像
$e : \SheafHom_\mathcal{O}(\mathcal{F}, \mathcal{O}) \to \mathcal{G}$
は同型である；
\item $\mathcal{F}$ と $\mathcal{G}$ の局所切断 $f$，$g$ に対して
$\epsilon(f, g) = e^{-1}(g)(f)$ である。
\end{enumerate}
\end{lemma}

\begin{proof}
仮定は，
$$
\mathcal{F} \xrightarrow{\eta \otimes 1}
\mathcal{F} \otimes_\mathcal{O} \mathcal{G}
\otimes_\mathcal{O} \mathcal{F}
\xrightarrow{1 \otimes \epsilon} \mathcal{F}
\quad\text{かつ}\quad
\mathcal{G} \xrightarrow{1 \otimes \eta}
\mathcal{G} \otimes_\mathcal{O} \mathcal{F}
\otimes_\mathcal{O} \mathcal{G}
\xrightarrow{\epsilon \otimes 1} \mathcal{G}
$$
が恒等写像であることを意味する。$U$ を $\mathcal{C}$ の対象とする。
$U$ をその被覆の各元で置き換えれば，$U$ 上の $\mathcal{F}$ と
$\mathcal{G}$ の有限個の切断 $f_1, \ldots, f_n$ と
$g_1, \ldots, g_n$ で，$\eta(1) = \sum f_i g_i$ を満たすものを取れる。
写像
$$
\mathcal{O}_U^{\oplus n} \to \mathcal{F}|_U
$$
を，第 $i$ 基底ベクトルを $f_i$ へ送る写像とする。このとき，写像
$\eta|_U$ は写像
$\tilde \eta : \mathcal{O}_U \to
\mathcal{O}_U^{\oplus n} \otimes_{\mathcal{O}_U} \mathcal{G}|_U$.
を経由して分解する。可換図式
$$
\xymatrix{
\mathcal{F}|_U
\ar[rr]_-{\eta \otimes 1} \ar[rrd]_-{\tilde \eta \otimes 1} & &
\mathcal{F}|_U \otimes \mathcal{G}|_U \otimes \mathcal{F}|_U
\ar[r]_-{1 \otimes \epsilon} &
\mathcal{F}|_U \\
& &
\mathcal{O}_U^{\oplus n} \otimes \mathcal{G}|_U \otimes \mathcal{F}|_U
\ar[u] \ar[r]^-{1 \otimes \epsilon} &
\mathcal{O}_U^{\oplus n} \ar[u]
}
$$
を得る。これにより，$\mathcal{F}|_U$ 上の恒等写像は有限自由
$\mathcal{O}_U$-加群を経由して分解することが分かる。これで (1) が示された。
(2) は《圏》の補題 \ref{categories-lemma-left-dual} とその証明から従う。
(3) は証明中の最初の等式から従う。また，左双対の一意性
（《圏》の注意 \ref{categories-remark-left-dual-adjoint}）と，
例 \ref{sites-modules-example-dual} における左双対の構成から (2) と (3) を導くこともできる。
\end{proof}

\begin{lemma}
\label{sites-modules-lemma-flat-locally-finite-presentation}
$(\mathcal{C}, \mathcal{O})$ を環付きサイトとする。
$\mathcal{F}$ が局所有限表示かつ平坦であるとする。このとき，
$\mathcal{C}$ の対象 $U$ が与えられると，被覆 $\{U_i \to U\}$ が存在して，
$\mathcal{F}|_{U_i}$ は有限自由 $\mathcal{O}_{U_i}$-加群の直和因子となる。
\end{lemma}

\begin{proof}
$\mathcal{C}$ の対象 $U$ を選ぶ。$U$ を被覆の各元で置き換えれば，表示
$$
\mathcal{O}_U^{\oplus r} \to
\mathcal{O}_U^{\oplus n} \to \mathcal{F}|_U \to 0
$$
が存在すると仮定してよい。補題 \ref{sites-modules-lemma-flat-eq} により，さらに $U$ を
被覆の各元で置き換えれば，分解
$$
\mathcal{O}_U^{\oplus n} \to
\mathcal{O}_U^{\oplus n_1} \to \mathcal{F}|_U
$$
が存在して，合成
$\mathcal{O}_U^{\oplus r} \to \mathcal{O}_U^{\oplus n}
\to \mathcal{O}_U^{\oplus n_r}$ が零であると仮定してよい。
これは全射 $\mathcal{O}_U^{\oplus n_r} \to \mathcal{F}|_U$ が切断をもつことを
意味し，結論が従う。
\end{proof}











\section{構成可能加群に向けて}
\label{sites-modules-section-constructible}

\noindent
サイトの準コンパクト対象とは，大まかには，その任意の被覆が有限被覆によって
細分されるような対象であったことを思い出そう（実際の定義はもう少し複雑である。
《サイト》の第 \ref{sites-section-quasi-compact} 節を参照）。サイトの各対象が
準コンパクト対象による被覆をもつならば，$U$ が準コンパクトである加群
$j_!\mathcal{O}_U$ は，すべての加群の圏に対して特に良い生成系をなす。

\begin{lemma}
\label{sites-modules-lemma-covering-gives-surjection}
$(\mathcal{C}, \mathcal{O})$ を環付きサイトとし，
$\{U_i \to U\}$ を $\mathcal{C}$ の被覆とする。このとき列
$$
\bigoplus j_{U_i \times_U U_j!}\mathcal{O}_{U_i \times_U U_j} \to
\bigoplus j_{U_i!}\mathcal{O}_{U_i} \to j_!\mathcal{O}_U \to 0
$$
は完全である。
\end{lemma}

\begin{proof}
任意の $\mathcal{O}$-加群 $\mathcal{F}$ に対し，関手
$\Hom_\mathcal{O}(-, \mathcal{F})$ はこの列を完全列
$0 \to \mathcal{F}(U) \to \prod \mathcal{F}(U_i) \to
\prod \mathcal{F}(U_i \times_U U_j)$ へ写す。
(\ref{sites-modules-equation-map-lower-shriek-OU-into-module}) を参照。これと
《ホモロジー》の補題 \ref{homology-lemma-check-exactness} から本補題が従う。
\end{proof}

\begin{lemma}
\label{sites-modules-lemma-silly-quasi-compact}
$(\mathcal{C}, \mathcal{O})$ を環付きサイトとする。
$\mathcal{U} = \{U_i \to U\}_{i \in I}$ を $\mathcal{C}$ の被覆とする。
$U$ が準コンパクトならば，有限部分集合 $I' \subset I$ が存在して，列
$$
\bigoplus\nolimits_{i, i' \in I'}
j_{U_i \times_U U_{i'}!}\mathcal{O}_{U_i \times_U U_{i'}} \to
\bigoplus\nolimits_{i \in I'}
j_{U_i!}\mathcal{O}_{U_i} \to
j_!\mathcal{O}_U \to 0
$$
は完全である。
\end{lemma}

\begin{proof}
補題 \ref{sites-modules-lemma-covering-gives-surjection} により，$U$ が《サイト》の補題
\ref{sites-lemma-conclude-quasi-compact} の条件 (3) を満たす場合，本補題は直ちに従う。
一般の場合の証明は読み飛ばすことを勧める。定義により，被覆
$\mathcal{V} = \{V_j \to U\}_{j \in J}$ と，固定した終域をもつ写像族の射
$\mathcal{V} \to \mathcal{U}$ が存在する。この射は
$\text{id} : U \to U$，$\alpha : J \to I$，および
$f_j : V_j \to U_{\alpha(j)}$ で与えられ
（《サイト》の定義 \ref{sites-definition-morphism-coverings} を参照），
$\alpha$ の像 $I' \subset I$ は有限である。
《ホモロジー》の補題 \ref{homology-lemma-check-exactness} により，任意の
$\mathcal{O}$-加群の層 $\mathcal{F}$ に対して関手
$\Hom_\mathcal{O}(-, \mathcal{F})$ が補題の列を完全列へ写すことを示せば十分である。
(\ref{sites-modules-equation-map-lower-shriek-OU-into-module}) により，通常の列
$$
0 \to
\mathcal{F}(U) \to
\prod\nolimits_{i \in I'} \mathcal{F}(U_i) \to
\prod\nolimits_{i, i' \in I'} \mathcal{F}(U_i \times_U U_{i'})
$$
を得る。被覆 $\mathcal{V}$ によって細分される写像族
$\{U_i \to U\}_{i \in I'}$ に《サイト》の補題
\ref{sites-lemma-compare-sheaf-condition} を適用すれば，これは完全列である。
\end{proof}

\begin{lemma}
\label{sites-modules-lemma-sections-over-quasi-compact}
$\mathcal{C}$ をサイトとし，$W$ を $\mathcal{C}$ の準コンパクト対象とする。
\begin{enumerate}
\item 関手 $\Sh(\mathcal{C}) \to \textit{Sets}$，
$\mathcal{F} \mapsto \mathcal{F}(W)$ は余積と可換である。
\item $\mathcal{O}$ を $\mathcal{C}$ 上の環の層とする。関手
$\textit{Mod}(\mathcal{O}) \to \textit{Ab}$,
$\mathcal{F} \mapsto \mathcal{F}(W)$
は直和と可換である。
\end{enumerate}
\end{lemma}

\begin{proof}
(1) を証明する。《サイト》の補題
\ref{sites-lemma-directed-colimits-sections} により，$W$ 上で切断をとる操作は，
単射な遷移写像をもつフィルター付き余極限と可換である。
$\mathcal{F}_i$ を集合 $I$ で添字付けられた集合の層の族とする。このとき
$\coprod \mathcal{F}_i$ は，有限部分集合 $E \subset I$ の半順序集合上の余積
$\mathcal{F}_E = \coprod_{i \in E} \mathcal{F}_i$ のフィルター付き余極限である。
遷移写像は単射なので結論が従う。

\medskip\noindent
(2) を証明する。$\mathcal{F}_i$ を集合 $I$ で添字付けられた
$\mathcal{O}$-加群の層の族とする。このとき $\bigoplus \mathcal{F}_i$ は，
有限部分集合 $E \subset I$ の半順序集合上の直和
$\mathcal{F}_E = \bigoplus_{i \in E} \mathcal{F}_i$ のフィルター付き余極限である。
アーベル群の層のフィルター付き余極限は集合の層の圏で計算できる。
さらに，$E \subset E'$ に対する遷移写像
$\mathcal{F}_E \to \mathcal{F}_{E'}$ は単射である
（層化は完全であり，台となる前層上では単射性が明らかだからである）。
したがって，《サイト》の補題 \ref{sites-lemma-directed-colimits-sections} により，
有限添字集合の場合を示せば十分である。有限の場合は補題
\ref{sites-modules-lemma-limits-colimits-abelian-sheaves} で扱った
（これは $\mathcal{C}$ の任意の対象上で成り立つ）。
\end{proof}

\begin{lemma}
\label{sites-modules-lemma-quasi-compact-hom-from}
$(\mathcal{C}, \mathcal{O})$ を環付きサイトとし，
$U$ を $\mathcal{C}$ の準コンパクト対象とする。このとき関手
$\Hom_\mathcal{O}(j_!\mathcal{O}_U, -)$ は直和と可換である。
\end{lemma}

\begin{proof}
これは
$\Hom_\mathcal{O}(j_!\mathcal{O}_U, \mathcal{F}) = \mathcal{F}(U)$
が (\ref{sites-modules-equation-map-lower-shriek-OU-into-module}) により成り立ち，
補題 \ref{sites-modules-lemma-sections-over-quasi-compact} により関手
$\mathcal{F} \mapsto \mathcal{F}(U)$ が直和と可換するからである。
\end{proof}

\noindent
以下で可能な限り鋭い結果を述べるため，いくつかの記法を導入する。

\begin{situation}
\label{sites-modules-situation-quasi-compact-objects}
$\mathcal{C}$ をサイトとし，
$\mathcal{B} \subset \text{Ob}(\mathcal{C})$ を対象の集合とする。
次の条件を考える：
\begin{enumerate}
\item
\label{sites-modules-item-enough}
$\mathcal{C}$ の各対象は $\mathcal{B}$ の元による被覆をもつ。
\item
\label{sites-modules-item-enough-qc}
すべての $U \in \mathcal{B}$ は準コンパクトである
（《サイト》の第 \ref{sites-section-quasi-compact} 節）。
\item
\label{sites-modules-item-enough-qc-qs}
$U_i, U \in \mathcal{B}$ である被覆 $\{U_i \to U\}$ に対し，
ファイバー積 $U_i \times_U U_j$ は準コンパクトである。
\end{enumerate}
\end{situation}

\begin{lemma}
\label{sites-modules-lemma-module-quotient-direct-sum}
状況 \ref{sites-modules-situation-quasi-compact-objects} において
(\ref{sites-modules-item-enough}) が成り立つと仮定する。
\begin{enumerate}
\item 任意の集合の層は，始域が $U_i \in \mathcal{B}$ である余積
$\coprod h_{U_i}^\#$ となる全射の終域である。
\item $\mathcal{O}$ が環の層ならば，任意の $\mathcal{O}$-加群は，
$U_i \in \mathcal{B}$ である直和
$\bigoplus\nolimits j_{U_i!}\mathcal{O}_{U_i}$ の商である。
\end{enumerate}
\end{lemma}

\begin{proof}
(1) は《サイト》の補題
\ref{sites-lemma-sheaf-coequalizer-representable} と
\ref{sites-lemma-covering-surjective-after-sheafification} から従う。
(2) は補題 \ref{sites-modules-lemma-module-quotient-flat} と
\ref{sites-modules-lemma-covering-gives-surjection} から従う。
\end{proof}

\begin{lemma}
\label{sites-modules-lemma-module-filtered-colimit-constructibles}
状況 \ref{sites-modules-situation-quasi-compact-objects} において
(\ref{sites-modules-item-enough}) と (\ref{sites-modules-item-enough-qc}) が成り立つと仮定する。
\begin{enumerate}
\item 任意の集合の層は，次の形の層のフィルター付き余極限である：
\begin{equation}
\label{sites-modules-equation-towards-constructible-sets}
\text{余等化子}\left(
\xymatrix{
\coprod\nolimits_{j = 1, \ldots, m} h_{V_j}^\#
\ar@<1ex>[r] \ar@<-1ex>[r] &
\coprod\nolimits_{i = 1, \ldots, n} h_{U_i}^\#
}
\right)
\end{equation}
ここで $U_i$ と $V_j$ は $\mathcal{B}$ に属する。
\item $\mathcal{O}$ が環の層ならば，任意の $\mathcal{O}$-加群は
次の形の層のフィルター付き余極限である：
\begin{equation}
\label{sites-modules-equation-towards-constructible}
\Coker\left(
\bigoplus\nolimits_{j = 1, \ldots, m} j_{V_j!}\mathcal{O}_{V_j}
\longrightarrow
\bigoplus\nolimits_{i = 1, \ldots, n} j_{U_i!}\mathcal{O}_{U_i}
\right)
\end{equation}
ここで $U_i$ と $V_j$ は $\mathcal{B}$ に属する。
\end{enumerate}
\end{lemma}

\begin{proof}
(1) を証明する。補題 \ref{sites-modules-lemma-module-quotient-direct-sum} により，
任意の集合の層 $\mathcal{F}$ は，$U \in \mathcal{B}$ である
$h_U^\#$ の形の層の余積 $\mathcal{F}_0$ を始域とする全射の終域である。
これを $\mathcal{F}_0 \times_\mathcal{F} \mathcal{F}_0$ に適用すると，
$\mathcal{F}$ は一対の写像
$$
\xymatrix{
\coprod\nolimits_{j \in J} h_{V_j}^\#
\ar@<1ex>[r] \ar@<-1ex>[r] &
\coprod\nolimits_{i \in I} h_{U_i}^\#
}
$$
の余等化子であることが分かる。ここで $I$，$J$ はある添字集合で，
$V_j$ と $U_i$ は $\mathcal{B}$ に属する。各有限部分集合 $J' \subset J$ に対し，
有限部分集合 $I' \subset I$ が存在して，$j \in J'$ 上の余積は，二つの写像の
いずれによっても $i \in I'$ 上の余積へ写る。《サイト》の補題
\ref{sites-lemma-directed-colimits-sections} を参照。
（細部は省略する。ヒント：無限余積は有限部分余積のフィルター付き余極限である。）
したがって，この層は，有限余積間のこれらの写像の余核の余極限である。

\medskip\noindent
(2) を証明する。補題 \ref{sites-modules-lemma-module-quotient-direct-sum} により，任意の加群は，
$U \in \mathcal{B}$ である $j_{U!}\mathcal{O}_U$ の形の加群の直和の商である。
したがって，任意の加群は余核
$$
\Coker\left(
\bigoplus\nolimits_{j \in J} j_{V_j!}\mathcal{O}_{V_j}
\longrightarrow
\bigoplus\nolimits_{i \in I} j_{U_i!}\mathcal{O}_{U_i}
\right)
$$
として書ける。ここで $I$，$J$ はある添字集合で，$V_j$ と $U_i$ は
$\mathcal{B}$ に属する。各有限部分集合 $J' \subset J$ に対し，有限部分集合
$I' \subset I$ が存在して，$j \in J'$ 上の直和は $i \in I'$ 上の直和へ写る。
補題 \ref{sites-modules-lemma-quasi-compact-hom-from} を参照。
したがって，この加群は，有限直和間のこれらの写像の余核の余極限である。
\end{proof}

\begin{lemma}
\label{sites-modules-lemma-cokernel-map-towards-constructibles}
状況 \ref{sites-modules-situation-quasi-compact-objects} において
(\ref{sites-modules-item-enough}) と (\ref{sites-modules-item-enough-qc}) が成り立つと仮定し，
$\mathcal{O}$ を環の層とする。このとき
(\ref{sites-modules-equation-towards-constructible}) の形の加群間の写像の余核も，
(\ref{sites-modules-equation-towards-constructible}) の形の加群である。
\end{lemma}

\begin{proof}
$\mathcal{F} = \Coker(\bigoplus j_{V_j!}\mathcal{O}_{V_j} \to
\bigoplus j_{U_i!}\mathcal{O}_{U_i})$
を (\ref{sites-modules-equation-towards-constructible}) のとおりとする。
$W \in \mathcal{B}$ である写像
$\varphi : j_{W!}\mathcal{O}_W \to \mathcal{F}$ の余核も同じ形の加群であることを
示せば十分である。写像 $\varphi$ は $s \in \mathcal{F}(W)$ に対応する。
$\bigoplus j_{U_i!}\mathcal{O}_{U_i} \to \mathcal{F}$ は全射なので，
(\ref{sites-modules-item-enough}) により，$W_k \in \mathcal{B}$ である被覆
$\{W_k \to W\}_{k \in K}$ を，$s|_{W_k}$ が
$\bigoplus j_{U_i!}\mathcal{O}_{U_i})$ のある切断 $s_k$ の像となるように選べる。
(\ref{sites-modules-item-enough-qc}) により対象 $W$ は準コンパクトである。
補題 \ref{sites-modules-lemma-silly-quasi-compact} により，有限部分集合
$K' \subset K$ が存在して，
$\bigoplus_{k \in K'} j_{W_k!}\mathcal{O}_{W_k} \to j_{W!}\mathcal{O}_W$
は全射である。したがって，$\Coker(\varphi)$ は
$$
\Coker\left(
\bigoplus\nolimits_{k \in K'} j_{W_k!}\mathcal{O}_{W_k} \oplus
\bigoplus j_{V_j!}\mathcal{O}_{V_j}
\longrightarrow
\bigoplus j_{U_i!}\mathcal{O}_{U_i}
\right)
$$
に等しい。ここで写像
$\bigoplus_{k \in K'} j_{W_k!}\mathcal{O}_{W_k}
\to \bigoplus j_{U_i!}\mathcal{O}_{U_i}$ は
$\sum_{k \in K'} s_k$ に対応する。これで証明が完了する。
\end{proof}

\begin{lemma}
\label{sites-modules-lemma-change-presentation-towards-constructibles}
状況 \ref{sites-modules-situation-quasi-compact-objects} において
(\ref{sites-modules-item-enough})，(\ref{sites-modules-item-enough-qc})，
(\ref{sites-modules-item-enough-qc-qs}) が成り立つと仮定する。
$\mathcal{O}$ を環の層とし，写像
$$
\bigoplus\nolimits_{j = 1, \ldots, m} j_{V_j!}\mathcal{O}_{V_j}
\longrightarrow
\bigoplus\nolimits_{i = 1, \ldots, n} j_{U_i!}\mathcal{O}_{U_i}
$$
が与えられたとする。ここで $U_i$ と $V_j$ は $\mathcal{B}$ に属する。
さらに，$U_{ik} \in \mathcal{B}$ である被覆
$\{U_{ik} \to U_i\}_{k \in K_i}$ が与えられているとする。
このとき，有限部分集合 $K'_i \subset K_i$，$W_l \in \mathcal{B}$ からなる
有限集合 $L$，および可換図式
$$
\xymatrix{
\bigoplus_{l \in L} j_{W_l!}\mathcal{O}_{W_l} \ar[d] \ar[r] &
\bigoplus_{i = 1, \ldots, n} \bigoplus_{k \in K'_i}
j_{U_{ik}!}\mathcal{O}_{U_{ik}} \ar[d] \\
\bigoplus_{j = 1, \ldots, m} j_{V_j!}\mathcal{O}_{V_j} \ar[r] &
\bigoplus_{i = 1, \ldots, n} j_{U_i!}\mathcal{O}_{U_i}
}
$$
が存在し，これは水平写像の余核の間に同型を誘導する。
\end{lemma}

\begin{proof}
$U_i$ は準コンパクトなので，補題 \ref{sites-modules-lemma-silly-quasi-compact} のように
有限部分集合 $K'_i \subset K_i$ を選べる。さらに，
$\bigoplus_{i = 1, \ldots, n}
\bigoplus_{k \in K'_i} j_{U_{ik}!}\mathcal{O}_{U_{ik}} \to
\bigoplus_{i = 1, \ldots, n} j_{U_i!}\mathcal{O}_{U_i}$ は全射なので，
$V_{jm} \in \mathcal{B}$ である被覆 $\{V_{jm} \to V_j\}_{m \in M_j}$ を選び，
可換図式
$$
\xymatrix{
\bigoplus_{j = 1, \ldots, m} \bigoplus_{m \in M_j}
j_{V_{jm}!}\mathcal{O}_{V_{jm}} \ar[d] \ar[r] &
\bigoplus_{i = 1, \ldots n} \bigoplus_{k \in K'_i}
j_{U_{ik}!}\mathcal{O}_{U_{ik}} \ar[d] \\
\bigoplus_{j = 1, \ldots, m} j_{V_j!}\mathcal{O}_{V_j} \ar[r] &
\bigoplus_{i = 1, \ldots, n} j_{U_i!}\mathcal{O}_{U_i}
}
$$
を得ることができる。$V_j$ は準コンパクトなので，補題
\ref{sites-modules-lemma-silly-quasi-compact} のように有限部分集合
$M'_j \subset M_j$ を選べる。さらに
$$
L = \left(\coprod\nolimits_{i = 1, \ldots, n} K'_i \times K'_i \right)
\coprod
\left(\coprod\nolimits_{j = 1, \ldots, m} M'_j\right)
$$
とおく。$l = (k, k') \in K'_i \times K'_i \subset L$ に対して
$W_l = U_{ik} \times_{U_i} U_{ik'}$ とおき，
$l = m \in M'_j \subset L$ に対して $W_l = V_{jm}$ とおく。
族 $\{U_{ik} \to U_i\}_{k \in K'_i}$ に対して補題
\ref{sites-modules-lemma-silly-quasi-compact} の完全列があるので，補題の主張にある図式を
得る（細部は省略する）。ただし，まだ $W_l \in \mathcal{B}$ であることは
明らかでない。しかし，すべての $l \in L$ に対して $W_l$ は準コンパクトなので，
補題 \ref{sites-modules-lemma-silly-quasi-compact} をもう一度適用し，
$W_{lt} \in \mathcal{B}$ である有限写像族
$\{W_{lt} \to W_l\}_{t \in T_l}$ を，
$\bigoplus j_{W_{lt}!}\mathcal{O}_{W_{lt}} \to j_{W_l!}\mathcal{O}_{W_l}$
が全射となるように取る。そして $L$ を $\coprod_{l \in L} T_l$ で置き換えれば，
すべて明らかである。
\end{proof}

\begin{lemma}
\label{sites-modules-lemma-extension-towards-constructibles}
状況 \ref{sites-modules-situation-quasi-compact-objects} において
(\ref{sites-modules-item-enough})，(\ref{sites-modules-item-enough-qc})，
(\ref{sites-modules-item-enough-qc-qs}) が成り立つと仮定し，$\mathcal{O}$ を環の層とする。
このとき，(\ref{sites-modules-equation-towards-constructible}) の形の加群の拡大も，
(\ref{sites-modules-equation-towards-constructible}) の形の加群である。
\end{lemma}

\begin{proof}
$0 \to \mathcal{F}_1 \to \mathcal{F}_2 \to \mathcal{F}_3 \to 0$ を
$\mathcal{O}$-加群の短完全列とし，$\mathcal{F}_1$ と $\mathcal{F}_3$ は
(\ref{sites-modules-equation-towards-constructible}) の形であるとする。表示
$$
\bigoplus A_{V_j} \to \bigoplus A_{U_i} \to \mathcal{F}_1 \to 0
\quad\text{かつ}\quad
\bigoplus A_{T_j} \to \bigoplus A_{W_i} \to \mathcal{F}_3 \to 0
$$
を選ぶ。この証明では直和は常に有限であり，$U \in \mathcal{B}$ に対して
$A_U = j_{U!}\mathcal{O}_U$ と書く。
$\mathcal{F}_2 \to \mathcal{F}_3$ は全射なので，$W_{ik} \in \mathcal{B}$ である
被覆 $\{W_{ik} \to W_i\}$ を，$A_{W_{ik}} \to \mathcal{F}_3$ が写像
$A_{W_{ik}} \to \mathcal{F}_2$ へ持ち上がるように選べる。
補題 \ref{sites-modules-lemma-change-presentation-towards-constructibles} により，族 $\{W_i\}$ を
族 $\{W_{ik}\}$ の有限部分族で置き換え，写像
$\bigoplus A_{W_i} \to \mathcal{F}_3$ が $\mathcal{F}_2$ への写像へ持ち上がると
仮定してよい。核
$$
\mathcal{K}_2 = \Ker(\bigoplus A_{U_i} \oplus \bigoplus A_{W_i}
\longrightarrow \mathcal{F}_2)
$$
を考える。蛇の補題により，この核は
$\mathcal{K}_3 = \Ker(\bigoplus A_{W_i} \to \mathcal{F}_3)$ へ全射する。
したがって，上と同様に論じ，各 $T_j$ を $\mathcal{B}$ の元からなる有限族で
置き換えれば（補題 \ref{sites-modules-lemma-silly-quasi-compact} により許される），
与えられた写像 $\bigoplus A_{T_j} \to \mathcal{K}_3$ を持ち上げる写像
$\bigoplus A_{T_j} \to \mathcal{K}_2$ が存在すると仮定してよい。
このとき $\bigoplus A_{V_j} \oplus \bigoplus A_{T_j} \to \mathcal{K}_2$ は
全射であり，証明が完了する。
\end{proof}

\begin{lemma}
\label{sites-modules-lemma-towards-constructible-when-serre-subcategory}
状況 \ref{sites-modules-situation-quasi-compact-objects} において
(\ref{sites-modules-item-enough})，(\ref{sites-modules-item-enough-qc})，
(\ref{sites-modules-item-enough-qc-qs}) が成り立つと仮定し，$\mathcal{O}$ を環の層とする。
$\mathcal{A} \subset \textit{Mod}(\mathcal{O})$ を，
(\ref{sites-modules-equation-towards-constructible}) の形の余核と同型な加群からなる充満部分圏とする。
もし，$U_i$ と $V_j$ が $\mathcal{B}$ に属する，次の形の任意の
$\mathcal{O}$-加群の写像
$$
\bigoplus\nolimits_{j = 1, \ldots, m} j_{V_j!}\mathcal{O}_{V_j}
\longrightarrow
\bigoplus\nolimits_{i = 1, \ldots, n} j_{U_i!}\mathcal{O}_{U_i}
$$
の核が $\mathcal{A}$ に属するならば，$\mathcal{A}$ は
$\textit{Mod}(\mathcal{O})$ の弱 Serre 部分圏である。
\end{lemma}

\begin{proof}
《ホモロジー》の補題 \ref{homology-lemma-characterize-weak-serre-subcategory} の
判定法を用いる。補題 \ref{sites-modules-lemma-cokernel-map-towards-constructibles} と
\ref{sites-modules-lemma-extension-towards-constructibles} の結果により，$\mathcal{A}$ の対象間の
写像 $\mathcal{F} \to \mathcal{G}$ の核が $\mathcal{A}$ に属することを示せば十分である。
これを示すため，表示
$$
\bigoplus A_{V_j} \to \bigoplus A_{U_i} \to \mathcal{F} \to 0
\quad\text{かつ}\quad
\bigoplus A_{T_j} \to \bigoplus A_{W_i} \to \mathcal{G} \to 0
$$
を選ぶ。この証明では直和は常に有限であり，$U \in \mathcal{B}$ に対して
$A_U = j_{U!}\mathcal{O}_U$ と書く。補題
\ref{sites-modules-lemma-covering-gives-surjection} と
\ref{sites-modules-lemma-change-presentation-towards-constructibles} を用い，補題
\ref{sites-modules-lemma-extension-towards-constructibles} の証明と同様に論じれば，写像
$\mathcal{F} \to \mathcal{G}$ は表示の写像
$$
\xymatrix{
\bigoplus A_{V_j} \ar[r] \ar[d] &
\bigoplus A_{U_i} \ar[r] \ar[d] &
\mathcal{F} \ar[r] \ar[d] & 0 \\
\bigoplus A_{T_j} \ar[r] &
\bigoplus A_{W_i} \ar[r] &
\mathcal{G} \ar[r] & 0
}
$$
へ持ち上がると仮定してよい。このとき
$$
\Ker(\mathcal{F} \to \mathcal{G}) =
\Coker\left(\bigoplus A_{V_j} \to
\Ker\left(
\bigoplus A_{T_j} \oplus \bigoplus A_{U_i} \to \bigoplus A_{W_i}\right)\right)
$$
であることが分かり，本補題は仮定と補題
\ref{sites-modules-lemma-cokernel-map-towards-constructibles} から従う。
\end{proof}








\section{平坦射}
\label{sites-modules-section-flat-morphisms}

\begin{definition}
\label{sites-modules-definition-flat-morphism}
環付きトポスの射
$(f, f^\sharp) :
(\Sh(\mathcal{C}), \mathcal{O})
\longrightarrow
(\Sh(\mathcal{C}'), \mathcal{O}')$
が {\it 平坦}であるとは，環写像
$f^\sharp : f^{-1}\mathcal{O}' \to \mathcal{O}$ が平坦であることをいう。
環付きサイトの射が {\it 平坦}であるとは，付随する環付きトポスの射が
平坦であることをいう。
\end{definition}

\begin{lemma}
\label{sites-modules-lemma-flat-pullback-exact}
$f : \Sh(\mathcal{C}) \to \Sh(\mathcal{C}')$ を環付きトポスの射とする。このとき
$$
f^{-1} : \textit{Ab}(\mathcal{C}') \longrightarrow \textit{Ab}(\mathcal{C}),
\quad
\mathcal{F} \longmapsto f^{-1}\mathcal{F}
$$
は完全である。また，
$(f, f^\sharp) :
(\Sh(\mathcal{C}), \mathcal{O})
\to
(\Sh(\mathcal{C}'), \mathcal{O}')$
が環付きトポスの平坦射ならば，
$$
f^* : \textit{Mod}(\mathcal{O}') \longrightarrow \textit{Mod}(\mathcal{O}),
\quad
\mathcal{F} \longmapsto f^*\mathcal{F}
$$
は完全である。
\end{lemma}

\begin{proof}
$\mathcal{C}'$ 上のアーベル群の層 $\mathcal{G}$ が与えられたとき，
$f^{-1}\mathcal{G}$ の台となる集合の層は，$\mathcal{G}$ の台となる集合の層の
$f^{-1}$ と同じである。《サイト》の第
\ref{sites-section-sheaves-algebraic-structures} 節を参照。
したがって，集合の層に対する $f^{-1}$ の完全性
（トポスの射の定義で要求される。《サイト》の定義
\ref{sites-definition-topos} を参照）から，アーベル群の層上の関手としての
$f^{-1}$ の完全性が従う。

\medskip\noindent
加群に関する主張を見るため，$f^*\mathcal{F}$ はテンソル積
$f^{-1}\mathcal{F} \otimes_{f^{-1}\mathcal{O}', f^\sharp} \mathcal{O}$ として
定義されたことを思い出そう。したがって，$f^*$ は二つの完全関手の合成である。
\end{proof}

\begin{definition}
\label{sites-modules-definition-flat-module}
$f : (\Sh(\mathcal{C}), \mathcal{O}) \to (\Sh(\mathcal{D}), \mathcal{O}')$
を環付きトポスの射とし，$\mathcal{F}$ を $\mathcal{O}$-加群の層とする。
$\mathcal{F}$ が {\it $(\Sh(\mathcal{D}), \mathcal{O}')$ 上平坦}であるとは，
$\mathcal{F}$ が $f^{-1}\mathcal{O}'$-加群として平坦であることをいう。
\end{definition}

\noindent
これは環付き空間の射に対して定義された概念と両立する。
《加群》の定義 \ref{modules-definition-flat-module} とそれに続く議論を参照。

\begin{lemma}
\label{sites-modules-lemma-pullback-internal-hom}
$f : (\mathcal{C}, \mathcal{O}_\mathcal{C}) \to
(\mathcal{D}, \mathcal{O}_\mathcal{D})$ を環付きサイトの射とし，
$\mathcal{F}$，$\mathcal{G}$ を $\mathcal{O}_\mathcal{D}$-加群とする。
$\mathcal{F}$ が有限表示で $f$ が平坦ならば，注意
\ref{sites-modules-remark-pullback-internal-hom} の標準写像
$$
f^*\SheafHom_{\mathcal{O}_\mathcal{D}}(\mathcal{F}, \mathcal{G})
\longrightarrow
\SheafHom_{\mathcal{O}_\mathcal{C}}(f^*\mathcal{F}, f^*\mathcal{G})
$$
は同型である。
\end{lemma}

\begin{proof}
$f$ は連続関手 $u : \mathcal{D} \to \mathcal{C}$ によって与えられるとする。
任意の $U \in \Ob(\mathcal{C})$ に対し，写像の $\mathcal{C}/U$ への制限が
同型であることを示せばよい。$U$ を $U$ の被覆の各元で置き換えてよい。
したがって，《サイト》の補題 \ref{sites-lemma-morphism-of-sites-covering} により，
ある $V \in \Ob(\mathcal{C})$ に対する射 $U \to u(V)$ が存在すると仮定してよい。
すると，もちろん $U$ を $u(V)$ で置き換えてよい。さらに $u$ は連続なので，
$V$ を被覆で置き換え，表示
$\mathcal{O}_V^{\oplus m} \to \mathcal{O}_V^{\oplus n} \to
\mathcal{F}|_V \to 0$ が $\mathcal{D}/V$ 上で存在すると仮定してよい。
$\SheafHom$ の構成は局所化と可換するので
（補題 \ref{sites-modules-lemma-internal-hom-restriction}），$f$ を，$f$ が誘導する射
$(\mathcal{C}/u(V), \mathcal{O}_{u(V)}) \to
(\mathcal{D}/V, \mathcal{O}_V)$ で置き換えてよい。
したがって，$\mathcal{F}$ が大域表示
$\mathcal{O}_\mathcal{D}^{\oplus m} \to \mathcal{O}_\mathcal{D}^{\oplus n} \to
\mathcal{F} \to 0$ をもつ場合に帰着する。
$f$ は平坦で $f^*\mathcal{O}_\mathcal{D} = \mathcal{O}_\mathcal{C}$ なので，
対応する表示
$\mathcal{O}_\mathcal{C}^{\oplus m} \to \mathcal{O}_\mathcal{C}^{\oplus n} \to
f^*\mathcal{F} \to 0$ を得る。補題 \ref{sites-modules-lemma-flat-pullback-exact} を参照。
$\SheafHom$ が第 1 変数の有限直和と可換すること，
$\SheafHom_{\mathcal{O}_\mathcal{C}}(\mathcal{O}_\mathcal{C}, -)$ と
$\SheafHom_{\mathcal{O}_\mathcal{D}}(\mathcal{O}_\mathcal{D}, -)$
がともに恒等関手であること，および注意
\ref{sites-modules-remark-pullback-internal-hom} の構成の関手性を用いると，可換図式
$$
\xymatrix{
0 \ar[r] &
f^*\SheafHom_{\mathcal{O}_\mathcal{D}}(\mathcal{F}, \mathcal{G})
\ar[d] \ar[r] &
f^*\mathcal{G}^{\oplus n} \ar[d] \ar[r] &
f^*\mathcal{G}^{\oplus n} \ar[d] \\
0 \ar[r] &
\SheafHom_{\mathcal{O}_\mathcal{C}}(f^*\mathcal{F}, f^*\mathcal{G})
\ar[r] &
f^*\mathcal{G}^{\oplus n} \ar[r] &
f^*\mathcal{G}^{\oplus n}
}
$$
を得る。ここで右側の二つの垂直矢印は同型である。
補題 \ref{sites-modules-lemma-internal-hom-exact} により行は完全である。
五項補題から結論が従う。
\end{proof}







\section{可逆加群}
\label{sites-modules-section-invertible}

\noindent
定義は次のとおりである。

\begin{definition}
\label{sites-modules-definition-invertible-sheaf}
$(\mathcal{C}, \mathcal{O})$ を環付きサイトとする。
\begin{enumerate}
\item 有限局所自由 $\mathcal{O}$-加群 $\mathcal{F}$ が {\it 階数 $r$} を
もつとは，$\mathcal{C}$ の各対象 $U$ に対して $U$ の被覆
$\{U_i \to U\}$ が存在し，$\mathcal{F}|_{U_i}$ が
$\mathcal{O}_{U_i}$-加群として $\mathcal{O}_{U_i}^{\oplus r}$ と同型であることをいう。
\item $\mathcal{O}$-加群 $\mathcal{L}$ が {\it 可逆}であるとは，関手
$$
\textit{Mod}(\mathcal{O}) \longrightarrow \textit{Mod}(\mathcal{O}),\quad
\mathcal{F}  \longmapsto \mathcal{F} \otimes_\mathcal{O} \mathcal{L}
$$
が圏同値であることをいう。
\item 層 {\it $\mathcal{O}^*$} とは，対応
$$
U \longmapsto \mathcal{O}^*(U) = \{f \in \mathcal{O}(U) \mid
\exists g \in \mathcal{O}(U)\text{ を満たす }fg = 1\}
$$
によって定まる $\mathcal{O}$ の部分層である。これは乗法を群演算とする
アーベル群の層である。
\end{enumerate}
\end{definition}

\noindent
下の補題 \ref{sites-modules-lemma-invertible-is-locally-free-rank-1} で，階数 $1$ の
局所自由加群との関係を説明する。

\begin{lemma}
\label{sites-modules-lemma-invertible}
$(\mathcal{C}, \mathcal{O})$ を環付きサイトとし，
$\mathcal{L}$ を $\mathcal{O}$-加群とする。次は同値である：
\begin{enumerate}
\item $\mathcal{L}$ は可逆である；
\item $\mathcal{O}$-加群 $\mathcal{N}$ が存在して，
$\mathcal{L} \otimes_\mathcal{O} \mathcal{N} \cong \mathcal{O}$.
\end{enumerate}
この場合，次が成り立つ：
\begin{enumerate}
\item[(a)] $\mathcal{L}$ は有限表示の平坦な $\mathcal{O}$-加群である；
\item[(b)] $\mathcal{C}$ の各対象 $U$ に対して被覆
$\{U_i \to U\}$ が存在し，$\mathcal{L}|_{U_i}$ は有限自由加群の直和因子となる；
\item[(c)] (2) の加群 $\mathcal{N}$ は
$\SheafHom_\mathcal{O}(\mathcal{L}, \mathcal{O})$ と同型である。
\end{enumerate}
\end{lemma}

\begin{proof}
(1) を仮定する。このとき関手 $- \otimes_\mathcal{O} \mathcal{L}$ は
本質的全射なので，(2) のような $\mathcal{O}$-加群 $\mathcal{N}$ が存在する。
(2) が成り立つならば，関手 $- \otimes_\mathcal{O} \mathcal{N}$ は関手
$- \otimes_\mathcal{O} \mathcal{L}$ の擬逆であり，(1) が成り立つ。

\medskip\noindent
(1) と (2) が成り立つと仮定する。$- \otimes_\mathcal{O} \mathcal{L}$ は
圏同値なので完全であり，したがって $\mathcal{L}$ は平坦である。
与えられた同型を
$\psi : \mathcal{L} \otimes_\mathcal{O} \mathcal{N} \to \mathcal{O}$ と書く。
$U$ を $\mathcal{C}$ の対象とする。$U$ のある被覆の各元への
$\mathcal{L}$ の制限が自由加群の直和因子であることを示す。これは確かに
$\mathcal{L}$ が有限表示であることを含意する。$\otimes$ の構成により，
$U$ を被覆の各元で置き換えれば，整数 $n \geq 1$ と切断
$x_i \in \mathcal{L}(U)$，$y_i \in \mathcal{N}(U)$ で
$\psi(\sum x_i \otimes y_i) = 1$ を満たすものが存在すると仮定してよい。
同型
$$
\mathcal{L}|_U \to
\mathcal{L}|_U \otimes_{\mathcal{O}_U}
\mathcal{L}|_U \otimes_{\mathcal{O}_U} \mathcal{N}|_U \to \mathcal{L}|_U
$$
を考える。最初の矢印は $x$ を $\sum x_i \otimes x \otimes y_i$ へ送り，
二つ目の矢印は $x \otimes x' \otimes y$ を $\psi(x' \otimes y)x$ へ送る。
したがって，$x \mapsto \sum \psi(x \otimes y_i)x_i$ は
$\mathcal{L}|_U$ の自己同型である。この自己同型は
$$
\mathcal{L}|_U \to \mathcal{O}_U^{\oplus n} \to \mathcal{L}|_U
$$
と分解する。ここで最初の矢印は
$x \mapsto (\psi(x \otimes y_1), \ldots, \psi(x \otimes y_n))$，
二つ目の矢印は $(a_1, \ldots, a_n) \mapsto \sum a_i x_i$ で与えられる。
このようにして，$\mathcal{L}|_U$ は有限自由
$\mathcal{O}_U$-加群の直和因子であることが分かる。

\medskip\noindent
(1) と (2) が成り立つと仮定する。評価写像
$$
\mathcal{L} \otimes_\mathcal{O}
\SheafHom_\mathcal{O}(\mathcal{L}, \mathcal{O}_X)
\longrightarrow \mathcal{O}_X
$$
を考える。本補題の証明を終えるため，これが同型であることを示す。
補題 \ref{sites-modules-lemma-internal-hom-adjoint-tensor} により，
$$
\Hom_\mathcal{O}(\mathcal{O}, \mathcal{O}) =
\Hom_\mathcal{O}
(\mathcal{N} \otimes_\mathcal{O} \mathcal{L}, \mathcal{O})
\longrightarrow
\Hom_\mathcal{O}
(\mathcal{N}, \SheafHom_\mathcal{O}(\mathcal{L}, \mathcal{O}))
$$
$1$ の像から射
$\mathcal{N} \to \SheafHom_\mathcal{O}(\mathcal{L}, \mathcal{O})$ を得る。
$\mathcal{L}$ とのテンソル積をとると，
$$
\mathcal{O} = \mathcal{L} \otimes_\mathcal{O} \mathcal{N}
\longrightarrow
\mathcal{L} \otimes_\mathcal{O} \SheafHom_\mathcal{O}(\mathcal{L}, \mathcal{O})
$$
を得る。この写像は評価写像の逆である。計算は省略する。
\end{proof}

\begin{lemma}
\label{sites-modules-lemma-pullback-invertible}
$f : (\Sh(\mathcal{C}), \mathcal{O}_\mathcal{C}) \to
(\Sh(\mathcal{D}), \mathcal{O}_\mathcal{D})$ を環付きトポスの射とする。
可逆 $\mathcal{O}_\mathcal{D}$-加群の引き戻し $f^*\mathcal{L}$ は可逆である。
\end{lemma}

\begin{proof}
補題 \ref{sites-modules-lemma-invertible} により，$\mathcal{O}_\mathcal{D}$-加群
$\mathcal{N}$ が存在して，
$\mathcal{L} \otimes_{\mathcal{O}_\mathcal{D}} \mathcal{N} \cong
\mathcal{O}_\mathcal{D}$ を満たす。引き戻すと，
$f^*\mathcal{L} \otimes_{\mathcal{O}_\mathcal{C}} f^*\mathcal{N} \cong
\mathcal{O}_\mathcal{C}$
を補題 \ref{sites-modules-lemma-tensor-product-pullback} により得る。
したがって，補題 \ref{sites-modules-lemma-invertible} により $f^*\mathcal{L}$ は可逆である。
\end{proof}

\begin{lemma}
\label{sites-modules-lemma-constructions-invertible}
$(\mathcal{C}, \mathcal{O})$ を環付きサイトとする。
\begin{enumerate}
\item $\mathcal{L}$，$\mathcal{N}$ が可逆 $\mathcal{O}$-加群ならば，
$\mathcal{L} \otimes_\mathcal{O} \mathcal{N}$ も可逆である。
\item $\mathcal{L}$ が可逆 $\mathcal{O}$-加群ならば，
$\SheafHom_\mathcal{O}(\mathcal{L}, \mathcal{O})$ も可逆であり，評価写像
$\mathcal{L} \otimes_\mathcal{O}
\SheafHom_\mathcal{O}(\mathcal{L}, \mathcal{O}) \to \mathcal{O}$
は同型である。
\end{enumerate}
\end{lemma}

\begin{proof}
(1) は定義から明らかであり，(2) は補題 \ref{sites-modules-lemma-invertible} と
その証明から従う。
\end{proof}

\begin{lemma}
\label{sites-modules-lemma-pic-set}
$(\mathcal{C}, \mathcal{O})$ を環付きサイトとする。
可逆加群の集合 $\{\mathcal{L}_i\}_{i \in I}$ で，
$(\mathcal{C}, \mathcal{O})$ 上の各可逆加群がちょうど一つの
$\mathcal{L}_i$ と同型になるものが存在する。
\end{lemma}

\begin{proof}
省略する。《加群の層》の補題 \ref{modules-lemma-pic-set} を参照。
\end{proof}

\noindent
補題 \ref{sites-modules-lemma-pic-set} は，可逆層の同型類全体が集合をなすことを述べている。
補題 \ref{sites-modules-lemma-constructions-invertible} は，テンソル積がこの集合に
アーベル群の構造を定め，$\mathcal{L}$ の逆元が
$\SheafHom_\mathcal{O}(\mathcal{L}, \mathcal{O})$ で与えられることを述べている。

\medskip\noindent
実際，可逆 $\mathcal{O}$-加群 $\mathcal{L}$ と $n \in \mathbf{Z}$ が
与えられたとき，$\mathcal{L}$ の第 $n$ {\it テンソル冪}
$\mathcal{L}^{\otimes n}$ を，圏同値
$\mathcal{F} \mapsto \mathcal{F} \otimes_\mathcal{O} \mathcal{L}$ を
ちょうど $n$ 回適用したときの $\mathcal{O}$ の像として定義する。
可逆 $\mathcal{O}$-加群を，それとのテンソル積をとる操作が圏同値となるものとして
定義したので，これは負の $n$ に対しても意味をもつ。より明示的には，
$$
\mathcal{L}^{\otimes n} =
\left\{
\begin{matrix}
\mathcal{O} & & n = 0\text{ のとき} \\
\SheafHom_\mathcal{O}(\mathcal{L}, \mathcal{O}) & & n = -1\text{ のとき}\\
\mathcal{L} \otimes_\mathcal{O} \ldots \otimes_\mathcal{O} \mathcal{L}
& & n > 0\text{ のとき} \\
\mathcal{L}^{\otimes -1} \otimes_\mathcal{O} \ldots
\otimes_\mathcal{O} \mathcal{L}^{\otimes -1}
& & n < -1\text{ のとき}
\end{matrix}
\right.
$$
である。補題 \ref{sites-modules-lemma-constructions-invertible} を参照。この定義のもとで標準同型
$\mathcal{L}^{\otimes n} \otimes_\mathcal{O}
\mathcal{L}^{\otimes m} \to
\mathcal{L}^{\otimes n + m}$ があり，これらの同型は可換性制約と結合性制約を
満たす（定式化は省略する）。

\begin{definition}
\label{sites-modules-definition-pic}
$(\mathcal{C}, \mathcal{O})$ を環付きサイトとする。
この環付きサイトの {\it Picard 群} $\Pic(\mathcal{O})$ とは，
可逆 $\mathcal{O}$-加群の同型類を元とし，テンソル積に対応する演算を加法とする
アーベル群である。
\end{definition}









\section{微分加群}
\label{sites-modules-section-differentials}

\noindent
本節では，環付きトポスの射に対する相対微分加群をどのように定義するかを
簡単に説明する。読者には，可換環論の章にある対応する節
（《代数》，\ref{algebra-section-differentials} 節）も参照することを勧める。

\begin{definition}
\label{sites-modules-definition-derivation}
$\mathcal{C}$ をサイトとする。$\varphi : \mathcal{O}_1 \to \mathcal{O}_2$
を環の層の準同型とし，$\mathcal{F}$ を $\mathcal{O}_2$-加群とする。
$\mathcal{F}$ に値をとる {\it $\mathcal{O}_1$-導分}，より正確には
{\it $\varphi$-導分}とは，加法的で，$\mathcal{O}_1 \to \mathcal{O}_2$
の像を零化し，{\it Leibniz 則}
$$
D(ab) = aD(b) + D(a)b
$$
を満たす写像 $D : \mathcal{O}_2 \to \mathcal{F}$ をいう。ここで $a,b$ は，
ともに定義されている任意の場所における $\mathcal{O}_2$ の局所切断である。
$\mathcal{F}$ に値をとる $\varphi$-導分全体の集合を
$\text{Der}_{\mathcal{O}_1}(\mathcal{O}_2, \mathcal{F})$
と書く。
\end{definition}

\noindent
これは《代数》の定義 \ref{sites-modules-definition-derivation} の層論的類似である。
定義のような導分 $D : \mathcal{O}_2 \to \mathcal{F}$ が与えられると，
大域切断上の写像
$$
D : \Gamma(\mathcal{O}_2) \longrightarrow \Gamma(\mathcal{F})
$$
は明らかに，代数の定義における $\Gamma(\mathcal{O}_1)$-導分である。
また，$\alpha : \mathcal{F} \to \mathcal{G}$ が
$\mathcal{O}_2$-加群の写像ならば，
$$
\text{Der}_{\mathcal{O}_1}(\mathcal{O}_2, \mathcal{F})
\longrightarrow
\text{Der}_{\mathcal{O}_1}(\mathcal{O}_2, \mathcal{G})
$$
が $D \mapsto \alpha \circ D$ によって誘導される。すなわち，関手を得る。

\begin{lemma}
\label{sites-modules-lemma-universal-module}
$\mathcal{C}$ をサイトとし，$\varphi : \mathcal{O}_1 \to \mathcal{O}_2$
を環の層の準同型とする。関手
$$
\textit{Mod}(\mathcal{O}_2) \longrightarrow \textit{Ab}, \quad
\mathcal{F} \longmapsto \text{Der}_{\mathcal{O}_1}(\mathcal{O}_2, \mathcal{F})
$$
は表現可能である。
\end{lemma}

\begin{proof}
代数における類似の主張とまったく同じ方法で証明する。
この証明では，$\mathcal{C}$ 上の任意の集合の層 $\mathcal{F}$ に対し，
前層 $U \mapsto \mathcal{O}_2(U)[\mathcal{F}(U)]$ の層化を
$\mathcal{O}_2[\mathcal{F}]$ と書く。ここで右辺は，集合
$\mathcal{F}(U)$ 上の自由 $\mathcal{O}_2(U)$-加群を表す。
$s \in \mathcal{F}(U)$ に対し，対応する
$\mathcal{O}_2[\mathcal{F}]$ の $U$ 上の切断を $[s]$ と書く。
$\mathcal{F}$ が $\mathcal{O}_2$-加群の層ならば，標準写像
$$
c : \mathcal{O}_2[\mathcal{F}] \longrightarrow \mathcal{F}
$$
があり，これは前層の段階では
$\sum f_s[s] \mapsto \sum f_s s$ によって与えられる。
この写像を記述するのに略記 $[s] \mapsto s$ を用い，以下の他の写像にも
同様の略記を用いる。$\mathcal{O}_2$-加群の写像
\begin{equation}
\label{sites-modules-equation-define-module-differentials}
\begin{matrix}
\mathcal{O}_2[\mathcal{O}_2 \times \mathcal{O}_2] \oplus
\mathcal{O}_2[\mathcal{O}_2 \times \mathcal{O}_2] \oplus
\mathcal{O}_2[\mathcal{O}_1] &
\longrightarrow &
\mathcal{O}_2[\mathcal{O}_2] \\
[(a, b)] \oplus [(f, g)] \oplus [h] & \longmapsto & [a + b] - [a] - [b] + \\
& & [fg] - g[f] - f[g] + \\
& & [\varphi(h)]
\end{matrix}
\end{equation}
を考える。ここでは上記の略記を用いた。
$\Omega_{\mathcal{O}_2/\mathcal{O}_1}$ をこの写像の余核とする。
すると，集合の層の写像
$$
\text{d} : \mathcal{O}_2 \longrightarrow \Omega_{\mathcal{O}_2/\mathcal{O}_1}
$$
で，局所切断 $f$ を $\Omega_{\mathcal{O}_2/\mathcal{O}_1}$ における
$[f]$ の像へ送るものが明らかに存在する。構成により $\text{d}$ は
$\mathcal{O}_1$-導分である。次に，$\mathcal{F}$ を
$\mathcal{O}_2$-加群の層とし，
$D : \mathcal{O}_2 \to \mathcal{F}$ を $\mathcal{O}_1$-導分とする。
$[g]$ を $D(g)$ へ送る $\mathcal{O}_2$-線形写像
$\mathcal{O}_2[\mathcal{O}_2] \to \mathcal{F}$ を考えられる。
導分の定義から，この写像は
(\ref{sites-modules-equation-define-module-differentials}) の像に属する切断を零化し，
したがって写像
$$
\alpha_D : \Omega_{\mathcal{O}_2/\mathcal{O}_1} \longrightarrow \mathcal{F}
$$
を定める。$D = \alpha_D \circ \text{d}$ は明らかなので，補題が従う。
\end{proof}

\begin{definition}
\label{sites-modules-definition-module-differentials}
$\mathcal{C}$ をサイトとし，$\varphi : \mathcal{O}_1 \to \mathcal{O}_2$
を環の層の準同型とする。環写像 $\varphi$ の {\it 微分加群}とは，
補題 \ref{sites-modules-lemma-universal-module} により存在する関手
$\mathcal{F} \mapsto \text{Der}_{\mathcal{O}_1}(\mathcal{O}_2, \mathcal{F})$
を表現する対象をいう。これを $\Omega_{\mathcal{O}_2/\mathcal{O}_1}$ と書き，
{\it 普遍 $\varphi$-導分}を
$\text{d} : \mathcal{O}_2 \to \Omega_{\mathcal{O}_2/\mathcal{O}_1}$
と書く。
\end{definition}

\noindent
この加群と導分は関手を表現する普遍対象をなすので，この概念は明らかに
内在的である（すなわち，環付きトポスの基礎となるサイトの選択に依存しない。
\ref{sites-modules-section-intrinsic} 節を参照）。
$\Omega_{\mathcal{O}_2/\mathcal{O}_1}$ は
$\mathcal{O}_2$-加群の写像
(\ref{sites-modules-equation-define-module-differentials}) の余核である。
さらに写像 $\text{d}$ は，$\text{d}f$ が局所切断 $[f]$ の像となるという
規則で記述される。

\begin{lemma}
\label{sites-modules-lemma-differentials-sheafify}
$\mathcal{C}$ をサイトとし，$\varphi : \mathcal{O}_1 \to \mathcal{O}_2$
を環の前層の準同型とする。このとき
$\Omega_{\mathcal{O}_2^\#/\mathcal{O}_1^\#}$ は，前層
$U \mapsto \Omega_{\mathcal{O}_2(U)/\mathcal{O}_1(U)}$ に付随する層である。
\end{lemma}

\begin{proof}
写像 (\ref{sites-modules-equation-define-module-differentials}) を考える。
これと同様の前層の写像があり，$U \in \Ob(\mathcal{C})$ における値は
$$
\mathcal{O}_2(U)[\mathcal{O}_2(U) \times \mathcal{O}_2(U)] \oplus
\mathcal{O}_2(U)[\mathcal{O}_2(U) \times \mathcal{O}_2(U)] \oplus
\mathcal{O}_2(U)[\mathcal{O}_1(U)]
\longrightarrow
\mathcal{O}_2(U)[\mathcal{O}_2(U)]
$$
である。《代数》の定義 \ref{algebra-definition-differentials} における
微分加群の構成により，この写像の余核の $U$ における値は
$\Omega_{\mathcal{O}_2(U)/\mathcal{O}_1(U)}$ である。
他方，(\ref{sites-modules-equation-define-module-differentials}) の各層は
上記の前層の層化である。層化は完全なので，結論が従う。
\end{proof}

\begin{lemma}
\label{sites-modules-lemma-pullback-differentials}
$f : \Sh(\mathcal{D}) \to \Sh(\mathcal{C})$ をトポスの射とし，
$\varphi : \mathcal{O}_1 \to \mathcal{O}_2$ を
$\mathcal{C}$ 上の環の層の準同型とする。このとき標準的な同一視
$f^{-1}\Omega_{\mathcal{O}_2/\mathcal{O}_1} =
\Omega_{f^{-1}\mathcal{O}_2/f^{-1}\mathcal{O}_1}$
があり，普遍導分と両立する。
\end{lemma}

\begin{proof}
層 $\Omega_{\mathcal{O}_2/\mathcal{O}_1}$ は
写像 (\ref{sites-modules-equation-define-module-differentials}) の余核であり，
$\Omega_{f^{-1}\mathcal{O}_2/f^{-1}\mathcal{O}_1}$ にも同様の主張が
成り立つこと，関手 $f^{-1}$ が完全であること，および
$f^{-1}(\mathcal{O}_2[\mathcal{O}_2]) =
f^{-1}\mathcal{O}_2[f^{-1}\mathcal{O}_2]$，
$f^{-1}(\mathcal{O}_2[\mathcal{O}_2 \times \mathcal{O}_2]) =
f^{-1}\mathcal{O}_2[f^{-1}\mathcal{O}_2 \times f^{-1}\mathcal{O}_2]$，および
$f^{-1}(\mathcal{O}_2[\mathcal{O}_1]) =
f^{-1}\mathcal{O}_2[f^{-1}\mathcal{O}_1]$
であることから従う。
\end{proof}

\begin{lemma}
\label{sites-modules-lemma-localize-differentials}
$\mathcal{C}$ をサイトとし，$\varphi : \mathcal{O}_1 \to \mathcal{O}_2$
を環の層の準同型とする。$\mathcal{C}$ の任意の対象 $U$ に対して，
標準同型
$$
\Omega_{\mathcal{O}_2/\mathcal{O}_1}|_U =
\Omega_{(\mathcal{O}_2|_U)/(\mathcal{O}_1|_U)}
$$
があり，普遍導分と両立する。
\end{lemma}

\begin{proof}
これは補題 \ref{sites-modules-lemma-pullback-differentials} の特別な場合である。
\end{proof}

\begin{lemma}
\label{sites-modules-lemma-functoriality-differentials}
$\mathcal{C}$ をサイトとする。
$$
\xymatrix{
\mathcal{O}_2 \ar[r]_\varphi & \mathcal{O}_2' \\
\mathcal{O}_1 \ar[r] \ar[u] & \mathcal{O}'_1 \ar[u]
}
$$
を $\mathcal{C}$ 上の環の層の可換図式とする。
写像 $\mathcal{O}_2 \to \mathcal{O}'_2$ と写像
$\text{d} : \mathcal{O}'_2 \to \Omega_{\mathcal{O}'_2/\mathcal{O}'_1}$
の合成は $\mathcal{O}_1$-導分である。したがって，
$\mathcal{O}_2$-加群の標準写像
$\Omega_{\mathcal{O}_2/\mathcal{O}_1} \to
\Omega_{\mathcal{O}'_2/\mathcal{O}'_1}$
を得る。この写像は，$\mathcal{O}_2$ の任意の局所切断 $f$ に対して
$\text{d}(f)$ が $\text{d}(\varphi(f))$ へ写るという性質によって
一意に特徴付けられる。このようにして $\Omega_{-/-}$ は
環の層の射の圏上の関手となる。
\end{lemma}

\begin{proof}
主張は自明である。
\end{proof}

\begin{lemma}
\label{sites-modules-lemma-differential-seq}
補題 \ref{sites-modules-lemma-functoriality-differentials} において，
$\mathcal{O}_2 \to \mathcal{O}'_2$ が全射で，その核が
$\mathcal{I} \subset \mathcal{O}_2$ であり，
$\mathcal{O}_1 = \mathcal{O}'_1$ であると仮定する。このとき
$\mathcal{O}'_2$-加群の標準完全列
$$
\mathcal{I}/\mathcal{I}^2
\longrightarrow
\Omega_{\mathcal{O}_2/\mathcal{O}_1} \otimes_{\mathcal{O}_2} \mathcal{O}'_2
\longrightarrow
\Omega_{\mathcal{O}'_2/\mathcal{O}_1}
\longrightarrow
0
$$
がある。最左の写像は，$\mathcal{I}$ の局所切断 $f$ が
$\text{d}f \otimes 1$ へ写るという規則で特徴付けられる。
\end{lemma}

\begin{proof}
$\mathcal{I}$ の局所切断 $f$ に対し，
$\mathcal{I}/\mathcal{I}^2$ における $f$ の像を $\overline{f}$ と書く。
写像 $\overline{f} \mapsto \text{d}f \otimes 1$ が良定義であることを示すには，
$f_1,f_2$ が $\mathcal{I}$ の局所切断ならば
$\text{d} f_1f_2 \otimes 1 = 0$ であることを確かめればよい。
これは Leibniz 則と
$\text{d} f_1f_2 \otimes 1 =
(f_1 \text{d}f_2 + f_2 \text{d} f_1 )\otimes 1 =
\text{d}f_2 \otimes f_1 + \text{d}f_2 \otimes f_1 = 0$
から明らかである。同様の計算により，この写像は
$\mathcal{O}'_2 = \mathcal{O}_2/\mathcal{I}$-線形である。
右側の写像は補題 \ref{sites-modules-lemma-functoriality-differentials} の写像である。

\medskip\noindent
この列が完全であることを示すため，次のように論じる。
$\mathcal{O}''_2 \subset \mathcal{O}'_2$ を，$U$ における値が
$\mathcal{O}_2(U) \to \mathcal{O}'_2(U)$ の像であるような
$\mathcal{O}_1$-代数の前層とする。
《代数》の補題 \ref{algebra-lemma-differential-seq} により，列
$$
\mathcal{I}(U)/\mathcal{I}(U)^2
\longrightarrow
\Omega_{\mathcal{O}_2(U)/\mathcal{O}_1(U)}
\otimes_{\mathcal{O}_2(U)} \mathcal{O}''_2(U)
\longrightarrow
\Omega_{\mathcal{O}''_2(U)/\mathcal{O}_1(U)}
\longrightarrow
0
$$
は $\mathcal{C}$ のすべての対象 $U$ に対して完全である。
層化は完全なので，これは $(\mathcal{O}'_2)^\#$-加群の層の完全列を与える。
補題 \ref{sites-modules-lemma-differentials-sheafify} と
$(\mathcal{O}''_2)^\# = \mathcal{O}'_2$ であることから結論が従う。
\end{proof}

\noindent
導分が自然に現れる一つの特別な状況を次に示す。

\begin{lemma}
\label{sites-modules-lemma-double-structure-gives-derivation}
$\mathcal{C}$ をサイトとし，$\varphi : \mathcal{O}_1 \to \mathcal{O}_2$
を環の層の準同型とする。短完全列
$$
0 \to \mathcal{F} \to \mathcal{A} \to \mathcal{O}_2 \to 0
$$
を考える。ここで $\mathcal{A}$ は $\mathcal{O}_1$-代数の層，
$\pi : \mathcal{A} \to \mathcal{O}_2$ は
$\mathcal{O}_1$-代数の層の全射であり，
$\mathcal{F} = \Ker(\pi)$ はその核である。
$\mathcal{F}$ は $\mathcal{A}$ の二乗零イデアル層であると仮定する。
したがって，$\mathcal{F}$ は自然に $\mathcal{O}_2$-加群の構造をもつ。
$\pi$ の切断 $s : \mathcal{O}_2 \to \mathcal{A}$ とは，
$\pi \circ s = \text{id}$ を満たす $\mathcal{O}_1$-代数の写像をいう。
$\pi$ の任意の切断 $s : \mathcal{O}_2 \to \mathcal{F}$ と，
任意の $\varphi$-導分 $D : \mathcal{O}_1 \to \mathcal{F}$ に対し，写像
$$
s + D : \mathcal{O}_1 \to \mathcal{A}
$$
は $\pi$ の切断であり，すべての切断 $s'$ は一意な
$\varphi$-導分 $D$ によって $s+D$ の形に書ける。
\end{lemma}

\begin{proof}
$\mathcal{F}$ 上の $\mathcal{O}_2$-加群構造は
$h \tau = \tilde h \tau$（$\mathcal{A}$ における積）で与えられることを
思い出そう。ここで $h$ は $\mathcal{O}_2$ の局所切断，
$\tilde h$ は $h$ の $\mathcal{A}$ の局所切断への局所的な持ち上げ，
$\tau$ は $\mathcal{F}$ の局所切断である。特に $s$ が与えられれば，
$\tilde h = s(h)$ ととれる。$s+D$ が環の層の準同型であることを
確かめるため，次を計算する：
\begin{eqnarray*}
(s + D)(ab) & = & s(ab) + D(ab) \\
& = & s(a)s(b) + aD(b) + D(a)b \\
& = & s(a) s(b) + s(a)D(b) + D(a)s(b) \\
& = & (s(a) + D(a))(s(b) + D(b))
\end{eqnarray*}
ここでは Leibniz 則を用いた。同様に，$D$ が $\mathcal{O}_1$-導分であることから
$s+D$ は $\mathcal{O}_1$-代数の写像であることが分かる。
逆に $s'$ が与えられたならば $D=s'-s$ とおく。詳細は省略する。
\end{proof}

\begin{definition}
\label{sites-modules-definition-sheaf-differentials}
$X = (\Sh(\mathcal{C}), \mathcal{O})$ と
$Y = (\Sh(\mathcal{C}'), \mathcal{O}')$ を環付きトポスとし，
$(f, f^\sharp) : X \to Y$ を環付きトポスの射とする。この状況で
\begin{enumerate}
\item $\mathcal{F}$ を $\mathcal{O}$-加群の層とするとき，
{\it $Y$-導分} $D : \mathcal{O} \to \mathcal{F}$ とは
$f^\sharp$-導分のことであり，
\item {\it $Y$ 上の $X$ の微分の層 $\Omega_{X/Y}$}とは，
$f^\sharp : f^{-1}\mathcal{O}' \to \mathcal{O}$
の微分加群の層である。定義 \ref{sites-modules-definition-module-differentials} を参照。
\end{enumerate}
したがって $\Omega_{X/Y}$ は {\it 普遍 $Y$-導分}
$\text{d}_{X/Y} : \mathcal{O} \longrightarrow \Omega_{X/Y}$ を備える。
$\Omega_{X/Y} = \Omega_f$ と書くこともある。
\end{definition}

\noindent
$f^\sharp : f^{-1}\mathcal{O}' \to \mathcal{O}$ であるから，
この定義には意味があることを思い出そう。

\begin{lemma}
\label{sites-modules-lemma-functoriality-differentials-ringed-topoi}
$X = (\Sh(\mathcal{C}_X), \mathcal{O}_X)$，
$Y = (\Sh(\mathcal{C}_Y), \mathcal{O}_Y)$，
$X' = (\Sh(\mathcal{C}_{X'}), \mathcal{O}_{X'})$，および
$Y' = (\Sh(\mathcal{C}_{Y'}), \mathcal{O}_{Y'})$ を環付きトポスとし，
$$
\xymatrix{
X' \ar[d] \ar[r]_f & X \ar[d] \\
Y' \ar[r] & Y
}
$$
を環付きトポスの射の可換図式とする。写像
$f^\sharp : \mathcal{O}_X \to f_*\mathcal{O}_{X'}$ と写像
$f_*\text{d}_{X'/Y'} : f_*\mathcal{O}_{X'} \to f_*\Omega_{X'/Y'}$
の合成は $Y$-導分である。したがって，
$\mathcal{O}_X$-加群の標準写像
$\Omega_{X/Y} \to f_*\Omega_{X'/Y'}$ を得る。さらに
$f_*$ と $f^*$ の随伴性により，標準的な
$\mathcal{O}_{X'}$-加群準同型を得る：
$$
c_f : f^*\Omega_{X/Y} \longrightarrow \Omega_{X'/Y'}.
$$
これは，$\mathcal{O}_X$ の任意の局所切断 $t$ に対して
$f^*\text{d}_{X/Y}(t)$ が $\text{d}_{X'/Y'}(f^* t)$ へ写るという性質で
一意に特徴付けられる。
\end{lemma}

\begin{proof}
最後の主張を除けば明らかである。その意味を説明しよう。
$U \in \Ob(\mathcal{C}_X)$，$t \in \mathcal{O}_X(U)$ とする。
これが $t$ を $\mathcal{O}_X$ の局所切断と呼ぶ意味である。
$t$ を集合の層の写像 $t : h_U^\# \to \mathcal{O}_X$ とみなせる。
すると $f^{-1}t : f^{-1}h_U^\# \to f^{-1}\mathcal{O}_X$ である。
$f^*t$ とは，合成
$$
\xymatrix{
f^{-1}h_U^\# \ar[rr]^{f^{-1}t} \ar@/^4ex/[rrrr]^{f^*t} & &
f^{-1}\mathcal{O}_X \ar[rr]^{f^\sharp} & &
\mathcal{O}_{X'}
}
$$
をいう。$\text{d}_{X/Y}(t) \in \Omega_{X/Y}(U)$ であることに注意する。
したがって $\text{d}_{X/Y}(t)$ を写像
$\text{d}_{X/Y}(t) : h_U^\# \to \Omega_{X/Y}$ とみなせる。
すると
$f^{-1}\text{d}_{X/Y}(t) : f^{-1}h_U^\# \to f^{-1}\Omega_{X/Y}$ である。
$f^*\text{d}_{X/Y}(t)$ とは，合成
$$
\xymatrix{
f^{-1}h_U^\#
\ar[rr]^{f^{-1}\text{d}_{X/Y}(t)}
\ar@/^4ex/[rrrr]^{f^*\text{d}_{X/Y}(t)} & &
f^{-1}\Omega_{X/Y} \ar[rr]^{1 \otimes \text{id}} & &
f^*\Omega_{X/Y}
}
$$
をいう。さて，補題の主張は，$f^{-1}h_U^\#$ から
$\Omega_{X'/Y'}$ への写像として
$$
c_f \circ f^*t = f^*\text{d}_{X/Y}(t)
$$
が成り立つという意味である。補題で定義した $c_f$ がこの性質をもつことの
確認は省略する。（ヒント：最初の写像
$c'_f : \Omega_{X/Y} \to f_*\Omega_{X'/Y'}$ は，$U$ 上の
$f_*\Omega_{X'/Y'}$ の切断として
$c'_f(\text{d}_{X/Y}(t)) = f_*\text{d}_{X'/Y'}(f^\sharp(t))$
を満たす。随伴性を用いて，これを上の等式に直さなければならない。）
これが $c_f$ を一意に特徴付ける理由は，局所切断
$\text{d}_{X/Y}(t)$ が $\mathcal{O}_X$-加群 $\Omega_{X/Y}$ を生成するので，
$f^*\text{d}_{X/Y}(t)$ の像が
$\mathcal{O}_{X'}$-加群 $f^*\Omega_{X/Y}$ を生成するからである。
\end{proof}





\section{有限階微分作用素}
\label{sites-modules-section-differential-operators}

\noindent
本節では有限階微分作用素を導入する。読者には，可換代数の章の対応する節
（《代数》\ref{algebra-section-differential-operators} 節）も
参照することを勧める。

\begin{definition}
\label{sites-modules-definition-differential-operators}
$\mathcal{C}$ をサイトとし，
$\varphi : \mathcal{O}_1 \to \mathcal{O}_2$ を環の層の準同型とする。
$k \geq 0$ を整数とし，$\mathcal{F},\mathcal{G}$ を
$\mathcal{O}_2$-加群の層とする。{\it $k$ 階微分作用素
$D : \mathcal{F} \to \mathcal{G}$}とは，$\mathcal{O}_1$-線形写像で，
$\mathcal{O}_2$ のすべての局所切断 $g$ に対して写像
$s \mapsto D(gs)-gD(s)$ が $k-1$ 階微分作用素となるものである。
基底の場合 $k=0$ については，$0$ 階微分作用素を
$\mathcal{O}_2$-線形写像と定義する。
\end{definition}

\noindent
$D : \mathcal{F} \to \mathcal{G}$ が $k$ 階微分作用素ならば，
$\mathcal{O}_2$ のすべての局所切断 $g$ に対し $gD$ も
$k$ 階微分作用素である。二つの $k$ 階微分作用素の和も同じ階の
微分作用素である。したがって，これら全体の集合
$$
\text{Diff}^k(\mathcal{F}, \mathcal{G}) =
\text{Diff}^k_{\mathcal{O}_2/\mathcal{O}_1}(\mathcal{F}, \mathcal{G})
$$
は $\Gamma(\mathcal{C}, \mathcal{O}_2)$-加群である。また，
$$
\text{Diff}^0(\mathcal{F}, \mathcal{G}) \subset
\text{Diff}^1(\mathcal{F}, \mathcal{G}) \subset
\text{Diff}^2(\mathcal{F}, \mathcal{G}) \subset \ldots
$$
である。$U \in \Ob(\mathcal{C})$ を $k$ 階微分作用素
$D : \mathcal{F}|_U \to \mathcal{G}|_U$ の加群へ写す規則は，
サイト $\mathcal{C}$ 上の $\mathcal{O}_2$-加群の層である。
したがって微分作用素の層を得る
（これが必要になれば，ここに定義を追加する）。

\begin{lemma}
\label{sites-modules-lemma-composition-differential-operators}
$\mathcal{C}$ をサイトとし，$\mathcal{O}_1 \to \mathcal{O}_2$ を
環の層の写像とする。$\mathcal{E},\mathcal{F},\mathcal{G}$ を
$\mathcal{O}_2$-加群の層とする。
$D : \mathcal{E} \to \mathcal{F}$ と
$D' : \mathcal{F} \to \mathcal{G}$ がそれぞれ $k$ 階および $k'$ 階の
微分作用素ならば，$D' \circ D$ は $k+k'$ 階微分作用素である。
\end{lemma}

\begin{proof}
$g$ を $\mathcal{O}_2$ の局所切断とする。このとき
$\mathcal{E}$ の局所切断 $x$ を
$$
D'(D(gx)) - gD'(D(x)) = D'(D(gx)) - D'(gD(x)) + D'(gD(x)) - gD'(D(x))
$$
へ写す写像は，より低い階の微分作用素の二つの合成の和である。
したがって，$k+k'$ に関する帰納法により補題が従う。
\end{proof}

\begin{lemma}
\label{sites-modules-lemma-module-principal-parts}
$\mathcal{C}$ をサイトとし，$\mathcal{O}_1 \to \mathcal{O}_2$ を
環の層の写像とする。$\mathcal{F}$ を
$\mathcal{O}_2$-加群の層とし，$k \geq 0$ とする。
$\mathcal{O}_2$-加群の層
$\mathcal{P}^k_{\mathcal{O}_2/\mathcal{O}_1}(\mathcal{F})$
と，$\mathcal{O}_2$-加群の層 $\mathcal{G}$ に関して関手的な標準同型
$$
\text{Diff}^k_{\mathcal{O}_2/\mathcal{O}_1}(\mathcal{F}, \mathcal{G}) =
\Hom_{\mathcal{O}_2}(
\mathcal{P}^k_{\mathcal{O}_2/\mathcal{O}_1}(\mathcal{F}), \mathcal{G})
$$
が存在する。
\end{lemma}

\begin{proof}
存在は一般的な圏論的議論から従う（将来ここに参照を挿入する）が，
この構成は後の証明で有用なので，直接的な構成も与える。
補題 \ref{sites-modules-lemma-universal-module} の証明で導入した記法を自由に用いる。
任意の微分作用素 $D : \mathcal{F} \to \mathcal{G}$ に対し，
$[m]$ を $D(m)$ へ写す $\mathcal{O}_2$-線形写像
$L_D : \mathcal{O}_2[\mathcal{F}] \to \mathcal{G}$
を得る。$D$ が $0$ 階ならば，$L_D$ は局所切断
$$
[m + m'] - [m] - [m'],\quad
g_0[m] - [g_0m]
$$
を零化する。ここで $g_0$ は $\mathcal{O}_2$ の局所切断，
$m,m'$ は $\mathcal{F}$ の局所切断である。$D$ が $1$ 階ならば，
$L_D$ は局所切断
$$
[m + m'] - [m] - [m'],\quad
f[m] - [fm], \quad
g_0g_1[m] - g_0[g_1m] - g_1[g_0m] + [g_1g_0m]
$$
を零化する。ここで $f$ は $\mathcal{O}_1$ の局所切断，
$g_0,g_1$ は $\mathcal{O}_2$ の局所切断，$m,m'$ は
$\mathcal{F}$ の局所切断である。$D$ が $k$ 階ならば，$L_D$ は
局所切断 $[m + m']-[m]-[m']$，$f[m]-[fm]$，および局所切断
$$
g_0g_1\ldots g_k[m] - \sum g_0 \ldots \hat g_i \ldots g_k[g_im] + \ldots
+(-1)^{k + 1}[g_0\ldots g_km]
$$
を零化する。逆に，$L : \mathcal{O}_2[\mathcal{F}] \to \mathcal{G}$ が
前文に列挙したすべての局所切断を零化する
$\mathcal{O}_2$-線形写像ならば，$m \mapsto L([m])$ は $k$ 階微分作用素で
ある。したがって，
$\mathcal{P}^k_{\mathcal{O}_2/\mathcal{O}_1}(\mathcal{F})$
は，$\mathcal{O}_2[\mathcal{F}]$ をこれらの局所切断が生成する
$\mathcal{O}_2$-部分加群の層で割った商である。
\end{proof}

\begin{definition}
\label{sites-modules-definition-module-principal-parts}
$\mathcal{C}$ をサイトとし，$\mathcal{O}_1 \to \mathcal{O}_2$ を
環の層の写像とする。$\mathcal{F}$ を
$\mathcal{O}_2$-加群の層とする。補題
\ref{sites-modules-lemma-module-principal-parts} で構成した加群の層
$\mathcal{P}^k_{\mathcal{O}_2/\mathcal{O}_1}(\mathcal{F})$ を，
$\mathcal{F}$ の {\it $k$ 階主部加群の層}という。
\end{definition}

\noindent
包含
$$
\text{Diff}^0(\mathcal{F}, \mathcal{G}) \subset
\text{Diff}^1(\mathcal{F}, \mathcal{G}) \subset
\text{Diff}^2(\mathcal{F}, \mathcal{G}) \subset \ldots
$$
は，Yoneda の補題（《圏》補題 \ref{categories-lemma-yoneda}）を通じて全射
$$
\ldots \to \mathcal{P}^2_{\mathcal{O}_2/\mathcal{O}_1}(\mathcal{F})
\to \mathcal{P}^1_{\mathcal{O}_2/\mathcal{O}_1}(\mathcal{F})
\to \mathcal{P}^0_{\mathcal{O}_2/\mathcal{O}_1}(\mathcal{F}) = \mathcal{F}
$$
に対応することに注意する。

\begin{lemma}
\label{sites-modules-lemma-differential-operators-sheafify}
$\mathcal{C}$ をサイトとし，$\mathcal{O}_1 \to \mathcal{O}_2$ を
環の前層の準同型とする。$\mathcal{F}$ を
$\mathcal{O}_2$-加群の前層とする。このとき
$\mathcal{P}^k_{\mathcal{O}_2^\#/\mathcal{O}_1^\#}(\mathcal{F}^\#)$
は前層
$U \mapsto P^k_{\mathcal{O}_2(U)/\mathcal{O}_1(U)}(\mathcal{F}(U))$
に付随する層である。
\end{lemma}

\begin{proof}
補題 \ref{sites-modules-lemma-differentials-sheafify} で微分の層について行ったのと
まったく同じ方法で証明できる。あるいは，補題
\ref{sites-modules-lemma-module-principal-parts} の普遍性を直接用いて等式を示す方が
自然かもしれない。詳細は省略する。
\end{proof}

\begin{lemma}
\label{sites-modules-lemma-sequence-of-principal-parts}
$\mathcal{C}$ をサイトとし，$\mathcal{O}_1 \to \mathcal{O}_2$ を
環の層の準同型とする。$\mathcal{F}$ を
$\mathcal{O}_2$-加群の層とする。$\mathcal{F}$ に関して関手的な標準短完全列
$$
0 \to
\Omega_{\mathcal{O}_2/\mathcal{O}_1} \otimes_{\mathcal{O}_2} \mathcal{F} \to
\mathcal{P}^1_{\mathcal{O}_2/\mathcal{O}_1}(\mathcal{F}) \to
\mathcal{F} \to 0
$$
があり，{\it 主部列}と呼ぶ。
\end{lemma}

\begin{proof}
可換代数における版（《代数》補題
\ref{algebra-lemma-sequence-of-principal-parts}）と，補題
\ref{sites-modules-lemma-differentials-sheafify} および
\ref{sites-modules-lemma-differential-operators-sheafify} から従う。
\end{proof}

\begin{remark}
\label{sites-modules-remark-functoriality-principal-parts}
$\mathcal{C}$ をサイトとする。環の層の可換図式
$$
\xymatrix{
\mathcal{B} \ar[r] & \mathcal{B}' \\
\mathcal{A} \ar[u] \ar[r] & \mathcal{A}' \ar[u]
}
$$
と，$\mathcal{B}$-加群 $\mathcal{F}$，
$\mathcal{B}'$-加群 $\mathcal{F}'$，および
$\mathcal{B}$-線形写像 $\mathcal{F} \to \mathcal{F}'$ が
与えられたとする。このとき，整合的な加群写像の系
$$
\xymatrix{
\ldots \ar[r] &
\mathcal{P}^2_{\mathcal{B}'/\mathcal{A}'}(\mathcal{F}') \ar[r] &
\mathcal{P}^1_{\mathcal{B}'/\mathcal{A}'}(\mathcal{F}') \ar[r] &
\mathcal{P}^0_{\mathcal{B}'/\mathcal{A}'}(\mathcal{F}') \\
\ldots \ar[r] &
\mathcal{P}^2_{\mathcal{B}/\mathcal{A}}(\mathcal{F}) \ar[r] \ar[u] &
\mathcal{P}^1_{\mathcal{B}/\mathcal{A}}(\mathcal{F}) \ar[r] \ar[u] &
\mathcal{P}^0_{\mathcal{B}/\mathcal{A}}(\mathcal{F}) \ar[u]
}
$$
を得る。これらの写像は，この種の写像のさらなる合成と整合する。
これを見る最も簡単な方法は，補題 \ref{sites-modules-lemma-module-principal-parts} の
証明における $\mathcal{P}^k_{\mathcal{B}/\mathcal{A}}(\mathcal{M})$ の
（局所）生成元と関係式による記述を用いることであるが，これらの加群の
普遍性から直接見ることもできる。さらに，これらの写像は補題
\ref{sites-modules-lemma-sequence-of-principal-parts} の短完全列と整合する。
\end{remark}








\section{素朴な余接複体}
\label{sites-modules-section-netherlander}

\noindent
本節は《代数》\ref{algebra-section-netherlander} 節および
《加群の層》\ref{modules-section-netherlander} 節の類似である。
読者には，まずこれらの節を読むことを勧める。

\medskip\noindent
$\mathcal{C}$ をサイトとし，$\mathcal{A} \to \mathcal{B}$ を
$\mathcal{C}$ 上の環の層の準同型とする。本節では，
$\mathcal{C}$ 上の任意の集合の層 $\mathcal{E}$ に対し，前層
$U \mapsto \mathcal{A}(U)[\mathcal{E}(U)]$ の層化を
$\mathcal{A}[\mathcal{E}]$ と書く。ここで
$\mathcal{A}(U)[\mathcal{E}(U)]$
は，$\mathcal{E}(U)$ の元に対応する変数をもつ
$\mathcal{A}(U)$ 上の多項式代数を表す。
$e \in \mathcal{E}(U)$ に対応する変数を
$[e] \in \mathcal{A}(U)[\mathcal{E}(U)]$ と書く。
$\mathcal{A}$-代数の標準全射
\begin{equation}
\label{sites-modules-equation-canonical-presentation}
\mathcal{A}[\mathcal{B}] \longrightarrow \mathcal{B},\quad [b] \longmapsto b
\end{equation}
があり，その核を $\mathcal{I} \subset \mathcal{A}[\mathcal{B}]$ と書く。
$\mathcal{I}$ は局所切断 $[b][b']-[bb']$ と $[a]-a$ で
生成されることが直ちに分かる。補題 \ref{sites-modules-lemma-differential-seq} により，
標準写像
\begin{equation}
\label{sites-modules-equation-naive-cotangent-complex}
\mathcal{I}/\mathcal{I}^2
\longrightarrow
\Omega_{\mathcal{A}[\mathcal{B}]/\mathcal{A}}
\otimes_{\mathcal{A}[\mathcal{B}]} \mathcal{B}
\end{equation}
があり，その余核は $\Omega_{\mathcal{B}/\mathcal{A}}$ と標準的に同型である。

\begin{definition}
\label{sites-modules-definition-naive-cotangent-complex}
$\mathcal{C}$ をサイトとし，$\mathcal{A} \to \mathcal{B}$ を
$\mathcal{C}$ 上の環の層の準同型とする。{\it 素朴な余接複体}
$\NL_{\mathcal{B}/\mathcal{A}}$ とは，鎖複体
(\ref{sites-modules-equation-naive-cotangent-complex})
$$
\NL_{\mathcal{B}/\mathcal{A}} =
\left(\mathcal{I}/\mathcal{I}^2
\longrightarrow
\Omega_{\mathcal{A}[\mathcal{B}]/\mathcal{A}}
\otimes_{\mathcal{A}[\mathcal{B}]} \mathcal{B}\right)
$$
であり，$\mathcal{I}/\mathcal{I}^2$ を次数 $-1$ に，
$\Omega_{\mathcal{A}[\mathcal{B}]/\mathcal{A}}
\otimes_{\mathcal{A}[\mathcal{B}]} \mathcal{B}$
を次数 $0$ に置いたものである。
\end{definition}

\noindent
この構成は，微分加群について補題
\ref{sites-modules-lemma-functoriality-differentials} で論じたものと同様の関手性を満たす。
すなわち，$\mathcal{C}$ 上の環の層の可換図式
\begin{equation}
\label{sites-modules-equation-commutative-square-sheaves}
\vcenter{
\xymatrix{
\mathcal{B} \ar[r] & \mathcal{B}' \\
\mathcal{A} \ar[u] \ar[r] & \mathcal{A}' \ar[u]
}
}
\end{equation}
が与えられると，複体の標準 $\mathcal{B}$-線形写像
$$
\NL_{\mathcal{B}/\mathcal{A}} \longrightarrow \NL_{\mathcal{B}'/\mathcal{A}'}
$$
がある。実際，可換図式の写像は標準写像
$\mathcal{A}[\mathcal{B}] \to \mathcal{A}'[\mathcal{B}']$
を与え，これは $\mathcal{I}$ を
$\mathcal{I}'=\Ker(\mathcal{A}'[\mathcal{B}']\to\mathcal{B}')$ へ写す。
したがって，写像
$\mathcal{I}/\mathcal{I}^2 \to \mathcal{I}'/(\mathcal{I}')^2$ と
微分加群の間の写像を得て，両者を合わせると素朴な余接複体の間の
所望の写像が得られる。

\medskip\noindent
$\mathcal{B}$ を $\mathcal{A}$ 上の多項式代数の商として別の表示に
取り替えても，$D(\mathcal{B})$ の同じ対象を得る。このことを説明しよう。
$\mathcal{E}$ を $\mathcal{C}$ 上の集合の層，
$\alpha : \mathcal{E} \to \mathcal{B}$ を集合の層の写像とする。
すると $\mathcal{A}$-代数の準同型
$\mathcal{A}[\mathcal{E}] \to \mathcal{B}$ を得る。
この写像は全射であると仮定し，
$\mathcal{J} \subset \mathcal{A}[\mathcal{E}]$ をその核とする。
$$
\NL(\alpha) = \left(
\mathcal{J}/\mathcal{J}^2
\longrightarrow
\Omega_{\mathcal{A}[\mathcal{E}]/\mathcal{A}}
\otimes_{\mathcal{A}[\mathcal{E}]} \mathcal{B}\right)
$$
と置く。次が成り立つ。

\begin{lemma}
\label{sites-modules-lemma-NL-up-to-qis}
上の状況で，$D(\mathcal{B})$ における標準同型
$\NL(\alpha) = \NL_{\mathcal{B}/\mathcal{A}}$ がある。
\end{lemma}

\begin{proof}
$\NL_{\mathcal{B}/\mathcal{A}}=\NL(\text{id}_\mathcal{B})$ であることに
注意する。したがって，上のような二つの写像
$\alpha_i : \mathcal{E}_i \to \mathcal{B}$ が与えられたとき，
$D(\mathcal{B})$ における標準擬同型
$\NL(\alpha_1)=\NL(\alpha_2)$ があることを示せば十分である。
$\mathcal{E}=\mathcal{E}_1\amalg\mathcal{E}_2$，
$\alpha=\alpha_1\amalg\alpha_2 : \mathcal{E}\to\mathcal{B}$ と置く。また，
$\mathcal{J}_i = \Ker(\mathcal{A}[\mathcal{E}_i] \to \mathcal{B})$
および
$\mathcal{J} = \Ker(\mathcal{A}[\mathcal{E}] \to \mathcal{B})$
と置く。$\mathcal{J}_i$ を $\mathcal{J}$ へ写す写像
$\mathcal{A}[\mathcal{E}_i] \to \mathcal{A}[\mathcal{E}]$ を得る。
したがって，複体の標準写像
$$
\NL(\alpha_i) \longrightarrow \NL(\alpha)
$$
を得るので，これらの写像が擬同型であることを示せば十分である。
そのため次のように論じる。まず，
$H^0(\NL(\alpha_i)) = \Omega_{\mathcal{B}/\mathcal{A}}$ および
$H^0(\NL(\alpha)) = \Omega_{\mathcal{B}/\mathcal{A}}$ が
補題 \ref{sites-modules-lemma-differential-seq} により成り立つことに注意する。
したがって，この写像は次数 $0$ のコホモロジー層上で同型である。
同様に，$H^{-1}(\NL(\alpha_i))$ と $H^{-1}(\NL(\alpha))$ は前層
$U \mapsto H_1(L_{\mathcal{B}(U)/\mathcal{A}(U)})$ に付随する層であると
主張する。ここで $H_1(L_{-/-})$ は
《代数》の定義 \ref{algebra-definition-naive-cotangent-complex} における
ものである。この主張が成り立てば証明は完了する。

\medskip\noindent
主張を証明する。$\alpha : \mathcal{E} \to \mathcal{B}$ を上のような
写像とする。$\mathcal{B}' \subset \mathcal{B}$ を，$U$ における値が
$\mathcal{A}(U)[\mathcal{E}(U)] \to \mathcal{B}(U)$ の像であるような
$\mathcal{A}$-代数の部分前層とする。$\mathcal{I}'$ を，$U$ における値が
$\mathcal{A}(U)[\mathcal{E}(U)] \to \mathcal{B}(U)$ の核である前層とする。
すると $\mathcal{I}$ は $\mathcal{I}'$ の層化であり，
$\mathcal{B}$ は $\mathcal{B}'$ の層化である。同様に，
$H^{-1}(\NL(\alpha))$ は前層
$$
U \longmapsto
\Ker(\mathcal{I}'(U)/\mathcal{I}'(U)^2 \to
\Omega_{\mathcal{A}(U)[\mathcal{E}(U)]/\mathcal{A}(U)}
\otimes_{\mathcal{A}(U)[\mathcal{E}(U)]} \mathcal{B}'(U))
$$
の層化である。これは補題 \ref{sites-modules-lemma-differentials-sheafify} による。
《代数》の補題 \ref{algebra-lemma-NL-homotopy} により，
$H^{-1}(\NL(\alpha))$ は前層
$U \mapsto H_1(L_{\mathcal{B}'(U)/\mathcal{A}(U)})$ に付随する層である。
したがって，写像
$H_1(L_{\mathcal{B}'(U)/\mathcal{A}(U)}) \to
H_1(L_{\mathcal{B}(U)/\mathcal{A}(U)})$ が層化の同型
$\mathcal{H}'_1 \to \mathcal{H}_1$ を誘導することを示さなければならない。

\medskip\noindent
$\mathcal{H}'_1 \to \mathcal{H}_1$ の単射性を示す。
$f \in H_1(L_{\mathcal{B}'(U)/\mathcal{A}(U)})$ が
$\mathcal{H}_1(U)$ において零へ写るとする。示すべきことは，
$f$ が $\mathcal{H}'_1(U)$ において零へ写ることである。
仮定は，被覆 $\{U_i \to U\}$ が存在し，すべての $i$ に対して $f$ が
$H_1(L_{\mathcal{B}(U_i)/\mathcal{A}(U_i)})$ において零へ写ることを意味する。
$U$ を $U_i$ で置き換えることにより，$f$ が
$H_1(L_{\mathcal{B}(U)/\mathcal{A}(U)})$ において零へ写る場合に帰着する。
《代数》の補題 \ref{algebra-lemma-colimits-NL} により，有限生成部分代数
$\mathcal{B}'(U) \subset B \subset \mathcal{B}(U)$ で，$f$ が
$H_1(L_{B/\mathcal{A}(U)})$ において零へ写るものを見つけられる。
$\mathcal{B} = (\mathcal{B}')^\#$ なので，被覆 $\{U_i \to U\}$ で，
$B \to \mathcal{B}(U_i)$ が $\mathcal{B}'(U_i)$ を経由するものを見つけられる。
したがって，望みどおり $f$ は
$H_1(L_{\mathcal{B}'(U_i)/\mathcal{A}(U_i)})$ において零へ写る。

\medskip\noindent
$\mathcal{H}'_1 \to \mathcal{H}_1$ の全射性もまったく同じ方法で証明される。
\end{proof}

\begin{lemma}
\label{sites-modules-lemma-pullback-NL}
$f : \Sh(\mathcal{C}) \to \Sh(\mathcal{D})$ をトポスの射とし，
$\mathcal{A} \to \mathcal{B}$ を $\mathcal{D}$ 上の環の層の準同型とする。
このとき $f^{-1}\NL_{\mathcal{B}/\mathcal{A}} =
\NL_{f^{-1}\mathcal{B}/f^{-1}\mathcal{A}}$。
\end{lemma}

\begin{proof}
省略する。ヒント：補題 \ref{sites-modules-lemma-pullback-differentials} を用いよ。
\end{proof}

\noindent
環付きトポスの射の余接複体を，上で定義した余接複体を用いて定義する。

\begin{definition}
\label{sites-modules-definition-cotangent-complex-morphism-ringed-topoi}
$X = (\Sh(\mathcal{C}), \mathcal{O})$ と
$Y = (\Sh(\mathcal{C}'), \mathcal{O}')$ を環付きトポスとし，
$(f, f^\sharp) : X \to Y$ を環付きトポスの射とする。
この環付きトポスの射の {\it 素朴な余接複体}
$\NL_f = \NL_{X/Y}$ とは
$\NL_{\mathcal{O}/f^{-1}\mathcal{O}'}$ のことである。
$\NL_{X/Y} = \NL_{\mathcal{O}/\mathcal{O}'}$ と書くこともある。
\end{definition}









\section{加群の茎}
\label{sites-modules-section-stalks}

\noindent
点において茎をとる際には少し注意が必要である。茎を定義する余極限
（《サイト》，式 \ref{sites-equation-stalk} を参照）は
フィルター付きとは限らないからである\footnote{もちろん，自然に現れる
ほとんどすべての場合に余極限はフィルター付きであり，本節の議論の一部は
簡略化できる。}。他方，サイトの点の定義により，茎関手は完全であり，
任意の余極限と可換する。言い換えれば，余極限がフィルター付きである場合と
まったく同様に振る舞う。

\begin{lemma}
\label{sites-modules-lemma-stalk-exact}
$\mathcal{C}$ をサイトとし，$p$ を $\mathcal{C}$ の点とする。
\begin{enumerate}
\item $\mathcal{C}$ 上の任意の集合の前層について
$(\mathcal{F}^\#)_p = \mathcal{F}_p$ が成り立つ。
\item 茎関手
$\Sh(\mathcal{C}) \to \textit{Sets}$,
$\mathcal{F} \mapsto \mathcal{F}_p$ は完全であり
（《圏》，定義 \ref{categories-definition-exact} を参照），
任意の余極限と可換する。
\item 茎関手
$\textit{PSh}(\mathcal{C}) \to \textit{Sets}$,
$\mathcal{F} \mapsto \mathcal{F}_p$ は完全であり
（《圏》，定義 \ref{categories-definition-exact} を参照），
任意の余極限と可換する。
\end{enumerate}
\end{lemma}

\begin{proof}
《サイト》の補題 \ref{sites-lemma-point-pushforward-sheaf} により (1) が成り立つ。
《サイト》の補題 \ref{sites-lemma-adjoint-point-push-stalk} により，
$\textit{PSh}(\mathcal{C}) \to \textit{Sets}$,
$\mathcal{F} \mapsto \mathcal{F}_p$ は左随伴である。また，
《サイト》の補題 \ref{sites-lemma-point-pushforward-sheaf} により，
$\Sh(\mathcal{C}) \to \textit{Sets}$,
$\mathcal{F} \mapsto \mathcal{F}_p$ についても同じことが成り立つ。
したがって，茎関手は任意の余極限と可換する
（《圏》，補題 \ref{categories-lemma-adjoint-exact} を参照）。
サイトの点の定義（《サイト》，定義 \ref{sites-definition-point} を参照）から，
$\Sh(\mathcal{C}) \to \textit{Sets}$,
$\mathcal{F} \mapsto \mathcal{F}_p$
は完全である。層化は完全なので
（《サイト》，補題 \ref{sites-lemma-sheafification-exact}），
$\textit{PSh}(\mathcal{C}) \to \textit{Sets}$,
$\mathcal{F} \mapsto \mathcal{F}_p$
も完全である。
\end{proof}

\noindent
特に，前層上の茎関手 $\mathcal{F} \mapsto \mathcal{F}_p$ は
すべての有限極限および有限余極限と可換するので，《サイト》の命題
\ref{sites-proposition-functoriality-algebraic-structures-topoi} の証明の
議論を適用できる。この議論から，$\mathcal{F}$ が《サイト》の命題
\ref{sites-proposition-functoriality-algebraic-structures-topoi} に列挙された
代数構造の（前）層ならば，茎 $\mathcal{F}_p$ は自然に同じ種類の
代数構造となる。$\mathcal{F}$ がアーベル群の前層である場合を
詳しく説明しよう。この場合，加法写像
$+ : \mathcal{F} \times \mathcal{F} \to \mathcal{F}$ は写像
$$
+ :
\mathcal{F}_p \times \mathcal{F}_p
=
(\mathcal{F} \times \mathcal{F})_p
\longrightarrow
\mathcal{F}_p
$$
を誘導する。ここで等号には，集合の前層上の茎関手が有限極限と可換することを
用いた。これは茎 $\mathcal{F}_p$ 上に群構造を定める。このようにして茎関手
$$
\textit{PAb}(\mathcal{C}) \longrightarrow \textit{Ab}, \quad
\mathcal{F} \longmapsto \mathcal{F}_p
$$
を得る。構成により，$\mathcal{F}_p$ の台集合は，台となる集合の前層の
茎である。さらに，包含
$\textit{Ab}(\mathcal{C}) \subset \textit{PAb}(\mathcal{C})$
を前合成することで，アーベル群の層に対する茎関手も定義される。

\begin{lemma}
\label{sites-modules-lemma-stalk-exact-abelian}
$\mathcal{C}$ をサイトとし，$p$ を $\mathcal{C}$ の点とする。
\begin{enumerate}
\item 関手
$\textit{Ab}(\mathcal{C}) \to \textit{Ab}$,
$\mathcal{F} \mapsto \mathcal{F}_p$ は完全である。
\item 茎関手
$\textit{PAb}(\mathcal{C}) \to \textit{Ab}$,
$\mathcal{F}  \mapsto  \mathcal{F}_p$
は完全である。
\item $\mathcal{F} \in \Ob(\textit{PAb}(\mathcal{C}))$ に対して
$\mathcal{F}_p = \mathcal{F}^\#_p$ が成り立つ。
\end{enumerate}
\end{lemma}

\begin{proof}
補題 \ref{sites-modules-lemma-stalk-exact} の結果と上の茎関手の構成から形式的に従う。
\end{proof}

\noindent
次に，加群の層の場合を扱う。$(\mathcal{C},\mathcal{O})$ を
環付きサイトとする。（この議論には，$\mathcal{O}$ が環の前層であれば
十分である。）$\mathcal{F}$ を $\mathcal{O}$-加群の前層とし，
$p$ を $\mathcal{C}$ の点とする。この場合，写像
$$
\cdot :
\mathcal{O}_p \times \mathcal{O}_p
=
(\mathcal{O} \times \mathcal{O})_p
\longrightarrow
\mathcal{O}_p
$$
を得る。これは乗法写像の茎である。また，写像
$$
\cdot :
\mathcal{O}_p \times \mathcal{F}_p
=
(\mathcal{O} \times \mathcal{F})_p
\longrightarrow
\mathcal{F}_p
$$
を得る。これは乗法写像の茎である。これにより
$\mathcal{O}_p$ 上に環構造が，$\mathcal{F}_p$ 上に
$\mathcal{O}_p$-加群構造が定まることの確認は省略する。
このようにして関手
$$
\textit{PMod}(\mathcal{O}) \longrightarrow \textit{Mod}(\mathcal{O}_p), \quad
\mathcal{F} \longmapsto \mathcal{F}_p
$$
を得る。構成により，$\mathcal{F}_p$ の台集合は，台となる集合の前層の
茎である。さらに，包含
$\textit{Mod}(\mathcal{O}) \subset \textit{PMod}(\mathcal{O})$
を前合成することで，$\mathcal{O}$-加群の層に対する茎関手も定義される。

\begin{lemma}
\label{sites-modules-lemma-stalk-exact-modules}
$(\mathcal{C}, \mathcal{O})$ を環付きサイトとし，
$p$ を $\mathcal{C}$ の点とする。
\begin{enumerate}
\item 関手
$\textit{Mod}(\mathcal{O}) \to \textit{Mod}(\mathcal{O}_p)$,
$\mathcal{F} \mapsto \mathcal{F}_p$ は完全である。
\item 茎関手
$\textit{PMod}(\mathcal{O}) \to \textit{Mod}(\mathcal{O}_p)$,
$\mathcal{F} \mapsto \mathcal{F}_p$
は完全である。
\item $\mathcal{F} \in \Ob(\textit{PMod}(\mathcal{O}))$ に対して
$\mathcal{F}_p = \mathcal{F}^\#_p$ が成り立つ。
\end{enumerate}
\end{lemma}

\begin{proof}
補題 \ref{sites-modules-lemma-stalk-exact-abelian}，上の茎関手の構成，および
補題 \ref{sites-modules-lemma-abelian} から形式的に従う。
\end{proof}

\begin{lemma}
\label{sites-modules-lemma-pullback-stalk}
$(f, f^\sharp) :
(\Sh(\mathcal{C}), \mathcal{O}_\mathcal{C})
\to
(\Sh(\mathcal{D}), \mathcal{O}_\mathcal{D})$
を環付きトポスまたは環付きサイトの射とする。
$p$ を $\mathcal{C}$ または $\Sh(\mathcal{C})$ の点とし，
$q = f \circ p$ と置く。このとき，任意の
$\mathcal{O}_\mathcal{D}$-加群 $\mathcal{F}$ に対して
$$
(f^*\mathcal{F})_p =
\mathcal{F}_q \otimes_{\mathcal{O}_{\mathcal{D}, q}}
\mathcal{O}_{\mathcal{C}, p}
$$
が成り立つ。
\end{lemma}

\begin{proof}
定義により
$$
f^*\mathcal{F} =
f^{-1}\mathcal{F} \otimes_{f^{-1}\mathcal{O}_\mathcal{D}}
\mathcal{O}_\mathcal{C}
$$
である。補題 \ref{sites-modules-lemma-tensor-product-pullback} により，
$p$ で茎をとること（すなわち $p^{-1}$ を適用すること）は
$\otimes$ と可換する。あとは，$p$ における引き戻しの茎と $q$ における
茎との関係を述べる《サイト》の補題
\ref{sites-lemma-point-morphism-sites} または補題
\ref{sites-lemma-point-morphism-topoi} から従う。
\end{proof}






\section{摩天楼層}
\label{sites-modules-section-skyscraper}

\noindent
$p$ をサイト $\mathcal{C}$ またはトポス $\Sh(\mathcal{C})$ の点とする。
本節では，アーベル群 $A$ に摩天楼層 $p_*A$ を対応させる関手の
完全性に関する性質を調べる。まず，
$p_* : \textit{Sets} \to \Sh(\mathcal{C})$ は多くの完全性をもつことを
思い出そう。《サイト》の補題 \ref{sites-lemma-stalk-skyscraper} および
\ref{sites-lemma-skyscraper-functor-exact} を参照。

\begin{lemma}
\label{sites-modules-lemma-skyscraper-exact}
$\mathcal{C}$ をサイトとし，$p$ を $\mathcal{C}$ または
それに付随するトポスの点とする。
\begin{enumerate}
\item 関手 $p_* : \textit{Ab} \to \textit{Ab}(\mathcal{C})$，
$A \mapsto p_*A$ は完全である。
\item 関手的な直和分解
$$
p^{-1}p_*A = A \oplus I(A)
$$
が $A \in \Ob(\textit{Ab})$ に対して存在する。
\end{enumerate}
\end{lemma}

\begin{proof}
《サイト》の補題 \ref{sites-lemma-stalk-skyscraper} により，合成が
$\text{id}_A$ に等しい関手的な写像 $A \to p^{-1}p_*A \to A$ がある。
したがって，$I(A)$ を随伴写像 $p^{-1}p_*A \to A$ の核とすれば，
(2) の関手的な直和分解を得る。補題
\ref{sites-modules-lemma-exactness-pushforward-pullback} により関手 $p_*$ は左完全である。
また，《サイト》の補題 \ref{sites-lemma-skyscraper-functor-exact} により，
関手 $p_*$ は全射を全射へ写す。よって (1) が成り立つ。
\end{proof}

\noindent
加群の層について同じことを行うため，
$p$ を環付きトポス $(\Sh(\mathcal{C}),\mathcal{O})$ の点とする。
$p^{-1}$ は茎関手にほかならないことを思い出そう。
したがって $p$ を環付きトポスの射
$$
(p, \text{id}_{\mathcal{O}_p}) :
(\Sh(pt), \mathcal{O}_p)
\longrightarrow
(\Sh(\mathcal{C}), \mathcal{O}).
$$
とみなせる。これにより引き戻し関手
$p^* : \textit{Mod}(\mathcal{O}) \to \textit{Mod}(\mathcal{O}_p)$
を得る。これは茎関手に等しく，補題
\ref{sites-modules-lemma-stalk-exact-modules} で論じたものである。
本節では関手
$p_* : \textit{Mod}(\mathcal{O}_p) \to \textit{Mod}(\mathcal{O})$
を考える。

\begin{lemma}
\label{sites-modules-lemma-skyscraper-modules-exact}
$(\Sh(\mathcal{C}), \mathcal{O})$ を環付きトポスとし，
$p$ をトポス $\Sh(\mathcal{C})$ の点とする。
\begin{enumerate}
\item 関手
$p_* : \textit{Mod}(\mathcal{O}_p) \to \textit{Mod}(\mathcal{O})$,
$M \mapsto p_*M$ は完全である。
\item 標準全射 $p^{-1}p_*M \to M$ は $\mathcal{O}_p$-線形である。
\item 補題 \ref{sites-modules-lemma-skyscraper-exact} の関手的な直和分解
$p^{-1}p_*M = M \oplus I(M)$ は，一般には $\mathcal{O}_p$-線形
{\bf ではない}。
\end{enumerate}
\end{lemma}

\begin{proof}
(1) と (2) の全射性は，補題 \ref{sites-modules-lemma-skyscraper-exact} における
アーベル群の層についての対応する結果から直ちに従う。
$p^{-1}\mathcal{O} = \mathcal{O}_p$ なので $p^{-1}=p^*$ であり，
したがって $p^{-1}p_*M \to M$ は加群に対する随伴の余単位
$p^*p_*M \to M$ と同じであるから線形である。

\medskip\noindent
(3) を証明する。$G$ を群とする。トポス
$G\textit{-Sets} = \Sh(\mathcal{T}_G)$
と点 $p : \textit{Sets} \to G\textit{-Sets}$ を考える。
《サイト》の \ref{sites-section-example-sheaf-G-sets} 節および例
\ref{sites-example-point-G-sets} を参照。ここで $p^{-1}$ は
$G$-作用を忘れる関手である。また $p_*$ は忘却関手の右随伴であり，
$M$ を $\text{Map}(G,M)$ へ写す。直和分解に現れる写像は
$$
M \to \text{Map}(G, M) \to M
$$
である。最初の写像は $m \in M$ を値 $m$ の定値写像へ送り，
二つ目の写像は $G$ の単位元 $1$ における評価である。
次に，$R$ を $G$ の作用を備えた環とする。これは
$\mathcal{T}_G$ 上の環の層 $\mathcal{O}$ を定める。
$\mathcal{O}$-加群の圏は，$R$ 上の作用と両立する $G$-作用を備えた
$R$-加群 $M$ の圏である。$\text{Map}(G,M)$ 上の $R$-加群構造は
$$
( r f ) (\sigma) = \sigma(r) f(\sigma)
$$
によって与えられる。ここで $r \in R$，$f \in \text{Map}(G,M)$ である。
これは $1$ における評価と両立する一意な $G$-不変
$R$-加群構造だからである。一般に $M \to \text{Map}(G,M)$ の像は
$R$-部分加群ではない（例えば $M=R$ とし，$G$-作用が非自明であると
仮定すればよい）ことが分かり，証明が完了する。
\end{proof}

\begin{example}
\label{sites-modules-example-ring-with-group-action}
$G$ を群とする。サイト $\mathcal{T}_G$ とその点 $p$ を考える。
《サイト》の例 \ref{sites-example-point-G-sets} を参照。
$R$ を $G$-作用を備えた環とし，これに対応する
$\mathcal{T}_G$ 上の環の層を $\mathcal{O}$ とする。
このとき $\mathcal{O}_p=R$ であり，ここでは $G$-作用を忘れる。
この場合 $p^{-1}p_*M=\text{Map}(G,M)$，
$I(M)=\{f:G\to M\mid f(1_G)=0\}$ であり，写像
$M\to\text{Map}(G,M)$ は $m\in M$ に値 $m$ の定値関数を対応させる。
\end{example}




\section{局所化と点}
\label{sites-modules-section-localize-points}

\begin{lemma}
\label{sites-modules-lemma-stalk-j-shriek}
$(\mathcal{C}, \mathcal{O})$ を環付きサイトとする。
$p$ を $\mathcal{C}$ の点とし，$U$ を $\mathcal{C}$ の対象とする。
$\mathcal{G} \in \textit{Mod}(\mathcal{O}_U)$ に対して
$$
(j_{U!}\mathcal{G})_p =
\bigoplus\nolimits_q \mathcal{G}_q
$$
が成り立つ。ここで余積は，$p$ 上にある $\mathcal{C}/U$ の点 $q$ に
わたって取る。《サイト》の補題
\ref{sites-lemma-points-above-point} を参照。
\end{lemma}

\begin{proof}
補題 \ref{sites-modules-lemma-extension-by-zero} による，
$j_{U!}\mathcal{G}$ を前層
$V \mapsto
\bigoplus\nolimits_{\varphi \in \Mor_\mathcal{C}(V, U)}
\mathcal{G}(V/_\varphi U)$
に付随する層として記述する方法を用いる。
$p$ における $j_{U!}\mathcal{G}$ の茎はこの前層の茎に等しい。
補題 \ref{sites-modules-lemma-stalk-exact-modules} を参照。
$u : \mathcal{C} \to \textit{Sets}$ を $p$ に対応する関手とする
（《サイト》の \ref{sites-section-points} 節参照）。
したがって
$$
(j_{U!}\mathcal{G})_p = \colim_{(V, y)}
\bigoplus\nolimits_{\varphi : V \to U} \mathcal{G}(V/_\varphi U)
$$
となる。ここで余極限はアーベル群の圏で取る。
この余極限に現れる四つ組 $(V, y, \varphi, s)$ に
$x = u(\varphi)(y) \in u(U)$ を対応させられる。よって
$$
(j_{U!}\mathcal{G})_p =
\bigoplus\nolimits_{x \in u(U)}
\colim_{(\varphi : V \to U, y), \ u(\varphi)(y) = x} \mathcal{G}(V/_\varphi U).
$$
《サイト》の補題 \ref{sites-lemma-points-above-point} における
$x$ 上の点 $q$ の記述により，これは補題の表示に等しい。
\end{proof}

\begin{remark}
\label{sites-modules-remark-not-pushforward}
注意：補題 \ref{sites-modules-lemma-stalk-j-shriek} の結果には，
$j_{U, *}$ に対する類似はない。
\end{remark}












\section{平坦加群の引き戻し}
\label{sites-modules-section-pullback-flat}

\noindent
環付きトポスの射に沿う平坦加群の引き戻しは平坦である。
これを証明するには少し工夫が要る。

\begin{lemma}
\label{sites-modules-lemma-pullback-flat}
\begin{reference}
\cite[Expos\'e V, Corollary 1.7.1]{SGA4}
\end{reference}
写像
$(f, f^\sharp) :
(\Sh(\mathcal{C}), \mathcal{O}_\mathcal{C})
\to
(\Sh(\mathcal{D}), \mathcal{O}_\mathcal{D})$
を環付きトポスまたは環付きサイトの射とする。
$\mathcal{F}$ が平坦な $\mathcal{O}_\mathcal{D}$-加群ならば，
$f^*\mathcal{F}$ は平坦な $\mathcal{O}_\mathcal{C}$-加群である。
\end{lemma}

\begin{proof}
補題 \ref{sites-modules-lemma-morphism-ringed-topoi-comes-from-morphism-ringed-sites}
のような図式を選ぶ。平坦加群であることは内在的である
（\ref{sites-modules-section-intrinsic} 節および定義 \ref{sites-modules-definition-flat} 参照）。
したがって，射 $(h, h^\sharp) :
(\Sh(\mathcal{C}'), \mathcal{O}_{\mathcal{C}'})
\to
(\Sh(\mathcal{D}'), \mathcal{O}_{\mathcal{D}'})$
について補題を証明すれば十分である。言い換えれば，
サイト $\mathcal{C}$ と $\mathcal{D}$ はすべての有限極限をもち，
$f$ は有限極限と可換する連続関手
$u : \mathcal{D} \to \mathcal{C}$ が誘導するサイトの射であると
仮定してよい。

\medskip\noindent
$f^*\mathcal{F} =
\mathcal{O}_\mathcal{C} \otimes_{f^{-1}\mathcal{O}_\mathcal{D}}
f^{-1}\mathcal{F}$ であることを思い出そう
（定義 \ref{sites-modules-definition-pushforward}）。
補題 \ref{sites-modules-lemma-flat-change-of-rings} により，
$f^{-1}\mathcal{F}$ が平坦な
$f^{-1}\mathcal{O}_\mathcal{D}$-加群であることを証明すれば十分である。
前段落と合わせると，次の段落の状況へ帰着する。

\medskip\noindent
$\mathcal{C}$ と $\mathcal{D}$ をすべての有限極限をもつサイトとし，
$u : \mathcal{D} \to \mathcal{C}$ を有限極限と可換する連続関手とする。
$\mathcal{O}$ を $\mathcal{D}$ 上の環の層，
$\mathcal{F}$ を平坦な $\mathcal{O}$-加群とする。
このとき $u$ はサイトの射 $f : \mathcal{C} \to \mathcal{D}$ を定める
（《サイト》の命題 \ref{sites-proposition-get-morphism}）。
示すべきことは，$f^{-1}\mathcal{F}$ が平坦な
$f^{-1}\mathcal{O}$-加群であることである。
$U$ を $\mathcal{C}$ の対象とし，
$$
f^{-1}\mathcal{O}|_U \xrightarrow{(f_1, \ldots, f_n)}
f^{-1}\mathcal{O}|_U^{\oplus n} \xrightarrow{(s_1, \ldots, s_n)}
f^{-1}\mathcal{F}|_U
$$
を $f^{-1}\mathcal{O}|_U$-加群の複体とする。
目標は，補題 \ref{sites-modules-lemma-flat-eq} の (2) のように，
$U$ のある被覆の各元上で $(s_1, \ldots, s_n)$ の分解を構成することである。
元 $s_a \in f^{-1}\mathcal{F}(U)$ と
$f_a \in f^{-1}\mathcal{O}(U)$ を考える。
$f^{-1}\mathcal{F}$ および $f^{-1}\mathcal{O}$ は
$u_p\mathcal{F}$ の層化なので，$U$ をある被覆の各元で置き換えれば，
$s_a$ は元 $s'_a \in u_p\mathcal{F}(U)$ の像であり，
$f_a$ は元 $f'_a \in u_p\mathcal{O}(U)$ の像であると仮定してよい。
さらにもう一度 $U$ をある被覆の各元で置き換えると，
$\sum f'_as'_a$ は $u_p\mathcal{F}(U)$ において零であると仮定できる。
圏 $(\mathcal{I}_U^u)^{opp}$ は有向であり
（《サイト》の補題 \ref{sites-lemma-directed}），
$u_p\mathcal{F}(U) = \colim_{(\mathcal{I}_U^u)^{opp}} \mathcal{F}(V)$
および $u_p\mathcal{O}(U) = \colim_{(\mathcal{I}_U^u)^{opp}} \mathcal{O}(V)$
であることを思い出そう。したがって，
$(V, \phi) \in \Ob(\mathcal{I}_U^u)$，
$\mathcal{D}$ の対象 $V$，$\mathcal{D}$ の射
$\phi : U \to u(V)$，および元
$s''_a \in \mathcal{F}(V)$，$f''_a \in \mathcal{O}(V)$ が存在し，
後二者の $u_p\mathcal{F}(U)$ と $u_p\mathcal{O}(U)$ における像は
それぞれ $s'_a$ と $f'_a$ に等しく，さらに
$\sum f''_a s''_a = 0$ が $\mathcal{F}(V)$ で成り立つと仮定してよい。
すると複体
$$
\mathcal{O}|_V \xrightarrow{(f''_1, \ldots, f''_n)}
\mathcal{O}|_V^{\oplus n} \xrightarrow{(s''_1, \ldots, s''_n)}
\mathcal{F}|_V
$$
を得る。補題 \ref{sites-modules-lemma-flat-eq} の逆方向を適用すると，
$\mathcal{D}$ における被覆 $\{V_i \to V\}$ が存在し，
各 $i$ に対して $(s''_1, \ldots, s''_n)|_{V_i}$ の分解
$$
\mathcal{O}|_{V_i}^{\oplus n}
\xrightarrow{B''_i}
\mathcal{O}|_{V_i}^{\oplus l_i} \xrightarrow{(t''_{i1}, \ldots, t''_{il_i})}
\mathcal{F}|_{V_i}
$$
が存在して，
$B_i \circ (f''_1, \ldots, f''_n)|_{V_i} = 0$ を満たす。
$U_i = U \times_{\phi, u(V)} u(V_i)$ と置き，
$B''_i$ の像を
$B_i \in \text{Mat}(l_i \times n, f^{-1}\mathcal{O}(U_i))$，
$t''_{ij}$ の像を $t_{ij} \in f^{-1}\mathcal{F}(U_i)$ と書く。
このとき $(s_1, \ldots, s_n)|_{U_i}$ の分解
$$
f^{-1}\mathcal{O}|_{U_i}^{\oplus n}
\xrightarrow{B_i}
f^{-1}\mathcal{O}|_{U_i}^{\oplus l_i}
\xrightarrow{(t_{i1}, \ldots, t_{il_i})}
\mathcal{F}|_{U_i}
$$
を得て，$B_i \circ (f_1, \ldots, f_n)|_{U_i} = 0$ を満たす。
これで証明が完了する。
\end{proof}

\begin{lemma}
\label{sites-modules-lemma-stalk-flat}
$(\mathcal{C}, \mathcal{O})$ を環付きサイトとし，
$p$ を $\mathcal{C}$ の点とする。
$\mathcal{F}$ が平坦な $\mathcal{O}$-加群ならば，
$\mathcal{F}_p$ は平坦な $\mathcal{O}_p$-加群である。
\end{lemma}

\begin{proof}
\ref{sites-modules-section-skyscraper} 節で見たように，$p$ は環付きトポスの射
$$
(p, \text{id}_{\mathcal{O}_p}) :
(\Sh(pt), \mathcal{O}_p)
\longrightarrow
(\Sh(\mathcal{C}), \mathcal{O}).
$$
とみなせ，その引き戻し関手
$p^* : \textit{Mod}(\mathcal{O}) \to \textit{Mod}(\mathcal{O}_p)$
は茎関手に等しい。したがって補題 \ref{sites-modules-lemma-pullback-flat} から従う。
\end{proof}

\begin{lemma}
\label{sites-modules-lemma-check-flat-stalks}
$(\mathcal{C}, \mathcal{O})$ を環付きサイトとし，
$\mathcal{F}$ を $\mathcal{O}$-加群の層とする。
$\{p_i\}_{i \in I}$ を $\mathcal{C}$ の点の保守的な族とする。
このとき $\mathcal{F}$ が平坦であることと，すべての $i \in I$ に対して
$\mathcal{F}_{p_i}$ が平坦な $\mathcal{O}_{p_i}$-加群であることは同値である。
\end{lemma}

\begin{proof}
一方の含意は補題 \ref{sites-modules-lemma-stalk-flat} から分かる。
逆は，補題 \ref{sites-modules-lemma-tensor-product-pullback} による等式
$(\mathcal{F} \otimes_\mathcal{O} \mathcal{G})_p =
\mathcal{F}_p \otimes_{\mathcal{O}_p} \mathcal{G}_p$
（$p$ における茎を取ることは $p^{-1}$ で与えられる）と
補題 \ref{sites-modules-lemma-check-exactness-stalks} を用いればよい。
\end{proof}

\begin{lemma}
\label{sites-modules-lemma-pullback-ses}
射
$f : (\Sh(\mathcal{C}'), \mathcal{O}') \to (\Sh(\mathcal{C}), \mathcal{O})$
を環付きトポスの射とする。
$0 \to \mathcal{F} \to \mathcal{G} \to \mathcal{H} \to 0$
を $\mathcal{O}$-加群の短完全列とし，
$\mathcal{H}$ は平坦な $\mathcal{O}$-加群であるとする。このとき列
$0 \to f^*\mathcal{F} \to f^*\mathcal{G} \to f^*\mathcal{H} \to 0$
も完全である。
\end{lemma}

\begin{proof}
$f^{-1}$ は完全なので，$f^{-1}\mathcal{O}$-加群の短完全列
$0 \to f^{-1}\mathcal{F} \to f^{-1}\mathcal{G} \to f^{-1}\mathcal{H} \to 0$
を得る。補題 \ref{sites-modules-lemma-pullback-flat} により，
$f^{-1}\mathcal{O}$-加群 $f^{-1}\mathcal{H}$ は平坦である。
補題 \ref{sites-modules-lemma-flat-tor-zero} により，この列を
$f^{-1}\mathcal{O}$ 上で $\mathcal{O}'$ とテンソルしても完全性は保たれる。
さらに
$f^*\mathcal{F} = f^{-1}\mathcal{F} \otimes_{f^{-1}\mathcal{O}} \mathcal{O}'$
であり，$\mathcal{G}$ と $\mathcal{H}$ についても同様なので結論を得る。
\end{proof}





\section{局所環付きトポス}
\label{sites-modules-section-locally-ringed}


\noindent
本節の参考文献は
\cite[Expos\'e IV, Exercice 13.9]{SGA4} である。

\begin{lemma}
\label{sites-modules-lemma-locally-ringed}
$(\mathcal{C}, \mathcal{O})$ を環付きサイトとする。
次の条件は同値である。
\begin{enumerate}
\item $\mathcal{C}$ の任意の対象 $U$ と $f \in \mathcal{O}(U)$ に対して，
被覆 $\{U_j \to U\}$ が存在し，各 $j$ について
$f|_{U_j}$ または $(1 - f)|_{U_j}$ のいずれかが可逆である。
\item $U \in \Ob(\mathcal{C})$，$n \geq 1$ とし，
$f_1, \ldots, f_n \in \mathcal{O}(U)$ が $\mathcal{O}(U)$ の単位イデアルを
生成するとする。このとき被覆 $\{U_j \to U\}$ が存在し，各 $j$ について
ある $i$ が存在して $f_i|_{U_j}$ は可逆である。
\item 集合の層の写像
$$
(\mathcal{O} \times \mathcal{O})
\amalg
(\mathcal{O} \times \mathcal{O})
\longrightarrow
\mathcal{O} \times \mathcal{O}
$$
は，第一成分の $(f, a)$ を $(f, af)$ へ，
第二成分の $(f, b)$ を $(f, b(1 - f))$ へ写すものとすると，全射である。
\end{enumerate}
\end{lemma}

\begin{proof}
(2) が (1) を含意することは明らかである。(1) が (2) を含意することを
$n$ に関する帰納法で示す。最初の場合は $n = 2$ である
（$n = 1$ は自明である）。この場合，ある
$a_1, a_2 \in \mathcal{O}(U)$ に対して
$a_1f_1 + a_2f_2 = 1$ が成り立つ。仮定により被覆
$\{U_j \to U\}$ を選んで，各 $j$ について
$a_1f_1|_{U_j}$ または $a_2f_2|_{U_j}$ のいずれかを可逆にできる。
したがって，所望のとおり $f_1|_{U_j}$ または $f_2|_{U_j}$ のいずれかが
可逆である。$n > 2$ のとき，ある
$a_1, \ldots, a_n \in \mathcal{O}(U)$ に対して
$a_1f_1 + \ldots + a_nf_n = 1$ が成り立つ。
$n = 2$ の場合により，被覆 $\{U_j \to U\}_{j \in J}$ が存在し，
各 $j$ について $f_n|_{U_j}$ または
$a_1f_1 + \ldots + a_{n - 1}f_{n - 1}|_{U_j}$ のいずれかが可逆である。
前者が成り立つ $j$ の集合を $J_n$ とし，$J' = J \setminus J_n$ と置く。
帰納法の仮定により，各 $j \in J'$ について被覆
$\{U_{jk} \to U_j\}_{k \in K_j}$ を選べ，各 $k \in K_j$ について
ある $i \in \{1, \ldots, n - 1\}$ が存在して
$f_i|_{U_{jk}}$ は可逆である。サイトの公理により，射の族
$\{U_j \to U\}_{j \in J_n} \cup \{U_{jk} \to U\}_{j \in J', k \in K_j}$
は所望の性質をもつ被覆である。

\medskip\noindent
(1) を仮定する。(3) の写像が全射であることを示すため，
$(f, c)$ を $U$ 上の $\mathcal{O} \times \mathcal{O}$ の切断とする。
仮定により，被覆 $\{U_j \to U\}$ が存在して，各 $j$ について
$f$ または $1 - f$ の制限が可逆な切断となる。
前者の場合は $a = c|_{U_j} (f|_{U_j})^{-1}$，
後者の場合は $b = c|_{U_j} (1 - f|_{U_j})^{-1}$ と取れる。
したがって $(f, c)$ は各被覆元上で写像の像に属する。
逆に (3) を仮定する。任意の $U$ と $f \in \mathcal{O}(U)$ に対して，
$U$ の被覆 $\{U_j \to U\}$ が存在し，各 $U_j$ 上の切断
$(f, 1)|_{U_j}$ は (3) の写像の切断上の像に属する。
これはまさに $f$ または $1 - f$ のいずれかの $U_j$ 上への制限が
可逆な切断であることを意味するので，(1) が成り立つ。
\end{proof}

\begin{lemma}
\label{sites-modules-lemma-locally-ringed-stalk}
$(\mathcal{C}, \mathcal{O})$ を環付きサイトとする。
次の条件を考える。
\begin{enumerate}
\item $\mathcal{C}$ の任意の対象 $U$ と $f \in \mathcal{O}(U)$ に対して，
被覆 $\{U_j \to U\}$ が存在し，各 $j$ について
$f|_{U_j}$ または $(1 - f)|_{U_j}$ のいずれかが可逆である。
\item $\mathcal{C}$ の任意の点 $p$ に対して，茎 $\mathcal{O}_p$ は
零環または局所環である。
\end{enumerate}
常に (1) $\Rightarrow$ (2) が成り立つ。
$\mathcal{C}$ が十分多くの点をもつならば，(1) と (2) は同値である。
\end{lemma}

\begin{proof}
(1) を仮定する。$p$ を関手
$u : \mathcal{C} \to \textit{Sets}$ が与える $\mathcal{C}$ の点とし，
$f_p \in \mathcal{O}_p$ とする。$\mathcal{O}_p$ は
《サイト》の式 (\ref{sites-equation-stalk}) によって計算されるので，
$f_p$ は $x \in U(U)$，$f \in \mathcal{O}(U)$ を満たす三つ組
$(U, x, f)$ で表せる。仮定により被覆 $\{U_i \to U\}$ が存在し，
各 $i$ について $f$ または $1 - f$ のいずれかが $U_i$ 上で可逆である。
$u$ はサイトの点を定めるので，ある $i$ に対して
$x \in u(U)$ へ写る $x_i \in u(U_i)$ が存在する。
《サイト》の式 (\ref{sites-equation-stalk}) の前後の議論により，
$(U, x, f)$ と $(U_i, x_i, f|_{U_i})$ は
$\mathcal{O}_p$ の同じ元を定める。したがって，
$f_p$ または $1 - f_p$ のいずれかが可逆である。
ゆえに $\mathcal{O}_p$ は，任意の元 $a$ について
$a$ または $1 - a$ のいずれかが可逆であるような環である。
これは $\mathcal{O}_p$ が零環または局所環であることを意味する。
《代数》の補題 \ref{algebra-lemma-characterize-local-ring} を参照。

\medskip\noindent
(2) を仮定し，$\mathcal{C}$ は十分多くの点をもつとする。
補題 \ref{sites-modules-lemma-locally-ringed} の (3) の集合の層の写像
$$
\mathcal{O} \times \mathcal{O} \amalg \mathcal{O} \times \mathcal{O}
\longrightarrow
\mathcal{O} \times \mathcal{O}
$$
を考える。任意の局所環 $R$ に対して，対応する写像
$(R \times R) \amalg (R \times R) \to R \times R$
は全射である。例えば《代数》の補題
\ref{algebra-lemma-characterize-local-ring} を参照。
各 $\mathcal{O}_p$ は局所環または零環なので，この写像は茎上で全射である。
したがって，$\mathcal{C}$ が十分多くの点をもつという仮定により，
この写像は全射であり，主張が従う。
\end{proof}

\noindent
《加群》の \ref{modules-section-pathology} 節では，環付き空間
$(X, \mathcal{O}_X)$ に，その上で構造層が零となる開部分空間が
存在しうることを指摘した。これを防ぐには，空間 $X$ のすべての茎で，
切断 $1$ と $0$ が異なる値をもつことを要求すればよい。
環付きトポスと環付きサイトの場合，条件は
\begin{equation}
\label{sites-modules-equation-one-is-never-zero}
\emptyset^\# \longrightarrow
\text{Equalizer}(0, 1 : * \longrightarrow \mathcal{O})
\end{equation}
が層の同型となることである。ここで $*$ は一点層，
$\emptyset^\#$ は「空層」，すなわち層の圏におけるそれぞれ終対象と始対象である。
《サイト》の例 \ref{sites-example-singleton-sheaf} および
\ref{sites-section-almost-cocontinuous} 節をそれぞれ参照。
言い換えれば，条件は $U \in \Ob(\mathcal{C})$ が層論的に空でないならば，
$1, 0 \in \mathcal{O}(U)$ は等しくないということである。
ここで当然の補題を述べておく。

\begin{lemma}
\label{sites-modules-lemma-ringed-stalk-not-zero}
$(\mathcal{C}, \mathcal{O})$ を環付きサイトとする。次の主張を考える。
\begin{enumerate}
\item (\ref{sites-modules-equation-one-is-never-zero}) は同型である。
\item $\mathcal{C}$ の任意の点 $p$ について，茎 $\mathcal{O}_p$ は
零環ではない。
\end{enumerate}
常に (1) $\Rightarrow$ (2) が成り立ち，
$\mathcal{C}$ が十分多くの点をもつならば (1) $\Leftrightarrow$ (2) である。
\end{lemma}

\begin{proof}
省略する。
\end{proof}

\noindent
補題 \ref{sites-modules-lemma-locally-ringed}，
\ref{sites-modules-lemma-locally-ringed-stalk}，および
\ref{sites-modules-lemma-ringed-stalk-not-zero} により，次の定義が自然に導かれる。

\begin{definition}
\label{sites-modules-definition-locally-ringed}
環付きサイト $(\mathcal{C}, \mathcal{O})$ が
{\it 局所環付きサイト}であるとは，
(\ref{sites-modules-equation-one-is-never-zero}) が同型であり，かつ補題
\ref{sites-modules-lemma-locally-ringed} の同値な性質を満たすことをいう。
\end{definition}

\noindent
\cite[Expos\'e IV, Exercice 13.9]{SGA4} では，
(\ref{sites-modules-equation-one-is-never-zero}) が同型であるという条件が欠けており，
局所環付きサイトと局所環付きトポスの概念がわずかに異なる。
ここでは局所環付き空間の概念を動機としているため，この条件を加えることにした
（上の説明を参照）。

\begin{lemma}
\label{sites-modules-lemma-locally-ringed-intrinsic}
局所環付きサイトであることは内在的性質である。より正確には，
\begin{enumerate}
\item $f : \Sh(\mathcal{C}') \to \Sh(\mathcal{C})$ がトポスの射で，
$(\mathcal{C}, \mathcal{O})$ が局所環付きサイトならば，
$(\mathcal{C}', f^{-1}\mathcal{O})$ も局所環付きサイトである。
\item
$(f, f^\sharp) : (\Sh(\mathcal{C}'), \mathcal{O}')
\to (\Sh(\mathcal{C}), \mathcal{O})$
が環付きトポスの同値ならば，
$(\mathcal{C}, \mathcal{O})$ が局所環付きであることと，
$(\mathcal{C}', \mathcal{O}')$ が局所環付きであることは同値である。
\end{enumerate}
\end{lemma}

\begin{proof}
(2) が (1) から従うことは明らかである。(1) を証明する。
$f^{-1}$ は完全なので $f^{-1}* = *$，
$f^{-1}\emptyset^\# = \emptyset^\#$ であり，$f^{-1}$ は積および等化子と
可換し，同型と全射をそれぞれ同型と全射へ写す。
したがって $f^{-1}$ は同型 (\ref{sites-modules-equation-one-is-never-zero}) を
$f^{-1}\mathcal{O}$ に対するその類似へ写し，補題
\ref{sites-modules-lemma-locally-ringed} の (3) の全射を
$f^{-1}\mathcal{O}$ に対する対応する全射へ写す。
\end{proof}

\noindent
実際，補題 \ref{sites-modules-lemma-locally-ringed-intrinsic} の (2) は
《スキーム》の補題 \ref{schemes-lemma-isomorphism-locally-ringed}
の類似である。これにより次の定義が意味をもつ。

\begin{definition}
\label{sites-modules-definition-locally-ringed-topos}
環付きトポス $(\Sh(\mathcal{C}), \mathcal{O})$ が
{\it 局所環付き}であるとは，その基礎にある環付きサイト
$(\mathcal{C}, \mathcal{O})$ が局所環付きであることをいう。
\end{definition}

\noindent
局所環付きであることの帰結の一例を挙げる。

\begin{lemma}
\label{sites-modules-lemma-invertible-is-locally-free-rank-1}
$(\Sh(\mathcal{C}), \mathcal{O})$ を環付きトポスとする。
階数 $1$ の任意の局所自由 $\mathcal{O}$-加群は可逆である。
$(\mathcal{C}, \mathcal{O})$ が局所環付きならば逆も成り立つ
（ただし一般には成り立たない）。
\end{lemma}

\begin{proof}
$\mathcal{L}$ が階数 $1$ の局所自由加群であると仮定し，評価写像
$$
\mathcal{L} \otimes_\mathcal{O}
\SheafHom_\mathcal{O}(\mathcal{L}, \mathcal{O})
\longrightarrow \mathcal{O}
$$
を考える。$\mathcal{C}$ の任意の対象 $U$ に対し，
$\mathcal{L}$ を自明化する被覆の各元へ制限すれば，
この写像が同型であることが分かる（詳細は省略する）。
したがって補題 \ref{sites-modules-lemma-invertible} により $\mathcal{L}$ は可逆である。

\medskip\noindent
$(\Sh(\mathcal{C}), \mathcal{O})$ が局所環付きであると仮定し，
$U$ を $\mathcal{C}$ の対象とする。補題 \ref{sites-modules-lemma-invertible} の証明で，
被覆 $\{U_i \to U\}$ が存在して，
$\mathcal{L}|_{\mathcal{C}/U_i}$ は有限自由
$\mathcal{O}_{U_i}$-加群の直和因子となることを見た。
$U$ を $U_i$ で置き換え，
$p : \mathcal{O}_U^{\oplus r} \to \mathcal{O}_U^{\oplus r}$
を，像が $\mathcal{L}|_{\mathcal{C}/U}$ と同型な射影子とする。
このとき $p$ は行列
$$
P = (p_{ij}) \in \text{Mat}(r \times r, \mathcal{O}(U))
$$
に対応し，これは射影子である：$P^2  = P$。
$A = \mathcal{O}(U)$ と置けば $P \in \text{Mat}(r \times r, A)$ である。
《代数》の補題 \ref{algebra-lemma-finite-projective} により，
$P$ の像は $A$ 上の有限局所自由加群 $M$ である。
したがって，単位イデアルを生成する $f_1, \ldots, f_t \in A$ が存在し，
$M_{f_i}$ は有限自由である。補題 \ref{sites-modules-lemma-locally-ringed} により，
$U$ をある開被覆の各元で置き換えれば，$M$ は自由であると仮定できる。
これは $\mathcal{L}|_U$ が自由であることを意味する（詳細は省略する）。
もちろん $\mathcal{L}$ は可逆なので，
これは $\mathcal{L}|_U$ の階数が $1$ の場合に限って可能であり，
証明が完了する。
\end{proof}

\noindent
次に，局所環付き空間の射をもつとは何を意味するかを明らかにしたい。
そのために次の補題を用いる。

\begin{lemma}
\label{sites-modules-lemma-locally-ringed-morphism}
射
$(f, f^\sharp) : (\Sh(\mathcal{C}), \mathcal{O}_\mathcal{C})
\to (\Sh(\mathcal{D}), \mathcal{O}_\mathcal{D})$
を環付きトポスの射とする。次の条件を考える。
\begin{enumerate}
\item 層の図式
$$
\xymatrix{
f^{-1}(\mathcal{O}^*_\mathcal{D}) \ar[r]_-{f^\sharp} \ar[d] &
\mathcal{O}^*_\mathcal{C} \ar[d] \\
f^{-1}(\mathcal{O}_\mathcal{D}) \ar[r]^-{f^\sharp} &
\mathcal{O}_\mathcal{C}
}
$$
はカルテジアンである。
\item $\mathcal{C}$ の任意の点 $p$ に対して $q = f \circ p$ と置くと，
集合の図式
$$
\xymatrix{
\mathcal{O}^*_{\mathcal{D}, q} \ar[r] \ar[d] &
\mathcal{O}^*_{\mathcal{C}, p} \ar[d] \\
\mathcal{O}_{\mathcal{D}, q} \ar[r] &
\mathcal{O}_{\mathcal{C}, p}
}
$$
はカルテジアンである。
\end{enumerate}
常に (1) $\Rightarrow$ (2) が成り立つ。
$\mathcal{C}$ が十分多くの点をもつならば，(1) と (2) は同値である。
$(\Sh(\mathcal{C}), \mathcal{O}_\mathcal{C})$
と
$(\Sh(\mathcal{D}), \mathcal{O}_\mathcal{D})$
が局所環付きトポスならば，(2) は次と同値である。
\begin{enumerate}
\item[(3)] $\mathcal{C}$ の任意の点 $p$ に対して $q = f \circ p$ と置くと，
環準同型 $\mathcal{O}_{\mathcal{D}, q} \to \mathcal{O}_{\mathcal{C}, p}$ は
局所環準同型である。
\end{enumerate}
実際，点の保守的な族に対して性質 (2) または (3) が成り立てば，
(1) が従う。
\end{lemma}

\begin{proof}
定義を展開すれば直ちに従う。
\end{proof}

\begin{definition}
\label{sites-modules-definition-morphism-locally-ringed-topoi}
$(f, f^\sharp) : (\Sh(\mathcal{C}), \mathcal{O}_\mathcal{C})
\to (\Sh(\mathcal{D}), \mathcal{O}_\mathcal{D})$
を環付きトポスの射とする。
$(\Sh(\mathcal{C}), \mathcal{O}_\mathcal{C})$
と
$(\Sh(\mathcal{D}), \mathcal{O}_\mathcal{D})$
は局所環付きトポスであると仮定する。層の図式
$$
\xymatrix{
f^{-1}(\mathcal{O}^*_\mathcal{D}) \ar[r]_-{f^\sharp} \ar[d] &
\mathcal{O}^*_\mathcal{C} \ar[d] \\
f^{-1}(\mathcal{O}_\mathcal{D}) \ar[r]^-{f^\sharp} &
\mathcal{O}_\mathcal{C}
}
$$
がカルテジアンであることと，$(f, f^\sharp)$ が
{\it 局所環付きトポスの射}であることは同値である
（補題 \ref{sites-modules-lemma-locally-ringed-morphism} 参照）。
$(f, f^\sharp)$ が環付きサイトの射であるとき，付随する環付きトポスの射が
局所環付きトポスの射であるならば，これを
{\it 局所環付きサイトの射}という。
\end{definition}

\noindent
局所環付きトポス間の環付きトポスの同型は，自動的に
局所環付きトポスの同型であることは明らかである。

\begin{lemma}
\label{sites-modules-lemma-composition-morphisms-locally-ringed-topoi}
射
$(f, f^\sharp) :
(\Sh(\mathcal{C}_1), \mathcal{O}_1)
\to (\Sh(\mathcal{C}_2), \mathcal{O}_2)$ と
$(g, g^\sharp) :
(\Sh(\mathcal{C}_2), \mathcal{O}_2) \to
(\Sh(\mathcal{C}_3), \mathcal{O}_3)$
を局所環付きトポスの射とする。このとき合成
$(g, g^\sharp) \circ (f, f^\sharp)$ も局所環付きトポスの射である
（定義 \ref{sites-modules-definition-ringed-topos} 参照）。
\end{lemma}

\begin{proof}
省略する。
\end{proof}

\begin{lemma}
\label{sites-modules-lemma-locally-ringed-intrinsic-morphism}
$f : \Sh(\mathcal{C}') \to \Sh(\mathcal{C})$ をトポスの射とする。
$\mathcal{O}$ が $\mathcal{C}$ 上の環の層ならば，
$$
f^{-1}(\mathcal{O}^*) = (f^{-1}\mathcal{O})^*.
$$
特に，$\mathcal{O}$ によって $\mathcal{C}$ が局所環付きサイトとなるならば，
$f^\sharp = \text{id}$ と置いた環付きトポスの射
$$
(f, f^\sharp) :
(\Sh(\mathcal{C}'), f^{-1}\mathcal{O})
\to
(\Sh(\mathcal{C}, \mathcal{O})
$$
は局所環付きトポスの射である。
\end{lemma}

\begin{proof}
図式
$$
\xymatrix{
\mathcal{O}^* \ar[rr] \ar[d]_{u \mapsto (u, u^{-1})} & &
{*} \ar[d]^{1} \\
\mathcal{O} \times \mathcal{O} \ar[rr]^-{(a, b) \mapsto ab} & &
\mathcal{O}
}
$$
はカルテジアンである。$f^{-1}$ は完全なので，図式
$$
\xymatrix{
f^{-1}(\mathcal{O}^*)
\ar[d]_{u \mapsto (u, u^{-1})} \ar[rr] & &
{*} \ar[d]^{1} \\
f^{-1}\mathcal{O} \times f^{-1}\mathcal{O} \ar[rr]^-{(a, b) \mapsto ab} & &
f^{-1}\mathcal{O}
}
$$
もカルテジアンであり，第一の主張が従う。第二の主張について，
補題 \ref{sites-modules-lemma-locally-ringed-intrinsic} により
$(\mathcal{C}', f^{-1}\mathcal{O})$ は局所環付きサイトなので，
主張は意味をもつ。ここで第一の主張により，この射は局所環付きトポスの射である。
\end{proof}

\begin{lemma}
\label{sites-modules-lemma-localize-morphism-locally-ringed-topoi}
局所環付きサイトと局所環付きトポスの局所化について，次が成り立つ。
\begin{enumerate}
\item $(\mathcal{C}, \mathcal{O})$ を局所環付きサイトとし，
$U$ を $\mathcal{C}$ の対象とする。このとき局所化
$(\mathcal{C}/U, \mathcal{O}_U)$ は局所環付きサイトであり，局所化射
$$
(j_U, j_U^\sharp) :
(\Sh(\mathcal{C}/U), \mathcal{O}_U)
\to
(\Sh(\mathcal{C}), \mathcal{O})
$$
は局所環付きトポスの射である。
\item $(\mathcal{C}, \mathcal{O})$ を局所環付きサイトとし，
$f : V \to U$ を $\mathcal{C}$ の射とする。このとき，補題
\ref{sites-modules-lemma-relocalize} の射
$$
(j, j^\sharp) :
(\Sh(\mathcal{C}/V), \mathcal{O}_V)
\to
(\Sh(\mathcal{C}/U), \mathcal{O}_U)
$$
は局所環付きトポスの射である。
\item
$(f, f^\sharp) :
(\mathcal{C}, \mathcal{O})
\longrightarrow
(\mathcal{D}, \mathcal{O}')$
を局所環付きサイトの射とし，$f$ は連続関手
$u : \mathcal{D} \to \mathcal{C}$ で与えられるとする。
$V$ を $\mathcal{D}$ の対象とし，$U = u(V)$ と置く。このとき，補題
\ref{sites-modules-lemma-localize-morphism-ringed-sites} の射
$$
(f', (f')^\sharp) :
(\Sh(\mathcal{C}/U), \mathcal{O}_U)
\to
(\Sh(\mathcal{D}/V), \mathcal{O}'_V)
$$
は局所環付きサイトの射である。
\item
$(f, f^\sharp) :
(\mathcal{C}, \mathcal{O})
\longrightarrow
(\mathcal{D}, \mathcal{O}')$
を局所環付きサイトの射とし，$f$ は連続関手
$u : \mathcal{D} \to \mathcal{C}$ で与えられるとする。
$V \in \Ob(\mathcal{D})$，$U \in \Ob(\mathcal{C})$，
$c : U \to u(V)$ とする。このとき，補題
\ref{sites-modules-lemma-relocalize-morphism-ringed-sites} の射
$$
(f_c, (f_c)^\sharp) :
(\Sh(\mathcal{C}/U), \mathcal{O}_U)
\to
(\Sh(\mathcal{D}/V), \mathcal{O}'_V)
$$
は局所環付きトポスの射である。
\item $(\Sh(\mathcal{C}), \mathcal{O})$ を局所環付きトポスとし，
$\mathcal{F}$ を $\mathcal{C}$ 上の層とする。このとき局所化
$(\Sh(\mathcal{C})/\mathcal{F}, \mathcal{O}_\mathcal{F})$ は
局所環付きトポスであり，局所化射
$$
(j_\mathcal{F}, j_\mathcal{F}^\sharp) :
(\Sh(\mathcal{C})/\mathcal{F}, \mathcal{O}_\mathcal{F})
\to
(\Sh(\mathcal{C}), \mathcal{O})
$$
は局所環付きトポスの射である。
\item $(\Sh(\mathcal{C}), \mathcal{O})$ を局所環付きトポスとし，
$s : \mathcal{G} \to \mathcal{F}$ を $\mathcal{C}$ 上の層の写像とする。
このとき，補題 \ref{sites-modules-lemma-relocalize-ringed-topos} の射
$$
(j, j^\sharp) :
(\Sh(\mathcal{C})/\mathcal{G}, \mathcal{O}_\mathcal{G})
\longrightarrow
(\Sh(\mathcal{C})/\mathcal{F}, \mathcal{O}_\mathcal{F})
$$
は局所環付きトポスの射である。
\item
$f :
(\Sh(\mathcal{C}), \mathcal{O})
\longrightarrow
(\Sh(\mathcal{D}), \mathcal{O}')$
を局所環付きトポスの射とする。$\mathcal{G}$ を $\mathcal{D}$ 上の層とし，
$\mathcal{F} = f^{-1}\mathcal{G}$ と置く。このとき，補題
\ref{sites-modules-lemma-localize-morphism-ringed-topoi} の射
$$
(f', (f')^\sharp) :
(\Sh(\mathcal{C})/\mathcal{F}, \mathcal{O}_\mathcal{F})
\longrightarrow
(\Sh(\mathcal{D})/\mathcal{G}, \mathcal{O}'_\mathcal{G})
$$
は局所環付きトポスの射である。
\item
$f :
(\Sh(\mathcal{C}), \mathcal{O})
\longrightarrow
(\Sh(\mathcal{D}), \mathcal{O}')$
を局所環付きトポスの射とする。$\mathcal{G}$ を $\mathcal{D}$ 上の層，
$\mathcal{F}$ を $\mathcal{C}$ 上の層とし，
$s : \mathcal{F} \to f^{-1}\mathcal{G}$ を層の射とする。このとき，補題
\ref{sites-modules-lemma-relocalize-morphism-ringed-topoi} の射
$$
(f_s, (f_s)^\sharp) :
(\Sh(\mathcal{C})/\mathcal{F}, \mathcal{O}_\mathcal{F})
\longrightarrow
(\Sh(\mathcal{D})/\mathcal{G}, \mathcal{O}'_\mathcal{G})
$$
は局所環付きトポスの射である。
\end{enumerate}
\end{lemma}

\begin{proof}
(1) は明らかである。実際，$\mathcal{O}_U$ は
$\mathcal{O}$ の制限にすぎないので，補題
\ref{sites-modules-lemma-locally-ringed-intrinsic} と
\ref{sites-modules-lemma-locally-ringed-intrinsic-morphism} を適用できる。
(2) も明らかである。射 $(j, j^\sharp)$ は実際に局所環付きサイトの
局所化なので，(1) を適用できる。
(3) も明らかである。実際，$(f')^\sharp$ は
$f^\sharp$ のトポス $\Sh(\mathcal{C})/\mathcal{F}$ への制限にすぎない。
補題 \ref{sites-modules-lemma-localize-morphism-ringed-topoi} の証明を参照
（したがって射 $f'$ に対する定義
\ref{sites-modules-definition-morphism-locally-ringed-topoi} の図式は，
$f$ に対する対応する図式の制限にすぎず，制限は完全関手である）。
(4) は (2) と (3) を合わせれば形式的に従う。
(5)，(6)，(7)，(8) は，補題
\ref{sites-modules-lemma-morphism-ringed-topoi-comes-from-morphism-ringed-sites}
のようにサイトを拡大し，補題
\ref{sites-modules-lemma-localize-compare}，
\ref{sites-modules-lemma-relocalize-compare}，
\ref{sites-modules-lemma-localize-morphism-compare}，および
\ref{sites-modules-lemma-relocalize-morphism-compare} の比較を用いて，
すべてをサイトとサイトの射の言葉に翻訳することにより，
それぞれ対応する (1)，(2)，(3)，(4) から従う。
（別法として，(1)，(2)，(3)，(4) の証明と同じ議論を用いて
(5)，(6)，(7)，(8) を直接証明してもよい。）
\end{proof}






\section{加群に対する関手 $g_!$}
\label{sites-modules-section-lower-shriek-modules}

\noindent
本節では，\ref{sites-modules-section-exactness-lower-shriek} 節で論じた
$g_!$ の構成を加群の層の場合へ拡張する。

\begin{lemma}
\label{sites-modules-lemma-lower-shriek-modules}
$u : \mathcal{C} \to \mathcal{D}$ をサイト間の連続かつ余連続な関手とする。
付随するトポスの射を
$g : \Sh(\mathcal{C}) \to \Sh(\mathcal{D})$ と書く。
$\mathcal{O}_\mathcal{D}$ を $\mathcal{D}$ 上の環の層とし，
$\mathcal{O}_\mathcal{C} = g^{-1}\mathcal{O}_\mathcal{D}$ と置く。
すると $g$ は $g^* = g^{-1}$ を満たす環付きトポスの射となる。
この場合，関手
$$
g_! :
\textit{Mod}(\mathcal{O}_\mathcal{C})
\longrightarrow
\textit{Mod}(\mathcal{O}_\mathcal{D})
$$
で $g^*$ の左随伴となるものが存在する。
\end{lemma}

\begin{proof}
$U$ を $\mathcal{C}$ の対象とする。任意の
$\mathcal{O}_\mathcal{D}$-加群 $\mathcal{G}$ に対して
\begin{align*}
\Hom_{\mathcal{O}_\mathcal{C}}(j_{U!}\mathcal{O}_U, g^{-1}\mathcal{G})
& =
g^{-1}\mathcal{G}(U) \\
& =
\mathcal{G}(u(U)) \\
& =
\Hom_{\mathcal{O}_\mathcal{D}}(j_{u(U)!}\mathcal{O}_{u(U)}, \mathcal{G})
\end{align*}
が成り立つ。これは $g^{-1}$ が制限によって記述されるからである。
《サイト》の補題 \ref{sites-lemma-when-shriek} を参照。
もちろん，同様の公式は $j_{U!}\mathcal{O}_U$ の形の加群の直和に対しても
成り立つ。《ホモロジー》の補題
\ref{homology-lemma-partially-defined-adjoint} と補題
\ref{sites-modules-lemma-module-quotient-flat} により，$g_!$ が存在することが分かる。
\end{proof}

\begin{remark}
\label{sites-modules-remark-when-shriek-equal}
注意！$u : \mathcal{C} \to \mathcal{D}$，$g$，
$\mathcal{O}_\mathcal{D}$，$\mathcal{O}_\mathcal{C}$ を補題
\ref{sites-modules-lemma-lower-shriek-modules} のものとする。一般には図式
$$
\xymatrix{
\textit{Mod}(\mathcal{O}_\mathcal{C}) \ar[r]_{g_!} \ar[d]_{forget} &
\textit{Mod}(\mathcal{O}_\mathcal{D}) \ar[d]^{forget} \\
\textit{Ab}(\mathcal{C}) \ar[r]^{g^{Ab}_!} &
\textit{Ab}(\mathcal{D})
}
$$
は可換{\bf ではない}（ここで $g^{Ab}_!$ は補題
\ref{sites-modules-lemma-g-shriek-adjoint} のものである）。関手の変換
$$
g_!^{Ab} \circ forget \longrightarrow forget \circ g_!
$$
が存在する。補題 \ref{sites-modules-lemma-lower-shriek-modules} の証明から，
これが同型であることと，$\mathcal{C}$ のすべての対象 $U$ に対して
$g^{Ab}_!j_{U!}\mathcal{O}_U \to g_!j_{U!}\mathcal{O}_U$
が同型であることは同値である。
$g_!j_{U!}\mathcal{O}_U = j_{u(U)!}\mathcal{O}_{u(U)}$ なので，
これはすべての対象 $U$ に対して
$$
g^{Ab}_!j_{U!}\mathcal{O}_U \longrightarrow j_{u(U)!}\mathcal{O}_{u(U)}
$$
が同型であることと同値である。このような $U$ に対して，
サイトの余連続関手の可換図式
$$
\xymatrix{
\mathcal{C}/U \ar[r]_-{u'} \ar[d]_{j_U} & \mathcal{D}/u(U) \ar[d]^{j_{u(U)}} \\
\mathcal{C} \ar[r]^u & \mathcal{D}
}
$$
を得る。《サイト》の補題 \ref{sites-lemma-localize-cocontinuous} を参照。
したがって $g^{Ab}_!j_{U!} = j_{u(U)!}(g')^{Ab}_!$ である。
ここで $g' : \Sh(\mathcal{C}/U) \to \Sh(\mathcal{D}/u(U))$ は
余連続関手 $u'$ が誘導するトポスの射である。
ゆえに，$\mathcal{C}$ のすべての対象 $U$ に対して標準写像
\begin{equation}
\label{sites-modules-equation-compare-on-localizations}
(g')^{Ab}_!\mathcal{O}_U \longrightarrow \mathcal{O}_{u(U)}
\end{equation}
が同型ならば $g_! = g^{Ab}_!$ である。
\end{remark}

\noindent
次の二つの結果は少し性質が異なる。

\begin{lemma}
\label{sites-modules-lemma-special-square-cocontinuous}
環付きトポスの可換図式
$$
\xymatrix{
(\Sh(\mathcal{C}'), \mathcal{O}_{\mathcal{C}'})
\ar[r]_{(g', (g')^\sharp)} \ar[d]_{(f', (f')^\sharp)} &
(\Sh(\mathcal{C}), \mathcal{O}_\mathcal{C}) \ar[d]^{(f, f^\sharp)} \\
(\Sh(\mathcal{D}'), \mathcal{O}_{\mathcal{D}'}) \ar[r]^{(g, g^\sharp)} &
(\Sh(\mathcal{D}), \mathcal{O}_\mathcal{D})
}
$$
が与えられたとし，次を仮定する。
\begin{enumerate}
\item $f$，$f'$，$g$，$g'$ は，《サイト》の補題
\ref{sites-lemma-cocontinuous-morphism-topoi} のように，
それぞれ余連続関手 $u$，$u'$，$v$，$v'$ に対応する。
\item $v \circ u' = u \circ v'$。
\item $v$ と $v'$ は余連続であると同時に連続である。
\item $\mathcal{D}'$ の任意の対象 $V'$ に対して，$v$ が与える関手
${}^{u'}_{V'}\mathcal{I} \to {}^{\ \ \ u}_{v(V')}\mathcal{I}$
は共終である。
\item $g^{-1}\mathcal{O}_{\mathcal{D}} = \mathcal{O}_{\mathcal{D}'}$
かつ $(g')^{-1}\mathcal{O}_{\mathcal{C}} = \mathcal{O}_{\mathcal{C}'}$。
\end{enumerate}
このとき加群上で $f'_* \circ (g')^* = g^* \circ f_*$ および
$g'_! \circ (f')^{-1} = f^{-1} \circ g_!$ が成り立つ。
\end{lemma}

\begin{proof}
条件 (5) により $(g')^*\mathcal{F} = (g')^{-1}\mathcal{F}$ および
$g^*\mathcal{G} = g^{-1}\mathcal{G}$ である。
したがって第一の等式は《サイト》の補題
\ref{sites-lemma-special-square-cocontinuous} の対応する等式から
直ちに従う。補題 \ref{sites-modules-lemma-lower-shriek-modules} により，
$g^*$ と $(g')^*$ の左随伴関手 $g_!$ と $g'_!$ が存在するので，
第二の等式は随伴関手の一意性から従う。
\end{proof}

\begin{lemma}
\label{sites-modules-lemma-special-square-continuous}
環付きトポスの可換図式
$$
\xymatrix{
(\Sh(\mathcal{C}'), \mathcal{O}_{\mathcal{C}'})
\ar[r]_{(g', (g')^\sharp)} \ar[d]_{(f', (f')^\sharp)} &
(\Sh(\mathcal{C}), \mathcal{O}_\mathcal{C}) \ar[d]^{(f, f^\sharp)} \\
(\Sh(\mathcal{D}'), \mathcal{O}_{\mathcal{D}'}) \ar[r]^{(g, g^\sharp)} &
(\Sh(\mathcal{D}), \mathcal{O}_\mathcal{D})
}
$$
と関手
$$
\xymatrix{
\mathcal{C}' \ar[r]_{v'} &
\mathcal{C} \\
\mathcal{D}' \ar[r]^v \ar[u]^{u'} &
\mathcal{D} \ar[u]_u
}
$$
が与えられ，次を満たすとする（記法は《サイト》の
\ref{sites-section-morphism-sites} 節および
\ref{sites-section-cocontinuous-morphism-topoi} 節のものとする）。
\begin{enumerate}
\item $u$ と $u'$ は連続で，射 $f$ と $f'$ を誘導する。
\item $v$ と $v'$ は余連続で，射 $g$ と $g'$ を誘導する。
\item $u \circ v = v' \circ u'$。
\item $v$ と $v'$ は余連続であると同時に連続である。
\item $g^{-1}\mathcal{O}_{\mathcal{D}} = \mathcal{O}_{\mathcal{D}'}$
かつ $(g')^{-1}\mathcal{O}_{\mathcal{C}} = \mathcal{O}_{\mathcal{C}'}$。
\end{enumerate}
このとき加群上で $f'_* \circ (g')^* = g^* \circ f_*$ および
$g'_! \circ (f')^{-1} = f^{-1} \circ g_!$ が成り立つ。
\end{lemma}

\begin{proof}
条件 (5) により $(g')^*\mathcal{F} = (g')^{-1}\mathcal{F}$ および
$g^*\mathcal{G} = g^{-1}\mathcal{G}$ である。
したがって第一の等式は《サイト》の補題
\ref{sites-lemma-special-square-continuous} の対応する等式から直ちに従う。
補題 \ref{sites-modules-lemma-lower-shriek-modules} により，
$g^*$ と $(g')^*$ の左随伴関手 $g_!$ と $g'_!$ が存在するので，
第二の等式は随伴関手の一意性から従う。
\end{proof}







\section{定数層}
\label{sites-modules-section-constant}

\noindent
$E$ を集合，$\mathcal{C}$ をサイトとする。値 $E$ の定数層を
$\underline{E}$ と書く。《サイト》の例
\ref{sites-example-constant-sheaf} を参照。
$E$ がアーベル群，環，加群などならば，
$\underline{E}$ はそれぞれアーベル群，環，加群などの層である。

\begin{lemma}
\label{sites-modules-lemma-constant-exact}
サイト $\mathcal{C}$ 上で，アーベル群の短完全列
$0 \to A \to B \to C \to 0$ が与えられたとする。このとき
$0 \to \underline{A} \to \underline{B} \to \underline{C} \to 0$
はアーベル層の完全列であり，実際にはアーベル前層の列としても完全である。
\end{lemma}

\begin{proof}
層化は完全なので，
$0 \to \underline{A} \to \underline{B} \to \underline{C} \to 0$
がアーベル層の完全列であることは明らかである。したがって
$0 \to \underline{A} \to \underline{B} \to \underline{C}$
はアーベル前層の完全列である。
$\underline{B} \to \underline{C}$ が全射であることを見るため，
集合論的切断 $s : C \to B$ を選ぶ。これは集合の層の切断
$\underline{s} : \underline{C} \to \underline{B}$ を誘導し，
全射 $\underline{B} \to \underline{C}$ の左逆となる。
\end{proof}

\begin{lemma}
\label{sites-modules-lemma-tensor-with-finitely-presented}
サイト $\mathcal{C}$，環 $\Lambda$，$\Lambda$-加群 $M$ と $Q$ を取る。
$Q$ が有限表示 $\Lambda$-加群ならば，すべての
$U \in \Ob(\mathcal{C})$ に対して
$\underline{M \otimes_\Lambda Q}(U) = \underline{M}(U) \otimes_\Lambda Q$
が成り立つ。
\end{lemma}

\begin{proof}
表示 $\Lambda^{\oplus m} \to \Lambda^{\oplus n} \to Q \to 0$ を選ぶ。
これは完全列
$M^{\oplus m} \to M^{\oplus n} \to M \otimes Q \to 0$ を与える。
補題 \ref{sites-modules-lemma-constant-exact} により完全列
$$
\underline{M}(U)^{\oplus m} \to
\underline{M}(U)^{\oplus n} \to
\underline{M \otimes Q}(U) \to 0
$$
を得るので，補題が従う
（$U$ 上の切断を取ることは常に有限直和と可換するが，
任意直和とは可換しないことに注意する）。
\end{proof}

\begin{lemma}
\label{sites-modules-lemma-flat-sections}
サイト $\mathcal{C}$，連接環 $\Lambda$，平坦 $\Lambda$-加群 $M$ を取る。
$U \in \Ob(\mathcal{C})$ に対して，加群 $\underline{M}(U)$ は
平坦な $\Lambda$-加群である。
\end{lemma}

\begin{proof}
有限生成イデアル $I \subset \Lambda$ を取る。《代数》の補題
\ref{algebra-lemma-flat} により，写像
$\underline{M}(U) \otimes_\Lambda I \to \underline{M}(U)$
が単射であることを示せば十分である。
$\Lambda$ は連接なので，$I$ は $\Lambda$-加群として有限表示である。
補題 \ref{sites-modules-lemma-tensor-with-finitely-presented} により
$\underline{M}(U) \otimes I = \underline{M \otimes I}$ である。
$M$ は平坦なので写像 $M \otimes I \to M$ は単射であり，
したがって $\underline{M \otimes I} \to \underline{M}$ も単射である。
\end{proof}

\begin{lemma}
\label{sites-modules-lemma-completion-flat}
サイト $\mathcal{C}$，ネーター環 $\Lambda$，
イデアル $I \subset \Lambda$ を取る。層
$\underline{\Lambda}^\wedge = \lim \underline{\Lambda/I^n}$
は平坦な $\underline{\Lambda}$-代数である。
さらに標準的な同一視
$$
\underline{\Lambda}/I\underline{\Lambda} =
\underline{\Lambda}/\underline{I} =
\underline{\Lambda}^\wedge/I\underline{\Lambda}^\wedge =
\underline{\Lambda}^\wedge/\underline{I} \cdot \underline{\Lambda}^\wedge =
\underline{\Lambda}^\wedge/\underline{I}^\wedge =
\underline{\Lambda/I}
$$
が成り立つ。ここで
$\underline{I}^\wedge = \lim \underline{I/I^n}$ である。
\end{lemma}

\begin{proof}
$\underline{\Lambda}^\wedge$ が平坦であることを証明するには，
各 $U \in \Ob(\mathcal{C})$ に対して
$\underline{\Lambda}^\wedge(U)$ が $\Lambda$-加群として平坦であることを
示せば十分である。補題 \ref{sites-modules-lemma-flatness-presheaves} と
\ref{sites-modules-lemma-flatness-sheafification} を参照。
補題 \ref{sites-modules-lemma-flat-sections} により，
$$
\underline{\Lambda}^\wedge(U) = \lim \underline{\Lambda/I^n}(U)
$$
は平坦な $\Lambda/I^n$-加群の系の極限である。
補題 \ref{sites-modules-lemma-constant-exact} により遷移写像は全射である。
《代数詳論》の補題 \ref{more-algebra-lemma-limit-flat} から結論を得る。

\medskip\noindent
等式を示すため，補題 \ref{sites-modules-lemma-tensor-with-finitely-presented} による等式
$\underline{\Lambda}(U)/I\underline{\Lambda}(U) = \underline{\Lambda/I}(U)$
に注意する。これから
$\underline{\Lambda}/I\underline{\Lambda} =
\underline{\Lambda}/\underline{I} = \underline{\Lambda/I}$ が従う。
短完全列の系
$$
0 \to \underline{I/I^n}(U) \to \underline{\Lambda/I^n}(U) \to
\underline{\Lambda/I}(U) \to 0
$$
は全射な遷移写像をもち，したがって短完全列
$$
0 \to \lim \underline{I/I^n}(U) \to \lim \underline{\Lambda/I^n}(U) \to
\lim \underline{\Lambda/I}(U) \to 0
$$
を与える。《ホモロジー》の補題
\ref{homology-lemma-Mittag-Leffler} を参照。したがって
$\underline{\Lambda}^\wedge/\underline{I}^\wedge =
\underline{\Lambda/I}$ である。
$$
I \underline{\Lambda}^\wedge \subset
\underline{I} \cdot \underline{\Lambda}^\wedge \subset
\underline{I}^\wedge
$$
なので，すべての $U$ に対して
$I \underline{\Lambda}^\wedge(U) = \underline{I}^\wedge(U)$
を示せば十分である。生成元 $I = (f_1, \ldots, f_r)$ を選ぶ。
各 $n$ に対して短完全列
$$
0 \to K_n/(I^n)^{\oplus r} \to (\Lambda/I^n)^{\oplus r}
\xrightarrow{(f_1, \ldots, f_r)}
I/I^{n + 1} \to 0
$$
を得る。ここで
$K_n =
\{(x_1, \ldots, x_r) \in \Lambda^{\oplus r} \mid \sum x_i f_i \in I^{n + 1}\}$
である。短完全列
$$
0 \to \underline{K_n/(I^n)^{\oplus r}}(U) \to
\underline{(\Lambda/I^n)^{\oplus r}}(U) \to
\underline{I/I^{n + 1}}(U) \to 0
$$
を得る。計算により $K_n = K + (I^n)^{\oplus r}$ なので，遷移写像
$K_{n + 1}/(I^{n + 1})^{\oplus r} \to K_n/(I^n)^{\oplus r}$ は全射である。
したがって左辺の加群の系は全射な遷移写像をもち，特に ML 条件を満たす。
ゆえに
$(f_1, \ldots, f_r) :
(\underline{\Lambda}^\wedge)^{\oplus r}(U) \to \underline{I}^\wedge(U)$
は《ホモロジー》の補題 \ref{homology-lemma-Mittag-Leffler} により全射である。
これが示すべきことであった。
\end{proof}

\begin{lemma}
\label{sites-modules-lemma-locally-constant-finite-type}
サイト $\mathcal{C}$，環 $\Lambda$，$\Lambda$-加群 $M$ を取る。
$\Sh(\mathcal{C})$ は空トポスでないと仮定する。このとき，
\begin{enumerate}
\item $\underline{M}$ が有限型 $\underline{\Lambda}$-加群の層であることと，
$M$ が有限 $\Lambda$-加群であることは同値である。
\item $\underline{M}$ が有限表示 $\underline{\Lambda}$-加群の層であることと，
$M$ が有限表示 $\Lambda$-加群であることは同値である。
\end{enumerate}
\end{lemma}

\begin{proof}
(1) を証明する。$M$ が $x_1, \ldots, x_r$ で生成されるならば，
$x_1, \ldots, x_r$ は $\underline{M}$ を生成する大域切断を定めるので，
$\underline{M}$ は有限型である。逆に $\underline{M}$ が有限型であると仮定する。
$U \in \mathcal{C}$ を層論的に空でない対象とする
（《サイト》の定義 \ref{sites-definition-empty}）。
$\Sh(\mathcal{C})$ が空トポスでないと仮定したので，このような対象は存在する。
このとき被覆 $\{U_i \to U\}$ と有限個の切断
$s_{ij} \in \underline{M}(U_i)$ が存在し，
$\underline{M}|_{U_i}$ を生成する。被覆を細分すれば，
$s_{ij}$ は $M$ の元 $x_{ij}$ から来ると仮定できる。
すると $x_{ij}$ は $\underline{M}$ の大域切断を定め，
その $U$ への制限は $\underline{M}$ を生成する。

\medskip\noindent
$M$ の元 $x_1, \ldots, x_r$ が $\underline{M}$ の大域切断を定め，
$\underline{\Lambda}$-加群の層として $\underline{M}$ を生成すると仮定する。
$x_1, \ldots, x_r$ が $\Lambda$-加群として $M$ を生成することを示す。
$x \in M$ とする。被覆 $\{U_i \to U\}_{i \in I}$ と
$f_{i, j} \in \underline{\Lambda}(U_i)$ を選んで，
$x|_{U_i} = \sum f_{i, j} x_j|_{U_i}$ とできる。
被覆を細分すれば $f_{i, j} \in \Lambda$ と仮定できる。
$U$ は層論的に空でないので，ある $i \in I$ に対して
$U_i$ は層論的に空でない。このとき写像
$M \to \underline{M}(U_i)$ は単射である（詳細は省略する）。
したがって，所望のとおり $M$ において
$x = \sum f_{i, j}x_j$ である。

\medskip\noindent
(2) を証明する。$\underline{M}$ が有限表示
$\underline{\Lambda}$-加群であると仮定する。(1) により $M$ は有限型である。
$M$ の $\Lambda$-加群としての生成元 $x_1, \ldots, x_r$ を選ぶ。
これは短完全列
$0 \to K \to \Lambda^{\oplus r} \to M \to 0$ を定め，補題
\ref{sites-modules-lemma-constant-exact} により短完全列
$$
0 \to \underline{K} \to \underline{\Lambda}^{\oplus r} \to \underline{M} \to 0
$$
へ移る。補題
\ref{sites-modules-lemma-kernel-surjection-finite-onto-finite-presentation} により，
$\underline{K}$ は有限型である。したがって (1) により
$K$ は有限 $\Lambda$-加群である。ゆえに $M$ は有限表示
$\Lambda$-加群である。
\end{proof}




\section{局所定数層}
\label{sites-modules-section-locally-constant}

\noindent
一般的な定義は次のとおりである。

\begin{definition}
\label{sites-modules-definition-locally-constant}
$\mathcal{C}$ をサイトとし，$\mathcal{F}$ を集合，群，アーベル群，環，
固定した環 $\Lambda$ 上の加群などの層とする。
\begin{enumerate}
\item $\mathcal{F}$ が，集合，群，アーベル群，環，固定した環
$\Lambda$ 上の加群などの層として，\ref{sites-modules-section-constant} 節の定数層
$\underline{E}$ と同型であるとき，$\mathcal{F}$ を集合，群，アーベル群，
環，固定した環 $\Lambda$ 上の加群などの{\it 定数層}という。
\item $\mathcal{C}$ の任意の対象 $U$ に対して被覆
$\{U_i \to U\}$ が存在し，$\mathcal{F}|_{U_i}$ が定数層となるとき，
$\mathcal{F}$ を{\it 局所定数}という。
\item $\mathcal{F}$ が集合または群の層であるとする。
その定数値が有限集合または有限群であるとき，$\mathcal{F}$ を
{\it 有限局所定数}という。
\end{enumerate}
\end{definition}

\begin{lemma}
\label{sites-modules-lemma-pullback-locally-constant}
$f : \Sh(\mathcal{C}) \to \Sh(\mathcal{D})$ をトポスの射とする。
$\mathcal{G}$ が $\mathcal{D}$ 上の集合，群，アーベル群，環，
固定した環 $\Lambda$ 上の加群などの局所定数層ならば，
$f^{-1}\mathcal{G}$ も $\mathcal{C}$ 上の同じ種類の局所定数層である。
\end{lemma}

\begin{proof}
省略する。
\end{proof}

\begin{lemma}
\label{sites-modules-lemma-morphism-locally-constant}
$\mathcal{C}$ を終対象 $X$ をもつサイトとする。
\begin{enumerate}
\item $\varphi : \mathcal{F} \to \mathcal{G}$ を
$\mathcal{C}$ 上の集合の局所定数層の写像とする。
$\mathcal{F}$ が有限局所定数ならば，被覆 $\{U_i \to X\}$ が存在し，
$\varphi|_{U_i}$ は集合の写像に付随する定数層の写像である。
\item $\varphi : \mathcal{F} \to \mathcal{G}$ を
$\mathcal{C}$ 上のアーベル群の局所定数層の写像とする。
$\mathcal{F}$ が有限局所定数ならば，被覆 $\{U_i \to X\}$ が存在し，
$\varphi|_{U_i}$ はアーベル群の写像に付随するアーベル定数層の写像である。
\item $\Lambda$ を環とし，$\varphi : \mathcal{F} \to \mathcal{G}$ を
$\mathcal{C}$ 上の $\Lambda$-加群の局所定数層の写像とする。
$\mathcal{F}$ が有限型ならば，被覆 $\{U_i \to X\}$ が存在し，
$\varphi|_{U_i}$ は $\Lambda$-加群の写像に付随する
$\Lambda$-加群の定数層の写像である。
\end{enumerate}
\end{lemma}

\begin{proof}
証明は省略する。
\end{proof}

\begin{lemma}
\label{sites-modules-lemma-locally-constant}
$\mathcal{C}$ をサイト，$\Lambda$ を環とする。
$M$，$N$ を $\Lambda$-加群とし，$\mathcal{F}$，$\mathcal{G}$ を
$\Lambda$-加群の局所定数層とする。
\begin{enumerate}
\item $M$ が有限表示ならば，
$$
\underline{\Hom_\Lambda(M, N)} =
\SheafHom_{\underline{\Lambda}}(\underline{M}, \underline{N})
$$
が成り立つ。
\item $M$ と $N$ がともに有限表示ならば，
$$
\underline{\text{Isom}_\Lambda(M, N)} =
\mathit{Isom}_{\underline{\Lambda}}(\underline{M}, \underline{N})
$$
が成り立つ。
\item $\mathcal{F}$ が有限表示ならば，
$\SheafHom_{\underline{\Lambda}}(\mathcal{F}, \mathcal{G})$
は $\Lambda$-加群の局所定数層である。
\item $\mathcal{F}$ と $\mathcal{G}$ がともに有限表示ならば，
$\mathit{Isom}_{\underline{\Lambda}}(\mathcal{F}, \mathcal{G})$
は集合の局所定数層である。
\end{enumerate}
\end{lemma}

\begin{proof}
(1) を証明する。$E = \Hom_\Lambda(M, N)$ と置く。標準写像
$$
\underline{E}
\longrightarrow
\SheafHom_{\underline{\Lambda}}(\underline{M}, \underline{N})
$$
が同型であることを示したい。加群 $M$ は表示
$\Lambda^{\oplus s} \to \Lambda^{\oplus t} \to M \to 0$ をもつ。
すると $E$ は完全列
$$
0 \to E \to \Hom_\Lambda(\Lambda^{\oplus t}, N) \to
\Hom_\Lambda(\Lambda^{\oplus s}, N)
$$
に入り，同様に完全列
$$
0 \to
\SheafHom_{\underline{\Lambda}}(\underline{M}, \underline{N}) \to
\SheafHom_{\underline{\Lambda}}(\underline{\Lambda^{\oplus t}}, \underline{N})
\to
\SheafHom_{\underline{\Lambda}}(\underline{\Lambda^{\oplus s}}, \underline{N})
$$
を得る。これにより問題は $M$ が有限自由加群の場合へ帰着し，
その場合は明らかである。

\medskip\noindent
(3) を証明する。問題は $\mathcal{C}$ 上で局所的なので，ある
$\Lambda$-加群 $M$ と $N$ に対して
$\mathcal{F} = \underline{M}$ および
$\mathcal{G} = \underline{N}$ と仮定してよい。
補題 \ref{sites-modules-lemma-locally-constant-finite-type} により
加群 $M$ は有限表示である。したがって (1) から従う。

\medskip\noindent
(2) と (4) は (1) と (3)，および $\mathit{Isom}$ が
$\SheafHom_{\underline{\Lambda}}(\mathcal{F}, \mathcal{G})$
の切断のうち
$\SheafHom_{\underline{\Lambda}}(\mathcal{G}, \mathcal{F})$
に逆をもつものからなる部分層とみなせることから従う。
\end{proof}

\begin{lemma}
\label{sites-modules-lemma-kernel-finite-locally-constant}
$\mathcal{C}$ をサイトとする。
\begin{enumerate}
\item 集合の有限局所定数層の圏は，$\Sh(\mathcal{C})$ の中で
有限極限と有限余極限について閉じている。
\item 有限局所定数アーベル層の圏は
$\textit{Ab}(\mathcal{C})$ の弱 Serre 部分圏である。
\item $\Lambda$ をネーター環とする。$\mathcal{C}$ 上の有限型な
$\Lambda$-加群の局所定数層の圏は
$\textit{Mod}(\mathcal{C}, \Lambda)$ の弱 Serre 部分圏である。
\end{enumerate}
\end{lemma}

\begin{proof}
(1) を証明する。$\mathcal{C}$ 上で局所的に作業してよい。
したがって補題 \ref{sites-modules-lemma-morphism-locally-constant} により，
有限集合の有限図式が与えられ，層の図式はそれに付随する
定数層の図式であると仮定できる。そこで集合の圏で極限または余極限を取り，
付随する定数層を取ればよい。細部の一部は省略する。

\medskip\noindent
(2) と (3) を証明するため，《ホモロジー》の補題
\ref{homology-lemma-characterize-weak-serre-subcategory} の判定法を用いる。
核と余核の存在は上と同様に論じられる。もちろん (3) でネーター環を
用いる理由は，有限 $\Lambda$-加群間の写像の核が有限
$\Lambda$-加群であることを保証するためである。
圏が拡大について閉じていることを見るため
（$\Lambda$-加群の層の場合），$\mathcal{C}$ 上の
$\Lambda$-加群の層の拡大
$$
0 \to \mathcal{F} \to \mathcal{E} \to \mathcal{G} \to 0
$$
が与えられ，$\mathcal{F}$ と $\mathcal{G}$ は有限型かつ局所定数で
あると仮定する。$\mathcal{C}$ 上で局所化すれば，
$\mathcal{F}$ と $\mathcal{G}$ は定数層であると仮定できる。すなわち，
ある $\Lambda$-加群 $M$ と $N$ に対して
$$
0 \to \underline{M} \to \mathcal{E} \to \underline{N} \to 0
$$
を得る。$N$ の生成元 $y_1, \ldots, y_m$ を選ぶと，
$\Lambda$-加群の短完全列
$0 \to K \to \Lambda^{\oplus m} \to N \to 0$ を得る。
さらに局所化すれば，各 $y_j$ は $\mathcal{E}$ の切断 $s_j$ へ
持ち上がると仮定できる。したがって，次の図式において
$\mathcal{E}$ は押し出しである。
$$
\xymatrix{
0 \ar[r] &
\underline{K} \ar[d] \ar[r] &
\underline{\Lambda^{\oplus m}} \ar[d] \ar[r] &
\underline{N} \ar[d] \ar[r] & 0 \\
0 \ar[r] &
\underline{M} \ar[r] &
\mathcal{E} \ar[r] &
\underline{N} \ar[r] & 0
}
$$
再び補題 \ref{sites-modules-lemma-morphism-locally-constant} により
（また，$\Lambda$ はネーターなので $K$ は有限 $\Lambda$-加群であるから），
写像 $\underline{K} \to \underline{M}$ は局所定数であることが分かり，
結論を得る。
\end{proof}

\begin{lemma}
\label{sites-modules-lemma-tensor-product-locally-constant}
$\mathcal{C}$ をサイト，$\Lambda$ を環とする。
$\mathcal{C}$ 上の二つの $\Lambda$-加群の局所定数層のテンソル積は，
$\Lambda$-加群の局所定数層である。
\end{lemma}

\begin{proof}
省略する。
\end{proof}











\section{環の層の局所化}
\label{sites-modules-section-localizing-sheaves-rings}

\noindent
$(\mathcal{C}, \mathcal{O})$ を環付きサイトとする。
$\mathcal{S} \subset \mathcal{O}$ を集合の部分前層とし，すべての
$U \in \Ob(\mathcal{C})$ に対して
$\mathcal{S}(U) \subset \mathcal{O}(U)$ は乗法的部分集合であるとする。
《代数》の定義 \ref{algebra-definition-multiplicative-subset} を参照。
この場合，環の前層
$$
\mathcal{S}^{-1}\mathcal{O} :
U \longmapsto \mathcal{S}(U)^{-1}\mathcal{O}(U).
$$
を考えられる。制限写像は，切断
$f/s$（$f \in \mathcal{O}(U)$，$s \in \mathcal{S}(U)$）を，
$\mathcal{C}$ の $V \to U$ に対して $(f|_V)/(s|_V)$ へ写す。

\begin{lemma}
\label{sites-modules-lemma-simple-invert}
上の状況で，層化への写像
$$
\mathcal{O} \longrightarrow (\mathcal{S}^{-1}\mathcal{O})^\#
$$
は環の層の準同型であり，次の普遍性をもつ：
$\mathcal{S}$ の各局所切断を $\mathcal{A}$ の可逆な切断へ写す
任意の環の層の準同型 $\mathcal{O} \to \mathcal{A}$ に対して，
一意な分解
$(\mathcal{S}^{-1}\mathcal{O})^\# \to \mathcal{A}$ が存在する。
\end{lemma}

\begin{proof}
省略する。
\end{proof}

\noindent
$(\mathcal{C}, \mathcal{O})$ を環付きサイトとする。
$\mathcal{S} \subset \mathcal{O}$ を集合の部分前層とし，すべての
$U \in \mathcal{C}$ に対して
$\mathcal{S}(U) \subset \mathcal{O}(U)$ は乗法的部分集合であるとする。
$\mathcal{F}$ を $\mathcal{O}$-加群の層とする。
この場合，$\mathcal{S}^{-1}\mathcal{O}$-加群の前層
$$
\mathcal{S}^{-1}\mathcal{F} :
U \longmapsto \mathcal{S}(U)^{-1}\mathcal{F}(U).
$$
を考えられる。制限写像は，切断 $t/s$（$t \in \mathcal{F}(U)$，
$s \in \mathcal{S}(U)$）を，$V \to U$ が $\mathcal{C}$ の射ならば
$(t|_V)/(s|_V)$ へ写す。このとき
$\mathcal{S}^{-1}\mathcal{F}$ は
$\mathcal{S}^{-1}\mathcal{O}$-加群の前層である。

\begin{lemma}
\label{sites-modules-lemma-simple-invert-module}
上の状況で，層化への写像
$$
\mathcal{F} \longrightarrow (\mathcal{S}^{-1}\mathcal{F})^\#
$$
は次の普遍性をもつ：$\mathcal{S}$ の各局所切断が
$\mathcal{G}$ 上で可逆に作用するような任意の
$\mathcal{O}$-加群準同型 $\mathcal{F} \to \mathcal{G}$ に対して，
一意な分解
$(\mathcal{S}^{-1}\mathcal{F})^\# \to \mathcal{G}$ が存在する。
さらに
$$
(\mathcal{S}^{-1}\mathcal{F})^\#
=
(\mathcal{S}^{-1}\mathcal{O})^\# \otimes_\mathcal{O} \mathcal{F}
$$
が $(\mathcal{S}^{-1}\mathcal{O})^\#$-加群の層として成り立つ。
\end{lemma}

\begin{proof}
省略する。
\end{proof}









\section{基点付き集合の層}
\label{sites-modules-section-pointed}

\noindent
本節では，これまでアーベル層についてのみ述べてきた事実のうち，
基点付き集合の層に関するものをまとめる。

\medskip\noindent
{\it 基点付き集合}とは，集合 $S$ とその元 $0 \in S$ の対 $(S,0)$ である。
基点付き集合の射 $(S, 0) \to (S', 0')$ とは，単に $0$ を $0'$ へ写す
集合の写像 $S \to S'$ である。記法を濫用し，
「$S$ を基点付き集合とする」とは，$S$ に指定された元
$0 \in S$ が備わっていることを意味するものとする。
基点付き集合の層とは，集合の層 $\mathcal{F}$ に「基点指定」
$0 : * \to \mathcal{F}$ を備えたものにほかならない。
ここで $*$ は終層である
（《サイト》の例 \ref{sites-example-singleton-sheaf}）。

\medskip\noindent
サイトまたはトポスの射が与えられると，基点付き集合の層の圏上に
直像関手と引き戻し関手がある。《サイト》の
\ref{sites-section-sheaves-algebraic-structures} 節を参照。
これらは，基礎にある集合の層の直像または引き戻しを取り，
適切に基点を指定することで構成される
（終層の引き戻しが終層であることを用いる）。

\medskip\noindent
$u : \mathcal{C} \to \mathcal{D}$ をサイト間の連続かつ余連続な関手とし，
$g : \Sh(\mathcal{C}) \to \Sh(\mathcal{D})$ を $u$ に付随する
トポスの射とする。《サイト》の補題
\ref{sites-lemma-cocontinuous-morphism-topoi} を参照。
このとき基点付き集合の層上の $g^{-1}$ も左随伴 $g_!$ をもつ。
この関手の構成は，\ref{sites-modules-section-exactness-lower-shriek} 節における
アーベル層上の $g_!$ の構成とまったく同様である。

\medskip\noindent
同様に，$j : \mathcal{C}/U \to \mathcal{C}$ が
\ref{sites-modules-section-localize} 節のものであるならば，
基点付き集合の層上の関手 $j^{-1}$ は左随伴 $j_!$ をもつ。

\medskip\noindent
これらの事実と構成が必要になった場合には，対応する補題をここに
正確に述べて証明することにする。
